\documentclass[11pt]{article}
\usepackage[margin=1in]{geometry}
\usepackage{hyperref}
\usepackage[T1]{fontenc}
\usepackage[utf8]{inputenc}
\title{Expert Conjectures}
\author{Abducer}
\date{}
\begin{document}
\maketitle

\noindent
\textbf{KB ref:} p-vs-np\\
\textbf{Executed discovery branches:} 25\\
\textbf{Expert conjecture packages:} 90\\
\textbf{Ready for prove:} 90\\
\textbf{Ready for discovery:} 0\\
\textbf{Restricted-regime programmes:} 7

\section*{Judgment}
\begin{itemize}
\item Discovery now compiles branch executions into theorem packages rather than leaving them as abstract branch prompts.
\item Next move: Promote the ready-for-prove expert conjectures into the proof queue and keep the restricted-regime and foundation-search programmes on the discovery frontier.
\end{itemize}

\section*{Ready For Prove}
\begin{itemize}
\item conjpkg.disclayer-p-vs-np-pattern-invariant-discovery.10: There exists an admissible 'MCSP obstruction family' whose realization internalizes 'A hardness magnification theorem has the form:"A barely-nontrivial lower bound for an explicit problem Q in a weak model implies a major lower bound/separation in a strong model."Typically, Q is a meta-problem (e.g. a Gap-MCSP variant) and"major"means consequences such as NP not in P/poly.' while preserving 'MCSP complexity profile in nonuniform-circuits' strongly enough to support 'In refined regimes, lower bounds for Gap-MCSP slightly beyond known frontiers (e.g. N\textasciicircum{} 1+varepsilon -type circuit lower bounds for suitable promise gaps and input length N=2\textasciicircum{}n ) imply major consequences such as NP not in P/poly.'.
\item conjpkg.disclayer-p-vs-np-pattern-invariant-discovery.09: No admissible 'MCSP obstruction family' can simultaneously preserve 'MCSP complexity profile in nonuniform-circuits' and realize 'In refined regimes, lower bounds for Gap-MCSP slightly beyond known frontiers (e.g. N\textasciicircum{} 1+varepsilon -type circuit lower bounds for suitable promise gaps and input length N=2\textasciicircum{}n ) imply major consequences such as NP not in P/poly.' under the current route commitments.
\item conjpkg.disclayer-p-vs-np-pattern-invariant-discovery.08: There exists an admissible 'MCSP witness family' whose realization internalizes 'A hardness magnification theorem has the form:"A barely-nontrivial lower bound for an explicit problem Q in a weak model implies a major lower bound/separation in a strong model."Typically, Q is a meta-problem (e.g. a Gap-MCSP variant) and"major"means consequences such as NP not in P/poly.' while preserving 'MCSP complexity profile in nonuniform-circuits' strongly enough to support 'There exist explicit"meta"problems (including MCSP variants) such that proving certain emph barely-superlinear lower bounds against restricted models (e.g. formulas) would imply major lower bounds for NP in much stronger models (including superpolynomial formula/circuit lower bounds, and in some parameter regimes consequences as strong as NP not in P/poly ).'.
\item conjpkg.disclayer-p-vs-np-pattern-invariant-discovery.07: No admissible 'MCSP witness family' can simultaneously preserve 'MCSP complexity profile in nonuniform-circuits' and realize 'There exist explicit"meta"problems (including MCSP variants) such that proving certain emph barely-superlinear lower bounds against restricted models (e.g. formulas) would imply major lower bounds for NP in much stronger models (including superpolynomial formula/circuit lower bounds, and in some parameter regimes consequences as strong as NP not in P/poly ).' under the current route commitments.
\item conjpkg.disclayer-p-vs-np-pattern-invariant-discovery.06: For every admissible 'MCSP witness family', there is a transport law from 'nonuniform-circuits' to 'formula-size' that preserves 'MCSP complexity profile in nonuniform-circuits' and retains enough strength to activate 'There exist explicit"meta"problems (including MCSP variants) such that proving certain emph barely-superlinear lower bounds against restricted models (e.g. formulas) would imply major lower bounds for NP in much stronger models (including superpolynomial formula/circuit lower bounds, and in some parameter regimes consequences as strong as NP not in P/poly ).'.
\item conjpkg.disclayer-p-vs-np-pattern-invariant-discovery.05: There exists an admissible 'MCSP witness family' whose realization internalizes 'A hardness magnification theorem has the form:"A barely-nontrivial lower bound for an explicit problem Q in a weak model implies a major lower bound/separation in a strong model."Typically, Q is a meta-problem (e.g. a Gap-MCSP variant) and"major"means consequences such as NP not in P/poly.' while preserving 'MCSP complexity profile in nonuniform-circuits' strongly enough to support 'In refined regimes, lower bounds for Gap-MCSP slightly beyond known frontiers (e.g. N\textasciicircum{} 1+varepsilon -type circuit lower bounds for suitable promise gaps and input length N=2\textasciicircum{}n ) imply major consequences such as NP not in P/poly.'.
\item conjpkg.disclayer-p-vs-np-pattern-invariant-discovery.04: No admissible 'MCSP witness family' can simultaneously preserve 'MCSP complexity profile in nonuniform-circuits' and realize 'In refined regimes, lower bounds for Gap-MCSP slightly beyond known frontiers (e.g. N\textasciicircum{} 1+varepsilon -type circuit lower bounds for suitable promise gaps and input length N=2\textasciicircum{}n ) imply major consequences such as NP not in P/poly.' under the current route commitments.
\item conjpkg.disclayer-p-vs-np-pattern-invariant-discovery.03: There exists an admissible 'MCSP obstruction family' such that 'A hardness magnification theorem has the form:"A barely-nontrivial lower bound for an explicit problem Q in a weak model implies a major lower bound/separation in a strong model."Typically, Q is a meta-problem (e.g. a Gap-MCSP variant) and"major"means consequences such as NP not in P/poly.' can be organized into a route preserving 'MCSP complexity profile in nonuniform-circuits' and thereby supports 'In refined regimes, lower bounds for Gap-MCSP slightly beyond known frontiers (e.g. N\textasciicircum{} 1+varepsilon -type circuit lower bounds for suitable promise gaps and input length N=2\textasciicircum{}n ) imply major consequences such as NP not in P/poly.'.
\item conjpkg.disclayer-p-vs-np-pattern-invariant-discovery.02: There exists an admissible 'MCSP witness family' such that 'A hardness magnification theorem has the form:"A barely-nontrivial lower bound for an explicit problem Q in a weak model implies a major lower bound/separation in a strong model."Typically, Q is a meta-problem (e.g. a Gap-MCSP variant) and"major"means consequences such as NP not in P/poly.' can be organized into a route preserving 'MCSP complexity profile in nonuniform-circuits' and thereby supports 'There exist explicit"meta"problems (including MCSP variants) such that proving certain emph barely-superlinear lower bounds against restricted models (e.g. formulas) would imply major lower bounds for NP in much stronger models (including superpolynomial formula/circuit lower bounds, and in some parameter regimes consequences as strong as NP not in P/poly ).'.
\item conjpkg.disclayer-p-vs-np-pattern-invariant-discovery.01: There exists an admissible 'MCSP witness family' such that 'A hardness magnification theorem has the form:"A barely-nontrivial lower bound for an explicit problem Q in a weak model implies a major lower bound/separation in a strong model."Typically, Q is a meta-problem (e.g. a Gap-MCSP variant) and"major"means consequences such as NP not in P/poly.' can be organized into a route preserving 'MCSP complexity profile in nonuniform-circuits' and thereby supports 'In refined regimes, lower bounds for Gap-MCSP slightly beyond known frontiers (e.g. N\textasciicircum{} 1+varepsilon -type circuit lower bounds for suitable promise gaps and input length N=2\textasciicircum{}n ) imply major consequences such as NP not in P/poly.'.
\item conjpkg.disclayer-p-vs-np-pattern-regime-localization.10: There exists an admissible 'MCSP obstruction family' whose realization internalizes 'A hardness magnification theorem has the form:"A barely-nontrivial lower bound for an explicit problem Q in a weak model implies a major lower bound/separation in a strong model."Typically, Q is a meta-problem (e.g. a Gap-MCSP variant) and"major"means consequences such as NP not in P/poly.' while preserving 'MCSP complexity profile in nonuniform-circuits' strongly enough to support 'In refined regimes, lower bounds for Gap-MCSP slightly beyond known frontiers (e.g. N\textasciicircum{} 1+varepsilon -type circuit lower bounds for suitable promise gaps and input length N=2\textasciicircum{}n ) imply major consequences such as NP not in P/poly.'.
\item conjpkg.disclayer-p-vs-np-pattern-regime-localization.09: No admissible 'MCSP obstruction family' can simultaneously preserve 'MCSP complexity profile in nonuniform-circuits' and realize 'In refined regimes, lower bounds for Gap-MCSP slightly beyond known frontiers (e.g. N\textasciicircum{} 1+varepsilon -type circuit lower bounds for suitable promise gaps and input length N=2\textasciicircum{}n ) imply major consequences such as NP not in P/poly.' under the current route commitments.
\item conjpkg.disclayer-p-vs-np-pattern-regime-localization.08: There exists an admissible 'MCSP witness family' whose realization internalizes 'A hardness magnification theorem has the form:"A barely-nontrivial lower bound for an explicit problem Q in a weak model implies a major lower bound/separation in a strong model."Typically, Q is a meta-problem (e.g. a Gap-MCSP variant) and"major"means consequences such as NP not in P/poly.' while preserving 'MCSP complexity profile in nonuniform-circuits' strongly enough to support 'There exist explicit"meta"problems (including MCSP variants) such that proving certain emph barely-superlinear lower bounds against restricted models (e.g. formulas) would imply major lower bounds for NP in much stronger models (including superpolynomial formula/circuit lower bounds, and in some parameter regimes consequences as strong as NP not in P/poly ).'.
\item conjpkg.disclayer-p-vs-np-pattern-regime-localization.07: No admissible 'MCSP witness family' can simultaneously preserve 'MCSP complexity profile in nonuniform-circuits' and realize 'There exist explicit"meta"problems (including MCSP variants) such that proving certain emph barely-superlinear lower bounds against restricted models (e.g. formulas) would imply major lower bounds for NP in much stronger models (including superpolynomial formula/circuit lower bounds, and in some parameter regimes consequences as strong as NP not in P/poly ).' under the current route commitments.
\item conjpkg.disclayer-p-vs-np-pattern-regime-localization.06: For every admissible 'MCSP witness family', there is a transport law from 'nonuniform-circuits' to 'formula-size' that preserves 'MCSP complexity profile in nonuniform-circuits' and retains enough strength to activate 'There exist explicit"meta"problems (including MCSP variants) such that proving certain emph barely-superlinear lower bounds against restricted models (e.g. formulas) would imply major lower bounds for NP in much stronger models (including superpolynomial formula/circuit lower bounds, and in some parameter regimes consequences as strong as NP not in P/poly ).'.
\item conjpkg.disclayer-p-vs-np-pattern-regime-localization.05: There exists an admissible 'MCSP witness family' whose realization internalizes 'A hardness magnification theorem has the form:"A barely-nontrivial lower bound for an explicit problem Q in a weak model implies a major lower bound/separation in a strong model."Typically, Q is a meta-problem (e.g. a Gap-MCSP variant) and"major"means consequences such as NP not in P/poly.' while preserving 'MCSP complexity profile in nonuniform-circuits' strongly enough to support 'In refined regimes, lower bounds for Gap-MCSP slightly beyond known frontiers (e.g. N\textasciicircum{} 1+varepsilon -type circuit lower bounds for suitable promise gaps and input length N=2\textasciicircum{}n ) imply major consequences such as NP not in P/poly.'.
\item conjpkg.disclayer-p-vs-np-pattern-regime-localization.04: No admissible 'MCSP witness family' can simultaneously preserve 'MCSP complexity profile in nonuniform-circuits' and realize 'In refined regimes, lower bounds for Gap-MCSP slightly beyond known frontiers (e.g. N\textasciicircum{} 1+varepsilon -type circuit lower bounds for suitable promise gaps and input length N=2\textasciicircum{}n ) imply major consequences such as NP not in P/poly.' under the current route commitments.
\item conjpkg.disclayer-p-vs-np-pattern-regime-localization.03: There exists an admissible 'MCSP obstruction family' such that 'A hardness magnification theorem has the form:"A barely-nontrivial lower bound for an explicit problem Q in a weak model implies a major lower bound/separation in a strong model."Typically, Q is a meta-problem (e.g. a Gap-MCSP variant) and"major"means consequences such as NP not in P/poly.' can be organized into a route preserving 'MCSP complexity profile in nonuniform-circuits' and thereby supports 'In refined regimes, lower bounds for Gap-MCSP slightly beyond known frontiers (e.g. N\textasciicircum{} 1+varepsilon -type circuit lower bounds for suitable promise gaps and input length N=2\textasciicircum{}n ) imply major consequences such as NP not in P/poly.'.
\item conjpkg.disclayer-p-vs-np-pattern-regime-localization.02: There exists an admissible 'MCSP witness family' such that 'A hardness magnification theorem has the form:"A barely-nontrivial lower bound for an explicit problem Q in a weak model implies a major lower bound/separation in a strong model."Typically, Q is a meta-problem (e.g. a Gap-MCSP variant) and"major"means consequences such as NP not in P/poly.' can be organized into a route preserving 'MCSP complexity profile in nonuniform-circuits' and thereby supports 'There exist explicit"meta"problems (including MCSP variants) such that proving certain emph barely-superlinear lower bounds against restricted models (e.g. formulas) would imply major lower bounds for NP in much stronger models (including superpolynomial formula/circuit lower bounds, and in some parameter regimes consequences as strong as NP not in P/poly ).'.
\item conjpkg.disclayer-p-vs-np-pattern-regime-localization.01: There exists an admissible 'MCSP witness family' such that 'A hardness magnification theorem has the form:"A barely-nontrivial lower bound for an explicit problem Q in a weak model implies a major lower bound/separation in a strong model."Typically, Q is a meta-problem (e.g. a Gap-MCSP variant) and"major"means consequences such as NP not in P/poly.' can be organized into a route preserving 'MCSP complexity profile in nonuniform-circuits' and thereby supports 'In refined regimes, lower bounds for Gap-MCSP slightly beyond known frontiers (e.g. N\textasciicircum{} 1+varepsilon -type circuit lower bounds for suitable promise gaps and input length N=2\textasciicircum{}n ) imply major consequences such as NP not in P/poly.'.
\item conjpkg.disclayer-p-vs-np-pattern-hybridization-of-partial-bridges.13: There exists an intermediate typed object built from 'MCSP obstruction family' such that the route from 'A hardness magnification theorem has the form:"A barely-nontrivial lower bound for an explicit problem Q in a weak model implies a major lower bound/separation in a strong model."Typically, Q is a meta-problem (e.g. a Gap-MCSP variant) and"major"means consequences such as NP not in P/poly.' to 'In refined regimes, lower bounds for Gap-MCSP slightly beyond known frontiers (e.g. N\textasciicircum{} 1+varepsilon -type circuit lower bounds for suitable promise gaps and input length N=2\textasciicircum{}n ) imply major consequences such as NP not in P/poly.' factors through theorem-sized maps preserving 'MCSP complexity profile in nonuniform-circuits'.
\item conjpkg.disclayer-p-vs-np-pattern-hybridization-of-partial-bridges.12: There exists an admissible 'MCSP obstruction family' whose realization internalizes 'A hardness magnification theorem has the form:"A barely-nontrivial lower bound for an explicit problem Q in a weak model implies a major lower bound/separation in a strong model."Typically, Q is a meta-problem (e.g. a Gap-MCSP variant) and"major"means consequences such as NP not in P/poly.' while preserving 'MCSP complexity profile in nonuniform-circuits' strongly enough to support 'In refined regimes, lower bounds for Gap-MCSP slightly beyond known frontiers (e.g. N\textasciicircum{} 1+varepsilon -type circuit lower bounds for suitable promise gaps and input length N=2\textasciicircum{}n ) imply major consequences such as NP not in P/poly.'.
\item conjpkg.disclayer-p-vs-np-pattern-hybridization-of-partial-bridges.11: No admissible 'MCSP obstruction family' can simultaneously preserve 'MCSP complexity profile in nonuniform-circuits' and realize 'In refined regimes, lower bounds for Gap-MCSP slightly beyond known frontiers (e.g. N\textasciicircum{} 1+varepsilon -type circuit lower bounds for suitable promise gaps and input length N=2\textasciicircum{}n ) imply major consequences such as NP not in P/poly.' under the current route commitments.
\item conjpkg.disclayer-p-vs-np-pattern-hybridization-of-partial-bridges.10: There exists an intermediate typed object built from 'MCSP witness family' such that the route from 'A hardness magnification theorem has the form:"A barely-nontrivial lower bound for an explicit problem Q in a weak model implies a major lower bound/separation in a strong model."Typically, Q is a meta-problem (e.g. a Gap-MCSP variant) and"major"means consequences such as NP not in P/poly.' to 'There exist explicit"meta"problems (including MCSP variants) such that proving certain emph barely-superlinear lower bounds against restricted models (e.g. formulas) would imply major lower bounds for NP in much stronger models (including superpolynomial formula/circuit lower bounds, and in some parameter regimes consequences as strong as NP not in P/poly ).' factors through theorem-sized maps preserving 'MCSP complexity profile in nonuniform-circuits'.
\item conjpkg.disclayer-p-vs-np-pattern-hybridization-of-partial-bridges.09: There exists an admissible 'MCSP witness family' whose realization internalizes 'A hardness magnification theorem has the form:"A barely-nontrivial lower bound for an explicit problem Q in a weak model implies a major lower bound/separation in a strong model."Typically, Q is a meta-problem (e.g. a Gap-MCSP variant) and"major"means consequences such as NP not in P/poly.' while preserving 'MCSP complexity profile in nonuniform-circuits' strongly enough to support 'There exist explicit"meta"problems (including MCSP variants) such that proving certain emph barely-superlinear lower bounds against restricted models (e.g. formulas) would imply major lower bounds for NP in much stronger models (including superpolynomial formula/circuit lower bounds, and in some parameter regimes consequences as strong as NP not in P/poly ).'.
\item conjpkg.disclayer-p-vs-np-pattern-hybridization-of-partial-bridges.08: No admissible 'MCSP witness family' can simultaneously preserve 'MCSP complexity profile in nonuniform-circuits' and realize 'There exist explicit"meta"problems (including MCSP variants) such that proving certain emph barely-superlinear lower bounds against restricted models (e.g. formulas) would imply major lower bounds for NP in much stronger models (including superpolynomial formula/circuit lower bounds, and in some parameter regimes consequences as strong as NP not in P/poly ).' under the current route commitments.
\item conjpkg.disclayer-p-vs-np-pattern-hybridization-of-partial-bridges.07: For every admissible 'MCSP witness family', there is a transport law from 'nonuniform-circuits' to 'formula-size' that preserves 'MCSP complexity profile in nonuniform-circuits' and retains enough strength to activate 'There exist explicit"meta"problems (including MCSP variants) such that proving certain emph barely-superlinear lower bounds against restricted models (e.g. formulas) would imply major lower bounds for NP in much stronger models (including superpolynomial formula/circuit lower bounds, and in some parameter regimes consequences as strong as NP not in P/poly ).'.
\item conjpkg.disclayer-p-vs-np-pattern-hybridization-of-partial-bridges.06: There exists an intermediate typed object built from 'MCSP witness family' such that the route from 'A hardness magnification theorem has the form:"A barely-nontrivial lower bound for an explicit problem Q in a weak model implies a major lower bound/separation in a strong model."Typically, Q is a meta-problem (e.g. a Gap-MCSP variant) and"major"means consequences such as NP not in P/poly.' to 'In refined regimes, lower bounds for Gap-MCSP slightly beyond known frontiers (e.g. N\textasciicircum{} 1+varepsilon -type circuit lower bounds for suitable promise gaps and input length N=2\textasciicircum{}n ) imply major consequences such as NP not in P/poly.' factors through theorem-sized maps preserving 'MCSP complexity profile in nonuniform-circuits'.
\item conjpkg.disclayer-p-vs-np-pattern-hybridization-of-partial-bridges.05: There exists an admissible 'MCSP witness family' whose realization internalizes 'A hardness magnification theorem has the form:"A barely-nontrivial lower bound for an explicit problem Q in a weak model implies a major lower bound/separation in a strong model."Typically, Q is a meta-problem (e.g. a Gap-MCSP variant) and"major"means consequences such as NP not in P/poly.' while preserving 'MCSP complexity profile in nonuniform-circuits' strongly enough to support 'In refined regimes, lower bounds for Gap-MCSP slightly beyond known frontiers (e.g. N\textasciicircum{} 1+varepsilon -type circuit lower bounds for suitable promise gaps and input length N=2\textasciicircum{}n ) imply major consequences such as NP not in P/poly.'.
\item conjpkg.disclayer-p-vs-np-pattern-hybridization-of-partial-bridges.04: No admissible 'MCSP witness family' can simultaneously preserve 'MCSP complexity profile in nonuniform-circuits' and realize 'In refined regimes, lower bounds for Gap-MCSP slightly beyond known frontiers (e.g. N\textasciicircum{} 1+varepsilon -type circuit lower bounds for suitable promise gaps and input length N=2\textasciicircum{}n ) imply major consequences such as NP not in P/poly.' under the current route commitments.
\item conjpkg.disclayer-p-vs-np-pattern-hybridization-of-partial-bridges.03: There exists an admissible 'MCSP obstruction family' such that 'A hardness magnification theorem has the form:"A barely-nontrivial lower bound for an explicit problem Q in a weak model implies a major lower bound/separation in a strong model."Typically, Q is a meta-problem (e.g. a Gap-MCSP variant) and"major"means consequences such as NP not in P/poly.' can be organized into a route preserving 'MCSP complexity profile in nonuniform-circuits' and thereby supports 'In refined regimes, lower bounds for Gap-MCSP slightly beyond known frontiers (e.g. N\textasciicircum{} 1+varepsilon -type circuit lower bounds for suitable promise gaps and input length N=2\textasciicircum{}n ) imply major consequences such as NP not in P/poly.'.
\item conjpkg.disclayer-p-vs-np-pattern-hybridization-of-partial-bridges.02: There exists an admissible 'MCSP witness family' such that 'A hardness magnification theorem has the form:"A barely-nontrivial lower bound for an explicit problem Q in a weak model implies a major lower bound/separation in a strong model."Typically, Q is a meta-problem (e.g. a Gap-MCSP variant) and"major"means consequences such as NP not in P/poly.' can be organized into a route preserving 'MCSP complexity profile in nonuniform-circuits' and thereby supports 'There exist explicit"meta"problems (including MCSP variants) such that proving certain emph barely-superlinear lower bounds against restricted models (e.g. formulas) would imply major lower bounds for NP in much stronger models (including superpolynomial formula/circuit lower bounds, and in some parameter regimes consequences as strong as NP not in P/poly ).'.
\item conjpkg.disclayer-p-vs-np-pattern-hybridization-of-partial-bridges.01: There exists an admissible 'MCSP witness family' such that 'A hardness magnification theorem has the form:"A barely-nontrivial lower bound for an explicit problem Q in a weak model implies a major lower bound/separation in a strong model."Typically, Q is a meta-problem (e.g. a Gap-MCSP variant) and"major"means consequences such as NP not in P/poly.' can be organized into a route preserving 'MCSP complexity profile in nonuniform-circuits' and thereby supports 'In refined regimes, lower bounds for Gap-MCSP slightly beyond known frontiers (e.g. N\textasciicircum{} 1+varepsilon -type circuit lower bounds for suitable promise gaps and input length N=2\textasciicircum{}n ) imply major consequences such as NP not in P/poly.'.
\item conjpkg.disclayer-p-vs-np-family-abstraction-search.10: There exists an admissible 'MCSP obstruction family' whose realization internalizes 'A hardness magnification theorem has the form:"A barely-nontrivial lower bound for an explicit problem Q in a weak model implies a major lower bound/separation in a strong model."Typically, Q is a meta-problem (e.g. a Gap-MCSP variant) and"major"means consequences such as NP not in P/poly.' while preserving 'MCSP complexity profile in nonuniform-circuits' strongly enough to support 'In refined regimes, lower bounds for Gap-MCSP slightly beyond known frontiers (e.g. N\textasciicircum{} 1+varepsilon -type circuit lower bounds for suitable promise gaps and input length N=2\textasciicircum{}n ) imply major consequences such as NP not in P/poly.'.
\item conjpkg.disclayer-p-vs-np-family-abstraction-search.09: No admissible 'MCSP obstruction family' can simultaneously preserve 'MCSP complexity profile in nonuniform-circuits' and realize 'In refined regimes, lower bounds for Gap-MCSP slightly beyond known frontiers (e.g. N\textasciicircum{} 1+varepsilon -type circuit lower bounds for suitable promise gaps and input length N=2\textasciicircum{}n ) imply major consequences such as NP not in P/poly.' under the current route commitments.
\item conjpkg.disclayer-p-vs-np-family-abstraction-search.08: There exists an admissible 'MCSP witness family' whose realization internalizes 'A hardness magnification theorem has the form:"A barely-nontrivial lower bound for an explicit problem Q in a weak model implies a major lower bound/separation in a strong model."Typically, Q is a meta-problem (e.g. a Gap-MCSP variant) and"major"means consequences such as NP not in P/poly.' while preserving 'MCSP complexity profile in nonuniform-circuits' strongly enough to support 'There exist explicit"meta"problems (including MCSP variants) such that proving certain emph barely-superlinear lower bounds against restricted models (e.g. formulas) would imply major lower bounds for NP in much stronger models (including superpolynomial formula/circuit lower bounds, and in some parameter regimes consequences as strong as NP not in P/poly ).'.
\item conjpkg.disclayer-p-vs-np-family-abstraction-search.07: No admissible 'MCSP witness family' can simultaneously preserve 'MCSP complexity profile in nonuniform-circuits' and realize 'There exist explicit"meta"problems (including MCSP variants) such that proving certain emph barely-superlinear lower bounds against restricted models (e.g. formulas) would imply major lower bounds for NP in much stronger models (including superpolynomial formula/circuit lower bounds, and in some parameter regimes consequences as strong as NP not in P/poly ).' under the current route commitments.
\item conjpkg.disclayer-p-vs-np-family-abstraction-search.06: For every admissible 'MCSP witness family', there is a transport law from 'nonuniform-circuits' to 'formula-size' that preserves 'MCSP complexity profile in nonuniform-circuits' and retains enough strength to activate 'There exist explicit"meta"problems (including MCSP variants) such that proving certain emph barely-superlinear lower bounds against restricted models (e.g. formulas) would imply major lower bounds for NP in much stronger models (including superpolynomial formula/circuit lower bounds, and in some parameter regimes consequences as strong as NP not in P/poly ).'.
\item conjpkg.disclayer-p-vs-np-family-abstraction-search.05: There exists an admissible 'MCSP witness family' whose realization internalizes 'A hardness magnification theorem has the form:"A barely-nontrivial lower bound for an explicit problem Q in a weak model implies a major lower bound/separation in a strong model."Typically, Q is a meta-problem (e.g. a Gap-MCSP variant) and"major"means consequences such as NP not in P/poly.' while preserving 'MCSP complexity profile in nonuniform-circuits' strongly enough to support 'In refined regimes, lower bounds for Gap-MCSP slightly beyond known frontiers (e.g. N\textasciicircum{} 1+varepsilon -type circuit lower bounds for suitable promise gaps and input length N=2\textasciicircum{}n ) imply major consequences such as NP not in P/poly.'.
\item conjpkg.disclayer-p-vs-np-family-abstraction-search.04: No admissible 'MCSP witness family' can simultaneously preserve 'MCSP complexity profile in nonuniform-circuits' and realize 'In refined regimes, lower bounds for Gap-MCSP slightly beyond known frontiers (e.g. N\textasciicircum{} 1+varepsilon -type circuit lower bounds for suitable promise gaps and input length N=2\textasciicircum{}n ) imply major consequences such as NP not in P/poly.' under the current route commitments.
\item conjpkg.disclayer-p-vs-np-family-abstraction-search.03: There exists an admissible 'MCSP obstruction family' such that 'A hardness magnification theorem has the form:"A barely-nontrivial lower bound for an explicit problem Q in a weak model implies a major lower bound/separation in a strong model."Typically, Q is a meta-problem (e.g. a Gap-MCSP variant) and"major"means consequences such as NP not in P/poly.' can be organized into a route preserving 'MCSP complexity profile in nonuniform-circuits' and thereby supports 'In refined regimes, lower bounds for Gap-MCSP slightly beyond known frontiers (e.g. N\textasciicircum{} 1+varepsilon -type circuit lower bounds for suitable promise gaps and input length N=2\textasciicircum{}n ) imply major consequences such as NP not in P/poly.'.
\item conjpkg.disclayer-p-vs-np-family-abstraction-search.02: There exists an admissible 'MCSP witness family' such that 'A hardness magnification theorem has the form:"A barely-nontrivial lower bound for an explicit problem Q in a weak model implies a major lower bound/separation in a strong model."Typically, Q is a meta-problem (e.g. a Gap-MCSP variant) and"major"means consequences such as NP not in P/poly.' can be organized into a route preserving 'MCSP complexity profile in nonuniform-circuits' and thereby supports 'There exist explicit"meta"problems (including MCSP variants) such that proving certain emph barely-superlinear lower bounds against restricted models (e.g. formulas) would imply major lower bounds for NP in much stronger models (including superpolynomial formula/circuit lower bounds, and in some parameter regimes consequences as strong as NP not in P/poly ).'.
\item conjpkg.disclayer-p-vs-np-family-abstraction-search.01: There exists an admissible 'MCSP witness family' such that 'A hardness magnification theorem has the form:"A barely-nontrivial lower bound for an explicit problem Q in a weak model implies a major lower bound/separation in a strong model."Typically, Q is a meta-problem (e.g. a Gap-MCSP variant) and"major"means consequences such as NP not in P/poly.' can be organized into a route preserving 'MCSP complexity profile in nonuniform-circuits' and thereby supports 'In refined regimes, lower bounds for Gap-MCSP slightly beyond known frontiers (e.g. N\textasciicircum{} 1+varepsilon -type circuit lower bounds for suitable promise gaps and input length N=2\textasciicircum{}n ) imply major consequences such as NP not in P/poly.'.
\item conjpkg.disclayer-p-vs-np-family-barrier-aware-transfer.06: There exists an admissible 'unspecified reduction gadget family' whose realization internalizes 'There exist many problems L in combinatorial optimisation/graph theory such that L is NP-complete under polynomial-time many-one reductions.' while preserving 'consequence schema NP-complete-classification' strongly enough to support 'Under"local"polynomial-time reductions (where each output bit is computable with very small time), MCSP is not NP-complete. More broadly, the paper develops evidence and conditional consequences suggesting NP-hardness of MCSP under various reductions would entail unexpected class separations.'.
\item conjpkg.disclayer-p-vs-np-family-barrier-aware-transfer.05: No admissible 'unspecified reduction gadget family' can simultaneously preserve 'consequence schema NP-complete-classification' and realize 'Under"local"polynomial-time reductions (where each output bit is computable with very small time), MCSP is not NP-complete. More broadly, the paper develops evidence and conditional consequences suggesting NP-hardness of MCSP under various reductions would entail unexpected class separations.' under the current route commitments.
\item conjpkg.disclayer-p-vs-np-family-barrier-aware-transfer.04: There exists an admissible 'unspecified witness family' whose realization internalizes 'There exist many problems L in combinatorial optimisation/graph theory such that L is NP-complete under polynomial-time many-one reductions.' while preserving 'consequence schema NP-complete-classification' strongly enough to support 'Under"local"polynomial-time reductions (where each output bit is computable with very small time), MCSP is not NP-complete. More broadly, the paper develops evidence and conditional consequences suggesting NP-hardness of MCSP under various reductions would entail unexpected class separations.'.
\item conjpkg.disclayer-p-vs-np-family-barrier-aware-transfer.03: No admissible 'unspecified witness family' can simultaneously preserve 'consequence schema NP-complete-classification' and realize 'Under"local"polynomial-time reductions (where each output bit is computable with very small time), MCSP is not NP-complete. More broadly, the paper develops evidence and conditional consequences suggesting NP-hardness of MCSP under various reductions would entail unexpected class separations.' under the current route commitments.
\item conjpkg.disclayer-p-vs-np-family-barrier-aware-transfer.02: There exists an admissible 'unspecified reduction gadget family' such that 'There exist many problems L in combinatorial optimisation/graph theory such that L is NP-complete under polynomial-time many-one reductions.' can be organized into a route preserving 'consequence schema NP-complete-classification' and thereby supports 'Under"local"polynomial-time reductions (where each output bit is computable with very small time), MCSP is not NP-complete. More broadly, the paper develops evidence and conditional consequences suggesting NP-hardness of MCSP under various reductions would entail unexpected class separations.'.
\item conjpkg.disclayer-p-vs-np-family-barrier-aware-transfer.01: There exists an admissible 'unspecified witness family' such that 'There exist many problems L in combinatorial optimisation/graph theory such that L is NP-complete under polynomial-time many-one reductions.' can be organized into a route preserving 'consequence schema NP-complete-classification' and thereby supports 'Under"local"polynomial-time reductions (where each output bit is computable with very small time), MCSP is not NP-complete. More broadly, the paper develops evidence and conditional consequences suggesting NP-hardness of MCSP under various reductions would entail unexpected class separations.'.
\item conjpkg.disclayer-p-vs-np-family-explicitness.08: There exists an admissible 'PIT witness family' whose realization internalizes 'A set Hsubseteq mathbb F \textasciicircum{}n is a hitting set for a class mathcal C of n -variate circuits if for every nonzero polynomial p computed by some circuit in mathcal C , there exists mathbf a in H with p(mathbf a )neq 0.' while preserving 'PIT semantic object' strongly enough to support 'If mathrm PIT in Pclass (or even suitable infinitely-often derandomisations in nondeterministic subexponential time, in standard variants), then at least one of the following holds: NEXP not in P/poly quad or quad Permanent does not have polynomial-size arithmetic circuits.'.
\item conjpkg.disclayer-p-vs-np-family-explicitness.07: No admissible 'PIT witness family' can simultaneously preserve 'PIT semantic object' and realize 'If mathrm PIT in Pclass (or even suitable infinitely-often derandomisations in nondeterministic subexponential time, in standard variants), then at least one of the following holds: NEXP not in P/poly quad or quad Permanent does not have polynomial-size arithmetic circuits.' under the current route commitments.
\item conjpkg.disclayer-p-vs-np-family-explicitness.06: For every admissible 'PIT witness family', there is a transport law from 'black-box-PIT' to 'arithmetic-circuits' that preserves 'PIT semantic object' and retains enough strength to activate 'If mathrm PIT in Pclass (or even suitable infinitely-often derandomisations in nondeterministic subexponential time, in standard variants), then at least one of the following holds: NEXP not in P/poly quad or quad Permanent does not have polynomial-size arithmetic circuits.'.
\item conjpkg.disclayer-p-vs-np-family-explicitness.05: There exists an admissible 'PIT certificate family' whose realization internalizes 'A set Hsubseteq mathbb F \textasciicircum{}n is a hitting set for a class mathcal C of n -variate circuits if for every nonzero polynomial p computed by some circuit in mathcal C , there exists mathbf a in H with p(mathbf a )neq 0.' while preserving 'PIT semantic object' strongly enough to support 'If mathrm PIT in Pclass (or even suitable infinitely-often derandomisations in nondeterministic subexponential time, in standard variants), then at least one of the following holds: NEXP not in P/poly quad or quad Permanent does not have polynomial-size arithmetic circuits.'.
\item conjpkg.disclayer-p-vs-np-family-explicitness.04: No admissible 'PIT certificate family' can simultaneously preserve 'PIT semantic object' and realize 'If mathrm PIT in Pclass (or even suitable infinitely-often derandomisations in nondeterministic subexponential time, in standard variants), then at least one of the following holds: NEXP not in P/poly quad or quad Permanent does not have polynomial-size arithmetic circuits.' under the current route commitments.
\item conjpkg.disclayer-p-vs-np-family-explicitness.03: For every admissible 'PIT certificate family', there is a transport law from 'black-box-PIT' to 'arithmetic-circuits' that preserves 'PIT semantic object' and retains enough strength to activate 'If mathrm PIT in Pclass (or even suitable infinitely-often derandomisations in nondeterministic subexponential time, in standard variants), then at least one of the following holds: NEXP not in P/poly quad or quad Permanent does not have polynomial-size arithmetic circuits.'.
\item conjpkg.disclayer-p-vs-np-family-explicitness.02: There exists an admissible 'PIT witness family' such that 'A set Hsubseteq mathbb F \textasciicircum{}n is a hitting set for a class mathcal C of n -variate circuits if for every nonzero polynomial p computed by some circuit in mathcal C , there exists mathbf a in H with p(mathbf a )neq 0.' can be organized into a route preserving 'PIT semantic object' and thereby supports 'If mathrm PIT in Pclass (or even suitable infinitely-often derandomisations in nondeterministic subexponential time, in standard variants), then at least one of the following holds: NEXP not in P/poly quad or quad Permanent does not have polynomial-size arithmetic circuits.'.
\item conjpkg.disclayer-p-vs-np-family-explicitness.01: There exists an admissible 'PIT certificate family' such that 'A set Hsubseteq mathbb F \textasciicircum{}n is a hitting set for a class mathcal C of n -variate circuits if for every nonzero polynomial p computed by some circuit in mathcal C , there exists mathbf a in H with p(mathbf a )neq 0.' can be organized into a route preserving 'PIT semantic object' and thereby supports 'If mathrm PIT in Pclass (or even suitable infinitely-often derandomisations in nondeterministic subexponential time, in standard variants), then at least one of the following holds: NEXP not in P/poly quad or quad Permanent does not have polynomial-size arithmetic circuits.'.
\item conjpkg.disclayer-p-vs-np-family-categorical-reframing.13: There exists an admissible 'MCSP obstruction family' whose realization internalizes 'A hardness magnification theorem has the form:"A barely-nontrivial lower bound for an explicit problem Q in a weak model implies a major lower bound/separation in a strong model."Typically, Q is a meta-problem (e.g. a Gap-MCSP variant) and"major"means consequences such as NP not in P/poly.' while preserving 'MCSP complexity profile in nonuniform-circuits' strongly enough to support 'In refined regimes, lower bounds for Gap-MCSP slightly beyond known frontiers (e.g. N\textasciicircum{} 1+varepsilon -type circuit lower bounds for suitable promise gaps and input length N=2\textasciicircum{}n ) imply major consequences such as NP not in P/poly.'.
\item conjpkg.disclayer-p-vs-np-family-categorical-reframing.12: There exists an intermediate typed object built from 'MCSP obstruction family' such that the route from 'A hardness magnification theorem has the form:"A barely-nontrivial lower bound for an explicit problem Q in a weak model implies a major lower bound/separation in a strong model."Typically, Q is a meta-problem (e.g. a Gap-MCSP variant) and"major"means consequences such as NP not in P/poly.' to 'In refined regimes, lower bounds for Gap-MCSP slightly beyond known frontiers (e.g. N\textasciicircum{} 1+varepsilon -type circuit lower bounds for suitable promise gaps and input length N=2\textasciicircum{}n ) imply major consequences such as NP not in P/poly.' factors through theorem-sized maps preserving 'MCSP complexity profile in nonuniform-circuits'.
\item conjpkg.disclayer-p-vs-np-family-categorical-reframing.11: No admissible 'MCSP obstruction family' can simultaneously preserve 'MCSP complexity profile in nonuniform-circuits' and realize 'In refined regimes, lower bounds for Gap-MCSP slightly beyond known frontiers (e.g. N\textasciicircum{} 1+varepsilon -type circuit lower bounds for suitable promise gaps and input length N=2\textasciicircum{}n ) imply major consequences such as NP not in P/poly.' under the current route commitments.
\item conjpkg.disclayer-p-vs-np-family-categorical-reframing.10: There exists an admissible 'MCSP witness family' whose realization internalizes 'A hardness magnification theorem has the form:"A barely-nontrivial lower bound for an explicit problem Q in a weak model implies a major lower bound/separation in a strong model."Typically, Q is a meta-problem (e.g. a Gap-MCSP variant) and"major"means consequences such as NP not in P/poly.' while preserving 'MCSP complexity profile in nonuniform-circuits' strongly enough to support 'There exist explicit"meta"problems (including MCSP variants) such that proving certain emph barely-superlinear lower bounds against restricted models (e.g. formulas) would imply major lower bounds for NP in much stronger models (including superpolynomial formula/circuit lower bounds, and in some parameter regimes consequences as strong as NP not in P/poly ).'.
\item conjpkg.disclayer-p-vs-np-family-categorical-reframing.09: There exists an intermediate typed object built from 'MCSP witness family' such that the route from 'A hardness magnification theorem has the form:"A barely-nontrivial lower bound for an explicit problem Q in a weak model implies a major lower bound/separation in a strong model."Typically, Q is a meta-problem (e.g. a Gap-MCSP variant) and"major"means consequences such as NP not in P/poly.' to 'There exist explicit"meta"problems (including MCSP variants) such that proving certain emph barely-superlinear lower bounds against restricted models (e.g. formulas) would imply major lower bounds for NP in much stronger models (including superpolynomial formula/circuit lower bounds, and in some parameter regimes consequences as strong as NP not in P/poly ).' factors through theorem-sized maps preserving 'MCSP complexity profile in nonuniform-circuits'.
\item conjpkg.disclayer-p-vs-np-family-categorical-reframing.08: No admissible 'MCSP witness family' can simultaneously preserve 'MCSP complexity profile in nonuniform-circuits' and realize 'There exist explicit"meta"problems (including MCSP variants) such that proving certain emph barely-superlinear lower bounds against restricted models (e.g. formulas) would imply major lower bounds for NP in much stronger models (including superpolynomial formula/circuit lower bounds, and in some parameter regimes consequences as strong as NP not in P/poly ).' under the current route commitments.
\item conjpkg.disclayer-p-vs-np-family-categorical-reframing.07: For every admissible 'MCSP witness family', there is a transport law from 'nonuniform-circuits' to 'formula-size' that preserves 'MCSP complexity profile in nonuniform-circuits' and retains enough strength to activate 'There exist explicit"meta"problems (including MCSP variants) such that proving certain emph barely-superlinear lower bounds against restricted models (e.g. formulas) would imply major lower bounds for NP in much stronger models (including superpolynomial formula/circuit lower bounds, and in some parameter regimes consequences as strong as NP not in P/poly ).'.
\item conjpkg.disclayer-p-vs-np-family-categorical-reframing.06: There exists an admissible 'MCSP witness family' whose realization internalizes 'A hardness magnification theorem has the form:"A barely-nontrivial lower bound for an explicit problem Q in a weak model implies a major lower bound/separation in a strong model."Typically, Q is a meta-problem (e.g. a Gap-MCSP variant) and"major"means consequences such as NP not in P/poly.' while preserving 'MCSP complexity profile in nonuniform-circuits' strongly enough to support 'In refined regimes, lower bounds for Gap-MCSP slightly beyond known frontiers (e.g. N\textasciicircum{} 1+varepsilon -type circuit lower bounds for suitable promise gaps and input length N=2\textasciicircum{}n ) imply major consequences such as NP not in P/poly.'.
\item conjpkg.disclayer-p-vs-np-family-categorical-reframing.05: There exists an intermediate typed object built from 'MCSP witness family' such that the route from 'A hardness magnification theorem has the form:"A barely-nontrivial lower bound for an explicit problem Q in a weak model implies a major lower bound/separation in a strong model."Typically, Q is a meta-problem (e.g. a Gap-MCSP variant) and"major"means consequences such as NP not in P/poly.' to 'In refined regimes, lower bounds for Gap-MCSP slightly beyond known frontiers (e.g. N\textasciicircum{} 1+varepsilon -type circuit lower bounds for suitable promise gaps and input length N=2\textasciicircum{}n ) imply major consequences such as NP not in P/poly.' factors through theorem-sized maps preserving 'MCSP complexity profile in nonuniform-circuits'.
\item conjpkg.disclayer-p-vs-np-family-categorical-reframing.04: No admissible 'MCSP witness family' can simultaneously preserve 'MCSP complexity profile in nonuniform-circuits' and realize 'In refined regimes, lower bounds for Gap-MCSP slightly beyond known frontiers (e.g. N\textasciicircum{} 1+varepsilon -type circuit lower bounds for suitable promise gaps and input length N=2\textasciicircum{}n ) imply major consequences such as NP not in P/poly.' under the current route commitments.
\item conjpkg.disclayer-p-vs-np-family-categorical-reframing.03: There exists a mediating object built from 'MCSP obstruction family' such that the route from 'A hardness magnification theorem has the form:"A barely-nontrivial lower bound for an explicit problem Q in a weak model implies a major lower bound/separation in a strong model."Typically, Q is a meta-problem (e.g. a Gap-MCSP variant) and"major"means consequences such as NP not in P/poly.' to 'In refined regimes, lower bounds for Gap-MCSP slightly beyond known frontiers (e.g. N\textasciicircum{} 1+varepsilon -type circuit lower bounds for suitable promise gaps and input length N=2\textasciicircum{}n ) imply major consequences such as NP not in P/poly.' factors through typed morphisms preserving 'MCSP complexity profile in nonuniform-circuits' and making the relevant square commute.
\item conjpkg.disclayer-p-vs-np-family-categorical-reframing.03: There exists a mediating object built from 'MCSP obstruction family' such that the route from 'A hardness magnification theorem has the form:"A barely-nontrivial lower bound for an explicit problem Q in a weak model implies a major lower bound/separation in a strong model."Typically, Q is a meta-problem (e.g. a Gap-MCSP variant) and"major"means consequences such as NP not in P/poly.' to 'In refined regimes, lower bounds for Gap-MCSP slightly beyond known frontiers (e.g. N\textasciicircum{} 1+varepsilon -type circuit lower bounds for suitable promise gaps and input length N=2\textasciicircum{}n ) imply major consequences such as NP not in P/poly.' factors through typed morphisms preserving 'MCSP complexity profile in nonuniform-circuits' and making the relevant square commute.
\item conjpkg.disclayer-p-vs-np-family-categorical-reframing.02: There exists a mediating object built from 'MCSP witness family' such that the route from 'A hardness magnification theorem has the form:"A barely-nontrivial lower bound for an explicit problem Q in a weak model implies a major lower bound/separation in a strong model."Typically, Q is a meta-problem (e.g. a Gap-MCSP variant) and"major"means consequences such as NP not in P/poly.' to 'There exist explicit"meta"problems (including MCSP variants) such that proving certain emph barely-superlinear lower bounds against restricted models (e.g. formulas) would imply major lower bounds for NP in much stronger models (including superpolynomial formula/circuit lower bounds, and in some parameter regimes consequences as strong as NP not in P/poly ).' factors through typed morphisms preserving 'MCSP complexity profile in nonuniform-circuits' and making the relevant square commute.
\item conjpkg.disclayer-p-vs-np-family-categorical-reframing.02: There exists a mediating object built from 'MCSP witness family' such that the route from 'A hardness magnification theorem has the form:"A barely-nontrivial lower bound for an explicit problem Q in a weak model implies a major lower bound/separation in a strong model."Typically, Q is a meta-problem (e.g. a Gap-MCSP variant) and"major"means consequences such as NP not in P/poly.' to 'There exist explicit"meta"problems (including MCSP variants) such that proving certain emph barely-superlinear lower bounds against restricted models (e.g. formulas) would imply major lower bounds for NP in much stronger models (including superpolynomial formula/circuit lower bounds, and in some parameter regimes consequences as strong as NP not in P/poly ).' factors through typed morphisms preserving 'MCSP complexity profile in nonuniform-circuits' and making the relevant square commute.
\item conjpkg.disclayer-p-vs-np-family-categorical-reframing.01: There exists a mediating object built from 'MCSP witness family' such that the route from 'A hardness magnification theorem has the form:"A barely-nontrivial lower bound for an explicit problem Q in a weak model implies a major lower bound/separation in a strong model."Typically, Q is a meta-problem (e.g. a Gap-MCSP variant) and"major"means consequences such as NP not in P/poly.' to 'In refined regimes, lower bounds for Gap-MCSP slightly beyond known frontiers (e.g. N\textasciicircum{} 1+varepsilon -type circuit lower bounds for suitable promise gaps and input length N=2\textasciicircum{}n ) imply major consequences such as NP not in P/poly.' factors through typed morphisms preserving 'MCSP complexity profile in nonuniform-circuits' and making the relevant square commute.
\item conjpkg.disclayer-p-vs-np-family-categorical-reframing.01: There exists a mediating object built from 'MCSP witness family' such that the route from 'A hardness magnification theorem has the form:"A barely-nontrivial lower bound for an explicit problem Q in a weak model implies a major lower bound/separation in a strong model."Typically, Q is a meta-problem (e.g. a Gap-MCSP variant) and"major"means consequences such as NP not in P/poly.' to 'In refined regimes, lower bounds for Gap-MCSP slightly beyond known frontiers (e.g. N\textasciicircum{} 1+varepsilon -type circuit lower bounds for suitable promise gaps and input length N=2\textasciicircum{}n ) imply major consequences such as NP not in P/poly.' factors through typed morphisms preserving 'MCSP complexity profile in nonuniform-circuits' and making the relevant square commute.
\item conjpkg.disclayer-p-vs-np-family-equivalence-coherence-reframing.14: There exists an admissible 'MCSP witness family' whose realization internalizes 'A hardness magnification theorem has the form:"A barely-nontrivial lower bound for an explicit problem Q in a weak model implies a major lower bound/separation in a strong model."Typically, Q is a meta-problem (e.g. a Gap-MCSP variant) and"major"means consequences such as NP not in P/poly.' while preserving 'exponential-lower-bound witness complexity' strongly enough to support 'There exist general-purpose magnification theorems (not problem-specific) showing that certain mild lower bounds for distinguishers/sparse-code objects can be amplified into stronger circuit/formula lower bounds, achieving magnification thresholds below previously known lower bounds in some settings.'.
\item conjpkg.disclayer-p-vs-np-family-equivalence-coherence-reframing.13: No admissible 'MCSP witness family' can simultaneously preserve 'exponential-lower-bound witness complexity' and realize 'There exist general-purpose magnification theorems (not problem-specific) showing that certain mild lower bounds for distinguishers/sparse-code objects can be amplified into stronger circuit/formula lower bounds, achieving magnification thresholds below previously known lower bounds in some settings.' under the current route commitments.
\item conjpkg.disclayer-p-vs-np-family-equivalence-coherence-reframing.12: For every admissible 'MCSP witness family', the route-side and target-side operationalizations are equivalent in the sense needed to transport 'exponential-lower-bound witness complexity' toward 'There exist general-purpose magnification theorems (not problem-specific) showing that certain mild lower bounds for distinguishers/sparse-code objects can be amplified into stronger circuit/formula lower bounds, achieving magnification thresholds below previously known lower bounds in some settings.'.
\item conjpkg.disclayer-p-vs-np-family-equivalence-coherence-reframing.11: For every admissible 'MCSP witness family', there is a transport law from 'nonuniform-circuits' to 'unspecified' that preserves 'exponential-lower-bound witness complexity' and retains enough strength to activate 'There exist general-purpose magnification theorems (not problem-specific) showing that certain mild lower bounds for distinguishers/sparse-code objects can be amplified into stronger circuit/formula lower bounds, achieving magnification thresholds below previously known lower bounds in some settings.'.
\item conjpkg.disclayer-p-vs-np-family-equivalence-coherence-reframing.10: There exists an admissible 'MCSP witness family' whose realization internalizes 'A hardness magnification theorem has the form:"A barely-nontrivial lower bound for an explicit problem Q in a weak model implies a major lower bound/separation in a strong model."Typically, Q is a meta-problem (e.g. a Gap-MCSP variant) and"major"means consequences such as NP not in P/poly.' while preserving 'MCSP complexity profile in nonuniform-circuits' strongly enough to support 'There exist explicit"meta"problems (including MCSP variants) such that proving certain emph barely-superlinear lower bounds against restricted models (e.g. formulas) would imply major lower bounds for NP in much stronger models (including superpolynomial formula/circuit lower bounds, and in some parameter regimes consequences as strong as NP not in P/poly ).'.
\item conjpkg.disclayer-p-vs-np-family-equivalence-coherence-reframing.09: No admissible 'MCSP witness family' can simultaneously preserve 'MCSP complexity profile in nonuniform-circuits' and realize 'There exist explicit"meta"problems (including MCSP variants) such that proving certain emph barely-superlinear lower bounds against restricted models (e.g. formulas) would imply major lower bounds for NP in much stronger models (including superpolynomial formula/circuit lower bounds, and in some parameter regimes consequences as strong as NP not in P/poly ).' under the current route commitments.
\item conjpkg.disclayer-p-vs-np-family-equivalence-coherence-reframing.08: For every admissible 'MCSP witness family', the route-side and target-side operationalizations are equivalent in the sense needed to transport 'MCSP complexity profile in nonuniform-circuits' toward 'There exist explicit"meta"problems (including MCSP variants) such that proving certain emph barely-superlinear lower bounds against restricted models (e.g. formulas) would imply major lower bounds for NP in much stronger models (including superpolynomial formula/circuit lower bounds, and in some parameter regimes consequences as strong as NP not in P/poly ).'.
\item conjpkg.disclayer-p-vs-np-family-equivalence-coherence-reframing.07: For every admissible 'MCSP witness family', there is a transport law from 'nonuniform-circuits' to 'formula-size' that preserves 'MCSP complexity profile in nonuniform-circuits' and retains enough strength to activate 'There exist explicit"meta"problems (including MCSP variants) such that proving certain emph barely-superlinear lower bounds against restricted models (e.g. formulas) would imply major lower bounds for NP in much stronger models (including superpolynomial formula/circuit lower bounds, and in some parameter regimes consequences as strong as NP not in P/poly ).'.
\item conjpkg.disclayer-p-vs-np-family-equivalence-coherence-reframing.06: There exists an admissible 'MCSP witness family' whose realization internalizes 'A hardness magnification theorem has the form:"A barely-nontrivial lower bound for an explicit problem Q in a weak model implies a major lower bound/separation in a strong model."Typically, Q is a meta-problem (e.g. a Gap-MCSP variant) and"major"means consequences such as NP not in P/poly.' while preserving 'MCSP complexity profile in nonuniform-circuits' strongly enough to support 'In refined regimes, lower bounds for Gap-MCSP slightly beyond known frontiers (e.g. N\textasciicircum{} 1+varepsilon -type circuit lower bounds for suitable promise gaps and input length N=2\textasciicircum{}n ) imply major consequences such as NP not in P/poly.'.
\item conjpkg.disclayer-p-vs-np-family-equivalence-coherence-reframing.05: No admissible 'MCSP witness family' can simultaneously preserve 'MCSP complexity profile in nonuniform-circuits' and realize 'In refined regimes, lower bounds for Gap-MCSP slightly beyond known frontiers (e.g. N\textasciicircum{} 1+varepsilon -type circuit lower bounds for suitable promise gaps and input length N=2\textasciicircum{}n ) imply major consequences such as NP not in P/poly.' under the current route commitments.
\item conjpkg.disclayer-p-vs-np-family-equivalence-coherence-reframing.04: For every admissible 'MCSP witness family', the route-side and target-side operationalizations are equivalent in the sense needed to transport 'MCSP complexity profile in nonuniform-circuits' toward 'In refined regimes, lower bounds for Gap-MCSP slightly beyond known frontiers (e.g. N\textasciicircum{} 1+varepsilon -type circuit lower bounds for suitable promise gaps and input length N=2\textasciicircum{}n ) imply major consequences such as NP not in P/poly.'.
\item conjpkg.disclayer-p-vs-np-family-equivalence-coherence-reframing.03: There exists an equivalence witness identifying 'MCSP witness family' with the target-side object class relevant to 'There exist general-purpose magnification theorems (not problem-specific) showing that certain mild lower bounds for distinguishers/sparse-code objects can be amplified into stronger circuit/formula lower bounds, achieving magnification thresholds below previously known lower bounds in some settings.', and 'exponential-lower-bound witness complexity' transports coherently along that witness.
\item conjpkg.disclayer-p-vs-np-family-equivalence-coherence-reframing.03: There exists an equivalence witness identifying 'MCSP witness family' with the target-side object class relevant to 'There exist general-purpose magnification theorems (not problem-specific) showing that certain mild lower bounds for distinguishers/sparse-code objects can be amplified into stronger circuit/formula lower bounds, achieving magnification thresholds below previously known lower bounds in some settings.', and 'exponential-lower-bound witness complexity' transports coherently along that witness.
\item conjpkg.disclayer-p-vs-np-family-equivalence-coherence-reframing.02: There exists an equivalence witness identifying 'MCSP witness family' with the target-side object class relevant to 'There exist explicit"meta"problems (including MCSP variants) such that proving certain emph barely-superlinear lower bounds against restricted models (e.g. formulas) would imply major lower bounds for NP in much stronger models (including superpolynomial formula/circuit lower bounds, and in some parameter regimes consequences as strong as NP not in P/poly ).', and 'MCSP complexity profile in nonuniform-circuits' transports coherently along that witness.
\item conjpkg.disclayer-p-vs-np-family-equivalence-coherence-reframing.02: There exists an equivalence witness identifying 'MCSP witness family' with the target-side object class relevant to 'There exist explicit"meta"problems (including MCSP variants) such that proving certain emph barely-superlinear lower bounds against restricted models (e.g. formulas) would imply major lower bounds for NP in much stronger models (including superpolynomial formula/circuit lower bounds, and in some parameter regimes consequences as strong as NP not in P/poly ).', and 'MCSP complexity profile in nonuniform-circuits' transports coherently along that witness.
\item conjpkg.disclayer-p-vs-np-family-equivalence-coherence-reframing.01: There exists an equivalence witness identifying 'MCSP witness family' with the target-side object class relevant to 'In refined regimes, lower bounds for Gap-MCSP slightly beyond known frontiers (e.g. N\textasciicircum{} 1+varepsilon -type circuit lower bounds for suitable promise gaps and input length N=2\textasciicircum{}n ) imply major consequences such as NP not in P/poly.', and 'MCSP complexity profile in nonuniform-circuits' transports coherently along that witness.
\item conjpkg.disclayer-p-vs-np-family-equivalence-coherence-reframing.01: There exists an equivalence witness identifying 'MCSP witness family' with the target-side object class relevant to 'In refined regimes, lower bounds for Gap-MCSP slightly beyond known frontiers (e.g. N\textasciicircum{} 1+varepsilon -type circuit lower bounds for suitable promise gaps and input length N=2\textasciicircum{}n ) imply major consequences such as NP not in P/poly.', and 'MCSP complexity profile in nonuniform-circuits' transports coherently along that witness.
\end{itemize}

\section*{Restricted-Regime Programmes}
\begin{itemize}
\item Special-Case Theorem: invariant-discovery: There exists an admissible MCSP witness family, realized as a family of instances equipped with a monotone witness measure, such that whenever A barely-nontrivial lower bound for an explicit problem Q in a weak model holds for MCSP in nonuniform-circuits, the object organizes the source-side mathematical content while preserving the invariant MCSP complexity profile in nonuniform-circuits and yields major consequences such as NP not in P/poly.
\item Threshold Version: regime-localization: There exists an admissible MCSP witness family, realized as a family of instances equipped with a monotone witness measure, such that whenever A barely-nontrivial lower bound for an explicit problem Q in a weak model holds for MCSP in nonuniform-circuits, the object organizes the source-side mathematical content while preserving the invariant MCSP complexity profile in nonuniform-circuits and yields major consequences such as NP not in P/poly.
\item Special-Case Theorem: hybridization-of-partial-bridges: There exists an admissible MCSP witness family, realized as a family of instances equipped with a monotone witness measure, such that whenever A barely-nontrivial lower bound for an explicit problem Q in a weak model holds for MCSP in nonuniform-circuits, the object organizes the source-side mathematical content while preserving the invariant MCSP complexity profile in nonuniform-circuits and yields major consequences such as NP not in P/poly.
\item Special-Case Theorem: route-as-commutative-diagram: There exists a mediating object M and typed morphisms making the route commute; otherwise the failure is certified as no morphism, non-functorial transport, broken naturality, or missing universal property.
\item Special-Case Theorem: route-as-commutative-diagram: There exists a mediating object M and typed morphisms making the route commute; otherwise the failure is certified as no morphism, non-functorial transport, broken naturality, or missing universal property.
\item Equivalence Theorem: route-as-equivalence-problem: Either exhibit an equivalence E between the live formulations and prove that the target-relevant invariant transports coherently along E, or prove that no such equivalence witness exists in the current regime.
\item Equivalence Theorem: route-as-equivalence-problem: Either exhibit an equivalence E between the live formulations and prove that the target-relevant invariant transports coherently along E, or prove that no such equivalence witness exists in the current regime.
\end{itemize}

\section*{Foundation-Search Programmes}
\begin{itemize}
\item conjpkg.disclayer-p-vs-np-pattern-regime-localization.10: threshold and regime-restriction frameworks, restricted-model lower-bound programmes, proof complexity, pseudorandomness / derandomization, coding-theoretic hardness surrogates
\item conjpkg.disclayer-p-vs-np-pattern-regime-localization.09: threshold and regime-restriction frameworks, restricted-model lower-bound programmes, proof complexity, pseudorandomness / derandomization, coding-theoretic hardness surrogates
\item conjpkg.disclayer-p-vs-np-pattern-regime-localization.08: threshold and regime-restriction frameworks, restricted-model lower-bound programmes, proof complexity, pseudorandomness / derandomization, coding-theoretic hardness surrogates
\item conjpkg.disclayer-p-vs-np-pattern-regime-localization.07: threshold and regime-restriction frameworks, restricted-model lower-bound programmes, proof complexity, pseudorandomness / derandomization, coding-theoretic hardness surrogates
\item conjpkg.disclayer-p-vs-np-pattern-regime-localization.06: threshold and regime-restriction frameworks, restricted-model lower-bound programmes, proof complexity, pseudorandomness / derandomization, coding-theoretic hardness surrogates
\item conjpkg.disclayer-p-vs-np-pattern-regime-localization.05: threshold and regime-restriction frameworks, restricted-model lower-bound programmes, proof complexity, pseudorandomness / derandomization, coding-theoretic hardness surrogates
\item conjpkg.disclayer-p-vs-np-pattern-regime-localization.04: threshold and regime-restriction frameworks, restricted-model lower-bound programmes, proof complexity, pseudorandomness / derandomization, coding-theoretic hardness surrogates
\item conjpkg.disclayer-p-vs-np-pattern-regime-localization.03: threshold and regime-restriction frameworks, restricted-model lower-bound programmes, proof complexity, pseudorandomness / derandomization, coding-theoretic hardness surrogates
\item conjpkg.disclayer-p-vs-np-pattern-regime-localization.02: threshold and regime-restriction frameworks, restricted-model lower-bound programmes, proof complexity, pseudorandomness / derandomization, coding-theoretic hardness surrogates
\item conjpkg.disclayer-p-vs-np-pattern-regime-localization.01: threshold and regime-restriction frameworks, restricted-model lower-bound programmes, proof complexity, pseudorandomness / derandomization, coding-theoretic hardness surrogates
\item conjpkg.disclayer-p-vs-np-pattern-invariant-discovery.10: potential / energy method frameworks, monotone measure frameworks, certificate-complexity and witness-measure theories, proof complexity, pseudorandomness / derandomization
\item conjpkg.disclayer-p-vs-np-pattern-invariant-discovery.09: potential / energy method frameworks, monotone measure frameworks, certificate-complexity and witness-measure theories, proof complexity, pseudorandomness / derandomization
\item conjpkg.disclayer-p-vs-np-pattern-invariant-discovery.08: potential / energy method frameworks, monotone measure frameworks, certificate-complexity and witness-measure theories, proof complexity, pseudorandomness / derandomization
\item conjpkg.disclayer-p-vs-np-pattern-invariant-discovery.07: potential / energy method frameworks, monotone measure frameworks, certificate-complexity and witness-measure theories, proof complexity, pseudorandomness / derandomization
\item conjpkg.disclayer-p-vs-np-pattern-invariant-discovery.06: potential / energy method frameworks, monotone measure frameworks, certificate-complexity and witness-measure theories, proof complexity, pseudorandomness / derandomization
\item conjpkg.disclayer-p-vs-np-pattern-invariant-discovery.05: potential / energy method frameworks, monotone measure frameworks, certificate-complexity and witness-measure theories, proof complexity, pseudorandomness / derandomization
\item conjpkg.disclayer-p-vs-np-pattern-invariant-discovery.04: potential / energy method frameworks, monotone measure frameworks, certificate-complexity and witness-measure theories, proof complexity, pseudorandomness / derandomization
\item conjpkg.disclayer-p-vs-np-pattern-invariant-discovery.03: potential / energy method frameworks, monotone measure frameworks, certificate-complexity and witness-measure theories, proof complexity, pseudorandomness / derandomization
\item conjpkg.disclayer-p-vs-np-pattern-invariant-discovery.02: potential / energy method frameworks, monotone measure frameworks, certificate-complexity and witness-measure theories, proof complexity, pseudorandomness / derandomization
\item conjpkg.disclayer-p-vs-np-pattern-invariant-discovery.01: potential / energy method frameworks, monotone measure frameworks, certificate-complexity and witness-measure theories, proof complexity, pseudorandomness / derandomization
\item conjpkg.disclayer-p-vs-np-pattern-hybridization-of-partial-bridges.13: composed transfer frameworks, hybrid argument methods, multi-interface reduction theories, proof complexity, pseudorandomness / derandomization
\item conjpkg.disclayer-p-vs-np-pattern-hybridization-of-partial-bridges.12: composed transfer frameworks, hybrid argument methods, multi-interface reduction theories, proof complexity, pseudorandomness / derandomization
\item conjpkg.disclayer-p-vs-np-pattern-hybridization-of-partial-bridges.11: composed transfer frameworks, hybrid argument methods, multi-interface reduction theories, proof complexity, pseudorandomness / derandomization
\item conjpkg.disclayer-p-vs-np-pattern-hybridization-of-partial-bridges.10: composed transfer frameworks, hybrid argument methods, multi-interface reduction theories, proof complexity, pseudorandomness / derandomization
\item conjpkg.disclayer-p-vs-np-pattern-hybridization-of-partial-bridges.09: composed transfer frameworks, hybrid argument methods, multi-interface reduction theories, proof complexity, pseudorandomness / derandomization
\item conjpkg.disclayer-p-vs-np-pattern-hybridization-of-partial-bridges.08: composed transfer frameworks, hybrid argument methods, multi-interface reduction theories, proof complexity, pseudorandomness / derandomization
\item conjpkg.disclayer-p-vs-np-pattern-hybridization-of-partial-bridges.07: composed transfer frameworks, hybrid argument methods, multi-interface reduction theories, proof complexity, pseudorandomness / derandomization
\item conjpkg.disclayer-p-vs-np-pattern-hybridization-of-partial-bridges.06: composed transfer frameworks, hybrid argument methods, multi-interface reduction theories, proof complexity, pseudorandomness / derandomization
\item conjpkg.disclayer-p-vs-np-pattern-hybridization-of-partial-bridges.05: composed transfer frameworks, hybrid argument methods, multi-interface reduction theories, proof complexity, pseudorandomness / derandomization
\item conjpkg.disclayer-p-vs-np-pattern-hybridization-of-partial-bridges.04: composed transfer frameworks, hybrid argument methods, multi-interface reduction theories, proof complexity, pseudorandomness / derandomization
\item conjpkg.disclayer-p-vs-np-pattern-hybridization-of-partial-bridges.03: composed transfer frameworks, hybrid argument methods, multi-interface reduction theories, proof complexity, pseudorandomness / derandomization
\item conjpkg.disclayer-p-vs-np-pattern-hybridization-of-partial-bridges.02: composed transfer frameworks, hybrid argument methods, multi-interface reduction theories, proof complexity, pseudorandomness / derandomization
\item conjpkg.disclayer-p-vs-np-pattern-hybridization-of-partial-bridges.01: composed transfer frameworks, hybrid argument methods, multi-interface reduction theories, proof complexity, pseudorandomness / derandomization
\item conjpkg.disclayer-p-vs-np-family-explicitness.08: derandomization and generator frameworks, algebraic certificate constructions, extractor / hitter frameworks
\item conjpkg.disclayer-p-vs-np-family-explicitness.07: derandomization and generator frameworks, algebraic certificate constructions, extractor / hitter frameworks
\item conjpkg.disclayer-p-vs-np-family-explicitness.06: derandomization and generator frameworks, algebraic certificate constructions, extractor / hitter frameworks
\item conjpkg.disclayer-p-vs-np-family-explicitness.05: derandomization and generator frameworks, algebraic certificate constructions, extractor / hitter frameworks
\item conjpkg.disclayer-p-vs-np-family-explicitness.04: derandomization and generator frameworks, algebraic certificate constructions, extractor / hitter frameworks
\item conjpkg.disclayer-p-vs-np-family-explicitness.03: derandomization and generator frameworks, algebraic certificate constructions, extractor / hitter frameworks
\item conjpkg.disclayer-p-vs-np-family-explicitness.02: derandomization and generator frameworks, algebraic certificate constructions, extractor / hitter frameworks
\item conjpkg.disclayer-p-vs-np-family-explicitness.01: derandomization and generator frameworks, algebraic certificate constructions, extractor / hitter frameworks
\item conjpkg.disclayer-p-vs-np-family-equivalence-coherence-reframing.14: sparse-object and distinguisher frameworks, duality or correspondence theorems, representation-theoretic normal forms, derandomization and generator frameworks, algebraic certificate constructions
\item conjpkg.disclayer-p-vs-np-family-equivalence-coherence-reframing.13: sparse-object and distinguisher frameworks, duality or correspondence theorems, representation-theoretic normal forms, derandomization and generator frameworks, algebraic certificate constructions
\item conjpkg.disclayer-p-vs-np-family-equivalence-coherence-reframing.12: sparse-object and distinguisher frameworks, duality or correspondence theorems, representation-theoretic normal forms, derandomization and generator frameworks, algebraic certificate constructions
\item conjpkg.disclayer-p-vs-np-family-equivalence-coherence-reframing.11: sparse-object and distinguisher frameworks, duality or correspondence theorems, representation-theoretic normal forms, derandomization and generator frameworks, algebraic certificate constructions
\item conjpkg.disclayer-p-vs-np-family-equivalence-coherence-reframing.10: sparse-object and distinguisher frameworks, duality or correspondence theorems, representation-theoretic normal forms, derandomization and generator frameworks, algebraic certificate constructions
\item conjpkg.disclayer-p-vs-np-family-equivalence-coherence-reframing.09: sparse-object and distinguisher frameworks, duality or correspondence theorems, representation-theoretic normal forms, derandomization and generator frameworks, algebraic certificate constructions
\item conjpkg.disclayer-p-vs-np-family-equivalence-coherence-reframing.08: sparse-object and distinguisher frameworks, duality or correspondence theorems, representation-theoretic normal forms, derandomization and generator frameworks, algebraic certificate constructions
\item conjpkg.disclayer-p-vs-np-family-equivalence-coherence-reframing.07: sparse-object and distinguisher frameworks, duality or correspondence theorems, representation-theoretic normal forms, derandomization and generator frameworks, algebraic certificate constructions
\item conjpkg.disclayer-p-vs-np-family-equivalence-coherence-reframing.06: sparse-object and distinguisher frameworks, duality or correspondence theorems, representation-theoretic normal forms, derandomization and generator frameworks, algebraic certificate constructions
\item conjpkg.disclayer-p-vs-np-family-equivalence-coherence-reframing.05: sparse-object and distinguisher frameworks, duality or correspondence theorems, representation-theoretic normal forms, derandomization and generator frameworks, algebraic certificate constructions
\item conjpkg.disclayer-p-vs-np-family-equivalence-coherence-reframing.04: sparse-object and distinguisher frameworks, duality or correspondence theorems, representation-theoretic normal forms, derandomization and generator frameworks, algebraic certificate constructions
\item conjpkg.disclayer-p-vs-np-family-equivalence-coherence-reframing.03: sparse-object and distinguisher frameworks, duality or correspondence theorems, representation-theoretic normal forms, derandomization and generator frameworks, algebraic certificate constructions
\item conjpkg.disclayer-p-vs-np-family-equivalence-coherence-reframing.03: sparse-object and distinguisher frameworks, duality or correspondence theorems, representation-theoretic normal forms, derandomization and generator frameworks, algebraic certificate constructions
\item conjpkg.disclayer-p-vs-np-family-equivalence-coherence-reframing.02: sparse-object and distinguisher frameworks, duality or correspondence theorems, representation-theoretic normal forms, derandomization and generator frameworks, algebraic certificate constructions
\item conjpkg.disclayer-p-vs-np-family-equivalence-coherence-reframing.02: sparse-object and distinguisher frameworks, duality or correspondence theorems, representation-theoretic normal forms, derandomization and generator frameworks, algebraic certificate constructions
\item conjpkg.disclayer-p-vs-np-family-equivalence-coherence-reframing.01: sparse-object and distinguisher frameworks, duality or correspondence theorems, representation-theoretic normal forms, derandomization and generator frameworks, algebraic certificate constructions
\item conjpkg.disclayer-p-vs-np-family-equivalence-coherence-reframing.01: sparse-object and distinguisher frameworks, duality or correspondence theorems, representation-theoretic normal forms, derandomization and generator frameworks, algebraic certificate constructions
\item conjpkg.disclayer-p-vs-np-family-categorical-reframing.13: sparse-object and distinguisher frameworks, duality or correspondence theorems, representation-theoretic normal forms, lower-bound barrier frameworks, local-reduction impossibility theory
\item conjpkg.disclayer-p-vs-np-family-categorical-reframing.12: sparse-object and distinguisher frameworks, duality or correspondence theorems, representation-theoretic normal forms, lower-bound barrier frameworks, local-reduction impossibility theory
\item conjpkg.disclayer-p-vs-np-family-categorical-reframing.11: sparse-object and distinguisher frameworks, duality or correspondence theorems, representation-theoretic normal forms, lower-bound barrier frameworks, local-reduction impossibility theory
\item conjpkg.disclayer-p-vs-np-family-categorical-reframing.10: sparse-object and distinguisher frameworks, duality or correspondence theorems, representation-theoretic normal forms, lower-bound barrier frameworks, local-reduction impossibility theory
\item conjpkg.disclayer-p-vs-np-family-categorical-reframing.09: sparse-object and distinguisher frameworks, duality or correspondence theorems, representation-theoretic normal forms, lower-bound barrier frameworks, local-reduction impossibility theory
\item conjpkg.disclayer-p-vs-np-family-categorical-reframing.08: sparse-object and distinguisher frameworks, duality or correspondence theorems, representation-theoretic normal forms, lower-bound barrier frameworks, local-reduction impossibility theory
\item conjpkg.disclayer-p-vs-np-family-categorical-reframing.07: sparse-object and distinguisher frameworks, duality or correspondence theorems, representation-theoretic normal forms, lower-bound barrier frameworks, local-reduction impossibility theory
\item conjpkg.disclayer-p-vs-np-family-categorical-reframing.06: sparse-object and distinguisher frameworks, duality or correspondence theorems, representation-theoretic normal forms, lower-bound barrier frameworks, local-reduction impossibility theory
\item conjpkg.disclayer-p-vs-np-family-categorical-reframing.05: sparse-object and distinguisher frameworks, duality or correspondence theorems, representation-theoretic normal forms, lower-bound barrier frameworks, local-reduction impossibility theory
\item conjpkg.disclayer-p-vs-np-family-categorical-reframing.04: sparse-object and distinguisher frameworks, duality or correspondence theorems, representation-theoretic normal forms, lower-bound barrier frameworks, local-reduction impossibility theory
\item conjpkg.disclayer-p-vs-np-family-categorical-reframing.03: sparse-object and distinguisher frameworks, duality or correspondence theorems, representation-theoretic normal forms, lower-bound barrier frameworks, local-reduction impossibility theory
\item conjpkg.disclayer-p-vs-np-family-categorical-reframing.03: sparse-object and distinguisher frameworks, duality or correspondence theorems, representation-theoretic normal forms, lower-bound barrier frameworks, local-reduction impossibility theory
\item conjpkg.disclayer-p-vs-np-family-categorical-reframing.02: sparse-object and distinguisher frameworks, duality or correspondence theorems, representation-theoretic normal forms, lower-bound barrier frameworks, local-reduction impossibility theory
\item conjpkg.disclayer-p-vs-np-family-categorical-reframing.02: sparse-object and distinguisher frameworks, duality or correspondence theorems, representation-theoretic normal forms, lower-bound barrier frameworks, local-reduction impossibility theory
\item conjpkg.disclayer-p-vs-np-family-categorical-reframing.01: sparse-object and distinguisher frameworks, duality or correspondence theorems, representation-theoretic normal forms, lower-bound barrier frameworks, local-reduction impossibility theory
\item conjpkg.disclayer-p-vs-np-family-categorical-reframing.01: sparse-object and distinguisher frameworks, duality or correspondence theorems, representation-theoretic normal forms, lower-bound barrier frameworks, local-reduction impossibility theory
\item conjpkg.disclayer-p-vs-np-family-barrier-aware-transfer.06: lower-bound barrier frameworks, local-reduction impossibility theory, gadget and locality-preservation frameworks
\item conjpkg.disclayer-p-vs-np-family-barrier-aware-transfer.05: lower-bound barrier frameworks, local-reduction impossibility theory, gadget and locality-preservation frameworks
\item conjpkg.disclayer-p-vs-np-family-barrier-aware-transfer.04: lower-bound barrier frameworks, local-reduction impossibility theory, gadget and locality-preservation frameworks
\item conjpkg.disclayer-p-vs-np-family-barrier-aware-transfer.03: lower-bound barrier frameworks, local-reduction impossibility theory, gadget and locality-preservation frameworks
\item conjpkg.disclayer-p-vs-np-family-barrier-aware-transfer.02: lower-bound barrier frameworks, local-reduction impossibility theory, gadget and locality-preservation frameworks
\item conjpkg.disclayer-p-vs-np-family-barrier-aware-transfer.01: lower-bound barrier frameworks, local-reduction impossibility theory, gadget and locality-preservation frameworks
\item conjpkg.disclayer-p-vs-np-family-abstraction-search.10: proof complexity, pseudorandomness / derandomization, coding-theoretic hardness surrogates
\item conjpkg.disclayer-p-vs-np-family-abstraction-search.09: proof complexity, pseudorandomness / derandomization, coding-theoretic hardness surrogates
\item conjpkg.disclayer-p-vs-np-family-abstraction-search.08: proof complexity, pseudorandomness / derandomization, coding-theoretic hardness surrogates
\item conjpkg.disclayer-p-vs-np-family-abstraction-search.07: proof complexity, pseudorandomness / derandomization, coding-theoretic hardness surrogates
\item conjpkg.disclayer-p-vs-np-family-abstraction-search.06: proof complexity, pseudorandomness / derandomization, coding-theoretic hardness surrogates
\item conjpkg.disclayer-p-vs-np-family-abstraction-search.05: proof complexity, pseudorandomness / derandomization, coding-theoretic hardness surrogates
\item conjpkg.disclayer-p-vs-np-family-abstraction-search.04: proof complexity, pseudorandomness / derandomization, coding-theoretic hardness surrogates
\item conjpkg.disclayer-p-vs-np-family-abstraction-search.03: proof complexity, pseudorandomness / derandomization, coding-theoretic hardness surrogates
\item conjpkg.disclayer-p-vs-np-family-abstraction-search.02: proof complexity, pseudorandomness / derandomization, coding-theoretic hardness surrogates
\item conjpkg.disclayer-p-vs-np-family-abstraction-search.01: proof complexity, pseudorandomness / derandomization, coding-theoretic hardness surrogates
\end{itemize}

\section*{Packages}
\subsection*{conjpkg.disclayer-p-vs-np-family-equivalence-coherence-reframing.01}
\begin{itemize}
\item Exact statement: There exists an equivalence witness identifying 'MCSP witness family' with the target-side object class relevant to 'In refined regimes, lower bounds for Gap-MCSP slightly beyond known frontiers (e.g. N\textasciicircum{} 1+varepsilon -type circuit lower bounds for suitable promise gaps and input length N=2\textasciicircum{}n ) imply major consequences such as NP not in P/poly.', and 'MCSP complexity profile in nonuniform-circuits' transports coherently along that witness.
\item Restricted-regime mode: equivalence-theorem
\item Bridge object: MCSP witness family (family of instances equipped with a monotone witness measure)
\item Bridge definition: 'MCSP witness family' denotes the structured intermediate object required by this branch. It is meant to organize the source-side content 'A hardness magnification theorem has the form:"A barely-nontrivial lower bound for an explicit problem Q in a weak model implies a major lower bound/separation in a strong model."Typically, Q is a meta-problem (e.g. a Gap-MCSP variant) and"major"means consequences such as NP not in P/poly.'. It must preserve the invariant 'MCSP complexity profile in nonuniform-circuits'. Its operational role is 'Recast the blocked frontier as an identity, equivalence, transport, or coherence problem rather than as a missing brute-force bridge.'.
\item First nontrivial lemma: For every admissible 'MCSP witness family', the target-side operationalization is equivalent to the route-side operationalization in the sense needed to transport 'MCSP complexity profile in nonuniform-circuits' toward 'In refined regimes, lower bounds for Gap-MCSP slightly beyond known frontiers (e.g. N\textasciicircum{} 1+varepsilon -type circuit lower bounds for suitable promise gaps and input length N=2\textasciicircum{}n ) imply major consequences such as NP not in P/poly.'.
\item Quantifier signature: For every admissible 'MCSP witness family' on the route from 'st.def.hardness-magnification-phenomenon-programme-leve' to 'st.thm.oliveira-pich-santhanam-magnification-near-state'.
\item LaTeX theorem: \textbackslash\{\}forall W,\textbackslash\{\} \textbackslash\{\}mathrm\{Adm\}(W) \textbackslash\{\}to \textbackslash\{\}exists w : \textbackslash\{\}mathsf\{BW\}(W),\textbackslash\{\} \textbackslash\{\}mathrm\{PresInv\}(W,w) \textbackslash\{\}wedge \textbackslash\{\}mathrm\{RealizesTarget\}(W,w).
\item Negative/barrier version: Absent an explicit equivalence witness for 'MCSP witness family', transport toward 'In refined regimes, lower bounds for Gap-MCSP slightly beyond known frontiers (e.g. N\textasciicircum{} 1+varepsilon -type circuit lower bounds for suitable promise gaps and input length N=2\textasciicircum{}n ) imply major consequences such as NP not in P/poly.' is not licensed. no\_equivalence\_witness
\item Proof strategy family: equivalence-transport
\item Negative route: Assume no equivalence witness for 'MCSP witness family' aligns the route-side and target-side objects.
\item Failure diagnosis: The route fails if no equivalence witness for 'MCSP witness family' transports 'MCSP complexity profile in nonuniform-circuits' to the target-side operationalization. Local negative probe: no\_equivalence\_witness.
\item Proof obligations: For every admissible 'MCSP witness family', the target-side operationalization is equivalent to the route-side operationalization in the sense needed to transport 'MCSP complexity profile in nonuniform-circuits' toward 'In refined regimes, lower bounds for Gap-MCSP slightly beyond known frontiers (e.g. N\textasciicircum{} 1+varepsilon -type circuit lower bounds for suitable promise gaps and input length N=2\textasciicircum{}n ) imply major consequences such as NP not in P/poly.'.; Transport obligation: move the invariant across the proposed equivalence for 'carrier-realization' without changing the target claim.; Transport obligation: move the invariant across the proposed equivalence for 'formula-size-transport' without changing the target claim.; Equivalence claim: determine whether the live carrier for 'carrier-realization' is identical to, equivalent to, or only analogous with the target-side object class.
\item Candidate lemmas: Equivalence and coherent transport theorem for MCSP witness family; Equivalence transport lemma for MCSP witness family
\item Independent value: An equivalence theorem for 'MCSP witness family' would independently matter because it identifies when two live operationalizations are mathematically interchangeable rather than merely analogous.
\item Expert gate: 6/6
\end{itemize}

\subsection*{conjpkg.disclayer-p-vs-np-family-equivalence-coherence-reframing.02}
\begin{itemize}
\item Exact statement: There exists an equivalence witness identifying 'MCSP witness family' with the target-side object class relevant to 'There exist explicit"meta"problems (including MCSP variants) such that proving certain emph barely-superlinear lower bounds against restricted models (e.g. formulas) would imply major lower bounds for NP in much stronger models (including superpolynomial formula/circuit lower bounds, and in some parameter regimes consequences as strong as NP not in P/poly ).', and 'MCSP complexity profile in nonuniform-circuits' transports coherently along that witness.
\item Restricted-regime mode: equivalence-theorem
\item Bridge object: MCSP witness family (family of instances equipped with a monotone witness measure)
\item Bridge definition: 'MCSP witness family' denotes the structured intermediate object required by this branch. It is meant to organize the source-side content 'A hardness magnification theorem has the form:"A barely-nontrivial lower bound for an explicit problem Q in a weak model implies a major lower bound/separation in a strong model."Typically, Q is a meta-problem (e.g. a Gap-MCSP variant) and"major"means consequences such as NP not in P/poly.'. It must preserve the invariant 'MCSP complexity profile in nonuniform-circuits'. Its operational role is 'Recast the blocked frontier as an identity, equivalence, transport, or coherence problem rather than as a missing brute-force bridge.'.
\item First nontrivial lemma: For every admissible 'MCSP witness family', the target-side operationalization is equivalent to the route-side operationalization in the sense needed to transport 'MCSP complexity profile in nonuniform-circuits' toward 'There exist explicit"meta"problems (including MCSP variants) such that proving certain emph barely-superlinear lower bounds against restricted models (e.g. formulas) would imply major lower bounds for NP in much stronger models (including superpolynomial formula/circuit lower bounds, and in some parameter regimes consequences as strong as NP not in P/poly ).'.
\item Quantifier signature: For every admissible 'MCSP witness family' on the route from 'st.def.hardness-magnification-phenomenon-programme-leve' to 'st.thm.oliveira-santhanam-hardness-magnification-for-na'.
\item LaTeX theorem: \textbackslash\{\}forall W,\textbackslash\{\} \textbackslash\{\}mathrm\{Adm\}(W) \textbackslash\{\}to \textbackslash\{\}exists w : \textbackslash\{\}mathsf\{BW\}(W),\textbackslash\{\} \textbackslash\{\}mathrm\{PresInv\}(W,w) \textbackslash\{\}wedge \textbackslash\{\}mathrm\{RealizesTarget\}(W,w).
\item Negative/barrier version: Absent an explicit equivalence witness for 'MCSP witness family', transport toward 'There exist explicit"meta"problems (including MCSP variants) such that proving certain emph barely-superlinear lower bounds against restricted models (e.g. formulas) would imply major lower bounds for NP in much stronger models (including superpolynomial formula/circuit lower bounds, and in some parameter regimes consequences as strong as NP not in P/poly ).' is not licensed. no\_equivalence\_witness
\item Proof strategy family: equivalence-transport
\item Negative route: Assume no equivalence witness for 'MCSP witness family' aligns the route-side and target-side objects.
\item Failure diagnosis: The route fails if no equivalence witness for 'MCSP witness family' transports 'MCSP complexity profile in nonuniform-circuits' to the target-side operationalization. Local negative probe: no\_equivalence\_witness.
\item Proof obligations: For every admissible 'MCSP witness family', the target-side operationalization is equivalent to the route-side operationalization in the sense needed to transport 'MCSP complexity profile in nonuniform-circuits' toward 'There exist explicit"meta"problems (including MCSP variants) such that proving certain emph barely-superlinear lower bounds against restricted models (e.g. formulas) would imply major lower bounds for NP in much stronger models (including superpolynomial formula/circuit lower bounds, and in some parameter regimes consequences as strong as NP not in P/poly ).'.; Transport obligation: move the invariant across the proposed equivalence for 'carrier-realization' without changing the target claim.; Transport obligation: move the invariant across the proposed equivalence for 'formula-size-transport' without changing the target claim.; Equivalence claim: determine whether the live carrier for 'carrier-realization' is identical to, equivalent to, or only analogous with the target-side object class.
\item Candidate lemmas: Equivalence and coherent transport theorem for MCSP witness family; Equivalence transport lemma for MCSP witness family
\item Independent value: An equivalence theorem for 'MCSP witness family' would independently matter because it identifies when two live operationalizations are mathematically interchangeable rather than merely analogous.
\item Expert gate: 6/6
\end{itemize}

\subsection*{conjpkg.disclayer-p-vs-np-family-equivalence-coherence-reframing.03}
\begin{itemize}
\item Exact statement: There exists an equivalence witness identifying 'MCSP witness family' with the target-side object class relevant to 'There exist general-purpose magnification theorems (not problem-specific) showing that certain mild lower bounds for distinguishers/sparse-code objects can be amplified into stronger circuit/formula lower bounds, achieving magnification thresholds below previously known lower bounds in some settings.', and 'exponential-lower-bound witness complexity' transports coherently along that witness.
\item Restricted-regime mode: equivalence-theorem
\item Bridge object: MCSP witness family (family of instances equipped with a monotone witness measure)
\item Bridge definition: 'MCSP witness family' denotes the structured intermediate object required by this branch. It is meant to organize the source-side content 'A hardness magnification theorem has the form:"A barely-nontrivial lower bound for an explicit problem Q in a weak model implies a major lower bound/separation in a strong model."Typically, Q is a meta-problem (e.g. a Gap-MCSP variant) and"major"means consequences such as NP not in P/poly.'. It must preserve the invariant 'exponential-lower-bound witness complexity'. Its operational role is 'Recast the blocked frontier as an identity, equivalence, transport, or coherence problem rather than as a missing brute-force bridge.'.
\item First nontrivial lemma: For every admissible 'MCSP witness family', the target-side operationalization is equivalent to the route-side operationalization in the sense needed to transport 'exponential-lower-bound witness complexity' toward 'There exist general-purpose magnification theorems (not problem-specific) showing that certain mild lower bounds for distinguishers/sparse-code objects can be amplified into stronger circuit/formula lower bounds, achieving magnification thresholds below previously known lower bounds in some settings.'.
\item Quantifier signature: For every admissible 'MCSP witness family' on the route from 'st.def.hardness-magnification-phenomenon-programme-leve' to 'st.thm.atserias-m-uller-simple-general-magnification-fr'.
\item LaTeX theorem: \textbackslash\{\}forall W,\textbackslash\{\} \textbackslash\{\}mathrm\{Adm\}(W) \textbackslash\{\}to \textbackslash\{\}exists w : \textbackslash\{\}mathsf\{BW\}(W),\textbackslash\{\} \textbackslash\{\}mathrm\{PresInv\}(W,w) \textbackslash\{\}wedge \textbackslash\{\}mathrm\{RealizesTarget\}(W,w).
\item Negative/barrier version: Absent an explicit equivalence witness for 'MCSP witness family', transport toward 'There exist general-purpose magnification theorems (not problem-specific) showing that certain mild lower bounds for distinguishers/sparse-code objects can be amplified into stronger circuit/formula lower bounds, achieving magnification thresholds below previously known lower bounds in some settings.' is not licensed. no\_equivalence\_witness
\item Proof strategy family: equivalence-transport
\item Negative route: Assume no equivalence witness for 'MCSP witness family' aligns the route-side and target-side objects.
\item Failure diagnosis: The route fails if no equivalence witness for 'MCSP witness family' transports 'exponential-lower-bound witness complexity' to the target-side operationalization. Local negative probe: no\_equivalence\_witness.
\item Proof obligations: For every admissible 'MCSP witness family', the target-side operationalization is equivalent to the route-side operationalization in the sense needed to transport 'exponential-lower-bound witness complexity' toward 'There exist general-purpose magnification theorems (not problem-specific) showing that certain mild lower bounds for distinguishers/sparse-code objects can be amplified into stronger circuit/formula lower bounds, achieving magnification thresholds below previously known lower bounds in some settings.'.; Transport obligation: move the invariant across the proposed equivalence for 'carrier-realization' without changing the target claim.; Transport obligation: move the invariant across the proposed equivalence for 'formula-size-transport' without changing the target claim.; Equivalence claim: determine whether the live carrier for 'carrier-realization' is identical to, equivalent to, or only analogous with the target-side object class.
\item Candidate lemmas: Equivalence and coherent transport theorem for MCSP witness family; Equivalence transport lemma for MCSP witness family
\item Independent value: An equivalence theorem for 'MCSP witness family' would independently matter because it identifies when two live operationalizations are mathematically interchangeable rather than merely analogous.
\item Expert gate: 6/6
\end{itemize}

\subsection*{conjpkg.disclayer-p-vs-np-family-equivalence-coherence-reframing.04}
\begin{itemize}
\item Exact statement: For every admissible 'MCSP witness family', the route-side and target-side operationalizations are equivalent in the sense needed to transport 'MCSP complexity profile in nonuniform-circuits' toward 'In refined regimes, lower bounds for Gap-MCSP slightly beyond known frontiers (e.g. N\textasciicircum{} 1+varepsilon -type circuit lower bounds for suitable promise gaps and input length N=2\textasciicircum{}n ) imply major consequences such as NP not in P/poly.'.
\item Restricted-regime mode: equivalence-theorem
\item Bridge object: MCSP witness family (family of instances equipped with a monotone witness measure)
\item Bridge definition: 'MCSP witness family' denotes the structured intermediate object required by this branch. It is meant to organize the source-side content 'A hardness magnification theorem has the form:"A barely-nontrivial lower bound for an explicit problem Q in a weak model implies a major lower bound/separation in a strong model."Typically, Q is a meta-problem (e.g. a Gap-MCSP variant) and"major"means consequences such as NP not in P/poly.'. It must preserve the invariant 'MCSP complexity profile in nonuniform-circuits'. Its operational role is 'Recast the blocked frontier as an identity, equivalence, transport, or coherence problem rather than as a missing brute-force bridge.'.
\item First nontrivial lemma: For every admissible 'MCSP witness family', the target-side operationalization is equivalent to the route-side operationalization in the sense needed to transport 'MCSP complexity profile in nonuniform-circuits' toward 'In refined regimes, lower bounds for Gap-MCSP slightly beyond known frontiers (e.g. N\textasciicircum{} 1+varepsilon -type circuit lower bounds for suitable promise gaps and input length N=2\textasciicircum{}n ) imply major consequences such as NP not in P/poly.'.
\item Quantifier signature: For every admissible 'MCSP witness family' on the route from 'st.def.hardness-magnification-phenomenon-programme-leve' to 'st.thm.oliveira-pich-santhanam-magnification-near-state'.
\item LaTeX theorem: \textbackslash\{\}forall W,\textbackslash\{\} \textbackslash\{\}mathrm\{Adm\}(W) \textbackslash\{\}to \textbackslash\{\}exists w : \textbackslash\{\}mathsf\{BW\}(W),\textbackslash\{\} \textbackslash\{\}mathrm\{PresInv\}(W,w) \textbackslash\{\}wedge \textbackslash\{\}mathrm\{RealizesTarget\}(W,w).
\item Negative/barrier version: Absent an explicit equivalence witness for 'MCSP witness family', transport toward 'In refined regimes, lower bounds for Gap-MCSP slightly beyond known frontiers (e.g. N\textasciicircum{} 1+varepsilon -type circuit lower bounds for suitable promise gaps and input length N=2\textasciicircum{}n ) imply major consequences such as NP not in P/poly.' is not licensed. no\_equivalence\_witness
\item Proof strategy family: equivalence-transport
\item Negative route: Assume no equivalence witness for 'MCSP witness family' aligns the route-side and target-side objects.
\item Failure diagnosis: The route fails if no equivalence witness for 'MCSP witness family' transports 'MCSP complexity profile in nonuniform-circuits' to the target-side operationalization. Local negative probe: no\_equivalence\_witness.
\item Proof obligations: For every admissible 'MCSP witness family', the target-side operationalization is equivalent to the route-side operationalization in the sense needed to transport 'MCSP complexity profile in nonuniform-circuits' toward 'In refined regimes, lower bounds for Gap-MCSP slightly beyond known frontiers (e.g. N\textasciicircum{} 1+varepsilon -type circuit lower bounds for suitable promise gaps and input length N=2\textasciicircum{}n ) imply major consequences such as NP not in P/poly.'.; Transport obligation: move the invariant across the proposed equivalence for 'carrier-realization' without changing the target claim.; Transport obligation: move the invariant across the proposed equivalence for 'formula-size-transport' without changing the target claim.; Equivalence claim: determine whether the live carrier for 'carrier-realization' is identical to, equivalent to, or only analogous with the target-side object class.
\item Candidate lemmas: Equivalence theorem for MCSP witness family; Equivalence transport lemma for MCSP witness family
\item Independent value: An equivalence theorem for 'MCSP witness family' would independently matter because it identifies when two live operationalizations are mathematically interchangeable rather than merely analogous.
\item Expert gate: 6/6
\end{itemize}

\subsection*{conjpkg.disclayer-p-vs-np-family-equivalence-coherence-reframing.05}
\begin{itemize}
\item Exact statement: No admissible 'MCSP witness family' can simultaneously preserve 'MCSP complexity profile in nonuniform-circuits' and realize 'In refined regimes, lower bounds for Gap-MCSP slightly beyond known frontiers (e.g. N\textasciicircum{} 1+varepsilon -type circuit lower bounds for suitable promise gaps and input length N=2\textasciicircum{}n ) imply major consequences such as NP not in P/poly.' under the current route commitments.
\item Restricted-regime mode: equivalence-theorem
\item Bridge object: MCSP witness family (family of instances equipped with a monotone witness measure)
\item Bridge definition: 'MCSP witness family' denotes the structured intermediate object required by this branch. It is meant to organize the source-side content 'A hardness magnification theorem has the form:"A barely-nontrivial lower bound for an explicit problem Q in a weak model implies a major lower bound/separation in a strong model."Typically, Q is a meta-problem (e.g. a Gap-MCSP variant) and"major"means consequences such as NP not in P/poly.'. It must preserve the invariant 'MCSP complexity profile in nonuniform-circuits'. Its operational role is 'Recast the blocked frontier as an identity, equivalence, transport, or coherence problem rather than as a missing brute-force bridge.'.
\item First nontrivial lemma: For every admissible 'MCSP witness family', the target-side operationalization is equivalent to the route-side operationalization in the sense needed to transport 'MCSP complexity profile in nonuniform-circuits' toward 'In refined regimes, lower bounds for Gap-MCSP slightly beyond known frontiers (e.g. N\textasciicircum{} 1+varepsilon -type circuit lower bounds for suitable promise gaps and input length N=2\textasciicircum{}n ) imply major consequences such as NP not in P/poly.'.
\item Quantifier signature: For every admissible 'MCSP witness family' on the route from 'st.def.hardness-magnification-phenomenon-programme-leve' to 'st.thm.oliveira-pich-santhanam-magnification-near-state'.
\item LaTeX theorem: \textbackslash\{\}forall W,\textbackslash\{\} \textbackslash\{\}mathrm\{Adm\}(W) \textbackslash\{\}to \textbackslash\{\}neg \textbackslash\{\}exists w : \textbackslash\{\}mathsf\{BW\}(W),\textbackslash\{\} \textbackslash\{\}mathrm\{PresInv\}(W,w) \textbackslash\{\}wedge \textbackslash\{\}mathrm\{RealizesTarget\}(W,w).
\item Negative/barrier version: Absent an explicit equivalence witness for 'MCSP witness family', transport toward 'In refined regimes, lower bounds for Gap-MCSP slightly beyond known frontiers (e.g. N\textasciicircum{} 1+varepsilon -type circuit lower bounds for suitable promise gaps and input length N=2\textasciicircum{}n ) imply major consequences such as NP not in P/poly.' is not licensed. no\_equivalence\_witness
\item Proof strategy family: equivalence-transport
\item Negative route: Assume no equivalence witness for 'MCSP witness family' aligns the route-side and target-side objects.
\item Failure diagnosis: The route fails if no equivalence witness for 'MCSP witness family' transports 'MCSP complexity profile in nonuniform-circuits' to the target-side operationalization. Local negative probe: no\_equivalence\_witness.
\item Proof obligations: For every admissible 'MCSP witness family', the target-side operationalization is equivalent to the route-side operationalization in the sense needed to transport 'MCSP complexity profile in nonuniform-circuits' toward 'In refined regimes, lower bounds for Gap-MCSP slightly beyond known frontiers (e.g. N\textasciicircum{} 1+varepsilon -type circuit lower bounds for suitable promise gaps and input length N=2\textasciicircum{}n ) imply major consequences such as NP not in P/poly.'.; Transport obligation: move the invariant across the proposed equivalence for 'carrier-realization' without changing the target claim.; Transport obligation: move the invariant across the proposed equivalence for 'formula-size-transport' without changing the target claim.; Equivalence claim: determine whether the live carrier for 'carrier-realization' is identical to, equivalent to, or only analogous with the target-side object class.
\item Candidate lemmas: Obstruction theorem for MCSP witness family; Equivalence transport lemma for MCSP witness family
\item Independent value: An equivalence theorem for 'MCSP witness family' would independently matter because it identifies when two live operationalizations are mathematically interchangeable rather than merely analogous.
\item Expert gate: 6/6
\end{itemize}

\subsection*{conjpkg.disclayer-p-vs-np-family-equivalence-coherence-reframing.06}
\begin{itemize}
\item Exact statement: There exists an admissible 'MCSP witness family' whose realization internalizes 'A hardness magnification theorem has the form:"A barely-nontrivial lower bound for an explicit problem Q in a weak model implies a major lower bound/separation in a strong model."Typically, Q is a meta-problem (e.g. a Gap-MCSP variant) and"major"means consequences such as NP not in P/poly.' while preserving 'MCSP complexity profile in nonuniform-circuits' strongly enough to support 'In refined regimes, lower bounds for Gap-MCSP slightly beyond known frontiers (e.g. N\textasciicircum{} 1+varepsilon -type circuit lower bounds for suitable promise gaps and input length N=2\textasciicircum{}n ) imply major consequences such as NP not in P/poly.'.
\item Restricted-regime mode: equivalence-theorem
\item Bridge object: MCSP witness family (family of instances equipped with a monotone witness measure)
\item Bridge definition: 'MCSP witness family' denotes the structured intermediate object required by this branch. It is meant to organize the source-side content 'A hardness magnification theorem has the form:"A barely-nontrivial lower bound for an explicit problem Q in a weak model implies a major lower bound/separation in a strong model."Typically, Q is a meta-problem (e.g. a Gap-MCSP variant) and"major"means consequences such as NP not in P/poly.'. It must preserve the invariant 'MCSP complexity profile in nonuniform-circuits'. Its operational role is 'Recast the blocked frontier as an identity, equivalence, transport, or coherence problem rather than as a missing brute-force bridge.'.
\item First nontrivial lemma: For every admissible 'MCSP witness family', the target-side operationalization is equivalent to the route-side operationalization in the sense needed to transport 'MCSP complexity profile in nonuniform-circuits' toward 'In refined regimes, lower bounds for Gap-MCSP slightly beyond known frontiers (e.g. N\textasciicircum{} 1+varepsilon -type circuit lower bounds for suitable promise gaps and input length N=2\textasciicircum{}n ) imply major consequences such as NP not in P/poly.'.
\item Quantifier signature: For every admissible 'MCSP witness family' on the route from 'st.def.hardness-magnification-phenomenon-programme-leve' to 'st.thm.oliveira-pich-santhanam-magnification-near-state'.
\item LaTeX theorem: \textbackslash\{\}forall W,\textbackslash\{\} \textbackslash\{\}mathrm\{Adm\}(W) \textbackslash\{\}to \textbackslash\{\}exists w : \textbackslash\{\}mathsf\{BW\}(W),\textbackslash\{\} \textbackslash\{\}mathrm\{PresInv\}(W,w) \textbackslash\{\}wedge \textbackslash\{\}mathrm\{RealizesTarget\}(W,w).
\item Negative/barrier version: Absent an explicit equivalence witness for 'MCSP witness family', transport toward 'In refined regimes, lower bounds for Gap-MCSP slightly beyond known frontiers (e.g. N\textasciicircum{} 1+varepsilon -type circuit lower bounds for suitable promise gaps and input length N=2\textasciicircum{}n ) imply major consequences such as NP not in P/poly.' is not licensed. no\_equivalence\_witness
\item Proof strategy family: equivalence-transport
\item Negative route: Assume no equivalence witness for 'MCSP witness family' aligns the route-side and target-side objects.
\item Failure diagnosis: The route fails if no equivalence witness for 'MCSP witness family' transports 'MCSP complexity profile in nonuniform-circuits' to the target-side operationalization. Local negative probe: no\_equivalence\_witness.
\item Proof obligations: For every admissible 'MCSP witness family', the target-side operationalization is equivalent to the route-side operationalization in the sense needed to transport 'MCSP complexity profile in nonuniform-circuits' toward 'In refined regimes, lower bounds for Gap-MCSP slightly beyond known frontiers (e.g. N\textasciicircum{} 1+varepsilon -type circuit lower bounds for suitable promise gaps and input length N=2\textasciicircum{}n ) imply major consequences such as NP not in P/poly.'.; Transport obligation: move the invariant across the proposed equivalence for 'carrier-realization' without changing the target claim.; Transport obligation: move the invariant across the proposed equivalence for 'formula-size-transport' without changing the target claim.; Equivalence claim: determine whether the live carrier for 'carrier-realization' is identical to, equivalent to, or only analogous with the target-side object class.
\item Candidate lemmas: Realization theorem for MCSP witness family; Equivalence transport lemma for MCSP witness family
\item Independent value: An equivalence theorem for 'MCSP witness family' would independently matter because it identifies when two live operationalizations are mathematically interchangeable rather than merely analogous.
\item Expert gate: 6/6
\end{itemize}

\subsection*{conjpkg.disclayer-p-vs-np-family-equivalence-coherence-reframing.07}
\begin{itemize}
\item Exact statement: For every admissible 'MCSP witness family', there is a transport law from 'nonuniform-circuits' to 'formula-size' that preserves 'MCSP complexity profile in nonuniform-circuits' and retains enough strength to activate 'There exist explicit"meta"problems (including MCSP variants) such that proving certain emph barely-superlinear lower bounds against restricted models (e.g. formulas) would imply major lower bounds for NP in much stronger models (including superpolynomial formula/circuit lower bounds, and in some parameter regimes consequences as strong as NP not in P/poly ).'.
\item Restricted-regime mode: equivalence-theorem
\item Bridge object: MCSP witness family (family of instances equipped with a monotone witness measure)
\item Bridge definition: 'MCSP witness family' denotes the structured intermediate object required by this branch. It is meant to organize the source-side content 'A hardness magnification theorem has the form:"A barely-nontrivial lower bound for an explicit problem Q in a weak model implies a major lower bound/separation in a strong model."Typically, Q is a meta-problem (e.g. a Gap-MCSP variant) and"major"means consequences such as NP not in P/poly.'. It must preserve the invariant 'MCSP complexity profile in nonuniform-circuits'. Its operational role is 'Recast the blocked frontier as an identity, equivalence, transport, or coherence problem rather than as a missing brute-force bridge.'.
\item First nontrivial lemma: For every admissible 'MCSP witness family', the target-side operationalization is equivalent to the route-side operationalization in the sense needed to transport 'MCSP complexity profile in nonuniform-circuits' toward 'There exist explicit"meta"problems (including MCSP variants) such that proving certain emph barely-superlinear lower bounds against restricted models (e.g. formulas) would imply major lower bounds for NP in much stronger models (including superpolynomial formula/circuit lower bounds, and in some parameter regimes consequences as strong as NP not in P/poly ).'.
\item Quantifier signature: For every admissible 'MCSP witness family' on the route from 'st.def.hardness-magnification-phenomenon-programme-leve' to 'st.thm.oliveira-santhanam-hardness-magnification-for-na'.
\item LaTeX theorem: \textbackslash\{\}forall W,\textbackslash\{\} \textbackslash\{\}mathrm\{Adm\}(W) \textbackslash\{\}to \textbackslash\{\}exists w : \textbackslash\{\}mathsf\{BW\}(W),\textbackslash\{\} \textbackslash\{\}mathrm\{PresInv\}(W,w) \textbackslash\{\}wedge \textbackslash\{\}mathrm\{RealizesTarget\}(W,w).
\item Negative/barrier version: Absent an explicit equivalence witness for 'MCSP witness family', transport toward 'There exist explicit"meta"problems (including MCSP variants) such that proving certain emph barely-superlinear lower bounds against restricted models (e.g. formulas) would imply major lower bounds for NP in much stronger models (including superpolynomial formula/circuit lower bounds, and in some parameter regimes consequences as strong as NP not in P/poly ).' is not licensed. no\_equivalence\_witness
\item Proof strategy family: equivalence-transport
\item Negative route: Assume no equivalence witness for 'MCSP witness family' aligns the route-side and target-side objects.
\item Failure diagnosis: The route fails if no equivalence witness for 'MCSP witness family' transports 'MCSP complexity profile in nonuniform-circuits' to the target-side operationalization. Local negative probe: no\_equivalence\_witness.
\item Proof obligations: For every admissible 'MCSP witness family', the target-side operationalization is equivalent to the route-side operationalization in the sense needed to transport 'MCSP complexity profile in nonuniform-circuits' toward 'There exist explicit"meta"problems (including MCSP variants) such that proving certain emph barely-superlinear lower bounds against restricted models (e.g. formulas) would imply major lower bounds for NP in much stronger models (including superpolynomial formula/circuit lower bounds, and in some parameter regimes consequences as strong as NP not in P/poly ).'.; Transport obligation: move the invariant across the proposed equivalence for 'carrier-realization' without changing the target claim.; Transport obligation: move the invariant across the proposed equivalence for 'formula-size-transport' without changing the target claim.; Equivalence claim: determine whether the live carrier for 'carrier-realization' is identical to, equivalent to, or only analogous with the target-side object class.
\item Candidate lemmas: Transport theorem for MCSP witness family; Equivalence transport lemma for MCSP witness family
\item Independent value: An equivalence theorem for 'MCSP witness family' would independently matter because it identifies when two live operationalizations are mathematically interchangeable rather than merely analogous.
\item Expert gate: 6/6
\end{itemize}

\subsection*{conjpkg.disclayer-p-vs-np-family-equivalence-coherence-reframing.08}
\begin{itemize}
\item Exact statement: For every admissible 'MCSP witness family', the route-side and target-side operationalizations are equivalent in the sense needed to transport 'MCSP complexity profile in nonuniform-circuits' toward 'There exist explicit"meta"problems (including MCSP variants) such that proving certain emph barely-superlinear lower bounds against restricted models (e.g. formulas) would imply major lower bounds for NP in much stronger models (including superpolynomial formula/circuit lower bounds, and in some parameter regimes consequences as strong as NP not in P/poly ).'.
\item Restricted-regime mode: equivalence-theorem
\item Bridge object: MCSP witness family (family of instances equipped with a monotone witness measure)
\item Bridge definition: 'MCSP witness family' denotes the structured intermediate object required by this branch. It is meant to organize the source-side content 'A hardness magnification theorem has the form:"A barely-nontrivial lower bound for an explicit problem Q in a weak model implies a major lower bound/separation in a strong model."Typically, Q is a meta-problem (e.g. a Gap-MCSP variant) and"major"means consequences such as NP not in P/poly.'. It must preserve the invariant 'MCSP complexity profile in nonuniform-circuits'. Its operational role is 'Recast the blocked frontier as an identity, equivalence, transport, or coherence problem rather than as a missing brute-force bridge.'.
\item First nontrivial lemma: For every admissible 'MCSP witness family', the target-side operationalization is equivalent to the route-side operationalization in the sense needed to transport 'MCSP complexity profile in nonuniform-circuits' toward 'There exist explicit"meta"problems (including MCSP variants) such that proving certain emph barely-superlinear lower bounds against restricted models (e.g. formulas) would imply major lower bounds for NP in much stronger models (including superpolynomial formula/circuit lower bounds, and in some parameter regimes consequences as strong as NP not in P/poly ).'.
\item Quantifier signature: For every admissible 'MCSP witness family' on the route from 'st.def.hardness-magnification-phenomenon-programme-leve' to 'st.thm.oliveira-santhanam-hardness-magnification-for-na'.
\item LaTeX theorem: \textbackslash\{\}forall W,\textbackslash\{\} \textbackslash\{\}mathrm\{Adm\}(W) \textbackslash\{\}to \textbackslash\{\}exists w : \textbackslash\{\}mathsf\{BW\}(W),\textbackslash\{\} \textbackslash\{\}mathrm\{PresInv\}(W,w) \textbackslash\{\}wedge \textbackslash\{\}mathrm\{RealizesTarget\}(W,w).
\item Negative/barrier version: Absent an explicit equivalence witness for 'MCSP witness family', transport toward 'There exist explicit"meta"problems (including MCSP variants) such that proving certain emph barely-superlinear lower bounds against restricted models (e.g. formulas) would imply major lower bounds for NP in much stronger models (including superpolynomial formula/circuit lower bounds, and in some parameter regimes consequences as strong as NP not in P/poly ).' is not licensed. no\_equivalence\_witness
\item Proof strategy family: equivalence-transport
\item Negative route: Assume no equivalence witness for 'MCSP witness family' aligns the route-side and target-side objects.
\item Failure diagnosis: The route fails if no equivalence witness for 'MCSP witness family' transports 'MCSP complexity profile in nonuniform-circuits' to the target-side operationalization. Local negative probe: no\_equivalence\_witness.
\item Proof obligations: For every admissible 'MCSP witness family', the target-side operationalization is equivalent to the route-side operationalization in the sense needed to transport 'MCSP complexity profile in nonuniform-circuits' toward 'There exist explicit"meta"problems (including MCSP variants) such that proving certain emph barely-superlinear lower bounds against restricted models (e.g. formulas) would imply major lower bounds for NP in much stronger models (including superpolynomial formula/circuit lower bounds, and in some parameter regimes consequences as strong as NP not in P/poly ).'.; Transport obligation: move the invariant across the proposed equivalence for 'carrier-realization' without changing the target claim.; Transport obligation: move the invariant across the proposed equivalence for 'formula-size-transport' without changing the target claim.; Equivalence claim: determine whether the live carrier for 'carrier-realization' is identical to, equivalent to, or only analogous with the target-side object class.
\item Candidate lemmas: Equivalence theorem for MCSP witness family; Equivalence transport lemma for MCSP witness family
\item Independent value: An equivalence theorem for 'MCSP witness family' would independently matter because it identifies when two live operationalizations are mathematically interchangeable rather than merely analogous.
\item Expert gate: 6/6
\end{itemize}

\subsection*{conjpkg.disclayer-p-vs-np-family-equivalence-coherence-reframing.09}
\begin{itemize}
\item Exact statement: No admissible 'MCSP witness family' can simultaneously preserve 'MCSP complexity profile in nonuniform-circuits' and realize 'There exist explicit"meta"problems (including MCSP variants) such that proving certain emph barely-superlinear lower bounds against restricted models (e.g. formulas) would imply major lower bounds for NP in much stronger models (including superpolynomial formula/circuit lower bounds, and in some parameter regimes consequences as strong as NP not in P/poly ).' under the current route commitments.
\item Restricted-regime mode: equivalence-theorem
\item Bridge object: MCSP witness family (family of instances equipped with a monotone witness measure)
\item Bridge definition: 'MCSP witness family' denotes the structured intermediate object required by this branch. It is meant to organize the source-side content 'A hardness magnification theorem has the form:"A barely-nontrivial lower bound for an explicit problem Q in a weak model implies a major lower bound/separation in a strong model."Typically, Q is a meta-problem (e.g. a Gap-MCSP variant) and"major"means consequences such as NP not in P/poly.'. It must preserve the invariant 'MCSP complexity profile in nonuniform-circuits'. Its operational role is 'Recast the blocked frontier as an identity, equivalence, transport, or coherence problem rather than as a missing brute-force bridge.'.
\item First nontrivial lemma: For every admissible 'MCSP witness family', the target-side operationalization is equivalent to the route-side operationalization in the sense needed to transport 'MCSP complexity profile in nonuniform-circuits' toward 'There exist explicit"meta"problems (including MCSP variants) such that proving certain emph barely-superlinear lower bounds against restricted models (e.g. formulas) would imply major lower bounds for NP in much stronger models (including superpolynomial formula/circuit lower bounds, and in some parameter regimes consequences as strong as NP not in P/poly ).'.
\item Quantifier signature: For every admissible 'MCSP witness family' on the route from 'st.def.hardness-magnification-phenomenon-programme-leve' to 'st.thm.oliveira-santhanam-hardness-magnification-for-na'.
\item LaTeX theorem: \textbackslash\{\}forall W,\textbackslash\{\} \textbackslash\{\}mathrm\{Adm\}(W) \textbackslash\{\}to \textbackslash\{\}neg \textbackslash\{\}exists w : \textbackslash\{\}mathsf\{BW\}(W),\textbackslash\{\} \textbackslash\{\}mathrm\{PresInv\}(W,w) \textbackslash\{\}wedge \textbackslash\{\}mathrm\{RealizesTarget\}(W,w).
\item Negative/barrier version: Absent an explicit equivalence witness for 'MCSP witness family', transport toward 'There exist explicit"meta"problems (including MCSP variants) such that proving certain emph barely-superlinear lower bounds against restricted models (e.g. formulas) would imply major lower bounds for NP in much stronger models (including superpolynomial formula/circuit lower bounds, and in some parameter regimes consequences as strong as NP not in P/poly ).' is not licensed. no\_equivalence\_witness
\item Proof strategy family: equivalence-transport
\item Negative route: Assume no equivalence witness for 'MCSP witness family' aligns the route-side and target-side objects.
\item Failure diagnosis: The route fails if no equivalence witness for 'MCSP witness family' transports 'MCSP complexity profile in nonuniform-circuits' to the target-side operationalization. Local negative probe: no\_equivalence\_witness.
\item Proof obligations: For every admissible 'MCSP witness family', the target-side operationalization is equivalent to the route-side operationalization in the sense needed to transport 'MCSP complexity profile in nonuniform-circuits' toward 'There exist explicit"meta"problems (including MCSP variants) such that proving certain emph barely-superlinear lower bounds against restricted models (e.g. formulas) would imply major lower bounds for NP in much stronger models (including superpolynomial formula/circuit lower bounds, and in some parameter regimes consequences as strong as NP not in P/poly ).'.; Transport obligation: move the invariant across the proposed equivalence for 'carrier-realization' without changing the target claim.; Transport obligation: move the invariant across the proposed equivalence for 'formula-size-transport' without changing the target claim.; Equivalence claim: determine whether the live carrier for 'carrier-realization' is identical to, equivalent to, or only analogous with the target-side object class.
\item Candidate lemmas: Obstruction theorem for MCSP witness family; Equivalence transport lemma for MCSP witness family
\item Independent value: An equivalence theorem for 'MCSP witness family' would independently matter because it identifies when two live operationalizations are mathematically interchangeable rather than merely analogous.
\item Expert gate: 6/6
\end{itemize}

\subsection*{conjpkg.disclayer-p-vs-np-family-equivalence-coherence-reframing.10}
\begin{itemize}
\item Exact statement: There exists an admissible 'MCSP witness family' whose realization internalizes 'A hardness magnification theorem has the form:"A barely-nontrivial lower bound for an explicit problem Q in a weak model implies a major lower bound/separation in a strong model."Typically, Q is a meta-problem (e.g. a Gap-MCSP variant) and"major"means consequences such as NP not in P/poly.' while preserving 'MCSP complexity profile in nonuniform-circuits' strongly enough to support 'There exist explicit"meta"problems (including MCSP variants) such that proving certain emph barely-superlinear lower bounds against restricted models (e.g. formulas) would imply major lower bounds for NP in much stronger models (including superpolynomial formula/circuit lower bounds, and in some parameter regimes consequences as strong as NP not in P/poly ).'.
\item Restricted-regime mode: equivalence-theorem
\item Bridge object: MCSP witness family (family of instances equipped with a monotone witness measure)
\item Bridge definition: 'MCSP witness family' denotes the structured intermediate object required by this branch. It is meant to organize the source-side content 'A hardness magnification theorem has the form:"A barely-nontrivial lower bound for an explicit problem Q in a weak model implies a major lower bound/separation in a strong model."Typically, Q is a meta-problem (e.g. a Gap-MCSP variant) and"major"means consequences such as NP not in P/poly.'. It must preserve the invariant 'MCSP complexity profile in nonuniform-circuits'. Its operational role is 'Recast the blocked frontier as an identity, equivalence, transport, or coherence problem rather than as a missing brute-force bridge.'.
\item First nontrivial lemma: For every admissible 'MCSP witness family', the target-side operationalization is equivalent to the route-side operationalization in the sense needed to transport 'MCSP complexity profile in nonuniform-circuits' toward 'There exist explicit"meta"problems (including MCSP variants) such that proving certain emph barely-superlinear lower bounds against restricted models (e.g. formulas) would imply major lower bounds for NP in much stronger models (including superpolynomial formula/circuit lower bounds, and in some parameter regimes consequences as strong as NP not in P/poly ).'.
\item Quantifier signature: For every admissible 'MCSP witness family' on the route from 'st.def.hardness-magnification-phenomenon-programme-leve' to 'st.thm.oliveira-santhanam-hardness-magnification-for-na'.
\item LaTeX theorem: \textbackslash\{\}forall W,\textbackslash\{\} \textbackslash\{\}mathrm\{Adm\}(W) \textbackslash\{\}to \textbackslash\{\}exists w : \textbackslash\{\}mathsf\{BW\}(W),\textbackslash\{\} \textbackslash\{\}mathrm\{PresInv\}(W,w) \textbackslash\{\}wedge \textbackslash\{\}mathrm\{RealizesTarget\}(W,w).
\item Negative/barrier version: Absent an explicit equivalence witness for 'MCSP witness family', transport toward 'There exist explicit"meta"problems (including MCSP variants) such that proving certain emph barely-superlinear lower bounds against restricted models (e.g. formulas) would imply major lower bounds for NP in much stronger models (including superpolynomial formula/circuit lower bounds, and in some parameter regimes consequences as strong as NP not in P/poly ).' is not licensed. no\_equivalence\_witness
\item Proof strategy family: equivalence-transport
\item Negative route: Assume no equivalence witness for 'MCSP witness family' aligns the route-side and target-side objects.
\item Failure diagnosis: The route fails if no equivalence witness for 'MCSP witness family' transports 'MCSP complexity profile in nonuniform-circuits' to the target-side operationalization. Local negative probe: no\_equivalence\_witness.
\item Proof obligations: For every admissible 'MCSP witness family', the target-side operationalization is equivalent to the route-side operationalization in the sense needed to transport 'MCSP complexity profile in nonuniform-circuits' toward 'There exist explicit"meta"problems (including MCSP variants) such that proving certain emph barely-superlinear lower bounds against restricted models (e.g. formulas) would imply major lower bounds for NP in much stronger models (including superpolynomial formula/circuit lower bounds, and in some parameter regimes consequences as strong as NP not in P/poly ).'.; Transport obligation: move the invariant across the proposed equivalence for 'carrier-realization' without changing the target claim.; Transport obligation: move the invariant across the proposed equivalence for 'formula-size-transport' without changing the target claim.; Equivalence claim: determine whether the live carrier for 'carrier-realization' is identical to, equivalent to, or only analogous with the target-side object class.
\item Candidate lemmas: Realization theorem for MCSP witness family; Equivalence transport lemma for MCSP witness family
\item Independent value: An equivalence theorem for 'MCSP witness family' would independently matter because it identifies when two live operationalizations are mathematically interchangeable rather than merely analogous.
\item Expert gate: 6/6
\end{itemize}

\subsection*{conjpkg.disclayer-p-vs-np-family-equivalence-coherence-reframing.11}
\begin{itemize}
\item Exact statement: For every admissible 'MCSP witness family', there is a transport law from 'nonuniform-circuits' to 'unspecified' that preserves 'exponential-lower-bound witness complexity' and retains enough strength to activate 'There exist general-purpose magnification theorems (not problem-specific) showing that certain mild lower bounds for distinguishers/sparse-code objects can be amplified into stronger circuit/formula lower bounds, achieving magnification thresholds below previously known lower bounds in some settings.'.
\item Restricted-regime mode: equivalence-theorem
\item Bridge object: MCSP witness family (family of instances equipped with a monotone witness measure)
\item Bridge definition: 'MCSP witness family' denotes the structured intermediate object required by this branch. It is meant to organize the source-side content 'A hardness magnification theorem has the form:"A barely-nontrivial lower bound for an explicit problem Q in a weak model implies a major lower bound/separation in a strong model."Typically, Q is a meta-problem (e.g. a Gap-MCSP variant) and"major"means consequences such as NP not in P/poly.'. It must preserve the invariant 'exponential-lower-bound witness complexity'. Its operational role is 'Recast the blocked frontier as an identity, equivalence, transport, or coherence problem rather than as a missing brute-force bridge.'.
\item First nontrivial lemma: For every admissible 'MCSP witness family', the target-side operationalization is equivalent to the route-side operationalization in the sense needed to transport 'exponential-lower-bound witness complexity' toward 'There exist general-purpose magnification theorems (not problem-specific) showing that certain mild lower bounds for distinguishers/sparse-code objects can be amplified into stronger circuit/formula lower bounds, achieving magnification thresholds below previously known lower bounds in some settings.'.
\item Quantifier signature: For every admissible 'MCSP witness family' on the route from 'st.def.hardness-magnification-phenomenon-programme-leve' to 'st.thm.atserias-m-uller-simple-general-magnification-fr'.
\item LaTeX theorem: \textbackslash\{\}forall W,\textbackslash\{\} \textbackslash\{\}mathrm\{Adm\}(W) \textbackslash\{\}to \textbackslash\{\}exists w : \textbackslash\{\}mathsf\{BW\}(W),\textbackslash\{\} \textbackslash\{\}mathrm\{PresInv\}(W,w) \textbackslash\{\}wedge \textbackslash\{\}mathrm\{RealizesTarget\}(W,w).
\item Negative/barrier version: Absent an explicit equivalence witness for 'MCSP witness family', transport toward 'There exist general-purpose magnification theorems (not problem-specific) showing that certain mild lower bounds for distinguishers/sparse-code objects can be amplified into stronger circuit/formula lower bounds, achieving magnification thresholds below previously known lower bounds in some settings.' is not licensed. no\_equivalence\_witness
\item Proof strategy family: equivalence-transport
\item Negative route: Assume no equivalence witness for 'MCSP witness family' aligns the route-side and target-side objects.
\item Failure diagnosis: The route fails if no equivalence witness for 'MCSP witness family' transports 'exponential-lower-bound witness complexity' to the target-side operationalization. Local negative probe: no\_equivalence\_witness.
\item Proof obligations: For every admissible 'MCSP witness family', the target-side operationalization is equivalent to the route-side operationalization in the sense needed to transport 'exponential-lower-bound witness complexity' toward 'There exist general-purpose magnification theorems (not problem-specific) showing that certain mild lower bounds for distinguishers/sparse-code objects can be amplified into stronger circuit/formula lower bounds, achieving magnification thresholds below previously known lower bounds in some settings.'.; Transport obligation: move the invariant across the proposed equivalence for 'carrier-realization' without changing the target claim.; Transport obligation: move the invariant across the proposed equivalence for 'formula-size-transport' without changing the target claim.; Equivalence claim: determine whether the live carrier for 'carrier-realization' is identical to, equivalent to, or only analogous with the target-side object class.
\item Candidate lemmas: Transport theorem for MCSP witness family; Equivalence transport lemma for MCSP witness family
\item Independent value: An equivalence theorem for 'MCSP witness family' would independently matter because it identifies when two live operationalizations are mathematically interchangeable rather than merely analogous.
\item Expert gate: 6/6
\end{itemize}

\subsection*{conjpkg.disclayer-p-vs-np-family-equivalence-coherence-reframing.12}
\begin{itemize}
\item Exact statement: For every admissible 'MCSP witness family', the route-side and target-side operationalizations are equivalent in the sense needed to transport 'exponential-lower-bound witness complexity' toward 'There exist general-purpose magnification theorems (not problem-specific) showing that certain mild lower bounds for distinguishers/sparse-code objects can be amplified into stronger circuit/formula lower bounds, achieving magnification thresholds below previously known lower bounds in some settings.'.
\item Restricted-regime mode: equivalence-theorem
\item Bridge object: MCSP witness family (family of instances equipped with a monotone witness measure)
\item Bridge definition: 'MCSP witness family' denotes the structured intermediate object required by this branch. It is meant to organize the source-side content 'A hardness magnification theorem has the form:"A barely-nontrivial lower bound for an explicit problem Q in a weak model implies a major lower bound/separation in a strong model."Typically, Q is a meta-problem (e.g. a Gap-MCSP variant) and"major"means consequences such as NP not in P/poly.'. It must preserve the invariant 'exponential-lower-bound witness complexity'. Its operational role is 'Recast the blocked frontier as an identity, equivalence, transport, or coherence problem rather than as a missing brute-force bridge.'.
\item First nontrivial lemma: For every admissible 'MCSP witness family', the target-side operationalization is equivalent to the route-side operationalization in the sense needed to transport 'exponential-lower-bound witness complexity' toward 'There exist general-purpose magnification theorems (not problem-specific) showing that certain mild lower bounds for distinguishers/sparse-code objects can be amplified into stronger circuit/formula lower bounds, achieving magnification thresholds below previously known lower bounds in some settings.'.
\item Quantifier signature: For every admissible 'MCSP witness family' on the route from 'st.def.hardness-magnification-phenomenon-programme-leve' to 'st.thm.atserias-m-uller-simple-general-magnification-fr'.
\item LaTeX theorem: \textbackslash\{\}forall W,\textbackslash\{\} \textbackslash\{\}mathrm\{Adm\}(W) \textbackslash\{\}to \textbackslash\{\}exists w : \textbackslash\{\}mathsf\{BW\}(W),\textbackslash\{\} \textbackslash\{\}mathrm\{PresInv\}(W,w) \textbackslash\{\}wedge \textbackslash\{\}mathrm\{RealizesTarget\}(W,w).
\item Negative/barrier version: Absent an explicit equivalence witness for 'MCSP witness family', transport toward 'There exist general-purpose magnification theorems (not problem-specific) showing that certain mild lower bounds for distinguishers/sparse-code objects can be amplified into stronger circuit/formula lower bounds, achieving magnification thresholds below previously known lower bounds in some settings.' is not licensed. no\_equivalence\_witness
\item Proof strategy family: equivalence-transport
\item Negative route: Assume no equivalence witness for 'MCSP witness family' aligns the route-side and target-side objects.
\item Failure diagnosis: The route fails if no equivalence witness for 'MCSP witness family' transports 'exponential-lower-bound witness complexity' to the target-side operationalization. Local negative probe: no\_equivalence\_witness.
\item Proof obligations: For every admissible 'MCSP witness family', the target-side operationalization is equivalent to the route-side operationalization in the sense needed to transport 'exponential-lower-bound witness complexity' toward 'There exist general-purpose magnification theorems (not problem-specific) showing that certain mild lower bounds for distinguishers/sparse-code objects can be amplified into stronger circuit/formula lower bounds, achieving magnification thresholds below previously known lower bounds in some settings.'.; Transport obligation: move the invariant across the proposed equivalence for 'carrier-realization' without changing the target claim.; Transport obligation: move the invariant across the proposed equivalence for 'formula-size-transport' without changing the target claim.; Equivalence claim: determine whether the live carrier for 'carrier-realization' is identical to, equivalent to, or only analogous with the target-side object class.
\item Candidate lemmas: Equivalence theorem for MCSP witness family; Equivalence transport lemma for MCSP witness family
\item Independent value: An equivalence theorem for 'MCSP witness family' would independently matter because it identifies when two live operationalizations are mathematically interchangeable rather than merely analogous.
\item Expert gate: 6/6
\end{itemize}

\subsection*{conjpkg.disclayer-p-vs-np-family-equivalence-coherence-reframing.13}
\begin{itemize}
\item Exact statement: No admissible 'MCSP witness family' can simultaneously preserve 'exponential-lower-bound witness complexity' and realize 'There exist general-purpose magnification theorems (not problem-specific) showing that certain mild lower bounds for distinguishers/sparse-code objects can be amplified into stronger circuit/formula lower bounds, achieving magnification thresholds below previously known lower bounds in some settings.' under the current route commitments.
\item Restricted-regime mode: equivalence-theorem
\item Bridge object: MCSP witness family (family of instances equipped with a monotone witness measure)
\item Bridge definition: 'MCSP witness family' denotes the structured intermediate object required by this branch. It is meant to organize the source-side content 'A hardness magnification theorem has the form:"A barely-nontrivial lower bound for an explicit problem Q in a weak model implies a major lower bound/separation in a strong model."Typically, Q is a meta-problem (e.g. a Gap-MCSP variant) and"major"means consequences such as NP not in P/poly.'. It must preserve the invariant 'exponential-lower-bound witness complexity'. Its operational role is 'Recast the blocked frontier as an identity, equivalence, transport, or coherence problem rather than as a missing brute-force bridge.'.
\item First nontrivial lemma: For every admissible 'MCSP witness family', the target-side operationalization is equivalent to the route-side operationalization in the sense needed to transport 'exponential-lower-bound witness complexity' toward 'There exist general-purpose magnification theorems (not problem-specific) showing that certain mild lower bounds for distinguishers/sparse-code objects can be amplified into stronger circuit/formula lower bounds, achieving magnification thresholds below previously known lower bounds in some settings.'.
\item Quantifier signature: For every admissible 'MCSP witness family' on the route from 'st.def.hardness-magnification-phenomenon-programme-leve' to 'st.thm.atserias-m-uller-simple-general-magnification-fr'.
\item LaTeX theorem: \textbackslash\{\}forall W,\textbackslash\{\} \textbackslash\{\}mathrm\{Adm\}(W) \textbackslash\{\}to \textbackslash\{\}neg \textbackslash\{\}exists w : \textbackslash\{\}mathsf\{BW\}(W),\textbackslash\{\} \textbackslash\{\}mathrm\{PresInv\}(W,w) \textbackslash\{\}wedge \textbackslash\{\}mathrm\{RealizesTarget\}(W,w).
\item Negative/barrier version: Absent an explicit equivalence witness for 'MCSP witness family', transport toward 'There exist general-purpose magnification theorems (not problem-specific) showing that certain mild lower bounds for distinguishers/sparse-code objects can be amplified into stronger circuit/formula lower bounds, achieving magnification thresholds below previously known lower bounds in some settings.' is not licensed. no\_equivalence\_witness
\item Proof strategy family: equivalence-transport
\item Negative route: Assume no equivalence witness for 'MCSP witness family' aligns the route-side and target-side objects.
\item Failure diagnosis: The route fails if no equivalence witness for 'MCSP witness family' transports 'exponential-lower-bound witness complexity' to the target-side operationalization. Local negative probe: no\_equivalence\_witness.
\item Proof obligations: For every admissible 'MCSP witness family', the target-side operationalization is equivalent to the route-side operationalization in the sense needed to transport 'exponential-lower-bound witness complexity' toward 'There exist general-purpose magnification theorems (not problem-specific) showing that certain mild lower bounds for distinguishers/sparse-code objects can be amplified into stronger circuit/formula lower bounds, achieving magnification thresholds below previously known lower bounds in some settings.'.; Transport obligation: move the invariant across the proposed equivalence for 'carrier-realization' without changing the target claim.; Transport obligation: move the invariant across the proposed equivalence for 'formula-size-transport' without changing the target claim.; Equivalence claim: determine whether the live carrier for 'carrier-realization' is identical to, equivalent to, or only analogous with the target-side object class.
\item Candidate lemmas: Obstruction theorem for MCSP witness family; Equivalence transport lemma for MCSP witness family
\item Independent value: An equivalence theorem for 'MCSP witness family' would independently matter because it identifies when two live operationalizations are mathematically interchangeable rather than merely analogous.
\item Expert gate: 6/6
\end{itemize}

\subsection*{conjpkg.disclayer-p-vs-np-family-equivalence-coherence-reframing.14}
\begin{itemize}
\item Exact statement: There exists an admissible 'MCSP witness family' whose realization internalizes 'A hardness magnification theorem has the form:"A barely-nontrivial lower bound for an explicit problem Q in a weak model implies a major lower bound/separation in a strong model."Typically, Q is a meta-problem (e.g. a Gap-MCSP variant) and"major"means consequences such as NP not in P/poly.' while preserving 'exponential-lower-bound witness complexity' strongly enough to support 'There exist general-purpose magnification theorems (not problem-specific) showing that certain mild lower bounds for distinguishers/sparse-code objects can be amplified into stronger circuit/formula lower bounds, achieving magnification thresholds below previously known lower bounds in some settings.'.
\item Restricted-regime mode: equivalence-theorem
\item Bridge object: MCSP witness family (family of instances equipped with a monotone witness measure)
\item Bridge definition: 'MCSP witness family' denotes the structured intermediate object required by this branch. It is meant to organize the source-side content 'A hardness magnification theorem has the form:"A barely-nontrivial lower bound for an explicit problem Q in a weak model implies a major lower bound/separation in a strong model."Typically, Q is a meta-problem (e.g. a Gap-MCSP variant) and"major"means consequences such as NP not in P/poly.'. It must preserve the invariant 'exponential-lower-bound witness complexity'. Its operational role is 'Recast the blocked frontier as an identity, equivalence, transport, or coherence problem rather than as a missing brute-force bridge.'.
\item First nontrivial lemma: For every admissible 'MCSP witness family', the target-side operationalization is equivalent to the route-side operationalization in the sense needed to transport 'exponential-lower-bound witness complexity' toward 'There exist general-purpose magnification theorems (not problem-specific) showing that certain mild lower bounds for distinguishers/sparse-code objects can be amplified into stronger circuit/formula lower bounds, achieving magnification thresholds below previously known lower bounds in some settings.'.
\item Quantifier signature: For every admissible 'MCSP witness family' on the route from 'st.def.hardness-magnification-phenomenon-programme-leve' to 'st.thm.atserias-m-uller-simple-general-magnification-fr'.
\item LaTeX theorem: \textbackslash\{\}forall W,\textbackslash\{\} \textbackslash\{\}mathrm\{Adm\}(W) \textbackslash\{\}to \textbackslash\{\}exists w : \textbackslash\{\}mathsf\{BW\}(W),\textbackslash\{\} \textbackslash\{\}mathrm\{PresInv\}(W,w) \textbackslash\{\}wedge \textbackslash\{\}mathrm\{RealizesTarget\}(W,w).
\item Negative/barrier version: Absent an explicit equivalence witness for 'MCSP witness family', transport toward 'There exist general-purpose magnification theorems (not problem-specific) showing that certain mild lower bounds for distinguishers/sparse-code objects can be amplified into stronger circuit/formula lower bounds, achieving magnification thresholds below previously known lower bounds in some settings.' is not licensed. no\_equivalence\_witness
\item Proof strategy family: equivalence-transport
\item Negative route: Assume no equivalence witness for 'MCSP witness family' aligns the route-side and target-side objects.
\item Failure diagnosis: The route fails if no equivalence witness for 'MCSP witness family' transports 'exponential-lower-bound witness complexity' to the target-side operationalization. Local negative probe: no\_equivalence\_witness.
\item Proof obligations: For every admissible 'MCSP witness family', the target-side operationalization is equivalent to the route-side operationalization in the sense needed to transport 'exponential-lower-bound witness complexity' toward 'There exist general-purpose magnification theorems (not problem-specific) showing that certain mild lower bounds for distinguishers/sparse-code objects can be amplified into stronger circuit/formula lower bounds, achieving magnification thresholds below previously known lower bounds in some settings.'.; Transport obligation: move the invariant across the proposed equivalence for 'carrier-realization' without changing the target claim.; Transport obligation: move the invariant across the proposed equivalence for 'formula-size-transport' without changing the target claim.; Equivalence claim: determine whether the live carrier for 'carrier-realization' is identical to, equivalent to, or only analogous with the target-side object class.
\item Candidate lemmas: Realization theorem for MCSP witness family; Equivalence transport lemma for MCSP witness family
\item Independent value: An equivalence theorem for 'MCSP witness family' would independently matter because it identifies when two live operationalizations are mathematically interchangeable rather than merely analogous.
\item Expert gate: 6/6
\end{itemize}

\subsection*{conjpkg.disclayer-p-vs-np-family-categorical-reframing.01}
\begin{itemize}
\item Exact statement: There exists a mediating object built from 'MCSP witness family' such that the route from 'A hardness magnification theorem has the form:"A barely-nontrivial lower bound for an explicit problem Q in a weak model implies a major lower bound/separation in a strong model."Typically, Q is a meta-problem (e.g. a Gap-MCSP variant) and"major"means consequences such as NP not in P/poly.' to 'In refined regimes, lower bounds for Gap-MCSP slightly beyond known frontiers (e.g. N\textasciicircum{} 1+varepsilon -type circuit lower bounds for suitable promise gaps and input length N=2\textasciicircum{}n ) imply major consequences such as NP not in P/poly.' factors through typed morphisms preserving 'MCSP complexity profile in nonuniform-circuits' and making the relevant square commute.
\item Restricted-regime mode: special-case
\item Bridge object: MCSP witness family (family of instances equipped with a monotone witness measure)
\item Bridge definition: 'MCSP witness family' denotes the structured intermediate object required by this branch. It is meant to organize the source-side content 'A hardness magnification theorem has the form:"A barely-nontrivial lower bound for an explicit problem Q in a weak model implies a major lower bound/separation in a strong model."Typically, Q is a meta-problem (e.g. a Gap-MCSP variant) and"major"means consequences such as NP not in P/poly.'. It must preserve the invariant 'MCSP complexity profile in nonuniform-circuits'. Its operational role is 'Recast the blocked frontier as a diagrammatic composition problem: identify objects, morphisms, mediating objects, and the exact commutativity or factorization obligation.'.
\item First nontrivial lemma: For every admissible 'MCSP witness family', the first route-specific realization step preserves 'MCSP complexity profile in nonuniform-circuits' and prepares the passage from 'A hardness magnification theorem has the form:"A barely-nontrivial lower bound for an explicit problem Q in a weak model implies a major lower bound/separation in a strong model."Typically, Q is a meta-problem (e.g. a Gap-MCSP variant) and"major"means consequences such as NP not in P/poly.' to 'In refined regimes, lower bounds for Gap-MCSP slightly beyond known frontiers (e.g. N\textasciicircum{} 1+varepsilon -type circuit lower bounds for suitable promise gaps and input length N=2\textasciicircum{}n ) imply major consequences such as NP not in P/poly.'.
\item Quantifier signature: For every admissible 'MCSP witness family' on the route from 'st.def.hardness-magnification-phenomenon-programme-leve' to 'st.thm.oliveira-pich-santhanam-magnification-near-state'.
\item LaTeX theorem: \textbackslash\{\}forall W,\textbackslash\{\} \textbackslash\{\}mathrm\{Adm\}(W) \textbackslash\{\}to \textbackslash\{\}exists w : \textbackslash\{\}mathsf\{BW\}(W),\textbackslash\{\} \textbackslash\{\}mathrm\{PresInv\}(W,w) \textbackslash\{\}wedge \textbackslash\{\}mathrm\{RealizesTarget\}(W,w).
\item Negative/barrier version: If 'MCSP witness family' cannot be realized with invariant 'MCSP complexity profile in nonuniform-circuits', the route cannot force 'In refined regimes, lower bounds for Gap-MCSP slightly beyond known frontiers (e.g. N\textasciicircum{} 1+varepsilon -type circuit lower bounds for suitable promise gaps and input length N=2\textasciicircum{}n ) imply major consequences such as NP not in P/poly.'. non\_commuting\_route
\item Proof strategy family: constructive-realization
\item Negative route: Assume 'MCSP witness family' cannot be realized while preserving 'MCSP complexity profile in nonuniform-circuits' strongly enough to activate the local target claim.
\item Failure diagnosis: The route fails if 'MCSP witness family' cannot be realized while preserving 'MCSP complexity profile in nonuniform-circuits'. Local negative probe: non\_commuting\_route.
\item Proof obligations: For every admissible 'MCSP witness family', the first route-specific realization step preserves 'MCSP complexity profile in nonuniform-circuits' and prepares the passage from 'A hardness magnification theorem has the form:"A barely-nontrivial lower bound for an explicit problem Q in a weak model implies a major lower bound/separation in a strong model."Typically, Q is a meta-problem (e.g. a Gap-MCSP variant) and"major"means consequences such as NP not in P/poly.' to 'In refined regimes, lower bounds for Gap-MCSP slightly beyond known frontiers (e.g. N\textasciicircum{} 1+varepsilon -type circuit lower bounds for suitable promise gaps and input length N=2\textasciicircum{}n ) imply major consequences such as NP not in P/poly.'.; Commutativity obligation: the route for 'carrier-realization' must commute between source theorem hypotheses, the mediating carrier, and the target theorem hypotheses.
\item Candidate lemmas: Commutative factorization theorem for MCSP witness family; Coherence diagram lemma for MCSP witness family
\item Independent value: A concrete special-case theorem for 'MCSP witness family' would be independently valuable because it converts a broad discovery programme into an inspectable mathematical subresult.
\item Expert gate: 6/6
\end{itemize}

\subsection*{conjpkg.disclayer-p-vs-np-family-categorical-reframing.02}
\begin{itemize}
\item Exact statement: There exists a mediating object built from 'MCSP witness family' such that the route from 'A hardness magnification theorem has the form:"A barely-nontrivial lower bound for an explicit problem Q in a weak model implies a major lower bound/separation in a strong model."Typically, Q is a meta-problem (e.g. a Gap-MCSP variant) and"major"means consequences such as NP not in P/poly.' to 'There exist explicit"meta"problems (including MCSP variants) such that proving certain emph barely-superlinear lower bounds against restricted models (e.g. formulas) would imply major lower bounds for NP in much stronger models (including superpolynomial formula/circuit lower bounds, and in some parameter regimes consequences as strong as NP not in P/poly ).' factors through typed morphisms preserving 'MCSP complexity profile in nonuniform-circuits' and making the relevant square commute.
\item Restricted-regime mode: special-case
\item Bridge object: MCSP witness family (family of instances equipped with a monotone witness measure)
\item Bridge definition: 'MCSP witness family' denotes the structured intermediate object required by this branch. It is meant to organize the source-side content 'A hardness magnification theorem has the form:"A barely-nontrivial lower bound for an explicit problem Q in a weak model implies a major lower bound/separation in a strong model."Typically, Q is a meta-problem (e.g. a Gap-MCSP variant) and"major"means consequences such as NP not in P/poly.'. It must preserve the invariant 'MCSP complexity profile in nonuniform-circuits'. Its operational role is 'Recast the blocked frontier as a diagrammatic composition problem: identify objects, morphisms, mediating objects, and the exact commutativity or factorization obligation.'.
\item First nontrivial lemma: For every admissible 'MCSP witness family', the first route-specific realization step preserves 'MCSP complexity profile in nonuniform-circuits' and prepares the passage from 'A hardness magnification theorem has the form:"A barely-nontrivial lower bound for an explicit problem Q in a weak model implies a major lower bound/separation in a strong model."Typically, Q is a meta-problem (e.g. a Gap-MCSP variant) and"major"means consequences such as NP not in P/poly.' to 'There exist explicit"meta"problems (including MCSP variants) such that proving certain emph barely-superlinear lower bounds against restricted models (e.g. formulas) would imply major lower bounds for NP in much stronger models (including superpolynomial formula/circuit lower bounds, and in some parameter regimes consequences as strong as NP not in P/poly ).'.
\item Quantifier signature: For every admissible 'MCSP witness family' on the route from 'st.def.hardness-magnification-phenomenon-programme-leve' to 'st.thm.oliveira-santhanam-hardness-magnification-for-na'.
\item LaTeX theorem: \textbackslash\{\}forall W,\textbackslash\{\} \textbackslash\{\}mathrm\{Adm\}(W) \textbackslash\{\}to \textbackslash\{\}exists w : \textbackslash\{\}mathsf\{BW\}(W),\textbackslash\{\} \textbackslash\{\}mathrm\{PresInv\}(W,w) \textbackslash\{\}wedge \textbackslash\{\}mathrm\{RealizesTarget\}(W,w).
\item Negative/barrier version: If 'MCSP witness family' cannot be realized with invariant 'MCSP complexity profile in nonuniform-circuits', the route cannot force 'There exist explicit"meta"problems (including MCSP variants) such that proving certain emph barely-superlinear lower bounds against restricted models (e.g. formulas) would imply major lower bounds for NP in much stronger models (including superpolynomial formula/circuit lower bounds, and in some parameter regimes consequences as strong as NP not in P/poly ).'. non\_commuting\_route
\item Proof strategy family: constructive-realization
\item Negative route: Assume 'MCSP witness family' cannot be realized while preserving 'MCSP complexity profile in nonuniform-circuits' strongly enough to activate the local target claim.
\item Failure diagnosis: The route fails if 'MCSP witness family' cannot be realized while preserving 'MCSP complexity profile in nonuniform-circuits'. Local negative probe: non\_commuting\_route.
\item Proof obligations: For every admissible 'MCSP witness family', the first route-specific realization step preserves 'MCSP complexity profile in nonuniform-circuits' and prepares the passage from 'A hardness magnification theorem has the form:"A barely-nontrivial lower bound for an explicit problem Q in a weak model implies a major lower bound/separation in a strong model."Typically, Q is a meta-problem (e.g. a Gap-MCSP variant) and"major"means consequences such as NP not in P/poly.' to 'There exist explicit"meta"problems (including MCSP variants) such that proving certain emph barely-superlinear lower bounds against restricted models (e.g. formulas) would imply major lower bounds for NP in much stronger models (including superpolynomial formula/circuit lower bounds, and in some parameter regimes consequences as strong as NP not in P/poly ).'.; Commutativity obligation: the route for 'carrier-realization' must commute between source theorem hypotheses, the mediating carrier, and the target theorem hypotheses.
\item Candidate lemmas: Commutative factorization theorem for MCSP witness family; Coherence diagram lemma for MCSP witness family
\item Independent value: A concrete special-case theorem for 'MCSP witness family' would be independently valuable because it converts a broad discovery programme into an inspectable mathematical subresult.
\item Expert gate: 6/6
\end{itemize}

\subsection*{conjpkg.disclayer-p-vs-np-family-categorical-reframing.03}
\begin{itemize}
\item Exact statement: There exists a mediating object built from 'MCSP obstruction family' such that the route from 'A hardness magnification theorem has the form:"A barely-nontrivial lower bound for an explicit problem Q in a weak model implies a major lower bound/separation in a strong model."Typically, Q is a meta-problem (e.g. a Gap-MCSP variant) and"major"means consequences such as NP not in P/poly.' to 'In refined regimes, lower bounds for Gap-MCSP slightly beyond known frontiers (e.g. N\textasciicircum{} 1+varepsilon -type circuit lower bounds for suitable promise gaps and input length N=2\textasciicircum{}n ) imply major consequences such as NP not in P/poly.' factors through typed morphisms preserving 'MCSP complexity profile in nonuniform-circuits' and making the relevant square commute.
\item Restricted-regime mode: special-case
\item Bridge object: MCSP obstruction family (family of hard instances or obstructions forcing the target payoff)
\item Bridge definition: 'MCSP obstruction family' denotes the structured intermediate object required by this branch. It is meant to organize the source-side content 'A hardness magnification theorem has the form:"A barely-nontrivial lower bound for an explicit problem Q in a weak model implies a major lower bound/separation in a strong model."Typically, Q is a meta-problem (e.g. a Gap-MCSP variant) and"major"means consequences such as NP not in P/poly.'. It must preserve the invariant 'MCSP complexity profile in nonuniform-circuits'. Its operational role is 'Recast the blocked frontier as a diagrammatic composition problem: identify objects, morphisms, mediating objects, and the exact commutativity or factorization obligation.'.
\item First nontrivial lemma: For every admissible 'MCSP obstruction family', the first route-specific realization step preserves 'MCSP complexity profile in nonuniform-circuits' and prepares the passage from 'A hardness magnification theorem has the form:"A barely-nontrivial lower bound for an explicit problem Q in a weak model implies a major lower bound/separation in a strong model."Typically, Q is a meta-problem (e.g. a Gap-MCSP variant) and"major"means consequences such as NP not in P/poly.' to 'In refined regimes, lower bounds for Gap-MCSP slightly beyond known frontiers (e.g. N\textasciicircum{} 1+varepsilon -type circuit lower bounds for suitable promise gaps and input length N=2\textasciicircum{}n ) imply major consequences such as NP not in P/poly.'.
\item Quantifier signature: For every admissible 'MCSP obstruction family' on the route from 'st.def.hardness-magnification-phenomenon-programme-leve' to 'st.thm.oliveira-pich-santhanam-magnification-near-state'.
\item LaTeX theorem: \textbackslash\{\}forall W,\textbackslash\{\} \textbackslash\{\}mathrm\{Adm\}(W) \textbackslash\{\}to \textbackslash\{\}exists w : \textbackslash\{\}mathsf\{BW\}(W),\textbackslash\{\} \textbackslash\{\}mathrm\{PresInv\}(W,w) \textbackslash\{\}wedge \textbackslash\{\}mathrm\{RealizesTarget\}(W,w).
\item Negative/barrier version: If 'MCSP obstruction family' cannot be realized with invariant 'MCSP complexity profile in nonuniform-circuits', the route cannot force 'In refined regimes, lower bounds for Gap-MCSP slightly beyond known frontiers (e.g. N\textasciicircum{} 1+varepsilon -type circuit lower bounds for suitable promise gaps and input length N=2\textasciicircum{}n ) imply major consequences such as NP not in P/poly.'. non\_commuting\_route
\item Proof strategy family: constructive-realization
\item Negative route: Assume 'MCSP obstruction family' cannot be realized while preserving 'MCSP complexity profile in nonuniform-circuits' strongly enough to activate the local target claim.
\item Failure diagnosis: The route fails if 'MCSP obstruction family' cannot be realized while preserving 'MCSP complexity profile in nonuniform-circuits'. Local negative probe: non\_commuting\_route.
\item Proof obligations: For every admissible 'MCSP obstruction family', the first route-specific realization step preserves 'MCSP complexity profile in nonuniform-circuits' and prepares the passage from 'A hardness magnification theorem has the form:"A barely-nontrivial lower bound for an explicit problem Q in a weak model implies a major lower bound/separation in a strong model."Typically, Q is a meta-problem (e.g. a Gap-MCSP variant) and"major"means consequences such as NP not in P/poly.' to 'In refined regimes, lower bounds for Gap-MCSP slightly beyond known frontiers (e.g. N\textasciicircum{} 1+varepsilon -type circuit lower bounds for suitable promise gaps and input length N=2\textasciicircum{}n ) imply major consequences such as NP not in P/poly.'.; Commutativity obligation: the route for 'carrier-realization' must commute between source theorem hypotheses, the mediating carrier, and the target theorem hypotheses.
\item Candidate lemmas: Commutative factorization theorem for MCSP obstruction family; Coherence diagram lemma for MCSP obstruction family
\item Independent value: A concrete special-case theorem for 'MCSP obstruction family' would be independently valuable because it converts a broad discovery programme into an inspectable mathematical subresult.
\item Expert gate: 6/6
\end{itemize}

\subsection*{conjpkg.disclayer-p-vs-np-family-categorical-reframing.04}
\begin{itemize}
\item Exact statement: No admissible 'MCSP witness family' can simultaneously preserve 'MCSP complexity profile in nonuniform-circuits' and realize 'In refined regimes, lower bounds for Gap-MCSP slightly beyond known frontiers (e.g. N\textasciicircum{} 1+varepsilon -type circuit lower bounds for suitable promise gaps and input length N=2\textasciicircum{}n ) imply major consequences such as NP not in P/poly.' under the current route commitments.
\item Restricted-regime mode: special-case
\item Bridge object: MCSP witness family (family of instances equipped with a monotone witness measure)
\item Bridge definition: 'MCSP witness family' denotes the structured intermediate object required by this branch. It is meant to organize the source-side content 'A hardness magnification theorem has the form:"A barely-nontrivial lower bound for an explicit problem Q in a weak model implies a major lower bound/separation in a strong model."Typically, Q is a meta-problem (e.g. a Gap-MCSP variant) and"major"means consequences such as NP not in P/poly.'. It must preserve the invariant 'MCSP complexity profile in nonuniform-circuits'. Its operational role is 'Recast the blocked frontier as a diagrammatic composition problem: identify objects, morphisms, mediating objects, and the exact commutativity or factorization obligation.'.
\item First nontrivial lemma: For every admissible 'MCSP witness family', the first route-specific realization step preserves 'MCSP complexity profile in nonuniform-circuits' and prepares the passage from 'A hardness magnification theorem has the form:"A barely-nontrivial lower bound for an explicit problem Q in a weak model implies a major lower bound/separation in a strong model."Typically, Q is a meta-problem (e.g. a Gap-MCSP variant) and"major"means consequences such as NP not in P/poly.' to 'In refined regimes, lower bounds for Gap-MCSP slightly beyond known frontiers (e.g. N\textasciicircum{} 1+varepsilon -type circuit lower bounds for suitable promise gaps and input length N=2\textasciicircum{}n ) imply major consequences such as NP not in P/poly.'.
\item Quantifier signature: For every admissible 'MCSP witness family' on the route from 'st.def.hardness-magnification-phenomenon-programme-leve' to 'st.thm.oliveira-pich-santhanam-magnification-near-state'.
\item LaTeX theorem: \textbackslash\{\}forall W,\textbackslash\{\} \textbackslash\{\}mathrm\{Adm\}(W) \textbackslash\{\}to \textbackslash\{\}neg \textbackslash\{\}exists w : \textbackslash\{\}mathsf\{BW\}(W),\textbackslash\{\} \textbackslash\{\}mathrm\{PresInv\}(W,w) \textbackslash\{\}wedge \textbackslash\{\}mathrm\{RealizesTarget\}(W,w).
\item Negative/barrier version: If 'MCSP witness family' cannot be realized with invariant 'MCSP complexity profile in nonuniform-circuits', the route cannot force 'In refined regimes, lower bounds for Gap-MCSP slightly beyond known frontiers (e.g. N\textasciicircum{} 1+varepsilon -type circuit lower bounds for suitable promise gaps and input length N=2\textasciicircum{}n ) imply major consequences such as NP not in P/poly.'. non\_commuting\_route
\item Proof strategy family: constructive-realization
\item Negative route: Assume 'MCSP witness family' cannot be realized while preserving 'MCSP complexity profile in nonuniform-circuits' strongly enough to activate the local target claim.
\item Failure diagnosis: The route fails if 'MCSP witness family' cannot be realized while preserving 'MCSP complexity profile in nonuniform-circuits'. Local negative probe: non\_commuting\_route.
\item Proof obligations: For every admissible 'MCSP witness family', the first route-specific realization step preserves 'MCSP complexity profile in nonuniform-circuits' and prepares the passage from 'A hardness magnification theorem has the form:"A barely-nontrivial lower bound for an explicit problem Q in a weak model implies a major lower bound/separation in a strong model."Typically, Q is a meta-problem (e.g. a Gap-MCSP variant) and"major"means consequences such as NP not in P/poly.' to 'In refined regimes, lower bounds for Gap-MCSP slightly beyond known frontiers (e.g. N\textasciicircum{} 1+varepsilon -type circuit lower bounds for suitable promise gaps and input length N=2\textasciicircum{}n ) imply major consequences such as NP not in P/poly.'.; Commutativity obligation: the route for 'carrier-realization' must commute between source theorem hypotheses, the mediating carrier, and the target theorem hypotheses.
\item Candidate lemmas: Obstruction theorem for MCSP witness family; Coherence diagram lemma for MCSP witness family
\item Independent value: A concrete special-case theorem for 'MCSP witness family' would be independently valuable because it converts a broad discovery programme into an inspectable mathematical subresult.
\item Expert gate: 6/6
\end{itemize}

\subsection*{conjpkg.disclayer-p-vs-np-family-categorical-reframing.05}
\begin{itemize}
\item Exact statement: There exists an intermediate typed object built from 'MCSP witness family' such that the route from 'A hardness magnification theorem has the form:"A barely-nontrivial lower bound for an explicit problem Q in a weak model implies a major lower bound/separation in a strong model."Typically, Q is a meta-problem (e.g. a Gap-MCSP variant) and"major"means consequences such as NP not in P/poly.' to 'In refined regimes, lower bounds for Gap-MCSP slightly beyond known frontiers (e.g. N\textasciicircum{} 1+varepsilon -type circuit lower bounds for suitable promise gaps and input length N=2\textasciicircum{}n ) imply major consequences such as NP not in P/poly.' factors through theorem-sized maps preserving 'MCSP complexity profile in nonuniform-circuits'.
\item Restricted-regime mode: special-case
\item Bridge object: MCSP witness family (family of instances equipped with a monotone witness measure)
\item Bridge definition: 'MCSP witness family' denotes the structured intermediate object required by this branch. It is meant to organize the source-side content 'A hardness magnification theorem has the form:"A barely-nontrivial lower bound for an explicit problem Q in a weak model implies a major lower bound/separation in a strong model."Typically, Q is a meta-problem (e.g. a Gap-MCSP variant) and"major"means consequences such as NP not in P/poly.'. It must preserve the invariant 'MCSP complexity profile in nonuniform-circuits'. Its operational role is 'Recast the blocked frontier as a diagrammatic composition problem: identify objects, morphisms, mediating objects, and the exact commutativity or factorization obligation.'.
\item First nontrivial lemma: For every admissible 'MCSP witness family', the first route-specific realization step preserves 'MCSP complexity profile in nonuniform-circuits' and prepares the passage from 'A hardness magnification theorem has the form:"A barely-nontrivial lower bound for an explicit problem Q in a weak model implies a major lower bound/separation in a strong model."Typically, Q is a meta-problem (e.g. a Gap-MCSP variant) and"major"means consequences such as NP not in P/poly.' to 'In refined regimes, lower bounds for Gap-MCSP slightly beyond known frontiers (e.g. N\textasciicircum{} 1+varepsilon -type circuit lower bounds for suitable promise gaps and input length N=2\textasciicircum{}n ) imply major consequences such as NP not in P/poly.'.
\item Quantifier signature: For every admissible 'MCSP witness family' on the route from 'st.def.hardness-magnification-phenomenon-programme-leve' to 'st.thm.oliveira-pich-santhanam-magnification-near-state'.
\item LaTeX theorem: \textbackslash\{\}forall W,\textbackslash\{\} \textbackslash\{\}mathrm\{Adm\}(W) \textbackslash\{\}to \textbackslash\{\}exists w : \textbackslash\{\}mathsf\{BW\}(W),\textbackslash\{\} \textbackslash\{\}mathrm\{PresInv\}(W,w) \textbackslash\{\}wedge \textbackslash\{\}mathrm\{RealizesTarget\}(W,w).
\item Negative/barrier version: If 'MCSP witness family' cannot be realized with invariant 'MCSP complexity profile in nonuniform-circuits', the route cannot force 'In refined regimes, lower bounds for Gap-MCSP slightly beyond known frontiers (e.g. N\textasciicircum{} 1+varepsilon -type circuit lower bounds for suitable promise gaps and input length N=2\textasciicircum{}n ) imply major consequences such as NP not in P/poly.'. non\_commuting\_route
\item Proof strategy family: constructive-realization
\item Negative route: Assume 'MCSP witness family' cannot be realized while preserving 'MCSP complexity profile in nonuniform-circuits' strongly enough to activate the local target claim.
\item Failure diagnosis: The route fails if 'MCSP witness family' cannot be realized while preserving 'MCSP complexity profile in nonuniform-circuits'. Local negative probe: non\_commuting\_route.
\item Proof obligations: For every admissible 'MCSP witness family', the first route-specific realization step preserves 'MCSP complexity profile in nonuniform-circuits' and prepares the passage from 'A hardness magnification theorem has the form:"A barely-nontrivial lower bound for an explicit problem Q in a weak model implies a major lower bound/separation in a strong model."Typically, Q is a meta-problem (e.g. a Gap-MCSP variant) and"major"means consequences such as NP not in P/poly.' to 'In refined regimes, lower bounds for Gap-MCSP slightly beyond known frontiers (e.g. N\textasciicircum{} 1+varepsilon -type circuit lower bounds for suitable promise gaps and input length N=2\textasciicircum{}n ) imply major consequences such as NP not in P/poly.'.; Commutativity obligation: the route for 'carrier-realization' must commute between source theorem hypotheses, the mediating carrier, and the target theorem hypotheses.
\item Candidate lemmas: Factorization theorem for MCSP witness family; Coherence diagram lemma for MCSP witness family
\item Independent value: A concrete special-case theorem for 'MCSP witness family' would be independently valuable because it converts a broad discovery programme into an inspectable mathematical subresult.
\item Expert gate: 6/6
\end{itemize}

\subsection*{conjpkg.disclayer-p-vs-np-family-categorical-reframing.06}
\begin{itemize}
\item Exact statement: There exists an admissible 'MCSP witness family' whose realization internalizes 'A hardness magnification theorem has the form:"A barely-nontrivial lower bound for an explicit problem Q in a weak model implies a major lower bound/separation in a strong model."Typically, Q is a meta-problem (e.g. a Gap-MCSP variant) and"major"means consequences such as NP not in P/poly.' while preserving 'MCSP complexity profile in nonuniform-circuits' strongly enough to support 'In refined regimes, lower bounds for Gap-MCSP slightly beyond known frontiers (e.g. N\textasciicircum{} 1+varepsilon -type circuit lower bounds for suitable promise gaps and input length N=2\textasciicircum{}n ) imply major consequences such as NP not in P/poly.'.
\item Restricted-regime mode: special-case
\item Bridge object: MCSP witness family (family of instances equipped with a monotone witness measure)
\item Bridge definition: 'MCSP witness family' denotes the structured intermediate object required by this branch. It is meant to organize the source-side content 'A hardness magnification theorem has the form:"A barely-nontrivial lower bound for an explicit problem Q in a weak model implies a major lower bound/separation in a strong model."Typically, Q is a meta-problem (e.g. a Gap-MCSP variant) and"major"means consequences such as NP not in P/poly.'. It must preserve the invariant 'MCSP complexity profile in nonuniform-circuits'. Its operational role is 'Recast the blocked frontier as a diagrammatic composition problem: identify objects, morphisms, mediating objects, and the exact commutativity or factorization obligation.'.
\item First nontrivial lemma: For every admissible 'MCSP witness family', the first route-specific realization step preserves 'MCSP complexity profile in nonuniform-circuits' and prepares the passage from 'A hardness magnification theorem has the form:"A barely-nontrivial lower bound for an explicit problem Q in a weak model implies a major lower bound/separation in a strong model."Typically, Q is a meta-problem (e.g. a Gap-MCSP variant) and"major"means consequences such as NP not in P/poly.' to 'In refined regimes, lower bounds for Gap-MCSP slightly beyond known frontiers (e.g. N\textasciicircum{} 1+varepsilon -type circuit lower bounds for suitable promise gaps and input length N=2\textasciicircum{}n ) imply major consequences such as NP not in P/poly.'.
\item Quantifier signature: For every admissible 'MCSP witness family' on the route from 'st.def.hardness-magnification-phenomenon-programme-leve' to 'st.thm.oliveira-pich-santhanam-magnification-near-state'.
\item LaTeX theorem: \textbackslash\{\}forall W,\textbackslash\{\} \textbackslash\{\}mathrm\{Adm\}(W) \textbackslash\{\}to \textbackslash\{\}exists w : \textbackslash\{\}mathsf\{BW\}(W),\textbackslash\{\} \textbackslash\{\}mathrm\{PresInv\}(W,w) \textbackslash\{\}wedge \textbackslash\{\}mathrm\{RealizesTarget\}(W,w).
\item Negative/barrier version: If 'MCSP witness family' cannot be realized with invariant 'MCSP complexity profile in nonuniform-circuits', the route cannot force 'In refined regimes, lower bounds for Gap-MCSP slightly beyond known frontiers (e.g. N\textasciicircum{} 1+varepsilon -type circuit lower bounds for suitable promise gaps and input length N=2\textasciicircum{}n ) imply major consequences such as NP not in P/poly.'. non\_commuting\_route
\item Proof strategy family: constructive-realization
\item Negative route: Assume 'MCSP witness family' cannot be realized while preserving 'MCSP complexity profile in nonuniform-circuits' strongly enough to activate the local target claim.
\item Failure diagnosis: The route fails if 'MCSP witness family' cannot be realized while preserving 'MCSP complexity profile in nonuniform-circuits'. Local negative probe: non\_commuting\_route.
\item Proof obligations: For every admissible 'MCSP witness family', the first route-specific realization step preserves 'MCSP complexity profile in nonuniform-circuits' and prepares the passage from 'A hardness magnification theorem has the form:"A barely-nontrivial lower bound for an explicit problem Q in a weak model implies a major lower bound/separation in a strong model."Typically, Q is a meta-problem (e.g. a Gap-MCSP variant) and"major"means consequences such as NP not in P/poly.' to 'In refined regimes, lower bounds for Gap-MCSP slightly beyond known frontiers (e.g. N\textasciicircum{} 1+varepsilon -type circuit lower bounds for suitable promise gaps and input length N=2\textasciicircum{}n ) imply major consequences such as NP not in P/poly.'.; Commutativity obligation: the route for 'carrier-realization' must commute between source theorem hypotheses, the mediating carrier, and the target theorem hypotheses.
\item Candidate lemmas: Realization theorem for MCSP witness family; Coherence diagram lemma for MCSP witness family
\item Independent value: A concrete special-case theorem for 'MCSP witness family' would be independently valuable because it converts a broad discovery programme into an inspectable mathematical subresult.
\item Expert gate: 6/6
\end{itemize}

\subsection*{conjpkg.disclayer-p-vs-np-family-categorical-reframing.07}
\begin{itemize}
\item Exact statement: For every admissible 'MCSP witness family', there is a transport law from 'nonuniform-circuits' to 'formula-size' that preserves 'MCSP complexity profile in nonuniform-circuits' and retains enough strength to activate 'There exist explicit"meta"problems (including MCSP variants) such that proving certain emph barely-superlinear lower bounds against restricted models (e.g. formulas) would imply major lower bounds for NP in much stronger models (including superpolynomial formula/circuit lower bounds, and in some parameter regimes consequences as strong as NP not in P/poly ).'.
\item Restricted-regime mode: special-case
\item Bridge object: MCSP witness family (family of instances equipped with a monotone witness measure)
\item Bridge definition: 'MCSP witness family' denotes the structured intermediate object required by this branch. It is meant to organize the source-side content 'A hardness magnification theorem has the form:"A barely-nontrivial lower bound for an explicit problem Q in a weak model implies a major lower bound/separation in a strong model."Typically, Q is a meta-problem (e.g. a Gap-MCSP variant) and"major"means consequences such as NP not in P/poly.'. It must preserve the invariant 'MCSP complexity profile in nonuniform-circuits'. Its operational role is 'Recast the blocked frontier as a diagrammatic composition problem: identify objects, morphisms, mediating objects, and the exact commutativity or factorization obligation.'.
\item First nontrivial lemma: For every admissible 'MCSP witness family', the first route-specific realization step preserves 'MCSP complexity profile in nonuniform-circuits' and prepares the passage from 'A hardness magnification theorem has the form:"A barely-nontrivial lower bound for an explicit problem Q in a weak model implies a major lower bound/separation in a strong model."Typically, Q is a meta-problem (e.g. a Gap-MCSP variant) and"major"means consequences such as NP not in P/poly.' to 'There exist explicit"meta"problems (including MCSP variants) such that proving certain emph barely-superlinear lower bounds against restricted models (e.g. formulas) would imply major lower bounds for NP in much stronger models (including superpolynomial formula/circuit lower bounds, and in some parameter regimes consequences as strong as NP not in P/poly ).'.
\item Quantifier signature: For every admissible 'MCSP witness family' on the route from 'st.def.hardness-magnification-phenomenon-programme-leve' to 'st.thm.oliveira-santhanam-hardness-magnification-for-na'.
\item LaTeX theorem: \textbackslash\{\}forall W,\textbackslash\{\} \textbackslash\{\}mathrm\{Adm\}(W) \textbackslash\{\}to \textbackslash\{\}exists w : \textbackslash\{\}mathsf\{BW\}(W),\textbackslash\{\} \textbackslash\{\}mathrm\{PresInv\}(W,w) \textbackslash\{\}wedge \textbackslash\{\}mathrm\{RealizesTarget\}(W,w).
\item Negative/barrier version: If 'MCSP witness family' cannot be realized with invariant 'MCSP complexity profile in nonuniform-circuits', the route cannot force 'There exist explicit"meta"problems (including MCSP variants) such that proving certain emph barely-superlinear lower bounds against restricted models (e.g. formulas) would imply major lower bounds for NP in much stronger models (including superpolynomial formula/circuit lower bounds, and in some parameter regimes consequences as strong as NP not in P/poly ).'. non\_commuting\_route
\item Proof strategy family: constructive-realization
\item Negative route: Assume 'MCSP witness family' cannot be realized while preserving 'MCSP complexity profile in nonuniform-circuits' strongly enough to activate the local target claim.
\item Failure diagnosis: The route fails if 'MCSP witness family' cannot be realized while preserving 'MCSP complexity profile in nonuniform-circuits'. Local negative probe: non\_commuting\_route.
\item Proof obligations: For every admissible 'MCSP witness family', the first route-specific realization step preserves 'MCSP complexity profile in nonuniform-circuits' and prepares the passage from 'A hardness magnification theorem has the form:"A barely-nontrivial lower bound for an explicit problem Q in a weak model implies a major lower bound/separation in a strong model."Typically, Q is a meta-problem (e.g. a Gap-MCSP variant) and"major"means consequences such as NP not in P/poly.' to 'There exist explicit"meta"problems (including MCSP variants) such that proving certain emph barely-superlinear lower bounds against restricted models (e.g. formulas) would imply major lower bounds for NP in much stronger models (including superpolynomial formula/circuit lower bounds, and in some parameter regimes consequences as strong as NP not in P/poly ).'.; Commutativity obligation: the route for 'carrier-realization' must commute between source theorem hypotheses, the mediating carrier, and the target theorem hypotheses.
\item Candidate lemmas: Transport theorem for MCSP witness family; Coherence diagram lemma for MCSP witness family
\item Independent value: A concrete special-case theorem for 'MCSP witness family' would be independently valuable because it converts a broad discovery programme into an inspectable mathematical subresult.
\item Expert gate: 6/6
\end{itemize}

\subsection*{conjpkg.disclayer-p-vs-np-family-categorical-reframing.08}
\begin{itemize}
\item Exact statement: No admissible 'MCSP witness family' can simultaneously preserve 'MCSP complexity profile in nonuniform-circuits' and realize 'There exist explicit"meta"problems (including MCSP variants) such that proving certain emph barely-superlinear lower bounds against restricted models (e.g. formulas) would imply major lower bounds for NP in much stronger models (including superpolynomial formula/circuit lower bounds, and in some parameter regimes consequences as strong as NP not in P/poly ).' under the current route commitments.
\item Restricted-regime mode: special-case
\item Bridge object: MCSP witness family (family of instances equipped with a monotone witness measure)
\item Bridge definition: 'MCSP witness family' denotes the structured intermediate object required by this branch. It is meant to organize the source-side content 'A hardness magnification theorem has the form:"A barely-nontrivial lower bound for an explicit problem Q in a weak model implies a major lower bound/separation in a strong model."Typically, Q is a meta-problem (e.g. a Gap-MCSP variant) and"major"means consequences such as NP not in P/poly.'. It must preserve the invariant 'MCSP complexity profile in nonuniform-circuits'. Its operational role is 'Recast the blocked frontier as a diagrammatic composition problem: identify objects, morphisms, mediating objects, and the exact commutativity or factorization obligation.'.
\item First nontrivial lemma: For every admissible 'MCSP witness family', the first route-specific realization step preserves 'MCSP complexity profile in nonuniform-circuits' and prepares the passage from 'A hardness magnification theorem has the form:"A barely-nontrivial lower bound for an explicit problem Q in a weak model implies a major lower bound/separation in a strong model."Typically, Q is a meta-problem (e.g. a Gap-MCSP variant) and"major"means consequences such as NP not in P/poly.' to 'There exist explicit"meta"problems (including MCSP variants) such that proving certain emph barely-superlinear lower bounds against restricted models (e.g. formulas) would imply major lower bounds for NP in much stronger models (including superpolynomial formula/circuit lower bounds, and in some parameter regimes consequences as strong as NP not in P/poly ).'.
\item Quantifier signature: For every admissible 'MCSP witness family' on the route from 'st.def.hardness-magnification-phenomenon-programme-leve' to 'st.thm.oliveira-santhanam-hardness-magnification-for-na'.
\item LaTeX theorem: \textbackslash\{\}forall W,\textbackslash\{\} \textbackslash\{\}mathrm\{Adm\}(W) \textbackslash\{\}to \textbackslash\{\}neg \textbackslash\{\}exists w : \textbackslash\{\}mathsf\{BW\}(W),\textbackslash\{\} \textbackslash\{\}mathrm\{PresInv\}(W,w) \textbackslash\{\}wedge \textbackslash\{\}mathrm\{RealizesTarget\}(W,w).
\item Negative/barrier version: If 'MCSP witness family' cannot be realized with invariant 'MCSP complexity profile in nonuniform-circuits', the route cannot force 'There exist explicit"meta"problems (including MCSP variants) such that proving certain emph barely-superlinear lower bounds against restricted models (e.g. formulas) would imply major lower bounds for NP in much stronger models (including superpolynomial formula/circuit lower bounds, and in some parameter regimes consequences as strong as NP not in P/poly ).'. non\_commuting\_route
\item Proof strategy family: constructive-realization
\item Negative route: Assume 'MCSP witness family' cannot be realized while preserving 'MCSP complexity profile in nonuniform-circuits' strongly enough to activate the local target claim.
\item Failure diagnosis: The route fails if 'MCSP witness family' cannot be realized while preserving 'MCSP complexity profile in nonuniform-circuits'. Local negative probe: non\_commuting\_route.
\item Proof obligations: For every admissible 'MCSP witness family', the first route-specific realization step preserves 'MCSP complexity profile in nonuniform-circuits' and prepares the passage from 'A hardness magnification theorem has the form:"A barely-nontrivial lower bound for an explicit problem Q in a weak model implies a major lower bound/separation in a strong model."Typically, Q is a meta-problem (e.g. a Gap-MCSP variant) and"major"means consequences such as NP not in P/poly.' to 'There exist explicit"meta"problems (including MCSP variants) such that proving certain emph barely-superlinear lower bounds against restricted models (e.g. formulas) would imply major lower bounds for NP in much stronger models (including superpolynomial formula/circuit lower bounds, and in some parameter regimes consequences as strong as NP not in P/poly ).'.; Commutativity obligation: the route for 'carrier-realization' must commute between source theorem hypotheses, the mediating carrier, and the target theorem hypotheses.
\item Candidate lemmas: Obstruction theorem for MCSP witness family; Coherence diagram lemma for MCSP witness family
\item Independent value: A concrete special-case theorem for 'MCSP witness family' would be independently valuable because it converts a broad discovery programme into an inspectable mathematical subresult.
\item Expert gate: 6/6
\end{itemize}

\subsection*{conjpkg.disclayer-p-vs-np-family-categorical-reframing.09}
\begin{itemize}
\item Exact statement: There exists an intermediate typed object built from 'MCSP witness family' such that the route from 'A hardness magnification theorem has the form:"A barely-nontrivial lower bound for an explicit problem Q in a weak model implies a major lower bound/separation in a strong model."Typically, Q is a meta-problem (e.g. a Gap-MCSP variant) and"major"means consequences such as NP not in P/poly.' to 'There exist explicit"meta"problems (including MCSP variants) such that proving certain emph barely-superlinear lower bounds against restricted models (e.g. formulas) would imply major lower bounds for NP in much stronger models (including superpolynomial formula/circuit lower bounds, and in some parameter regimes consequences as strong as NP not in P/poly ).' factors through theorem-sized maps preserving 'MCSP complexity profile in nonuniform-circuits'.
\item Restricted-regime mode: special-case
\item Bridge object: MCSP witness family (family of instances equipped with a monotone witness measure)
\item Bridge definition: 'MCSP witness family' denotes the structured intermediate object required by this branch. It is meant to organize the source-side content 'A hardness magnification theorem has the form:"A barely-nontrivial lower bound for an explicit problem Q in a weak model implies a major lower bound/separation in a strong model."Typically, Q is a meta-problem (e.g. a Gap-MCSP variant) and"major"means consequences such as NP not in P/poly.'. It must preserve the invariant 'MCSP complexity profile in nonuniform-circuits'. Its operational role is 'Recast the blocked frontier as a diagrammatic composition problem: identify objects, morphisms, mediating objects, and the exact commutativity or factorization obligation.'.
\item First nontrivial lemma: For every admissible 'MCSP witness family', the first route-specific realization step preserves 'MCSP complexity profile in nonuniform-circuits' and prepares the passage from 'A hardness magnification theorem has the form:"A barely-nontrivial lower bound for an explicit problem Q in a weak model implies a major lower bound/separation in a strong model."Typically, Q is a meta-problem (e.g. a Gap-MCSP variant) and"major"means consequences such as NP not in P/poly.' to 'There exist explicit"meta"problems (including MCSP variants) such that proving certain emph barely-superlinear lower bounds against restricted models (e.g. formulas) would imply major lower bounds for NP in much stronger models (including superpolynomial formula/circuit lower bounds, and in some parameter regimes consequences as strong as NP not in P/poly ).'.
\item Quantifier signature: For every admissible 'MCSP witness family' on the route from 'st.def.hardness-magnification-phenomenon-programme-leve' to 'st.thm.oliveira-santhanam-hardness-magnification-for-na'.
\item LaTeX theorem: \textbackslash\{\}forall W,\textbackslash\{\} \textbackslash\{\}mathrm\{Adm\}(W) \textbackslash\{\}to \textbackslash\{\}exists w : \textbackslash\{\}mathsf\{BW\}(W),\textbackslash\{\} \textbackslash\{\}mathrm\{PresInv\}(W,w) \textbackslash\{\}wedge \textbackslash\{\}mathrm\{RealizesTarget\}(W,w).
\item Negative/barrier version: If 'MCSP witness family' cannot be realized with invariant 'MCSP complexity profile in nonuniform-circuits', the route cannot force 'There exist explicit"meta"problems (including MCSP variants) such that proving certain emph barely-superlinear lower bounds against restricted models (e.g. formulas) would imply major lower bounds for NP in much stronger models (including superpolynomial formula/circuit lower bounds, and in some parameter regimes consequences as strong as NP not in P/poly ).'. non\_commuting\_route
\item Proof strategy family: constructive-realization
\item Negative route: Assume 'MCSP witness family' cannot be realized while preserving 'MCSP complexity profile in nonuniform-circuits' strongly enough to activate the local target claim.
\item Failure diagnosis: The route fails if 'MCSP witness family' cannot be realized while preserving 'MCSP complexity profile in nonuniform-circuits'. Local negative probe: non\_commuting\_route.
\item Proof obligations: For every admissible 'MCSP witness family', the first route-specific realization step preserves 'MCSP complexity profile in nonuniform-circuits' and prepares the passage from 'A hardness magnification theorem has the form:"A barely-nontrivial lower bound for an explicit problem Q in a weak model implies a major lower bound/separation in a strong model."Typically, Q is a meta-problem (e.g. a Gap-MCSP variant) and"major"means consequences such as NP not in P/poly.' to 'There exist explicit"meta"problems (including MCSP variants) such that proving certain emph barely-superlinear lower bounds against restricted models (e.g. formulas) would imply major lower bounds for NP in much stronger models (including superpolynomial formula/circuit lower bounds, and in some parameter regimes consequences as strong as NP not in P/poly ).'.; Commutativity obligation: the route for 'carrier-realization' must commute between source theorem hypotheses, the mediating carrier, and the target theorem hypotheses.
\item Candidate lemmas: Factorization theorem for MCSP witness family; Coherence diagram lemma for MCSP witness family
\item Independent value: A concrete special-case theorem for 'MCSP witness family' would be independently valuable because it converts a broad discovery programme into an inspectable mathematical subresult.
\item Expert gate: 6/6
\end{itemize}

\subsection*{conjpkg.disclayer-p-vs-np-family-categorical-reframing.10}
\begin{itemize}
\item Exact statement: There exists an admissible 'MCSP witness family' whose realization internalizes 'A hardness magnification theorem has the form:"A barely-nontrivial lower bound for an explicit problem Q in a weak model implies a major lower bound/separation in a strong model."Typically, Q is a meta-problem (e.g. a Gap-MCSP variant) and"major"means consequences such as NP not in P/poly.' while preserving 'MCSP complexity profile in nonuniform-circuits' strongly enough to support 'There exist explicit"meta"problems (including MCSP variants) such that proving certain emph barely-superlinear lower bounds against restricted models (e.g. formulas) would imply major lower bounds for NP in much stronger models (including superpolynomial formula/circuit lower bounds, and in some parameter regimes consequences as strong as NP not in P/poly ).'.
\item Restricted-regime mode: special-case
\item Bridge object: MCSP witness family (family of instances equipped with a monotone witness measure)
\item Bridge definition: 'MCSP witness family' denotes the structured intermediate object required by this branch. It is meant to organize the source-side content 'A hardness magnification theorem has the form:"A barely-nontrivial lower bound for an explicit problem Q in a weak model implies a major lower bound/separation in a strong model."Typically, Q is a meta-problem (e.g. a Gap-MCSP variant) and"major"means consequences such as NP not in P/poly.'. It must preserve the invariant 'MCSP complexity profile in nonuniform-circuits'. Its operational role is 'Recast the blocked frontier as a diagrammatic composition problem: identify objects, morphisms, mediating objects, and the exact commutativity or factorization obligation.'.
\item First nontrivial lemma: For every admissible 'MCSP witness family', the first route-specific realization step preserves 'MCSP complexity profile in nonuniform-circuits' and prepares the passage from 'A hardness magnification theorem has the form:"A barely-nontrivial lower bound for an explicit problem Q in a weak model implies a major lower bound/separation in a strong model."Typically, Q is a meta-problem (e.g. a Gap-MCSP variant) and"major"means consequences such as NP not in P/poly.' to 'There exist explicit"meta"problems (including MCSP variants) such that proving certain emph barely-superlinear lower bounds against restricted models (e.g. formulas) would imply major lower bounds for NP in much stronger models (including superpolynomial formula/circuit lower bounds, and in some parameter regimes consequences as strong as NP not in P/poly ).'.
\item Quantifier signature: For every admissible 'MCSP witness family' on the route from 'st.def.hardness-magnification-phenomenon-programme-leve' to 'st.thm.oliveira-santhanam-hardness-magnification-for-na'.
\item LaTeX theorem: \textbackslash\{\}forall W,\textbackslash\{\} \textbackslash\{\}mathrm\{Adm\}(W) \textbackslash\{\}to \textbackslash\{\}exists w : \textbackslash\{\}mathsf\{BW\}(W),\textbackslash\{\} \textbackslash\{\}mathrm\{PresInv\}(W,w) \textbackslash\{\}wedge \textbackslash\{\}mathrm\{RealizesTarget\}(W,w).
\item Negative/barrier version: If 'MCSP witness family' cannot be realized with invariant 'MCSP complexity profile in nonuniform-circuits', the route cannot force 'There exist explicit"meta"problems (including MCSP variants) such that proving certain emph barely-superlinear lower bounds against restricted models (e.g. formulas) would imply major lower bounds for NP in much stronger models (including superpolynomial formula/circuit lower bounds, and in some parameter regimes consequences as strong as NP not in P/poly ).'. non\_commuting\_route
\item Proof strategy family: constructive-realization
\item Negative route: Assume 'MCSP witness family' cannot be realized while preserving 'MCSP complexity profile in nonuniform-circuits' strongly enough to activate the local target claim.
\item Failure diagnosis: The route fails if 'MCSP witness family' cannot be realized while preserving 'MCSP complexity profile in nonuniform-circuits'. Local negative probe: non\_commuting\_route.
\item Proof obligations: For every admissible 'MCSP witness family', the first route-specific realization step preserves 'MCSP complexity profile in nonuniform-circuits' and prepares the passage from 'A hardness magnification theorem has the form:"A barely-nontrivial lower bound for an explicit problem Q in a weak model implies a major lower bound/separation in a strong model."Typically, Q is a meta-problem (e.g. a Gap-MCSP variant) and"major"means consequences such as NP not in P/poly.' to 'There exist explicit"meta"problems (including MCSP variants) such that proving certain emph barely-superlinear lower bounds against restricted models (e.g. formulas) would imply major lower bounds for NP in much stronger models (including superpolynomial formula/circuit lower bounds, and in some parameter regimes consequences as strong as NP not in P/poly ).'.; Commutativity obligation: the route for 'carrier-realization' must commute between source theorem hypotheses, the mediating carrier, and the target theorem hypotheses.
\item Candidate lemmas: Realization theorem for MCSP witness family; Coherence diagram lemma for MCSP witness family
\item Independent value: A concrete special-case theorem for 'MCSP witness family' would be independently valuable because it converts a broad discovery programme into an inspectable mathematical subresult.
\item Expert gate: 6/6
\end{itemize}

\subsection*{conjpkg.disclayer-p-vs-np-family-categorical-reframing.11}
\begin{itemize}
\item Exact statement: No admissible 'MCSP obstruction family' can simultaneously preserve 'MCSP complexity profile in nonuniform-circuits' and realize 'In refined regimes, lower bounds for Gap-MCSP slightly beyond known frontiers (e.g. N\textasciicircum{} 1+varepsilon -type circuit lower bounds for suitable promise gaps and input length N=2\textasciicircum{}n ) imply major consequences such as NP not in P/poly.' under the current route commitments.
\item Restricted-regime mode: special-case
\item Bridge object: MCSP obstruction family (family of hard instances or obstructions forcing the target payoff)
\item Bridge definition: 'MCSP obstruction family' denotes the structured intermediate object required by this branch. It is meant to organize the source-side content 'A hardness magnification theorem has the form:"A barely-nontrivial lower bound for an explicit problem Q in a weak model implies a major lower bound/separation in a strong model."Typically, Q is a meta-problem (e.g. a Gap-MCSP variant) and"major"means consequences such as NP not in P/poly.'. It must preserve the invariant 'MCSP complexity profile in nonuniform-circuits'. Its operational role is 'Recast the blocked frontier as a diagrammatic composition problem: identify objects, morphisms, mediating objects, and the exact commutativity or factorization obligation.'.
\item First nontrivial lemma: For every admissible 'MCSP obstruction family', the first route-specific realization step preserves 'MCSP complexity profile in nonuniform-circuits' and prepares the passage from 'A hardness magnification theorem has the form:"A barely-nontrivial lower bound for an explicit problem Q in a weak model implies a major lower bound/separation in a strong model."Typically, Q is a meta-problem (e.g. a Gap-MCSP variant) and"major"means consequences such as NP not in P/poly.' to 'In refined regimes, lower bounds for Gap-MCSP slightly beyond known frontiers (e.g. N\textasciicircum{} 1+varepsilon -type circuit lower bounds for suitable promise gaps and input length N=2\textasciicircum{}n ) imply major consequences such as NP not in P/poly.'.
\item Quantifier signature: For every admissible 'MCSP obstruction family' on the route from 'st.def.hardness-magnification-phenomenon-programme-leve' to 'st.thm.oliveira-pich-santhanam-magnification-near-state'.
\item LaTeX theorem: \textbackslash\{\}forall W,\textbackslash\{\} \textbackslash\{\}mathrm\{Adm\}(W) \textbackslash\{\}to \textbackslash\{\}neg \textbackslash\{\}exists w : \textbackslash\{\}mathsf\{BW\}(W),\textbackslash\{\} \textbackslash\{\}mathrm\{PresInv\}(W,w) \textbackslash\{\}wedge \textbackslash\{\}mathrm\{RealizesTarget\}(W,w).
\item Negative/barrier version: If 'MCSP obstruction family' cannot be realized with invariant 'MCSP complexity profile in nonuniform-circuits', the route cannot force 'In refined regimes, lower bounds for Gap-MCSP slightly beyond known frontiers (e.g. N\textasciicircum{} 1+varepsilon -type circuit lower bounds for suitable promise gaps and input length N=2\textasciicircum{}n ) imply major consequences such as NP not in P/poly.'. non\_commuting\_route
\item Proof strategy family: constructive-realization
\item Negative route: Assume 'MCSP obstruction family' cannot be realized while preserving 'MCSP complexity profile in nonuniform-circuits' strongly enough to activate the local target claim.
\item Failure diagnosis: The route fails if 'MCSP obstruction family' cannot be realized while preserving 'MCSP complexity profile in nonuniform-circuits'. Local negative probe: non\_commuting\_route.
\item Proof obligations: For every admissible 'MCSP obstruction family', the first route-specific realization step preserves 'MCSP complexity profile in nonuniform-circuits' and prepares the passage from 'A hardness magnification theorem has the form:"A barely-nontrivial lower bound for an explicit problem Q in a weak model implies a major lower bound/separation in a strong model."Typically, Q is a meta-problem (e.g. a Gap-MCSP variant) and"major"means consequences such as NP not in P/poly.' to 'In refined regimes, lower bounds for Gap-MCSP slightly beyond known frontiers (e.g. N\textasciicircum{} 1+varepsilon -type circuit lower bounds for suitable promise gaps and input length N=2\textasciicircum{}n ) imply major consequences such as NP not in P/poly.'.; Commutativity obligation: the route for 'carrier-realization' must commute between source theorem hypotheses, the mediating carrier, and the target theorem hypotheses.
\item Candidate lemmas: Obstruction theorem for MCSP obstruction family; Coherence diagram lemma for MCSP obstruction family
\item Independent value: A concrete special-case theorem for 'MCSP obstruction family' would be independently valuable because it converts a broad discovery programme into an inspectable mathematical subresult.
\item Expert gate: 6/6
\end{itemize}

\subsection*{conjpkg.disclayer-p-vs-np-family-categorical-reframing.12}
\begin{itemize}
\item Exact statement: There exists an intermediate typed object built from 'MCSP obstruction family' such that the route from 'A hardness magnification theorem has the form:"A barely-nontrivial lower bound for an explicit problem Q in a weak model implies a major lower bound/separation in a strong model."Typically, Q is a meta-problem (e.g. a Gap-MCSP variant) and"major"means consequences such as NP not in P/poly.' to 'In refined regimes, lower bounds for Gap-MCSP slightly beyond known frontiers (e.g. N\textasciicircum{} 1+varepsilon -type circuit lower bounds for suitable promise gaps and input length N=2\textasciicircum{}n ) imply major consequences such as NP not in P/poly.' factors through theorem-sized maps preserving 'MCSP complexity profile in nonuniform-circuits'.
\item Restricted-regime mode: special-case
\item Bridge object: MCSP obstruction family (family of hard instances or obstructions forcing the target payoff)
\item Bridge definition: 'MCSP obstruction family' denotes the structured intermediate object required by this branch. It is meant to organize the source-side content 'A hardness magnification theorem has the form:"A barely-nontrivial lower bound for an explicit problem Q in a weak model implies a major lower bound/separation in a strong model."Typically, Q is a meta-problem (e.g. a Gap-MCSP variant) and"major"means consequences such as NP not in P/poly.'. It must preserve the invariant 'MCSP complexity profile in nonuniform-circuits'. Its operational role is 'Recast the blocked frontier as a diagrammatic composition problem: identify objects, morphisms, mediating objects, and the exact commutativity or factorization obligation.'.
\item First nontrivial lemma: For every admissible 'MCSP obstruction family', the first route-specific realization step preserves 'MCSP complexity profile in nonuniform-circuits' and prepares the passage from 'A hardness magnification theorem has the form:"A barely-nontrivial lower bound for an explicit problem Q in a weak model implies a major lower bound/separation in a strong model."Typically, Q is a meta-problem (e.g. a Gap-MCSP variant) and"major"means consequences such as NP not in P/poly.' to 'In refined regimes, lower bounds for Gap-MCSP slightly beyond known frontiers (e.g. N\textasciicircum{} 1+varepsilon -type circuit lower bounds for suitable promise gaps and input length N=2\textasciicircum{}n ) imply major consequences such as NP not in P/poly.'.
\item Quantifier signature: For every admissible 'MCSP obstruction family' on the route from 'st.def.hardness-magnification-phenomenon-programme-leve' to 'st.thm.oliveira-pich-santhanam-magnification-near-state'.
\item LaTeX theorem: \textbackslash\{\}forall W,\textbackslash\{\} \textbackslash\{\}mathrm\{Adm\}(W) \textbackslash\{\}to \textbackslash\{\}exists w : \textbackslash\{\}mathsf\{BW\}(W),\textbackslash\{\} \textbackslash\{\}mathrm\{PresInv\}(W,w) \textbackslash\{\}wedge \textbackslash\{\}mathrm\{RealizesTarget\}(W,w).
\item Negative/barrier version: If 'MCSP obstruction family' cannot be realized with invariant 'MCSP complexity profile in nonuniform-circuits', the route cannot force 'In refined regimes, lower bounds for Gap-MCSP slightly beyond known frontiers (e.g. N\textasciicircum{} 1+varepsilon -type circuit lower bounds for suitable promise gaps and input length N=2\textasciicircum{}n ) imply major consequences such as NP not in P/poly.'. non\_commuting\_route
\item Proof strategy family: constructive-realization
\item Negative route: Assume 'MCSP obstruction family' cannot be realized while preserving 'MCSP complexity profile in nonuniform-circuits' strongly enough to activate the local target claim.
\item Failure diagnosis: The route fails if 'MCSP obstruction family' cannot be realized while preserving 'MCSP complexity profile in nonuniform-circuits'. Local negative probe: non\_commuting\_route.
\item Proof obligations: For every admissible 'MCSP obstruction family', the first route-specific realization step preserves 'MCSP complexity profile in nonuniform-circuits' and prepares the passage from 'A hardness magnification theorem has the form:"A barely-nontrivial lower bound for an explicit problem Q in a weak model implies a major lower bound/separation in a strong model."Typically, Q is a meta-problem (e.g. a Gap-MCSP variant) and"major"means consequences such as NP not in P/poly.' to 'In refined regimes, lower bounds for Gap-MCSP slightly beyond known frontiers (e.g. N\textasciicircum{} 1+varepsilon -type circuit lower bounds for suitable promise gaps and input length N=2\textasciicircum{}n ) imply major consequences such as NP not in P/poly.'.; Commutativity obligation: the route for 'carrier-realization' must commute between source theorem hypotheses, the mediating carrier, and the target theorem hypotheses.
\item Candidate lemmas: Factorization theorem for MCSP obstruction family; Coherence diagram lemma for MCSP obstruction family
\item Independent value: A concrete special-case theorem for 'MCSP obstruction family' would be independently valuable because it converts a broad discovery programme into an inspectable mathematical subresult.
\item Expert gate: 6/6
\end{itemize}

\subsection*{conjpkg.disclayer-p-vs-np-family-categorical-reframing.13}
\begin{itemize}
\item Exact statement: There exists an admissible 'MCSP obstruction family' whose realization internalizes 'A hardness magnification theorem has the form:"A barely-nontrivial lower bound for an explicit problem Q in a weak model implies a major lower bound/separation in a strong model."Typically, Q is a meta-problem (e.g. a Gap-MCSP variant) and"major"means consequences such as NP not in P/poly.' while preserving 'MCSP complexity profile in nonuniform-circuits' strongly enough to support 'In refined regimes, lower bounds for Gap-MCSP slightly beyond known frontiers (e.g. N\textasciicircum{} 1+varepsilon -type circuit lower bounds for suitable promise gaps and input length N=2\textasciicircum{}n ) imply major consequences such as NP not in P/poly.'.
\item Restricted-regime mode: special-case
\item Bridge object: MCSP obstruction family (family of hard instances or obstructions forcing the target payoff)
\item Bridge definition: 'MCSP obstruction family' denotes the structured intermediate object required by this branch. It is meant to organize the source-side content 'A hardness magnification theorem has the form:"A barely-nontrivial lower bound for an explicit problem Q in a weak model implies a major lower bound/separation in a strong model."Typically, Q is a meta-problem (e.g. a Gap-MCSP variant) and"major"means consequences such as NP not in P/poly.'. It must preserve the invariant 'MCSP complexity profile in nonuniform-circuits'. Its operational role is 'Recast the blocked frontier as a diagrammatic composition problem: identify objects, morphisms, mediating objects, and the exact commutativity or factorization obligation.'.
\item First nontrivial lemma: For every admissible 'MCSP obstruction family', the first route-specific realization step preserves 'MCSP complexity profile in nonuniform-circuits' and prepares the passage from 'A hardness magnification theorem has the form:"A barely-nontrivial lower bound for an explicit problem Q in a weak model implies a major lower bound/separation in a strong model."Typically, Q is a meta-problem (e.g. a Gap-MCSP variant) and"major"means consequences such as NP not in P/poly.' to 'In refined regimes, lower bounds for Gap-MCSP slightly beyond known frontiers (e.g. N\textasciicircum{} 1+varepsilon -type circuit lower bounds for suitable promise gaps and input length N=2\textasciicircum{}n ) imply major consequences such as NP not in P/poly.'.
\item Quantifier signature: For every admissible 'MCSP obstruction family' on the route from 'st.def.hardness-magnification-phenomenon-programme-leve' to 'st.thm.oliveira-pich-santhanam-magnification-near-state'.
\item LaTeX theorem: \textbackslash\{\}forall W,\textbackslash\{\} \textbackslash\{\}mathrm\{Adm\}(W) \textbackslash\{\}to \textbackslash\{\}exists w : \textbackslash\{\}mathsf\{BW\}(W),\textbackslash\{\} \textbackslash\{\}mathrm\{PresInv\}(W,w) \textbackslash\{\}wedge \textbackslash\{\}mathrm\{RealizesTarget\}(W,w).
\item Negative/barrier version: If 'MCSP obstruction family' cannot be realized with invariant 'MCSP complexity profile in nonuniform-circuits', the route cannot force 'In refined regimes, lower bounds for Gap-MCSP slightly beyond known frontiers (e.g. N\textasciicircum{} 1+varepsilon -type circuit lower bounds for suitable promise gaps and input length N=2\textasciicircum{}n ) imply major consequences such as NP not in P/poly.'. non\_commuting\_route
\item Proof strategy family: constructive-realization
\item Negative route: Assume 'MCSP obstruction family' cannot be realized while preserving 'MCSP complexity profile in nonuniform-circuits' strongly enough to activate the local target claim.
\item Failure diagnosis: The route fails if 'MCSP obstruction family' cannot be realized while preserving 'MCSP complexity profile in nonuniform-circuits'. Local negative probe: non\_commuting\_route.
\item Proof obligations: For every admissible 'MCSP obstruction family', the first route-specific realization step preserves 'MCSP complexity profile in nonuniform-circuits' and prepares the passage from 'A hardness magnification theorem has the form:"A barely-nontrivial lower bound for an explicit problem Q in a weak model implies a major lower bound/separation in a strong model."Typically, Q is a meta-problem (e.g. a Gap-MCSP variant) and"major"means consequences such as NP not in P/poly.' to 'In refined regimes, lower bounds for Gap-MCSP slightly beyond known frontiers (e.g. N\textasciicircum{} 1+varepsilon -type circuit lower bounds for suitable promise gaps and input length N=2\textasciicircum{}n ) imply major consequences such as NP not in P/poly.'.; Commutativity obligation: the route for 'carrier-realization' must commute between source theorem hypotheses, the mediating carrier, and the target theorem hypotheses.
\item Candidate lemmas: Realization theorem for MCSP obstruction family; Coherence diagram lemma for MCSP obstruction family
\item Independent value: A concrete special-case theorem for 'MCSP obstruction family' would be independently valuable because it converts a broad discovery programme into an inspectable mathematical subresult.
\item Expert gate: 6/6
\end{itemize}

\subsection*{conjpkg.disclayer-p-vs-np-family-explicitness.01}
\begin{itemize}
\item Exact statement: There exists an admissible 'PIT certificate family' such that 'A set Hsubseteq mathbb F \textasciicircum{}n is a hitting set for a class mathcal C of n -variate circuits if for every nonzero polynomial p computed by some circuit in mathcal C , there exists mathbf a in H with p(mathbf a )neq 0.' can be organized into a route preserving 'PIT semantic object' and thereby supports 'If mathrm PIT in Pclass (or even suitable infinitely-often derandomisations in nondeterministic subexponential time, in standard variants), then at least one of the following holds: NEXP not in P/poly quad or quad Permanent does not have polynomial-size arithmetic circuits.'.
\item Restricted-regime mode: special-case
\item Bridge object: PIT certificate family (checkable certificate map or algebraic witness relation)
\item Bridge definition: 'PIT certificate family' denotes the structured intermediate object required by this branch. It is meant to organize the source-side content 'A set Hsubseteq mathbb F \textasciicircum{}n is a hitting set for a class mathcal C of n -variate circuits if for every nonzero polynomial p computed by some circuit in mathcal C , there exists mathbf a in H with p(mathbf a )neq 0.'. It must preserve the invariant 'PIT semantic object'. Its operational role is 'Upgrade a semantic or existential algebraic object into an explicit algorithmic derandomization premise.'.
\item First nontrivial lemma: For every admissible 'PIT certificate family', the first route-specific realization step preserves 'PIT semantic object' and prepares the passage from 'A set Hsubseteq mathbb F \textasciicircum{}n is a hitting set for a class mathcal C of n -variate circuits if for every nonzero polynomial p computed by some circuit in mathcal C , there exists mathbf a in H with p(mathbf a )neq 0.' to 'If mathrm PIT in Pclass (or even suitable infinitely-often derandomisations in nondeterministic subexponential time, in standard variants), then at least one of the following holds: NEXP not in P/poly quad or quad Permanent does not have polynomial-size arithmetic circuits.'.
\item Quantifier signature: For every admissible 'PIT certificate family' on the route from 'st.def.hitting-set-for-pit-black-box' to 'st.thm.kabanets-impagliazzo-pit-derandomisation-implies'.
\item LaTeX theorem: \textbackslash\{\}forall W,\textbackslash\{\} \textbackslash\{\}mathrm\{Adm\}(W) \textbackslash\{\}to \textbackslash\{\}exists w : \textbackslash\{\}mathsf\{BW\}(W),\textbackslash\{\} \textbackslash\{\}mathrm\{PresInv\}(W,w) \textbackslash\{\}wedge \textbackslash\{\}mathrm\{RealizesTarget\}(W,w).
\item Negative/barrier version: If 'PIT certificate family' cannot be realized with invariant 'PIT semantic object', the route cannot force 'If mathrm PIT in Pclass (or even suitable infinitely-often derandomisations in nondeterministic subexponential time, in standard variants), then at least one of the following holds: NEXP not in P/poly quad or quad Permanent does not have polynomial-size arithmetic circuits.'.
\item Proof strategy family: constructive-realization
\item Negative route: Assume 'PIT certificate family' cannot be realized while preserving 'PIT semantic object' strongly enough to activate the local target claim.
\item Failure diagnosis: The route fails if 'PIT certificate family' cannot be realized while preserving 'PIT semantic object'.
\item Proof obligations: For every admissible 'PIT certificate family', the first route-specific realization step preserves 'PIT semantic object' and prepares the passage from 'A set Hsubseteq mathbb F \textasciicircum{}n is a hitting set for a class mathcal C of n -variate circuits if for every nonzero polynomial p computed by some circuit in mathcal C , there exists mathbf a in H with p(mathbf a )neq 0.' to 'If mathrm PIT in Pclass (or even suitable infinitely-often derandomisations in nondeterministic subexponential time, in standard variants), then at least one of the following holds: NEXP not in P/poly quad or quad Permanent does not have polynomial-size arithmetic circuits.'.
\item Candidate lemmas: Prove that a hitting-set or pseudorandom-generator family preserves 'PIT semantic object' while transporting the route from the black-box-PIT regime to the arithmetic-circuits regime strongly enough to derive 'If mathrm PIT in Pclass (or even suitable infinitely-often derandomisations in nondeterministic subexponential time, in standard...'.
\item Independent value: A concrete special-case theorem for 'PIT certificate family' would be independently valuable because it converts a broad discovery programme into an inspectable mathematical subresult.
\item Expert gate: 6/6
\end{itemize}

\subsection*{conjpkg.disclayer-p-vs-np-family-explicitness.02}
\begin{itemize}
\item Exact statement: There exists an admissible 'PIT witness family' such that 'A set Hsubseteq mathbb F \textasciicircum{}n is a hitting set for a class mathcal C of n -variate circuits if for every nonzero polynomial p computed by some circuit in mathcal C , there exists mathbf a in H with p(mathbf a )neq 0.' can be organized into a route preserving 'PIT semantic object' and thereby supports 'If mathrm PIT in Pclass (or even suitable infinitely-often derandomisations in nondeterministic subexponential time, in standard variants), then at least one of the following holds: NEXP not in P/poly quad or quad Permanent does not have polynomial-size arithmetic circuits.'.
\item Restricted-regime mode: special-case
\item Bridge object: PIT witness family (family of instances equipped with a monotone witness measure)
\item Bridge definition: 'PIT witness family' denotes the structured intermediate object required by this branch. It is meant to organize the source-side content 'A set Hsubseteq mathbb F \textasciicircum{}n is a hitting set for a class mathcal C of n -variate circuits if for every nonzero polynomial p computed by some circuit in mathcal C , there exists mathbf a in H with p(mathbf a )neq 0.'. It must preserve the invariant 'PIT semantic object'. Its operational role is 'Upgrade a semantic or existential algebraic object into an explicit algorithmic derandomization premise.'.
\item First nontrivial lemma: For every admissible 'PIT witness family', the first route-specific realization step preserves 'PIT semantic object' and prepares the passage from 'A set Hsubseteq mathbb F \textasciicircum{}n is a hitting set for a class mathcal C of n -variate circuits if for every nonzero polynomial p computed by some circuit in mathcal C , there exists mathbf a in H with p(mathbf a )neq 0.' to 'If mathrm PIT in Pclass (or even suitable infinitely-often derandomisations in nondeterministic subexponential time, in standard variants), then at least one of the following holds: NEXP not in P/poly quad or quad Permanent does not have polynomial-size arithmetic circuits.'.
\item Quantifier signature: For every admissible 'PIT witness family' on the route from 'st.def.hitting-set-for-pit-black-box' to 'st.thm.kabanets-impagliazzo-pit-derandomisation-implies'.
\item LaTeX theorem: \textbackslash\{\}forall W,\textbackslash\{\} \textbackslash\{\}mathrm\{Adm\}(W) \textbackslash\{\}to \textbackslash\{\}exists w : \textbackslash\{\}mathsf\{BW\}(W),\textbackslash\{\} \textbackslash\{\}mathrm\{PresInv\}(W,w) \textbackslash\{\}wedge \textbackslash\{\}mathrm\{RealizesTarget\}(W,w).
\item Negative/barrier version: If 'PIT witness family' cannot be realized with invariant 'PIT semantic object', the route cannot force 'If mathrm PIT in Pclass (or even suitable infinitely-often derandomisations in nondeterministic subexponential time, in standard variants), then at least one of the following holds: NEXP not in P/poly quad or quad Permanent does not have polynomial-size arithmetic circuits.'.
\item Proof strategy family: constructive-realization
\item Negative route: Assume 'PIT witness family' cannot be realized while preserving 'PIT semantic object' strongly enough to activate the local target claim.
\item Failure diagnosis: The route fails if 'PIT witness family' cannot be realized while preserving 'PIT semantic object'.
\item Proof obligations: For every admissible 'PIT witness family', the first route-specific realization step preserves 'PIT semantic object' and prepares the passage from 'A set Hsubseteq mathbb F \textasciicircum{}n is a hitting set for a class mathcal C of n -variate circuits if for every nonzero polynomial p computed by some circuit in mathcal C , there exists mathbf a in H with p(mathbf a )neq 0.' to 'If mathrm PIT in Pclass (or even suitable infinitely-often derandomisations in nondeterministic subexponential time, in standard variants), then at least one of the following holds: NEXP not in P/poly quad or quad Permanent does not have polynomial-size arithmetic circuits.'.
\item Candidate lemmas: Prove that a hitting-set or pseudorandom-generator family preserves 'PIT semantic object' while transporting the route from the black-box-PIT regime to the arithmetic-circuits regime strongly enough to derive 'If mathrm PIT in Pclass (or even suitable infinitely-often derandomisations in nondeterministic subexponential time, in standard...'.
\item Independent value: A concrete special-case theorem for 'PIT witness family' would be independently valuable because it converts a broad discovery programme into an inspectable mathematical subresult.
\item Expert gate: 6/6
\end{itemize}

\subsection*{conjpkg.disclayer-p-vs-np-family-explicitness.03}
\begin{itemize}
\item Exact statement: For every admissible 'PIT certificate family', there is a transport law from 'black-box-PIT' to 'arithmetic-circuits' that preserves 'PIT semantic object' and retains enough strength to activate 'If mathrm PIT in Pclass (or even suitable infinitely-often derandomisations in nondeterministic subexponential time, in standard variants), then at least one of the following holds: NEXP not in P/poly quad or quad Permanent does not have polynomial-size arithmetic circuits.'.
\item Restricted-regime mode: special-case
\item Bridge object: PIT certificate family (checkable certificate map or algebraic witness relation)
\item Bridge definition: 'PIT certificate family' denotes the structured intermediate object required by this branch. It is meant to organize the source-side content 'A set Hsubseteq mathbb F \textasciicircum{}n is a hitting set for a class mathcal C of n -variate circuits if for every nonzero polynomial p computed by some circuit in mathcal C , there exists mathbf a in H with p(mathbf a )neq 0.'. It must preserve the invariant 'PIT semantic object'. Its operational role is 'Upgrade a semantic or existential algebraic object into an explicit algorithmic derandomization premise.'.
\item First nontrivial lemma: For every admissible 'PIT certificate family', the first route-specific realization step preserves 'PIT semantic object' and prepares the passage from 'A set Hsubseteq mathbb F \textasciicircum{}n is a hitting set for a class mathcal C of n -variate circuits if for every nonzero polynomial p computed by some circuit in mathcal C , there exists mathbf a in H with p(mathbf a )neq 0.' to 'If mathrm PIT in Pclass (or even suitable infinitely-often derandomisations in nondeterministic subexponential time, in standard variants), then at least one of the following holds: NEXP not in P/poly quad or quad Permanent does not have polynomial-size arithmetic circuits.'.
\item Quantifier signature: For every admissible 'PIT certificate family' on the route from 'st.def.hitting-set-for-pit-black-box' to 'st.thm.kabanets-impagliazzo-pit-derandomisation-implies'.
\item LaTeX theorem: \textbackslash\{\}forall W,\textbackslash\{\} \textbackslash\{\}mathrm\{Adm\}(W) \textbackslash\{\}to \textbackslash\{\}exists w : \textbackslash\{\}mathsf\{BW\}(W),\textbackslash\{\} \textbackslash\{\}mathrm\{PresInv\}(W,w) \textbackslash\{\}wedge \textbackslash\{\}mathrm\{RealizesTarget\}(W,w).
\item Negative/barrier version: If 'PIT certificate family' cannot be realized with invariant 'PIT semantic object', the route cannot force 'If mathrm PIT in Pclass (or even suitable infinitely-often derandomisations in nondeterministic subexponential time, in standard variants), then at least one of the following holds: NEXP not in P/poly quad or quad Permanent does not have polynomial-size arithmetic circuits.'.
\item Proof strategy family: constructive-realization
\item Negative route: Assume 'PIT certificate family' cannot be realized while preserving 'PIT semantic object' strongly enough to activate the local target claim.
\item Failure diagnosis: The route fails if 'PIT certificate family' cannot be realized while preserving 'PIT semantic object'.
\item Proof obligations: For every admissible 'PIT certificate family', the first route-specific realization step preserves 'PIT semantic object' and prepares the passage from 'A set Hsubseteq mathbb F \textasciicircum{}n is a hitting set for a class mathcal C of n -variate circuits if for every nonzero polynomial p computed by some circuit in mathcal C , there exists mathbf a in H with p(mathbf a )neq 0.' to 'If mathrm PIT in Pclass (or even suitable infinitely-often derandomisations in nondeterministic subexponential time, in standard variants), then at least one of the following holds: NEXP not in P/poly quad or quad Permanent does not have polynomial-size arithmetic circuits.'.
\item Candidate lemmas: Transport theorem for PIT certificate family
\item Independent value: A concrete special-case theorem for 'PIT certificate family' would be independently valuable because it converts a broad discovery programme into an inspectable mathematical subresult.
\item Expert gate: 6/6
\end{itemize}

\subsection*{conjpkg.disclayer-p-vs-np-family-explicitness.04}
\begin{itemize}
\item Exact statement: No admissible 'PIT certificate family' can simultaneously preserve 'PIT semantic object' and realize 'If mathrm PIT in Pclass (or even suitable infinitely-often derandomisations in nondeterministic subexponential time, in standard variants), then at least one of the following holds: NEXP not in P/poly quad or quad Permanent does not have polynomial-size arithmetic circuits.' under the current route commitments.
\item Restricted-regime mode: special-case
\item Bridge object: PIT certificate family (checkable certificate map or algebraic witness relation)
\item Bridge definition: 'PIT certificate family' denotes the structured intermediate object required by this branch. It is meant to organize the source-side content 'A set Hsubseteq mathbb F \textasciicircum{}n is a hitting set for a class mathcal C of n -variate circuits if for every nonzero polynomial p computed by some circuit in mathcal C , there exists mathbf a in H with p(mathbf a )neq 0.'. It must preserve the invariant 'PIT semantic object'. Its operational role is 'Upgrade a semantic or existential algebraic object into an explicit algorithmic derandomization premise.'.
\item First nontrivial lemma: For every admissible 'PIT certificate family', the first route-specific realization step preserves 'PIT semantic object' and prepares the passage from 'A set Hsubseteq mathbb F \textasciicircum{}n is a hitting set for a class mathcal C of n -variate circuits if for every nonzero polynomial p computed by some circuit in mathcal C , there exists mathbf a in H with p(mathbf a )neq 0.' to 'If mathrm PIT in Pclass (or even suitable infinitely-often derandomisations in nondeterministic subexponential time, in standard variants), then at least one of the following holds: NEXP not in P/poly quad or quad Permanent does not have polynomial-size arithmetic circuits.'.
\item Quantifier signature: For every admissible 'PIT certificate family' on the route from 'st.def.hitting-set-for-pit-black-box' to 'st.thm.kabanets-impagliazzo-pit-derandomisation-implies'.
\item LaTeX theorem: \textbackslash\{\}forall W,\textbackslash\{\} \textbackslash\{\}mathrm\{Adm\}(W) \textbackslash\{\}to \textbackslash\{\}neg \textbackslash\{\}exists w : \textbackslash\{\}mathsf\{BW\}(W),\textbackslash\{\} \textbackslash\{\}mathrm\{PresInv\}(W,w) \textbackslash\{\}wedge \textbackslash\{\}mathrm\{RealizesTarget\}(W,w).
\item Negative/barrier version: If 'PIT certificate family' cannot be realized with invariant 'PIT semantic object', the route cannot force 'If mathrm PIT in Pclass (or even suitable infinitely-often derandomisations in nondeterministic subexponential time, in standard variants), then at least one of the following holds: NEXP not in P/poly quad or quad Permanent does not have polynomial-size arithmetic circuits.'.
\item Proof strategy family: constructive-realization
\item Negative route: Assume 'PIT certificate family' cannot be realized while preserving 'PIT semantic object' strongly enough to activate the local target claim.
\item Failure diagnosis: The route fails if 'PIT certificate family' cannot be realized while preserving 'PIT semantic object'.
\item Proof obligations: For every admissible 'PIT certificate family', the first route-specific realization step preserves 'PIT semantic object' and prepares the passage from 'A set Hsubseteq mathbb F \textasciicircum{}n is a hitting set for a class mathcal C of n -variate circuits if for every nonzero polynomial p computed by some circuit in mathcal C , there exists mathbf a in H with p(mathbf a )neq 0.' to 'If mathrm PIT in Pclass (or even suitable infinitely-often derandomisations in nondeterministic subexponential time, in standard variants), then at least one of the following holds: NEXP not in P/poly quad or quad Permanent does not have polynomial-size arithmetic circuits.'.
\item Candidate lemmas: Obstruction theorem for PIT certificate family
\item Independent value: A concrete special-case theorem for 'PIT certificate family' would be independently valuable because it converts a broad discovery programme into an inspectable mathematical subresult.
\item Expert gate: 6/6
\end{itemize}

\subsection*{conjpkg.disclayer-p-vs-np-family-explicitness.05}
\begin{itemize}
\item Exact statement: There exists an admissible 'PIT certificate family' whose realization internalizes 'A set Hsubseteq mathbb F \textasciicircum{}n is a hitting set for a class mathcal C of n -variate circuits if for every nonzero polynomial p computed by some circuit in mathcal C , there exists mathbf a in H with p(mathbf a )neq 0.' while preserving 'PIT semantic object' strongly enough to support 'If mathrm PIT in Pclass (or even suitable infinitely-often derandomisations in nondeterministic subexponential time, in standard variants), then at least one of the following holds: NEXP not in P/poly quad or quad Permanent does not have polynomial-size arithmetic circuits.'.
\item Restricted-regime mode: special-case
\item Bridge object: PIT certificate family (checkable certificate map or algebraic witness relation)
\item Bridge definition: 'PIT certificate family' denotes the structured intermediate object required by this branch. It is meant to organize the source-side content 'A set Hsubseteq mathbb F \textasciicircum{}n is a hitting set for a class mathcal C of n -variate circuits if for every nonzero polynomial p computed by some circuit in mathcal C , there exists mathbf a in H with p(mathbf a )neq 0.'. It must preserve the invariant 'PIT semantic object'. Its operational role is 'Upgrade a semantic or existential algebraic object into an explicit algorithmic derandomization premise.'.
\item First nontrivial lemma: For every admissible 'PIT certificate family', the first route-specific realization step preserves 'PIT semantic object' and prepares the passage from 'A set Hsubseteq mathbb F \textasciicircum{}n is a hitting set for a class mathcal C of n -variate circuits if for every nonzero polynomial p computed by some circuit in mathcal C , there exists mathbf a in H with p(mathbf a )neq 0.' to 'If mathrm PIT in Pclass (or even suitable infinitely-often derandomisations in nondeterministic subexponential time, in standard variants), then at least one of the following holds: NEXP not in P/poly quad or quad Permanent does not have polynomial-size arithmetic circuits.'.
\item Quantifier signature: For every admissible 'PIT certificate family' on the route from 'st.def.hitting-set-for-pit-black-box' to 'st.thm.kabanets-impagliazzo-pit-derandomisation-implies'.
\item LaTeX theorem: \textbackslash\{\}forall W,\textbackslash\{\} \textbackslash\{\}mathrm\{Adm\}(W) \textbackslash\{\}to \textbackslash\{\}exists w : \textbackslash\{\}mathsf\{BW\}(W),\textbackslash\{\} \textbackslash\{\}mathrm\{PresInv\}(W,w) \textbackslash\{\}wedge \textbackslash\{\}mathrm\{RealizesTarget\}(W,w).
\item Negative/barrier version: If 'PIT certificate family' cannot be realized with invariant 'PIT semantic object', the route cannot force 'If mathrm PIT in Pclass (or even suitable infinitely-often derandomisations in nondeterministic subexponential time, in standard variants), then at least one of the following holds: NEXP not in P/poly quad or quad Permanent does not have polynomial-size arithmetic circuits.'.
\item Proof strategy family: constructive-realization
\item Negative route: Assume 'PIT certificate family' cannot be realized while preserving 'PIT semantic object' strongly enough to activate the local target claim.
\item Failure diagnosis: The route fails if 'PIT certificate family' cannot be realized while preserving 'PIT semantic object'.
\item Proof obligations: For every admissible 'PIT certificate family', the first route-specific realization step preserves 'PIT semantic object' and prepares the passage from 'A set Hsubseteq mathbb F \textasciicircum{}n is a hitting set for a class mathcal C of n -variate circuits if for every nonzero polynomial p computed by some circuit in mathcal C , there exists mathbf a in H with p(mathbf a )neq 0.' to 'If mathrm PIT in Pclass (or even suitable infinitely-often derandomisations in nondeterministic subexponential time, in standard variants), then at least one of the following holds: NEXP not in P/poly quad or quad Permanent does not have polynomial-size arithmetic circuits.'.
\item Candidate lemmas: Realization theorem for PIT certificate family
\item Independent value: A concrete special-case theorem for 'PIT certificate family' would be independently valuable because it converts a broad discovery programme into an inspectable mathematical subresult.
\item Expert gate: 6/6
\end{itemize}

\subsection*{conjpkg.disclayer-p-vs-np-family-explicitness.06}
\begin{itemize}
\item Exact statement: For every admissible 'PIT witness family', there is a transport law from 'black-box-PIT' to 'arithmetic-circuits' that preserves 'PIT semantic object' and retains enough strength to activate 'If mathrm PIT in Pclass (or even suitable infinitely-often derandomisations in nondeterministic subexponential time, in standard variants), then at least one of the following holds: NEXP not in P/poly quad or quad Permanent does not have polynomial-size arithmetic circuits.'.
\item Restricted-regime mode: special-case
\item Bridge object: PIT witness family (family of instances equipped with a monotone witness measure)
\item Bridge definition: 'PIT witness family' denotes the structured intermediate object required by this branch. It is meant to organize the source-side content 'A set Hsubseteq mathbb F \textasciicircum{}n is a hitting set for a class mathcal C of n -variate circuits if for every nonzero polynomial p computed by some circuit in mathcal C , there exists mathbf a in H with p(mathbf a )neq 0.'. It must preserve the invariant 'PIT semantic object'. Its operational role is 'Upgrade a semantic or existential algebraic object into an explicit algorithmic derandomization premise.'.
\item First nontrivial lemma: For every admissible 'PIT witness family', the first route-specific realization step preserves 'PIT semantic object' and prepares the passage from 'A set Hsubseteq mathbb F \textasciicircum{}n is a hitting set for a class mathcal C of n -variate circuits if for every nonzero polynomial p computed by some circuit in mathcal C , there exists mathbf a in H with p(mathbf a )neq 0.' to 'If mathrm PIT in Pclass (or even suitable infinitely-often derandomisations in nondeterministic subexponential time, in standard variants), then at least one of the following holds: NEXP not in P/poly quad or quad Permanent does not have polynomial-size arithmetic circuits.'.
\item Quantifier signature: For every admissible 'PIT witness family' on the route from 'st.def.hitting-set-for-pit-black-box' to 'st.thm.kabanets-impagliazzo-pit-derandomisation-implies'.
\item LaTeX theorem: \textbackslash\{\}forall W,\textbackslash\{\} \textbackslash\{\}mathrm\{Adm\}(W) \textbackslash\{\}to \textbackslash\{\}exists w : \textbackslash\{\}mathsf\{BW\}(W),\textbackslash\{\} \textbackslash\{\}mathrm\{PresInv\}(W,w) \textbackslash\{\}wedge \textbackslash\{\}mathrm\{RealizesTarget\}(W,w).
\item Negative/barrier version: If 'PIT witness family' cannot be realized with invariant 'PIT semantic object', the route cannot force 'If mathrm PIT in Pclass (or even suitable infinitely-often derandomisations in nondeterministic subexponential time, in standard variants), then at least one of the following holds: NEXP not in P/poly quad or quad Permanent does not have polynomial-size arithmetic circuits.'.
\item Proof strategy family: constructive-realization
\item Negative route: Assume 'PIT witness family' cannot be realized while preserving 'PIT semantic object' strongly enough to activate the local target claim.
\item Failure diagnosis: The route fails if 'PIT witness family' cannot be realized while preserving 'PIT semantic object'.
\item Proof obligations: For every admissible 'PIT witness family', the first route-specific realization step preserves 'PIT semantic object' and prepares the passage from 'A set Hsubseteq mathbb F \textasciicircum{}n is a hitting set for a class mathcal C of n -variate circuits if for every nonzero polynomial p computed by some circuit in mathcal C , there exists mathbf a in H with p(mathbf a )neq 0.' to 'If mathrm PIT in Pclass (or even suitable infinitely-often derandomisations in nondeterministic subexponential time, in standard variants), then at least one of the following holds: NEXP not in P/poly quad or quad Permanent does not have polynomial-size arithmetic circuits.'.
\item Candidate lemmas: Transport theorem for PIT witness family
\item Independent value: A concrete special-case theorem for 'PIT witness family' would be independently valuable because it converts a broad discovery programme into an inspectable mathematical subresult.
\item Expert gate: 6/6
\end{itemize}

\subsection*{conjpkg.disclayer-p-vs-np-family-explicitness.07}
\begin{itemize}
\item Exact statement: No admissible 'PIT witness family' can simultaneously preserve 'PIT semantic object' and realize 'If mathrm PIT in Pclass (or even suitable infinitely-often derandomisations in nondeterministic subexponential time, in standard variants), then at least one of the following holds: NEXP not in P/poly quad or quad Permanent does not have polynomial-size arithmetic circuits.' under the current route commitments.
\item Restricted-regime mode: special-case
\item Bridge object: PIT witness family (family of instances equipped with a monotone witness measure)
\item Bridge definition: 'PIT witness family' denotes the structured intermediate object required by this branch. It is meant to organize the source-side content 'A set Hsubseteq mathbb F \textasciicircum{}n is a hitting set for a class mathcal C of n -variate circuits if for every nonzero polynomial p computed by some circuit in mathcal C , there exists mathbf a in H with p(mathbf a )neq 0.'. It must preserve the invariant 'PIT semantic object'. Its operational role is 'Upgrade a semantic or existential algebraic object into an explicit algorithmic derandomization premise.'.
\item First nontrivial lemma: For every admissible 'PIT witness family', the first route-specific realization step preserves 'PIT semantic object' and prepares the passage from 'A set Hsubseteq mathbb F \textasciicircum{}n is a hitting set for a class mathcal C of n -variate circuits if for every nonzero polynomial p computed by some circuit in mathcal C , there exists mathbf a in H with p(mathbf a )neq 0.' to 'If mathrm PIT in Pclass (or even suitable infinitely-often derandomisations in nondeterministic subexponential time, in standard variants), then at least one of the following holds: NEXP not in P/poly quad or quad Permanent does not have polynomial-size arithmetic circuits.'.
\item Quantifier signature: For every admissible 'PIT witness family' on the route from 'st.def.hitting-set-for-pit-black-box' to 'st.thm.kabanets-impagliazzo-pit-derandomisation-implies'.
\item LaTeX theorem: \textbackslash\{\}forall W,\textbackslash\{\} \textbackslash\{\}mathrm\{Adm\}(W) \textbackslash\{\}to \textbackslash\{\}neg \textbackslash\{\}exists w : \textbackslash\{\}mathsf\{BW\}(W),\textbackslash\{\} \textbackslash\{\}mathrm\{PresInv\}(W,w) \textbackslash\{\}wedge \textbackslash\{\}mathrm\{RealizesTarget\}(W,w).
\item Negative/barrier version: If 'PIT witness family' cannot be realized with invariant 'PIT semantic object', the route cannot force 'If mathrm PIT in Pclass (or even suitable infinitely-often derandomisations in nondeterministic subexponential time, in standard variants), then at least one of the following holds: NEXP not in P/poly quad or quad Permanent does not have polynomial-size arithmetic circuits.'.
\item Proof strategy family: constructive-realization
\item Negative route: Assume 'PIT witness family' cannot be realized while preserving 'PIT semantic object' strongly enough to activate the local target claim.
\item Failure diagnosis: The route fails if 'PIT witness family' cannot be realized while preserving 'PIT semantic object'.
\item Proof obligations: For every admissible 'PIT witness family', the first route-specific realization step preserves 'PIT semantic object' and prepares the passage from 'A set Hsubseteq mathbb F \textasciicircum{}n is a hitting set for a class mathcal C of n -variate circuits if for every nonzero polynomial p computed by some circuit in mathcal C , there exists mathbf a in H with p(mathbf a )neq 0.' to 'If mathrm PIT in Pclass (or even suitable infinitely-often derandomisations in nondeterministic subexponential time, in standard variants), then at least one of the following holds: NEXP not in P/poly quad or quad Permanent does not have polynomial-size arithmetic circuits.'.
\item Candidate lemmas: Obstruction theorem for PIT witness family
\item Independent value: A concrete special-case theorem for 'PIT witness family' would be independently valuable because it converts a broad discovery programme into an inspectable mathematical subresult.
\item Expert gate: 6/6
\end{itemize}

\subsection*{conjpkg.disclayer-p-vs-np-family-explicitness.08}
\begin{itemize}
\item Exact statement: There exists an admissible 'PIT witness family' whose realization internalizes 'A set Hsubseteq mathbb F \textasciicircum{}n is a hitting set for a class mathcal C of n -variate circuits if for every nonzero polynomial p computed by some circuit in mathcal C , there exists mathbf a in H with p(mathbf a )neq 0.' while preserving 'PIT semantic object' strongly enough to support 'If mathrm PIT in Pclass (or even suitable infinitely-often derandomisations in nondeterministic subexponential time, in standard variants), then at least one of the following holds: NEXP not in P/poly quad or quad Permanent does not have polynomial-size arithmetic circuits.'.
\item Restricted-regime mode: special-case
\item Bridge object: PIT witness family (family of instances equipped with a monotone witness measure)
\item Bridge definition: 'PIT witness family' denotes the structured intermediate object required by this branch. It is meant to organize the source-side content 'A set Hsubseteq mathbb F \textasciicircum{}n is a hitting set for a class mathcal C of n -variate circuits if for every nonzero polynomial p computed by some circuit in mathcal C , there exists mathbf a in H with p(mathbf a )neq 0.'. It must preserve the invariant 'PIT semantic object'. Its operational role is 'Upgrade a semantic or existential algebraic object into an explicit algorithmic derandomization premise.'.
\item First nontrivial lemma: For every admissible 'PIT witness family', the first route-specific realization step preserves 'PIT semantic object' and prepares the passage from 'A set Hsubseteq mathbb F \textasciicircum{}n is a hitting set for a class mathcal C of n -variate circuits if for every nonzero polynomial p computed by some circuit in mathcal C , there exists mathbf a in H with p(mathbf a )neq 0.' to 'If mathrm PIT in Pclass (or even suitable infinitely-often derandomisations in nondeterministic subexponential time, in standard variants), then at least one of the following holds: NEXP not in P/poly quad or quad Permanent does not have polynomial-size arithmetic circuits.'.
\item Quantifier signature: For every admissible 'PIT witness family' on the route from 'st.def.hitting-set-for-pit-black-box' to 'st.thm.kabanets-impagliazzo-pit-derandomisation-implies'.
\item LaTeX theorem: \textbackslash\{\}forall W,\textbackslash\{\} \textbackslash\{\}mathrm\{Adm\}(W) \textbackslash\{\}to \textbackslash\{\}exists w : \textbackslash\{\}mathsf\{BW\}(W),\textbackslash\{\} \textbackslash\{\}mathrm\{PresInv\}(W,w) \textbackslash\{\}wedge \textbackslash\{\}mathrm\{RealizesTarget\}(W,w).
\item Negative/barrier version: If 'PIT witness family' cannot be realized with invariant 'PIT semantic object', the route cannot force 'If mathrm PIT in Pclass (or even suitable infinitely-often derandomisations in nondeterministic subexponential time, in standard variants), then at least one of the following holds: NEXP not in P/poly quad or quad Permanent does not have polynomial-size arithmetic circuits.'.
\item Proof strategy family: constructive-realization
\item Negative route: Assume 'PIT witness family' cannot be realized while preserving 'PIT semantic object' strongly enough to activate the local target claim.
\item Failure diagnosis: The route fails if 'PIT witness family' cannot be realized while preserving 'PIT semantic object'.
\item Proof obligations: For every admissible 'PIT witness family', the first route-specific realization step preserves 'PIT semantic object' and prepares the passage from 'A set Hsubseteq mathbb F \textasciicircum{}n is a hitting set for a class mathcal C of n -variate circuits if for every nonzero polynomial p computed by some circuit in mathcal C , there exists mathbf a in H with p(mathbf a )neq 0.' to 'If mathrm PIT in Pclass (or even suitable infinitely-often derandomisations in nondeterministic subexponential time, in standard variants), then at least one of the following holds: NEXP not in P/poly quad or quad Permanent does not have polynomial-size arithmetic circuits.'.
\item Candidate lemmas: Realization theorem for PIT witness family
\item Independent value: A concrete special-case theorem for 'PIT witness family' would be independently valuable because it converts a broad discovery programme into an inspectable mathematical subresult.
\item Expert gate: 6/6
\end{itemize}

\subsection*{conjpkg.disclayer-p-vs-np-family-barrier-aware-transfer.01}
\begin{itemize}
\item Exact statement: There exists an admissible 'unspecified witness family' such that 'There exist many problems L in combinatorial optimisation/graph theory such that L is NP-complete under polynomial-time many-one reductions.' can be organized into a route preserving 'consequence schema NP-complete-classification' and thereby supports 'Under"local"polynomial-time reductions (where each output bit is computable with very small time), MCSP is not NP-complete. More broadly, the paper develops evidence and conditional consequences suggesting NP-hardness of MCSP under various reductions would entail unexpected class separations.'.
\item Restricted-regime mode: no-go-theorem
\item Bridge object: unspecified witness family (family of instances equipped with a monotone witness measure)
\item Bridge definition: 'unspecified witness family' denotes the structured intermediate object required by this branch. It is meant to organize the source-side content 'There exist many problems L in combinatorial optimisation/graph theory such that L is NP-complete under polynomial-time many-one reductions.'. It must preserve the invariant 'consequence schema NP-complete-classification'. Its operational role is 'Design a transfer mechanism that survives known local-reduction barriers, or prove that no such mechanism exists under explicit assumptions.'.
\item First nontrivial lemma: If 'Design a transfer mechanism that survives known local-reduction barriers, or prove that no such mechanism exists under explicit assumptions.' fails for 'unspecified witness family', then no route preserving 'consequence schema NP-complete-classification' can force 'Under"local"polynomial-time reductions (where each output bit is computable with very small time), MCSP is not NP-complete. More broadly, the paper develops evidence and conditional consequences suggesting NP-hardness of MCSP under various reductions would entail unexpected class separations.' from 'There exist many problems L in combinatorial optimisation/graph theory such that L is NP-complete under polynomial-time many-one reductions.'.
\item Quantifier signature: For every admissible 'unspecified witness family' on the route from 'st.thm.karp-many-natural-problems-are-np-complete-karp1' to 'st.thm.murray-williams-non-np-hardness-of-mcsp-under-lo'.
\item LaTeX theorem: \textbackslash\{\}forall W,\textbackslash\{\} \textbackslash\{\}mathrm\{Adm\}(W) \textbackslash\{\}to \textbackslash\{\}exists w : \textbackslash\{\}mathsf\{BW\}(W),\textbackslash\{\} \textbackslash\{\}mathrm\{PresInv\}(W,w) \textbackslash\{\}wedge \textbackslash\{\}mathrm\{RealizesTarget\}(W,w).
\item Negative/barrier version: Every attempted realization of 'unspecified witness family' remains blocked by the current barrier, so 'Under"local"polynomial-time reductions (where each output bit is computable with very small time), MCSP is not NP-complete. More broadly, the paper develops evidence and conditional consequences suggesting NP-hardness of MCSP under various reductions would entail unexpected class separations.' does not follow.
\item Proof strategy family: barrier
\item Negative route: Assume every attempted realization of 'unspecified witness family' collapses into the current barrier.
\item Failure diagnosis: The route fails if every admissible realization of 'unspecified witness family' remains trapped behind the current barrier.
\item Proof obligations: If 'Design a transfer mechanism that survives known local-reduction barriers, or prove that no such mechanism exists under explicit assumptions.' fails for 'unspecified witness family', then no route preserving 'consequence schema NP-complete-classification' can force 'Under"local"polynomial-time reductions (where each output bit is computable with very small time), MCSP is not NP-complete. More broadly, the paper develops evidence and conditional consequences suggesting NP-hardness of MCSP under various reductions would entail unexpected class separations.' from 'There exist many problems L in combinatorial optimisation/graph theory such that L is NP-complete under polynomial-time many-one reductions.'.
\item Candidate lemmas: Prove that a hardness-magnification witness family turns 'NP-complete under polynomial-time many-one reductions' into 'is not NP-complete. More broadly, the paper develops evidence and conditional consequences suggesting NP-hardness of MCSP under various...' while preserving 'consequence schema NP-complete-classification'.; Barrier obstruction lemma for unspecified witness family
\item Independent value: A no-go theorem for 'unspecified witness family' would be independently valuable because it collapses an entire discovery subspace and exposes a genuine barrier theorem.
\item Expert gate: 6/6
\end{itemize}

\subsection*{conjpkg.disclayer-p-vs-np-family-barrier-aware-transfer.02}
\begin{itemize}
\item Exact statement: There exists an admissible 'unspecified reduction gadget family' such that 'There exist many problems L in combinatorial optimisation/graph theory such that L is NP-complete under polynomial-time many-one reductions.' can be organized into a route preserving 'consequence schema NP-complete-classification' and thereby supports 'Under"local"polynomial-time reductions (where each output bit is computable with very small time), MCSP is not NP-complete. More broadly, the paper develops evidence and conditional consequences suggesting NP-hardness of MCSP under various reductions would entail unexpected class separations.'.
\item Restricted-regime mode: no-go-theorem
\item Bridge object: unspecified reduction gadget family (family of reductions preserving the route witness structure)
\item Bridge definition: 'unspecified reduction gadget family' denotes the structured intermediate object required by this branch. It is meant to organize the source-side content 'There exist many problems L in combinatorial optimisation/graph theory such that L is NP-complete under polynomial-time many-one reductions.'. It must preserve the invariant 'consequence schema NP-complete-classification'. Its operational role is 'Design a transfer mechanism that survives known local-reduction barriers, or prove that no such mechanism exists under explicit assumptions.'.
\item First nontrivial lemma: If 'Design a transfer mechanism that survives known local-reduction barriers, or prove that no such mechanism exists under explicit assumptions.' fails for 'unspecified reduction gadget family', then no route preserving 'consequence schema NP-complete-classification' can force 'Under"local"polynomial-time reductions (where each output bit is computable with very small time), MCSP is not NP-complete. More broadly, the paper develops evidence and conditional consequences suggesting NP-hardness of MCSP under various reductions would entail unexpected class separations.' from 'There exist many problems L in combinatorial optimisation/graph theory such that L is NP-complete under polynomial-time many-one reductions.'.
\item Quantifier signature: For every admissible 'unspecified reduction gadget family' on the route from 'st.thm.karp-many-natural-problems-are-np-complete-karp1' to 'st.thm.murray-williams-non-np-hardness-of-mcsp-under-lo'.
\item LaTeX theorem: \textbackslash\{\}forall W,\textbackslash\{\} \textbackslash\{\}mathrm\{Adm\}(W) \textbackslash\{\}to \textbackslash\{\}exists w : \textbackslash\{\}mathsf\{BW\}(W),\textbackslash\{\} \textbackslash\{\}mathrm\{PresInv\}(W,w) \textbackslash\{\}wedge \textbackslash\{\}mathrm\{RealizesTarget\}(W,w).
\item Negative/barrier version: Every attempted realization of 'unspecified reduction gadget family' remains blocked by the current barrier, so 'Under"local"polynomial-time reductions (where each output bit is computable with very small time), MCSP is not NP-complete. More broadly, the paper develops evidence and conditional consequences suggesting NP-hardness of MCSP under various reductions would entail unexpected class separations.' does not follow.
\item Proof strategy family: barrier
\item Negative route: Assume every attempted realization of 'unspecified reduction gadget family' collapses into the current barrier.
\item Failure diagnosis: The route fails if every admissible realization of 'unspecified reduction gadget family' remains trapped behind the current barrier.
\item Proof obligations: If 'Design a transfer mechanism that survives known local-reduction barriers, or prove that no such mechanism exists under explicit assumptions.' fails for 'unspecified reduction gadget family', then no route preserving 'consequence schema NP-complete-classification' can force 'Under"local"polynomial-time reductions (where each output bit is computable with very small time), MCSP is not NP-complete. More broadly, the paper develops evidence and conditional consequences suggesting NP-hardness of MCSP under various reductions would entail unexpected class separations.' from 'There exist many problems L in combinatorial optimisation/graph theory such that L is NP-complete under polynomial-time many-one reductions.'.
\item Candidate lemmas: Prove that a hardness-magnification witness family turns 'NP-complete under polynomial-time many-one reductions' into 'is not NP-complete. More broadly, the paper develops evidence and conditional consequences suggesting NP-hardness of MCSP under various...' while preserving 'consequence schema NP-complete-classification'.; Barrier obstruction lemma for unspecified reduction gadget family
\item Independent value: A no-go theorem for 'unspecified reduction gadget family' would be independently valuable because it collapses an entire discovery subspace and exposes a genuine barrier theorem.
\item Expert gate: 6/6
\end{itemize}

\subsection*{conjpkg.disclayer-p-vs-np-family-barrier-aware-transfer.03}
\begin{itemize}
\item Exact statement: No admissible 'unspecified witness family' can simultaneously preserve 'consequence schema NP-complete-classification' and realize 'Under"local"polynomial-time reductions (where each output bit is computable with very small time), MCSP is not NP-complete. More broadly, the paper develops evidence and conditional consequences suggesting NP-hardness of MCSP under various reductions would entail unexpected class separations.' under the current route commitments.
\item Restricted-regime mode: no-go-theorem
\item Bridge object: unspecified witness family (family of instances equipped with a monotone witness measure)
\item Bridge definition: 'unspecified witness family' denotes the structured intermediate object required by this branch. It is meant to organize the source-side content 'There exist many problems L in combinatorial optimisation/graph theory such that L is NP-complete under polynomial-time many-one reductions.'. It must preserve the invariant 'consequence schema NP-complete-classification'. Its operational role is 'Design a transfer mechanism that survives known local-reduction barriers, or prove that no such mechanism exists under explicit assumptions.'.
\item First nontrivial lemma: If 'Design a transfer mechanism that survives known local-reduction barriers, or prove that no such mechanism exists under explicit assumptions.' fails for 'unspecified witness family', then no route preserving 'consequence schema NP-complete-classification' can force 'Under"local"polynomial-time reductions (where each output bit is computable with very small time), MCSP is not NP-complete. More broadly, the paper develops evidence and conditional consequences suggesting NP-hardness of MCSP under various reductions would entail unexpected class separations.' from 'There exist many problems L in combinatorial optimisation/graph theory such that L is NP-complete under polynomial-time many-one reductions.'.
\item Quantifier signature: For every admissible 'unspecified witness family' on the route from 'st.thm.karp-many-natural-problems-are-np-complete-karp1' to 'st.thm.murray-williams-non-np-hardness-of-mcsp-under-lo'.
\item LaTeX theorem: \textbackslash\{\}forall W,\textbackslash\{\} \textbackslash\{\}mathrm\{Adm\}(W) \textbackslash\{\}to \textbackslash\{\}neg \textbackslash\{\}exists w : \textbackslash\{\}mathsf\{BW\}(W),\textbackslash\{\} \textbackslash\{\}mathrm\{PresInv\}(W,w) \textbackslash\{\}wedge \textbackslash\{\}mathrm\{RealizesTarget\}(W,w).
\item Negative/barrier version: Every attempted realization of 'unspecified witness family' remains blocked by the current barrier, so 'Under"local"polynomial-time reductions (where each output bit is computable with very small time), MCSP is not NP-complete. More broadly, the paper develops evidence and conditional consequences suggesting NP-hardness of MCSP under various reductions would entail unexpected class separations.' does not follow.
\item Proof strategy family: barrier
\item Negative route: Assume every attempted realization of 'unspecified witness family' collapses into the current barrier.
\item Failure diagnosis: The route fails if every admissible realization of 'unspecified witness family' remains trapped behind the current barrier.
\item Proof obligations: If 'Design a transfer mechanism that survives known local-reduction barriers, or prove that no such mechanism exists under explicit assumptions.' fails for 'unspecified witness family', then no route preserving 'consequence schema NP-complete-classification' can force 'Under"local"polynomial-time reductions (where each output bit is computable with very small time), MCSP is not NP-complete. More broadly, the paper develops evidence and conditional consequences suggesting NP-hardness of MCSP under various reductions would entail unexpected class separations.' from 'There exist many problems L in combinatorial optimisation/graph theory such that L is NP-complete under polynomial-time many-one reductions.'.
\item Candidate lemmas: Obstruction theorem for unspecified witness family; Barrier obstruction lemma for unspecified witness family
\item Independent value: A no-go theorem for 'unspecified witness family' would be independently valuable because it collapses an entire discovery subspace and exposes a genuine barrier theorem.
\item Expert gate: 6/6
\end{itemize}

\subsection*{conjpkg.disclayer-p-vs-np-family-barrier-aware-transfer.04}
\begin{itemize}
\item Exact statement: There exists an admissible 'unspecified witness family' whose realization internalizes 'There exist many problems L in combinatorial optimisation/graph theory such that L is NP-complete under polynomial-time many-one reductions.' while preserving 'consequence schema NP-complete-classification' strongly enough to support 'Under"local"polynomial-time reductions (where each output bit is computable with very small time), MCSP is not NP-complete. More broadly, the paper develops evidence and conditional consequences suggesting NP-hardness of MCSP under various reductions would entail unexpected class separations.'.
\item Restricted-regime mode: no-go-theorem
\item Bridge object: unspecified witness family (family of instances equipped with a monotone witness measure)
\item Bridge definition: 'unspecified witness family' denotes the structured intermediate object required by this branch. It is meant to organize the source-side content 'There exist many problems L in combinatorial optimisation/graph theory such that L is NP-complete under polynomial-time many-one reductions.'. It must preserve the invariant 'consequence schema NP-complete-classification'. Its operational role is 'Design a transfer mechanism that survives known local-reduction barriers, or prove that no such mechanism exists under explicit assumptions.'.
\item First nontrivial lemma: If 'Design a transfer mechanism that survives known local-reduction barriers, or prove that no such mechanism exists under explicit assumptions.' fails for 'unspecified witness family', then no route preserving 'consequence schema NP-complete-classification' can force 'Under"local"polynomial-time reductions (where each output bit is computable with very small time), MCSP is not NP-complete. More broadly, the paper develops evidence and conditional consequences suggesting NP-hardness of MCSP under various reductions would entail unexpected class separations.' from 'There exist many problems L in combinatorial optimisation/graph theory such that L is NP-complete under polynomial-time many-one reductions.'.
\item Quantifier signature: For every admissible 'unspecified witness family' on the route from 'st.thm.karp-many-natural-problems-are-np-complete-karp1' to 'st.thm.murray-williams-non-np-hardness-of-mcsp-under-lo'.
\item LaTeX theorem: \textbackslash\{\}forall W,\textbackslash\{\} \textbackslash\{\}mathrm\{Adm\}(W) \textbackslash\{\}to \textbackslash\{\}exists w : \textbackslash\{\}mathsf\{BW\}(W),\textbackslash\{\} \textbackslash\{\}mathrm\{PresInv\}(W,w) \textbackslash\{\}wedge \textbackslash\{\}mathrm\{RealizesTarget\}(W,w).
\item Negative/barrier version: Every attempted realization of 'unspecified witness family' remains blocked by the current barrier, so 'Under"local"polynomial-time reductions (where each output bit is computable with very small time), MCSP is not NP-complete. More broadly, the paper develops evidence and conditional consequences suggesting NP-hardness of MCSP under various reductions would entail unexpected class separations.' does not follow.
\item Proof strategy family: barrier
\item Negative route: Assume every attempted realization of 'unspecified witness family' collapses into the current barrier.
\item Failure diagnosis: The route fails if every admissible realization of 'unspecified witness family' remains trapped behind the current barrier.
\item Proof obligations: If 'Design a transfer mechanism that survives known local-reduction barriers, or prove that no such mechanism exists under explicit assumptions.' fails for 'unspecified witness family', then no route preserving 'consequence schema NP-complete-classification' can force 'Under"local"polynomial-time reductions (where each output bit is computable with very small time), MCSP is not NP-complete. More broadly, the paper develops evidence and conditional consequences suggesting NP-hardness of MCSP under various reductions would entail unexpected class separations.' from 'There exist many problems L in combinatorial optimisation/graph theory such that L is NP-complete under polynomial-time many-one reductions.'.
\item Candidate lemmas: Realization theorem for unspecified witness family; Barrier obstruction lemma for unspecified witness family
\item Independent value: A no-go theorem for 'unspecified witness family' would be independently valuable because it collapses an entire discovery subspace and exposes a genuine barrier theorem.
\item Expert gate: 6/6
\end{itemize}

\subsection*{conjpkg.disclayer-p-vs-np-family-barrier-aware-transfer.05}
\begin{itemize}
\item Exact statement: No admissible 'unspecified reduction gadget family' can simultaneously preserve 'consequence schema NP-complete-classification' and realize 'Under"local"polynomial-time reductions (where each output bit is computable with very small time), MCSP is not NP-complete. More broadly, the paper develops evidence and conditional consequences suggesting NP-hardness of MCSP under various reductions would entail unexpected class separations.' under the current route commitments.
\item Restricted-regime mode: no-go-theorem
\item Bridge object: unspecified reduction gadget family (family of reductions preserving the route witness structure)
\item Bridge definition: 'unspecified reduction gadget family' denotes the structured intermediate object required by this branch. It is meant to organize the source-side content 'There exist many problems L in combinatorial optimisation/graph theory such that L is NP-complete under polynomial-time many-one reductions.'. It must preserve the invariant 'consequence schema NP-complete-classification'. Its operational role is 'Design a transfer mechanism that survives known local-reduction barriers, or prove that no such mechanism exists under explicit assumptions.'.
\item First nontrivial lemma: If 'Design a transfer mechanism that survives known local-reduction barriers, or prove that no such mechanism exists under explicit assumptions.' fails for 'unspecified reduction gadget family', then no route preserving 'consequence schema NP-complete-classification' can force 'Under"local"polynomial-time reductions (where each output bit is computable with very small time), MCSP is not NP-complete. More broadly, the paper develops evidence and conditional consequences suggesting NP-hardness of MCSP under various reductions would entail unexpected class separations.' from 'There exist many problems L in combinatorial optimisation/graph theory such that L is NP-complete under polynomial-time many-one reductions.'.
\item Quantifier signature: For every admissible 'unspecified reduction gadget family' on the route from 'st.thm.karp-many-natural-problems-are-np-complete-karp1' to 'st.thm.murray-williams-non-np-hardness-of-mcsp-under-lo'.
\item LaTeX theorem: \textbackslash\{\}forall W,\textbackslash\{\} \textbackslash\{\}mathrm\{Adm\}(W) \textbackslash\{\}to \textbackslash\{\}neg \textbackslash\{\}exists w : \textbackslash\{\}mathsf\{BW\}(W),\textbackslash\{\} \textbackslash\{\}mathrm\{PresInv\}(W,w) \textbackslash\{\}wedge \textbackslash\{\}mathrm\{RealizesTarget\}(W,w).
\item Negative/barrier version: Every attempted realization of 'unspecified reduction gadget family' remains blocked by the current barrier, so 'Under"local"polynomial-time reductions (where each output bit is computable with very small time), MCSP is not NP-complete. More broadly, the paper develops evidence and conditional consequences suggesting NP-hardness of MCSP under various reductions would entail unexpected class separations.' does not follow.
\item Proof strategy family: barrier
\item Negative route: Assume every attempted realization of 'unspecified reduction gadget family' collapses into the current barrier.
\item Failure diagnosis: The route fails if every admissible realization of 'unspecified reduction gadget family' remains trapped behind the current barrier.
\item Proof obligations: If 'Design a transfer mechanism that survives known local-reduction barriers, or prove that no such mechanism exists under explicit assumptions.' fails for 'unspecified reduction gadget family', then no route preserving 'consequence schema NP-complete-classification' can force 'Under"local"polynomial-time reductions (where each output bit is computable with very small time), MCSP is not NP-complete. More broadly, the paper develops evidence and conditional consequences suggesting NP-hardness of MCSP under various reductions would entail unexpected class separations.' from 'There exist many problems L in combinatorial optimisation/graph theory such that L is NP-complete under polynomial-time many-one reductions.'.
\item Candidate lemmas: Obstruction theorem for unspecified reduction gadget family; Barrier obstruction lemma for unspecified reduction gadget family
\item Independent value: A no-go theorem for 'unspecified reduction gadget family' would be independently valuable because it collapses an entire discovery subspace and exposes a genuine barrier theorem.
\item Expert gate: 6/6
\end{itemize}

\subsection*{conjpkg.disclayer-p-vs-np-family-barrier-aware-transfer.06}
\begin{itemize}
\item Exact statement: There exists an admissible 'unspecified reduction gadget family' whose realization internalizes 'There exist many problems L in combinatorial optimisation/graph theory such that L is NP-complete under polynomial-time many-one reductions.' while preserving 'consequence schema NP-complete-classification' strongly enough to support 'Under"local"polynomial-time reductions (where each output bit is computable with very small time), MCSP is not NP-complete. More broadly, the paper develops evidence and conditional consequences suggesting NP-hardness of MCSP under various reductions would entail unexpected class separations.'.
\item Restricted-regime mode: no-go-theorem
\item Bridge object: unspecified reduction gadget family (family of reductions preserving the route witness structure)
\item Bridge definition: 'unspecified reduction gadget family' denotes the structured intermediate object required by this branch. It is meant to organize the source-side content 'There exist many problems L in combinatorial optimisation/graph theory such that L is NP-complete under polynomial-time many-one reductions.'. It must preserve the invariant 'consequence schema NP-complete-classification'. Its operational role is 'Design a transfer mechanism that survives known local-reduction barriers, or prove that no such mechanism exists under explicit assumptions.'.
\item First nontrivial lemma: If 'Design a transfer mechanism that survives known local-reduction barriers, or prove that no such mechanism exists under explicit assumptions.' fails for 'unspecified reduction gadget family', then no route preserving 'consequence schema NP-complete-classification' can force 'Under"local"polynomial-time reductions (where each output bit is computable with very small time), MCSP is not NP-complete. More broadly, the paper develops evidence and conditional consequences suggesting NP-hardness of MCSP under various reductions would entail unexpected class separations.' from 'There exist many problems L in combinatorial optimisation/graph theory such that L is NP-complete under polynomial-time many-one reductions.'.
\item Quantifier signature: For every admissible 'unspecified reduction gadget family' on the route from 'st.thm.karp-many-natural-problems-are-np-complete-karp1' to 'st.thm.murray-williams-non-np-hardness-of-mcsp-under-lo'.
\item LaTeX theorem: \textbackslash\{\}forall W,\textbackslash\{\} \textbackslash\{\}mathrm\{Adm\}(W) \textbackslash\{\}to \textbackslash\{\}exists w : \textbackslash\{\}mathsf\{BW\}(W),\textbackslash\{\} \textbackslash\{\}mathrm\{PresInv\}(W,w) \textbackslash\{\}wedge \textbackslash\{\}mathrm\{RealizesTarget\}(W,w).
\item Negative/barrier version: Every attempted realization of 'unspecified reduction gadget family' remains blocked by the current barrier, so 'Under"local"polynomial-time reductions (where each output bit is computable with very small time), MCSP is not NP-complete. More broadly, the paper develops evidence and conditional consequences suggesting NP-hardness of MCSP under various reductions would entail unexpected class separations.' does not follow.
\item Proof strategy family: barrier
\item Negative route: Assume every attempted realization of 'unspecified reduction gadget family' collapses into the current barrier.
\item Failure diagnosis: The route fails if every admissible realization of 'unspecified reduction gadget family' remains trapped behind the current barrier.
\item Proof obligations: If 'Design a transfer mechanism that survives known local-reduction barriers, or prove that no such mechanism exists under explicit assumptions.' fails for 'unspecified reduction gadget family', then no route preserving 'consequence schema NP-complete-classification' can force 'Under"local"polynomial-time reductions (where each output bit is computable with very small time), MCSP is not NP-complete. More broadly, the paper develops evidence and conditional consequences suggesting NP-hardness of MCSP under various reductions would entail unexpected class separations.' from 'There exist many problems L in combinatorial optimisation/graph theory such that L is NP-complete under polynomial-time many-one reductions.'.
\item Candidate lemmas: Realization theorem for unspecified reduction gadget family; Barrier obstruction lemma for unspecified reduction gadget family
\item Independent value: A no-go theorem for 'unspecified reduction gadget family' would be independently valuable because it collapses an entire discovery subspace and exposes a genuine barrier theorem.
\item Expert gate: 6/6
\end{itemize}

\subsection*{conjpkg.disclayer-p-vs-np-pattern-hybridization-of-partial-bridges.01}
\begin{itemize}
\item Exact statement: There exists an admissible 'MCSP witness family' such that 'A hardness magnification theorem has the form:"A barely-nontrivial lower bound for an explicit problem Q in a weak model implies a major lower bound/separation in a strong model."Typically, Q is a meta-problem (e.g. a Gap-MCSP variant) and"major"means consequences such as NP not in P/poly.' can be organized into a route preserving 'MCSP complexity profile in nonuniform-circuits' and thereby supports 'In refined regimes, lower bounds for Gap-MCSP slightly beyond known frontiers (e.g. N\textasciicircum{} 1+varepsilon -type circuit lower bounds for suitable promise gaps and input length N=2\textasciicircum{}n ) imply major consequences such as NP not in P/poly.'.
\item Restricted-regime mode: special-case
\item Bridge object: MCSP witness family (family of instances equipped with a monotone witness measure)
\item Bridge definition: 'MCSP witness family' denotes the structured intermediate object required by this branch. It is meant to organize the source-side content 'A hardness magnification theorem has the form:"A barely-nontrivial lower bound for an explicit problem Q in a weak model implies a major lower bound/separation in a strong model."Typically, Q is a meta-problem (e.g. a Gap-MCSP variant) and"major"means consequences such as NP not in P/poly.'. It must preserve the invariant 'MCSP complexity profile in nonuniform-circuits'. Its operational role is 'Some blocked routes sit at the intersection of multiple breakthrough families, which is evidence that one bridge object may be insufficient and a composed route should be explored.'.
\item First nontrivial lemma: For every admissible 'MCSP witness family', the first route-specific realization step preserves 'MCSP complexity profile in nonuniform-circuits' and prepares the passage from 'A hardness magnification theorem has the form:"A barely-nontrivial lower bound for an explicit problem Q in a weak model implies a major lower bound/separation in a strong model."Typically, Q is a meta-problem (e.g. a Gap-MCSP variant) and"major"means consequences such as NP not in P/poly.' to 'In refined regimes, lower bounds for Gap-MCSP slightly beyond known frontiers (e.g. N\textasciicircum{} 1+varepsilon -type circuit lower bounds for suitable promise gaps and input length N=2\textasciicircum{}n ) imply major consequences such as NP not in P/poly.'.
\item Quantifier signature: For every admissible 'MCSP witness family' on the route from 'st.def.hardness-magnification-phenomenon-programme-leve' to 'st.thm.oliveira-pich-santhanam-magnification-near-state'.
\item LaTeX theorem: \textbackslash\{\}forall W,\textbackslash\{\} \textbackslash\{\}mathrm\{Adm\}(W) \textbackslash\{\}to \textbackslash\{\}exists w : \textbackslash\{\}mathsf\{BW\}(W),\textbackslash\{\} \textbackslash\{\}mathrm\{PresInv\}(W,w) \textbackslash\{\}wedge \textbackslash\{\}mathrm\{RealizesTarget\}(W,w).
\item Negative/barrier version: If 'MCSP witness family' cannot be realized with invariant 'MCSP complexity profile in nonuniform-circuits', the route cannot force 'In refined regimes, lower bounds for Gap-MCSP slightly beyond known frontiers (e.g. N\textasciicircum{} 1+varepsilon -type circuit lower bounds for suitable promise gaps and input length N=2\textasciicircum{}n ) imply major consequences such as NP not in P/poly.'.
\item Proof strategy family: constructive-realization
\item Negative route: Assume 'MCSP witness family' cannot be realized while preserving 'MCSP complexity profile in nonuniform-circuits' strongly enough to activate the local target claim.
\item Failure diagnosis: The route fails if 'MCSP witness family' cannot be realized while preserving 'MCSP complexity profile in nonuniform-circuits'.
\item Proof obligations: For every admissible 'MCSP witness family', the first route-specific realization step preserves 'MCSP complexity profile in nonuniform-circuits' and prepares the passage from 'A hardness magnification theorem has the form:"A barely-nontrivial lower bound for an explicit problem Q in a weak model implies a major lower bound/separation in a strong model."Typically, Q is a meta-problem (e.g. a Gap-MCSP variant) and"major"means consequences such as NP not in P/poly.' to 'In refined regimes, lower bounds for Gap-MCSP slightly beyond known frontiers (e.g. N\textasciicircum{} 1+varepsilon -type circuit lower bounds for suitable promise gaps and input length N=2\textasciicircum{}n ) imply major consequences such as NP not in P/poly.'.
\item Candidate lemmas: Prove that a hardness-magnification witness family turns 'A barely-nontrivial lower bound for an explicit problem Q in a weak model' into 'major consequences such as NP not in P/poly' while preserving 'MCSP complexity profile in nonuniform-circuits'.
\item Independent value: A concrete special-case theorem for 'MCSP witness family' would be independently valuable because it converts a broad discovery programme into an inspectable mathematical subresult.
\item Expert gate: 6/6
\end{itemize}

\subsection*{conjpkg.disclayer-p-vs-np-pattern-hybridization-of-partial-bridges.02}
\begin{itemize}
\item Exact statement: There exists an admissible 'MCSP witness family' such that 'A hardness magnification theorem has the form:"A barely-nontrivial lower bound for an explicit problem Q in a weak model implies a major lower bound/separation in a strong model."Typically, Q is a meta-problem (e.g. a Gap-MCSP variant) and"major"means consequences such as NP not in P/poly.' can be organized into a route preserving 'MCSP complexity profile in nonuniform-circuits' and thereby supports 'There exist explicit"meta"problems (including MCSP variants) such that proving certain emph barely-superlinear lower bounds against restricted models (e.g. formulas) would imply major lower bounds for NP in much stronger models (including superpolynomial formula/circuit lower bounds, and in some parameter regimes consequences as strong as NP not in P/poly ).'.
\item Restricted-regime mode: special-case
\item Bridge object: MCSP witness family (family of instances equipped with a monotone witness measure)
\item Bridge definition: 'MCSP witness family' denotes the structured intermediate object required by this branch. It is meant to organize the source-side content 'A hardness magnification theorem has the form:"A barely-nontrivial lower bound for an explicit problem Q in a weak model implies a major lower bound/separation in a strong model."Typically, Q is a meta-problem (e.g. a Gap-MCSP variant) and"major"means consequences such as NP not in P/poly.'. It must preserve the invariant 'MCSP complexity profile in nonuniform-circuits'. Its operational role is 'Some blocked routes sit at the intersection of multiple breakthrough families, which is evidence that one bridge object may be insufficient and a composed route should be explored.'.
\item First nontrivial lemma: For every admissible 'MCSP witness family', the first route-specific realization step preserves 'MCSP complexity profile in nonuniform-circuits' and prepares the passage from 'A hardness magnification theorem has the form:"A barely-nontrivial lower bound for an explicit problem Q in a weak model implies a major lower bound/separation in a strong model."Typically, Q is a meta-problem (e.g. a Gap-MCSP variant) and"major"means consequences such as NP not in P/poly.' to 'There exist explicit"meta"problems (including MCSP variants) such that proving certain emph barely-superlinear lower bounds against restricted models (e.g. formulas) would imply major lower bounds for NP in much stronger models (including superpolynomial formula/circuit lower bounds, and in some parameter regimes consequences as strong as NP not in P/poly ).'.
\item Quantifier signature: For every admissible 'MCSP witness family' on the route from 'st.def.hardness-magnification-phenomenon-programme-leve' to 'st.thm.oliveira-santhanam-hardness-magnification-for-na'.
\item LaTeX theorem: \textbackslash\{\}forall W,\textbackslash\{\} \textbackslash\{\}mathrm\{Adm\}(W) \textbackslash\{\}to \textbackslash\{\}exists w : \textbackslash\{\}mathsf\{BW\}(W),\textbackslash\{\} \textbackslash\{\}mathrm\{PresInv\}(W,w) \textbackslash\{\}wedge \textbackslash\{\}mathrm\{RealizesTarget\}(W,w).
\item Negative/barrier version: If 'MCSP witness family' cannot be realized with invariant 'MCSP complexity profile in nonuniform-circuits', the route cannot force 'There exist explicit"meta"problems (including MCSP variants) such that proving certain emph barely-superlinear lower bounds against restricted models (e.g. formulas) would imply major lower bounds for NP in much stronger models (including superpolynomial formula/circuit lower bounds, and in some parameter regimes consequences as strong as NP not in P/poly ).'.
\item Proof strategy family: constructive-realization
\item Negative route: Assume 'MCSP witness family' cannot be realized while preserving 'MCSP complexity profile in nonuniform-circuits' strongly enough to activate the local target claim.
\item Failure diagnosis: The route fails if 'MCSP witness family' cannot be realized while preserving 'MCSP complexity profile in nonuniform-circuits'.
\item Proof obligations: For every admissible 'MCSP witness family', the first route-specific realization step preserves 'MCSP complexity profile in nonuniform-circuits' and prepares the passage from 'A hardness magnification theorem has the form:"A barely-nontrivial lower bound for an explicit problem Q in a weak model implies a major lower bound/separation in a strong model."Typically, Q is a meta-problem (e.g. a Gap-MCSP variant) and"major"means consequences such as NP not in P/poly.' to 'There exist explicit"meta"problems (including MCSP variants) such that proving certain emph barely-superlinear lower bounds against restricted models (e.g. formulas) would imply major lower bounds for NP in much stronger models (including superpolynomial formula/circuit lower bounds, and in some parameter regimes consequences as strong as NP not in P/poly ).'.
\item Candidate lemmas: Prove that a hardness-magnification witness family preserves 'MCSP complexity profile in nonuniform-circuits' while transporting the route from the nonuniform-circuits regime to the formula-size regime strongly enough to derive 'major lower bounds for NP in much stronger models (including superpolynomial formula/circuit lower bounds, and in some parameter...'.
\item Independent value: A concrete special-case theorem for 'MCSP witness family' would be independently valuable because it converts a broad discovery programme into an inspectable mathematical subresult.
\item Expert gate: 6/6
\end{itemize}

\subsection*{conjpkg.disclayer-p-vs-np-pattern-hybridization-of-partial-bridges.03}
\begin{itemize}
\item Exact statement: There exists an admissible 'MCSP obstruction family' such that 'A hardness magnification theorem has the form:"A barely-nontrivial lower bound for an explicit problem Q in a weak model implies a major lower bound/separation in a strong model."Typically, Q is a meta-problem (e.g. a Gap-MCSP variant) and"major"means consequences such as NP not in P/poly.' can be organized into a route preserving 'MCSP complexity profile in nonuniform-circuits' and thereby supports 'In refined regimes, lower bounds for Gap-MCSP slightly beyond known frontiers (e.g. N\textasciicircum{} 1+varepsilon -type circuit lower bounds for suitable promise gaps and input length N=2\textasciicircum{}n ) imply major consequences such as NP not in P/poly.'.
\item Restricted-regime mode: special-case
\item Bridge object: MCSP obstruction family (family of hard instances or obstructions forcing the target payoff)
\item Bridge definition: 'MCSP obstruction family' denotes the structured intermediate object required by this branch. It is meant to organize the source-side content 'A hardness magnification theorem has the form:"A barely-nontrivial lower bound for an explicit problem Q in a weak model implies a major lower bound/separation in a strong model."Typically, Q is a meta-problem (e.g. a Gap-MCSP variant) and"major"means consequences such as NP not in P/poly.'. It must preserve the invariant 'MCSP complexity profile in nonuniform-circuits'. Its operational role is 'Some blocked routes sit at the intersection of multiple breakthrough families, which is evidence that one bridge object may be insufficient and a composed route should be explored.'.
\item First nontrivial lemma: For every admissible 'MCSP obstruction family', the first route-specific realization step preserves 'MCSP complexity profile in nonuniform-circuits' and prepares the passage from 'A hardness magnification theorem has the form:"A barely-nontrivial lower bound for an explicit problem Q in a weak model implies a major lower bound/separation in a strong model."Typically, Q is a meta-problem (e.g. a Gap-MCSP variant) and"major"means consequences such as NP not in P/poly.' to 'In refined regimes, lower bounds for Gap-MCSP slightly beyond known frontiers (e.g. N\textasciicircum{} 1+varepsilon -type circuit lower bounds for suitable promise gaps and input length N=2\textasciicircum{}n ) imply major consequences such as NP not in P/poly.'.
\item Quantifier signature: For every admissible 'MCSP obstruction family' on the route from 'st.def.hardness-magnification-phenomenon-programme-leve' to 'st.thm.oliveira-pich-santhanam-magnification-near-state'.
\item LaTeX theorem: \textbackslash\{\}forall W,\textbackslash\{\} \textbackslash\{\}mathrm\{Adm\}(W) \textbackslash\{\}to \textbackslash\{\}exists w : \textbackslash\{\}mathsf\{BW\}(W),\textbackslash\{\} \textbackslash\{\}mathrm\{PresInv\}(W,w) \textbackslash\{\}wedge \textbackslash\{\}mathrm\{RealizesTarget\}(W,w).
\item Negative/barrier version: If 'MCSP obstruction family' cannot be realized with invariant 'MCSP complexity profile in nonuniform-circuits', the route cannot force 'In refined regimes, lower bounds for Gap-MCSP slightly beyond known frontiers (e.g. N\textasciicircum{} 1+varepsilon -type circuit lower bounds for suitable promise gaps and input length N=2\textasciicircum{}n ) imply major consequences such as NP not in P/poly.'.
\item Proof strategy family: constructive-realization
\item Negative route: Assume 'MCSP obstruction family' cannot be realized while preserving 'MCSP complexity profile in nonuniform-circuits' strongly enough to activate the local target claim.
\item Failure diagnosis: The route fails if 'MCSP obstruction family' cannot be realized while preserving 'MCSP complexity profile in nonuniform-circuits'.
\item Proof obligations: For every admissible 'MCSP obstruction family', the first route-specific realization step preserves 'MCSP complexity profile in nonuniform-circuits' and prepares the passage from 'A hardness magnification theorem has the form:"A barely-nontrivial lower bound for an explicit problem Q in a weak model implies a major lower bound/separation in a strong model."Typically, Q is a meta-problem (e.g. a Gap-MCSP variant) and"major"means consequences such as NP not in P/poly.' to 'In refined regimes, lower bounds for Gap-MCSP slightly beyond known frontiers (e.g. N\textasciicircum{} 1+varepsilon -type circuit lower bounds for suitable promise gaps and input length N=2\textasciicircum{}n ) imply major consequences such as NP not in P/poly.'.
\item Candidate lemmas: Prove that a hardness-magnification witness family turns 'A barely-nontrivial lower bound for an explicit problem Q in a weak model' into 'major consequences such as NP not in P/poly' while preserving 'MCSP complexity profile in nonuniform-circuits'.
\item Independent value: A concrete special-case theorem for 'MCSP obstruction family' would be independently valuable because it converts a broad discovery programme into an inspectable mathematical subresult.
\item Expert gate: 6/6
\end{itemize}

\subsection*{conjpkg.disclayer-p-vs-np-pattern-hybridization-of-partial-bridges.04}
\begin{itemize}
\item Exact statement: No admissible 'MCSP witness family' can simultaneously preserve 'MCSP complexity profile in nonuniform-circuits' and realize 'In refined regimes, lower bounds for Gap-MCSP slightly beyond known frontiers (e.g. N\textasciicircum{} 1+varepsilon -type circuit lower bounds for suitable promise gaps and input length N=2\textasciicircum{}n ) imply major consequences such as NP not in P/poly.' under the current route commitments.
\item Restricted-regime mode: special-case
\item Bridge object: MCSP witness family (family of instances equipped with a monotone witness measure)
\item Bridge definition: 'MCSP witness family' denotes the structured intermediate object required by this branch. It is meant to organize the source-side content 'A hardness magnification theorem has the form:"A barely-nontrivial lower bound for an explicit problem Q in a weak model implies a major lower bound/separation in a strong model."Typically, Q is a meta-problem (e.g. a Gap-MCSP variant) and"major"means consequences such as NP not in P/poly.'. It must preserve the invariant 'MCSP complexity profile in nonuniform-circuits'. Its operational role is 'Some blocked routes sit at the intersection of multiple breakthrough families, which is evidence that one bridge object may be insufficient and a composed route should be explored.'.
\item First nontrivial lemma: For every admissible 'MCSP witness family', the first route-specific realization step preserves 'MCSP complexity profile in nonuniform-circuits' and prepares the passage from 'A hardness magnification theorem has the form:"A barely-nontrivial lower bound for an explicit problem Q in a weak model implies a major lower bound/separation in a strong model."Typically, Q is a meta-problem (e.g. a Gap-MCSP variant) and"major"means consequences such as NP not in P/poly.' to 'In refined regimes, lower bounds for Gap-MCSP slightly beyond known frontiers (e.g. N\textasciicircum{} 1+varepsilon -type circuit lower bounds for suitable promise gaps and input length N=2\textasciicircum{}n ) imply major consequences such as NP not in P/poly.'.
\item Quantifier signature: For every admissible 'MCSP witness family' on the route from 'st.def.hardness-magnification-phenomenon-programme-leve' to 'st.thm.oliveira-pich-santhanam-magnification-near-state'.
\item LaTeX theorem: \textbackslash\{\}forall W,\textbackslash\{\} \textbackslash\{\}mathrm\{Adm\}(W) \textbackslash\{\}to \textbackslash\{\}neg \textbackslash\{\}exists w : \textbackslash\{\}mathsf\{BW\}(W),\textbackslash\{\} \textbackslash\{\}mathrm\{PresInv\}(W,w) \textbackslash\{\}wedge \textbackslash\{\}mathrm\{RealizesTarget\}(W,w).
\item Negative/barrier version: If 'MCSP witness family' cannot be realized with invariant 'MCSP complexity profile in nonuniform-circuits', the route cannot force 'In refined regimes, lower bounds for Gap-MCSP slightly beyond known frontiers (e.g. N\textasciicircum{} 1+varepsilon -type circuit lower bounds for suitable promise gaps and input length N=2\textasciicircum{}n ) imply major consequences such as NP not in P/poly.'.
\item Proof strategy family: constructive-realization
\item Negative route: Assume 'MCSP witness family' cannot be realized while preserving 'MCSP complexity profile in nonuniform-circuits' strongly enough to activate the local target claim.
\item Failure diagnosis: The route fails if 'MCSP witness family' cannot be realized while preserving 'MCSP complexity profile in nonuniform-circuits'.
\item Proof obligations: For every admissible 'MCSP witness family', the first route-specific realization step preserves 'MCSP complexity profile in nonuniform-circuits' and prepares the passage from 'A hardness magnification theorem has the form:"A barely-nontrivial lower bound for an explicit problem Q in a weak model implies a major lower bound/separation in a strong model."Typically, Q is a meta-problem (e.g. a Gap-MCSP variant) and"major"means consequences such as NP not in P/poly.' to 'In refined regimes, lower bounds for Gap-MCSP slightly beyond known frontiers (e.g. N\textasciicircum{} 1+varepsilon -type circuit lower bounds for suitable promise gaps and input length N=2\textasciicircum{}n ) imply major consequences such as NP not in P/poly.'.
\item Candidate lemmas: Obstruction theorem for MCSP witness family
\item Independent value: A concrete special-case theorem for 'MCSP witness family' would be independently valuable because it converts a broad discovery programme into an inspectable mathematical subresult.
\item Expert gate: 6/6
\end{itemize}

\subsection*{conjpkg.disclayer-p-vs-np-pattern-hybridization-of-partial-bridges.05}
\begin{itemize}
\item Exact statement: There exists an admissible 'MCSP witness family' whose realization internalizes 'A hardness magnification theorem has the form:"A barely-nontrivial lower bound for an explicit problem Q in a weak model implies a major lower bound/separation in a strong model."Typically, Q is a meta-problem (e.g. a Gap-MCSP variant) and"major"means consequences such as NP not in P/poly.' while preserving 'MCSP complexity profile in nonuniform-circuits' strongly enough to support 'In refined regimes, lower bounds for Gap-MCSP slightly beyond known frontiers (e.g. N\textasciicircum{} 1+varepsilon -type circuit lower bounds for suitable promise gaps and input length N=2\textasciicircum{}n ) imply major consequences such as NP not in P/poly.'.
\item Restricted-regime mode: special-case
\item Bridge object: MCSP witness family (family of instances equipped with a monotone witness measure)
\item Bridge definition: 'MCSP witness family' denotes the structured intermediate object required by this branch. It is meant to organize the source-side content 'A hardness magnification theorem has the form:"A barely-nontrivial lower bound for an explicit problem Q in a weak model implies a major lower bound/separation in a strong model."Typically, Q is a meta-problem (e.g. a Gap-MCSP variant) and"major"means consequences such as NP not in P/poly.'. It must preserve the invariant 'MCSP complexity profile in nonuniform-circuits'. Its operational role is 'Some blocked routes sit at the intersection of multiple breakthrough families, which is evidence that one bridge object may be insufficient and a composed route should be explored.'.
\item First nontrivial lemma: For every admissible 'MCSP witness family', the first route-specific realization step preserves 'MCSP complexity profile in nonuniform-circuits' and prepares the passage from 'A hardness magnification theorem has the form:"A barely-nontrivial lower bound for an explicit problem Q in a weak model implies a major lower bound/separation in a strong model."Typically, Q is a meta-problem (e.g. a Gap-MCSP variant) and"major"means consequences such as NP not in P/poly.' to 'In refined regimes, lower bounds for Gap-MCSP slightly beyond known frontiers (e.g. N\textasciicircum{} 1+varepsilon -type circuit lower bounds for suitable promise gaps and input length N=2\textasciicircum{}n ) imply major consequences such as NP not in P/poly.'.
\item Quantifier signature: For every admissible 'MCSP witness family' on the route from 'st.def.hardness-magnification-phenomenon-programme-leve' to 'st.thm.oliveira-pich-santhanam-magnification-near-state'.
\item LaTeX theorem: \textbackslash\{\}forall W,\textbackslash\{\} \textbackslash\{\}mathrm\{Adm\}(W) \textbackslash\{\}to \textbackslash\{\}exists w : \textbackslash\{\}mathsf\{BW\}(W),\textbackslash\{\} \textbackslash\{\}mathrm\{PresInv\}(W,w) \textbackslash\{\}wedge \textbackslash\{\}mathrm\{RealizesTarget\}(W,w).
\item Negative/barrier version: If 'MCSP witness family' cannot be realized with invariant 'MCSP complexity profile in nonuniform-circuits', the route cannot force 'In refined regimes, lower bounds for Gap-MCSP slightly beyond known frontiers (e.g. N\textasciicircum{} 1+varepsilon -type circuit lower bounds for suitable promise gaps and input length N=2\textasciicircum{}n ) imply major consequences such as NP not in P/poly.'.
\item Proof strategy family: constructive-realization
\item Negative route: Assume 'MCSP witness family' cannot be realized while preserving 'MCSP complexity profile in nonuniform-circuits' strongly enough to activate the local target claim.
\item Failure diagnosis: The route fails if 'MCSP witness family' cannot be realized while preserving 'MCSP complexity profile in nonuniform-circuits'.
\item Proof obligations: For every admissible 'MCSP witness family', the first route-specific realization step preserves 'MCSP complexity profile in nonuniform-circuits' and prepares the passage from 'A hardness magnification theorem has the form:"A barely-nontrivial lower bound for an explicit problem Q in a weak model implies a major lower bound/separation in a strong model."Typically, Q is a meta-problem (e.g. a Gap-MCSP variant) and"major"means consequences such as NP not in P/poly.' to 'In refined regimes, lower bounds for Gap-MCSP slightly beyond known frontiers (e.g. N\textasciicircum{} 1+varepsilon -type circuit lower bounds for suitable promise gaps and input length N=2\textasciicircum{}n ) imply major consequences such as NP not in P/poly.'.
\item Candidate lemmas: Realization theorem for MCSP witness family
\item Independent value: A concrete special-case theorem for 'MCSP witness family' would be independently valuable because it converts a broad discovery programme into an inspectable mathematical subresult.
\item Expert gate: 6/6
\end{itemize}

\subsection*{conjpkg.disclayer-p-vs-np-pattern-hybridization-of-partial-bridges.06}
\begin{itemize}
\item Exact statement: There exists an intermediate typed object built from 'MCSP witness family' such that the route from 'A hardness magnification theorem has the form:"A barely-nontrivial lower bound for an explicit problem Q in a weak model implies a major lower bound/separation in a strong model."Typically, Q is a meta-problem (e.g. a Gap-MCSP variant) and"major"means consequences such as NP not in P/poly.' to 'In refined regimes, lower bounds for Gap-MCSP slightly beyond known frontiers (e.g. N\textasciicircum{} 1+varepsilon -type circuit lower bounds for suitable promise gaps and input length N=2\textasciicircum{}n ) imply major consequences such as NP not in P/poly.' factors through theorem-sized maps preserving 'MCSP complexity profile in nonuniform-circuits'.
\item Restricted-regime mode: special-case
\item Bridge object: MCSP witness family (family of instances equipped with a monotone witness measure)
\item Bridge definition: 'MCSP witness family' denotes the structured intermediate object required by this branch. It is meant to organize the source-side content 'A hardness magnification theorem has the form:"A barely-nontrivial lower bound for an explicit problem Q in a weak model implies a major lower bound/separation in a strong model."Typically, Q is a meta-problem (e.g. a Gap-MCSP variant) and"major"means consequences such as NP not in P/poly.'. It must preserve the invariant 'MCSP complexity profile in nonuniform-circuits'. Its operational role is 'Some blocked routes sit at the intersection of multiple breakthrough families, which is evidence that one bridge object may be insufficient and a composed route should be explored.'.
\item First nontrivial lemma: For every admissible 'MCSP witness family', the first route-specific realization step preserves 'MCSP complexity profile in nonuniform-circuits' and prepares the passage from 'A hardness magnification theorem has the form:"A barely-nontrivial lower bound for an explicit problem Q in a weak model implies a major lower bound/separation in a strong model."Typically, Q is a meta-problem (e.g. a Gap-MCSP variant) and"major"means consequences such as NP not in P/poly.' to 'In refined regimes, lower bounds for Gap-MCSP slightly beyond known frontiers (e.g. N\textasciicircum{} 1+varepsilon -type circuit lower bounds for suitable promise gaps and input length N=2\textasciicircum{}n ) imply major consequences such as NP not in P/poly.'.
\item Quantifier signature: For every admissible 'MCSP witness family' on the route from 'st.def.hardness-magnification-phenomenon-programme-leve' to 'st.thm.oliveira-pich-santhanam-magnification-near-state'.
\item LaTeX theorem: \textbackslash\{\}forall W,\textbackslash\{\} \textbackslash\{\}mathrm\{Adm\}(W) \textbackslash\{\}to \textbackslash\{\}exists w : \textbackslash\{\}mathsf\{BW\}(W),\textbackslash\{\} \textbackslash\{\}mathrm\{PresInv\}(W,w) \textbackslash\{\}wedge \textbackslash\{\}mathrm\{RealizesTarget\}(W,w).
\item Negative/barrier version: If 'MCSP witness family' cannot be realized with invariant 'MCSP complexity profile in nonuniform-circuits', the route cannot force 'In refined regimes, lower bounds for Gap-MCSP slightly beyond known frontiers (e.g. N\textasciicircum{} 1+varepsilon -type circuit lower bounds for suitable promise gaps and input length N=2\textasciicircum{}n ) imply major consequences such as NP not in P/poly.'.
\item Proof strategy family: constructive-realization
\item Negative route: Assume 'MCSP witness family' cannot be realized while preserving 'MCSP complexity profile in nonuniform-circuits' strongly enough to activate the local target claim.
\item Failure diagnosis: The route fails if 'MCSP witness family' cannot be realized while preserving 'MCSP complexity profile in nonuniform-circuits'.
\item Proof obligations: For every admissible 'MCSP witness family', the first route-specific realization step preserves 'MCSP complexity profile in nonuniform-circuits' and prepares the passage from 'A hardness magnification theorem has the form:"A barely-nontrivial lower bound for an explicit problem Q in a weak model implies a major lower bound/separation in a strong model."Typically, Q is a meta-problem (e.g. a Gap-MCSP variant) and"major"means consequences such as NP not in P/poly.' to 'In refined regimes, lower bounds for Gap-MCSP slightly beyond known frontiers (e.g. N\textasciicircum{} 1+varepsilon -type circuit lower bounds for suitable promise gaps and input length N=2\textasciicircum{}n ) imply major consequences such as NP not in P/poly.'.
\item Candidate lemmas: Factorization theorem for MCSP witness family
\item Independent value: A concrete special-case theorem for 'MCSP witness family' would be independently valuable because it converts a broad discovery programme into an inspectable mathematical subresult.
\item Expert gate: 6/6
\end{itemize}

\subsection*{conjpkg.disclayer-p-vs-np-pattern-hybridization-of-partial-bridges.07}
\begin{itemize}
\item Exact statement: For every admissible 'MCSP witness family', there is a transport law from 'nonuniform-circuits' to 'formula-size' that preserves 'MCSP complexity profile in nonuniform-circuits' and retains enough strength to activate 'There exist explicit"meta"problems (including MCSP variants) such that proving certain emph barely-superlinear lower bounds against restricted models (e.g. formulas) would imply major lower bounds for NP in much stronger models (including superpolynomial formula/circuit lower bounds, and in some parameter regimes consequences as strong as NP not in P/poly ).'.
\item Restricted-regime mode: special-case
\item Bridge object: MCSP witness family (family of instances equipped with a monotone witness measure)
\item Bridge definition: 'MCSP witness family' denotes the structured intermediate object required by this branch. It is meant to organize the source-side content 'A hardness magnification theorem has the form:"A barely-nontrivial lower bound for an explicit problem Q in a weak model implies a major lower bound/separation in a strong model."Typically, Q is a meta-problem (e.g. a Gap-MCSP variant) and"major"means consequences such as NP not in P/poly.'. It must preserve the invariant 'MCSP complexity profile in nonuniform-circuits'. Its operational role is 'Some blocked routes sit at the intersection of multiple breakthrough families, which is evidence that one bridge object may be insufficient and a composed route should be explored.'.
\item First nontrivial lemma: For every admissible 'MCSP witness family', the first route-specific realization step preserves 'MCSP complexity profile in nonuniform-circuits' and prepares the passage from 'A hardness magnification theorem has the form:"A barely-nontrivial lower bound for an explicit problem Q in a weak model implies a major lower bound/separation in a strong model."Typically, Q is a meta-problem (e.g. a Gap-MCSP variant) and"major"means consequences such as NP not in P/poly.' to 'There exist explicit"meta"problems (including MCSP variants) such that proving certain emph barely-superlinear lower bounds against restricted models (e.g. formulas) would imply major lower bounds for NP in much stronger models (including superpolynomial formula/circuit lower bounds, and in some parameter regimes consequences as strong as NP not in P/poly ).'.
\item Quantifier signature: For every admissible 'MCSP witness family' on the route from 'st.def.hardness-magnification-phenomenon-programme-leve' to 'st.thm.oliveira-santhanam-hardness-magnification-for-na'.
\item LaTeX theorem: \textbackslash\{\}forall W,\textbackslash\{\} \textbackslash\{\}mathrm\{Adm\}(W) \textbackslash\{\}to \textbackslash\{\}exists w : \textbackslash\{\}mathsf\{BW\}(W),\textbackslash\{\} \textbackslash\{\}mathrm\{PresInv\}(W,w) \textbackslash\{\}wedge \textbackslash\{\}mathrm\{RealizesTarget\}(W,w).
\item Negative/barrier version: If 'MCSP witness family' cannot be realized with invariant 'MCSP complexity profile in nonuniform-circuits', the route cannot force 'There exist explicit"meta"problems (including MCSP variants) such that proving certain emph barely-superlinear lower bounds against restricted models (e.g. formulas) would imply major lower bounds for NP in much stronger models (including superpolynomial formula/circuit lower bounds, and in some parameter regimes consequences as strong as NP not in P/poly ).'.
\item Proof strategy family: constructive-realization
\item Negative route: Assume 'MCSP witness family' cannot be realized while preserving 'MCSP complexity profile in nonuniform-circuits' strongly enough to activate the local target claim.
\item Failure diagnosis: The route fails if 'MCSP witness family' cannot be realized while preserving 'MCSP complexity profile in nonuniform-circuits'.
\item Proof obligations: For every admissible 'MCSP witness family', the first route-specific realization step preserves 'MCSP complexity profile in nonuniform-circuits' and prepares the passage from 'A hardness magnification theorem has the form:"A barely-nontrivial lower bound for an explicit problem Q in a weak model implies a major lower bound/separation in a strong model."Typically, Q is a meta-problem (e.g. a Gap-MCSP variant) and"major"means consequences such as NP not in P/poly.' to 'There exist explicit"meta"problems (including MCSP variants) such that proving certain emph barely-superlinear lower bounds against restricted models (e.g. formulas) would imply major lower bounds for NP in much stronger models (including superpolynomial formula/circuit lower bounds, and in some parameter regimes consequences as strong as NP not in P/poly ).'.
\item Candidate lemmas: Transport theorem for MCSP witness family
\item Independent value: A concrete special-case theorem for 'MCSP witness family' would be independently valuable because it converts a broad discovery programme into an inspectable mathematical subresult.
\item Expert gate: 6/6
\end{itemize}

\subsection*{conjpkg.disclayer-p-vs-np-pattern-hybridization-of-partial-bridges.08}
\begin{itemize}
\item Exact statement: No admissible 'MCSP witness family' can simultaneously preserve 'MCSP complexity profile in nonuniform-circuits' and realize 'There exist explicit"meta"problems (including MCSP variants) such that proving certain emph barely-superlinear lower bounds against restricted models (e.g. formulas) would imply major lower bounds for NP in much stronger models (including superpolynomial formula/circuit lower bounds, and in some parameter regimes consequences as strong as NP not in P/poly ).' under the current route commitments.
\item Restricted-regime mode: special-case
\item Bridge object: MCSP witness family (family of instances equipped with a monotone witness measure)
\item Bridge definition: 'MCSP witness family' denotes the structured intermediate object required by this branch. It is meant to organize the source-side content 'A hardness magnification theorem has the form:"A barely-nontrivial lower bound for an explicit problem Q in a weak model implies a major lower bound/separation in a strong model."Typically, Q is a meta-problem (e.g. a Gap-MCSP variant) and"major"means consequences such as NP not in P/poly.'. It must preserve the invariant 'MCSP complexity profile in nonuniform-circuits'. Its operational role is 'Some blocked routes sit at the intersection of multiple breakthrough families, which is evidence that one bridge object may be insufficient and a composed route should be explored.'.
\item First nontrivial lemma: For every admissible 'MCSP witness family', the first route-specific realization step preserves 'MCSP complexity profile in nonuniform-circuits' and prepares the passage from 'A hardness magnification theorem has the form:"A barely-nontrivial lower bound for an explicit problem Q in a weak model implies a major lower bound/separation in a strong model."Typically, Q is a meta-problem (e.g. a Gap-MCSP variant) and"major"means consequences such as NP not in P/poly.' to 'There exist explicit"meta"problems (including MCSP variants) such that proving certain emph barely-superlinear lower bounds against restricted models (e.g. formulas) would imply major lower bounds for NP in much stronger models (including superpolynomial formula/circuit lower bounds, and in some parameter regimes consequences as strong as NP not in P/poly ).'.
\item Quantifier signature: For every admissible 'MCSP witness family' on the route from 'st.def.hardness-magnification-phenomenon-programme-leve' to 'st.thm.oliveira-santhanam-hardness-magnification-for-na'.
\item LaTeX theorem: \textbackslash\{\}forall W,\textbackslash\{\} \textbackslash\{\}mathrm\{Adm\}(W) \textbackslash\{\}to \textbackslash\{\}neg \textbackslash\{\}exists w : \textbackslash\{\}mathsf\{BW\}(W),\textbackslash\{\} \textbackslash\{\}mathrm\{PresInv\}(W,w) \textbackslash\{\}wedge \textbackslash\{\}mathrm\{RealizesTarget\}(W,w).
\item Negative/barrier version: If 'MCSP witness family' cannot be realized with invariant 'MCSP complexity profile in nonuniform-circuits', the route cannot force 'There exist explicit"meta"problems (including MCSP variants) such that proving certain emph barely-superlinear lower bounds against restricted models (e.g. formulas) would imply major lower bounds for NP in much stronger models (including superpolynomial formula/circuit lower bounds, and in some parameter regimes consequences as strong as NP not in P/poly ).'.
\item Proof strategy family: constructive-realization
\item Negative route: Assume 'MCSP witness family' cannot be realized while preserving 'MCSP complexity profile in nonuniform-circuits' strongly enough to activate the local target claim.
\item Failure diagnosis: The route fails if 'MCSP witness family' cannot be realized while preserving 'MCSP complexity profile in nonuniform-circuits'.
\item Proof obligations: For every admissible 'MCSP witness family', the first route-specific realization step preserves 'MCSP complexity profile in nonuniform-circuits' and prepares the passage from 'A hardness magnification theorem has the form:"A barely-nontrivial lower bound for an explicit problem Q in a weak model implies a major lower bound/separation in a strong model."Typically, Q is a meta-problem (e.g. a Gap-MCSP variant) and"major"means consequences such as NP not in P/poly.' to 'There exist explicit"meta"problems (including MCSP variants) such that proving certain emph barely-superlinear lower bounds against restricted models (e.g. formulas) would imply major lower bounds for NP in much stronger models (including superpolynomial formula/circuit lower bounds, and in some parameter regimes consequences as strong as NP not in P/poly ).'.
\item Candidate lemmas: Obstruction theorem for MCSP witness family
\item Independent value: A concrete special-case theorem for 'MCSP witness family' would be independently valuable because it converts a broad discovery programme into an inspectable mathematical subresult.
\item Expert gate: 6/6
\end{itemize}

\subsection*{conjpkg.disclayer-p-vs-np-pattern-hybridization-of-partial-bridges.09}
\begin{itemize}
\item Exact statement: There exists an admissible 'MCSP witness family' whose realization internalizes 'A hardness magnification theorem has the form:"A barely-nontrivial lower bound for an explicit problem Q in a weak model implies a major lower bound/separation in a strong model."Typically, Q is a meta-problem (e.g. a Gap-MCSP variant) and"major"means consequences such as NP not in P/poly.' while preserving 'MCSP complexity profile in nonuniform-circuits' strongly enough to support 'There exist explicit"meta"problems (including MCSP variants) such that proving certain emph barely-superlinear lower bounds against restricted models (e.g. formulas) would imply major lower bounds for NP in much stronger models (including superpolynomial formula/circuit lower bounds, and in some parameter regimes consequences as strong as NP not in P/poly ).'.
\item Restricted-regime mode: special-case
\item Bridge object: MCSP witness family (family of instances equipped with a monotone witness measure)
\item Bridge definition: 'MCSP witness family' denotes the structured intermediate object required by this branch. It is meant to organize the source-side content 'A hardness magnification theorem has the form:"A barely-nontrivial lower bound for an explicit problem Q in a weak model implies a major lower bound/separation in a strong model."Typically, Q is a meta-problem (e.g. a Gap-MCSP variant) and"major"means consequences such as NP not in P/poly.'. It must preserve the invariant 'MCSP complexity profile in nonuniform-circuits'. Its operational role is 'Some blocked routes sit at the intersection of multiple breakthrough families, which is evidence that one bridge object may be insufficient and a composed route should be explored.'.
\item First nontrivial lemma: For every admissible 'MCSP witness family', the first route-specific realization step preserves 'MCSP complexity profile in nonuniform-circuits' and prepares the passage from 'A hardness magnification theorem has the form:"A barely-nontrivial lower bound for an explicit problem Q in a weak model implies a major lower bound/separation in a strong model."Typically, Q is a meta-problem (e.g. a Gap-MCSP variant) and"major"means consequences such as NP not in P/poly.' to 'There exist explicit"meta"problems (including MCSP variants) such that proving certain emph barely-superlinear lower bounds against restricted models (e.g. formulas) would imply major lower bounds for NP in much stronger models (including superpolynomial formula/circuit lower bounds, and in some parameter regimes consequences as strong as NP not in P/poly ).'.
\item Quantifier signature: For every admissible 'MCSP witness family' on the route from 'st.def.hardness-magnification-phenomenon-programme-leve' to 'st.thm.oliveira-santhanam-hardness-magnification-for-na'.
\item LaTeX theorem: \textbackslash\{\}forall W,\textbackslash\{\} \textbackslash\{\}mathrm\{Adm\}(W) \textbackslash\{\}to \textbackslash\{\}exists w : \textbackslash\{\}mathsf\{BW\}(W),\textbackslash\{\} \textbackslash\{\}mathrm\{PresInv\}(W,w) \textbackslash\{\}wedge \textbackslash\{\}mathrm\{RealizesTarget\}(W,w).
\item Negative/barrier version: If 'MCSP witness family' cannot be realized with invariant 'MCSP complexity profile in nonuniform-circuits', the route cannot force 'There exist explicit"meta"problems (including MCSP variants) such that proving certain emph barely-superlinear lower bounds against restricted models (e.g. formulas) would imply major lower bounds for NP in much stronger models (including superpolynomial formula/circuit lower bounds, and in some parameter regimes consequences as strong as NP not in P/poly ).'.
\item Proof strategy family: constructive-realization
\item Negative route: Assume 'MCSP witness family' cannot be realized while preserving 'MCSP complexity profile in nonuniform-circuits' strongly enough to activate the local target claim.
\item Failure diagnosis: The route fails if 'MCSP witness family' cannot be realized while preserving 'MCSP complexity profile in nonuniform-circuits'.
\item Proof obligations: For every admissible 'MCSP witness family', the first route-specific realization step preserves 'MCSP complexity profile in nonuniform-circuits' and prepares the passage from 'A hardness magnification theorem has the form:"A barely-nontrivial lower bound for an explicit problem Q in a weak model implies a major lower bound/separation in a strong model."Typically, Q is a meta-problem (e.g. a Gap-MCSP variant) and"major"means consequences such as NP not in P/poly.' to 'There exist explicit"meta"problems (including MCSP variants) such that proving certain emph barely-superlinear lower bounds against restricted models (e.g. formulas) would imply major lower bounds for NP in much stronger models (including superpolynomial formula/circuit lower bounds, and in some parameter regimes consequences as strong as NP not in P/poly ).'.
\item Candidate lemmas: Realization theorem for MCSP witness family
\item Independent value: A concrete special-case theorem for 'MCSP witness family' would be independently valuable because it converts a broad discovery programme into an inspectable mathematical subresult.
\item Expert gate: 6/6
\end{itemize}

\subsection*{conjpkg.disclayer-p-vs-np-pattern-hybridization-of-partial-bridges.10}
\begin{itemize}
\item Exact statement: There exists an intermediate typed object built from 'MCSP witness family' such that the route from 'A hardness magnification theorem has the form:"A barely-nontrivial lower bound for an explicit problem Q in a weak model implies a major lower bound/separation in a strong model."Typically, Q is a meta-problem (e.g. a Gap-MCSP variant) and"major"means consequences such as NP not in P/poly.' to 'There exist explicit"meta"problems (including MCSP variants) such that proving certain emph barely-superlinear lower bounds against restricted models (e.g. formulas) would imply major lower bounds for NP in much stronger models (including superpolynomial formula/circuit lower bounds, and in some parameter regimes consequences as strong as NP not in P/poly ).' factors through theorem-sized maps preserving 'MCSP complexity profile in nonuniform-circuits'.
\item Restricted-regime mode: special-case
\item Bridge object: MCSP witness family (family of instances equipped with a monotone witness measure)
\item Bridge definition: 'MCSP witness family' denotes the structured intermediate object required by this branch. It is meant to organize the source-side content 'A hardness magnification theorem has the form:"A barely-nontrivial lower bound for an explicit problem Q in a weak model implies a major lower bound/separation in a strong model."Typically, Q is a meta-problem (e.g. a Gap-MCSP variant) and"major"means consequences such as NP not in P/poly.'. It must preserve the invariant 'MCSP complexity profile in nonuniform-circuits'. Its operational role is 'Some blocked routes sit at the intersection of multiple breakthrough families, which is evidence that one bridge object may be insufficient and a composed route should be explored.'.
\item First nontrivial lemma: For every admissible 'MCSP witness family', the first route-specific realization step preserves 'MCSP complexity profile in nonuniform-circuits' and prepares the passage from 'A hardness magnification theorem has the form:"A barely-nontrivial lower bound for an explicit problem Q in a weak model implies a major lower bound/separation in a strong model."Typically, Q is a meta-problem (e.g. a Gap-MCSP variant) and"major"means consequences such as NP not in P/poly.' to 'There exist explicit"meta"problems (including MCSP variants) such that proving certain emph barely-superlinear lower bounds against restricted models (e.g. formulas) would imply major lower bounds for NP in much stronger models (including superpolynomial formula/circuit lower bounds, and in some parameter regimes consequences as strong as NP not in P/poly ).'.
\item Quantifier signature: For every admissible 'MCSP witness family' on the route from 'st.def.hardness-magnification-phenomenon-programme-leve' to 'st.thm.oliveira-santhanam-hardness-magnification-for-na'.
\item LaTeX theorem: \textbackslash\{\}forall W,\textbackslash\{\} \textbackslash\{\}mathrm\{Adm\}(W) \textbackslash\{\}to \textbackslash\{\}exists w : \textbackslash\{\}mathsf\{BW\}(W),\textbackslash\{\} \textbackslash\{\}mathrm\{PresInv\}(W,w) \textbackslash\{\}wedge \textbackslash\{\}mathrm\{RealizesTarget\}(W,w).
\item Negative/barrier version: If 'MCSP witness family' cannot be realized with invariant 'MCSP complexity profile in nonuniform-circuits', the route cannot force 'There exist explicit"meta"problems (including MCSP variants) such that proving certain emph barely-superlinear lower bounds against restricted models (e.g. formulas) would imply major lower bounds for NP in much stronger models (including superpolynomial formula/circuit lower bounds, and in some parameter regimes consequences as strong as NP not in P/poly ).'.
\item Proof strategy family: constructive-realization
\item Negative route: Assume 'MCSP witness family' cannot be realized while preserving 'MCSP complexity profile in nonuniform-circuits' strongly enough to activate the local target claim.
\item Failure diagnosis: The route fails if 'MCSP witness family' cannot be realized while preserving 'MCSP complexity profile in nonuniform-circuits'.
\item Proof obligations: For every admissible 'MCSP witness family', the first route-specific realization step preserves 'MCSP complexity profile in nonuniform-circuits' and prepares the passage from 'A hardness magnification theorem has the form:"A barely-nontrivial lower bound for an explicit problem Q in a weak model implies a major lower bound/separation in a strong model."Typically, Q is a meta-problem (e.g. a Gap-MCSP variant) and"major"means consequences such as NP not in P/poly.' to 'There exist explicit"meta"problems (including MCSP variants) such that proving certain emph barely-superlinear lower bounds against restricted models (e.g. formulas) would imply major lower bounds for NP in much stronger models (including superpolynomial formula/circuit lower bounds, and in some parameter regimes consequences as strong as NP not in P/poly ).'.
\item Candidate lemmas: Factorization theorem for MCSP witness family
\item Independent value: A concrete special-case theorem for 'MCSP witness family' would be independently valuable because it converts a broad discovery programme into an inspectable mathematical subresult.
\item Expert gate: 6/6
\end{itemize}

\subsection*{conjpkg.disclayer-p-vs-np-pattern-hybridization-of-partial-bridges.11}
\begin{itemize}
\item Exact statement: No admissible 'MCSP obstruction family' can simultaneously preserve 'MCSP complexity profile in nonuniform-circuits' and realize 'In refined regimes, lower bounds for Gap-MCSP slightly beyond known frontiers (e.g. N\textasciicircum{} 1+varepsilon -type circuit lower bounds for suitable promise gaps and input length N=2\textasciicircum{}n ) imply major consequences such as NP not in P/poly.' under the current route commitments.
\item Restricted-regime mode: special-case
\item Bridge object: MCSP obstruction family (family of hard instances or obstructions forcing the target payoff)
\item Bridge definition: 'MCSP obstruction family' denotes the structured intermediate object required by this branch. It is meant to organize the source-side content 'A hardness magnification theorem has the form:"A barely-nontrivial lower bound for an explicit problem Q in a weak model implies a major lower bound/separation in a strong model."Typically, Q is a meta-problem (e.g. a Gap-MCSP variant) and"major"means consequences such as NP not in P/poly.'. It must preserve the invariant 'MCSP complexity profile in nonuniform-circuits'. Its operational role is 'Some blocked routes sit at the intersection of multiple breakthrough families, which is evidence that one bridge object may be insufficient and a composed route should be explored.'.
\item First nontrivial lemma: For every admissible 'MCSP obstruction family', the first route-specific realization step preserves 'MCSP complexity profile in nonuniform-circuits' and prepares the passage from 'A hardness magnification theorem has the form:"A barely-nontrivial lower bound for an explicit problem Q in a weak model implies a major lower bound/separation in a strong model."Typically, Q is a meta-problem (e.g. a Gap-MCSP variant) and"major"means consequences such as NP not in P/poly.' to 'In refined regimes, lower bounds for Gap-MCSP slightly beyond known frontiers (e.g. N\textasciicircum{} 1+varepsilon -type circuit lower bounds for suitable promise gaps and input length N=2\textasciicircum{}n ) imply major consequences such as NP not in P/poly.'.
\item Quantifier signature: For every admissible 'MCSP obstruction family' on the route from 'st.def.hardness-magnification-phenomenon-programme-leve' to 'st.thm.oliveira-pich-santhanam-magnification-near-state'.
\item LaTeX theorem: \textbackslash\{\}forall W,\textbackslash\{\} \textbackslash\{\}mathrm\{Adm\}(W) \textbackslash\{\}to \textbackslash\{\}neg \textbackslash\{\}exists w : \textbackslash\{\}mathsf\{BW\}(W),\textbackslash\{\} \textbackslash\{\}mathrm\{PresInv\}(W,w) \textbackslash\{\}wedge \textbackslash\{\}mathrm\{RealizesTarget\}(W,w).
\item Negative/barrier version: If 'MCSP obstruction family' cannot be realized with invariant 'MCSP complexity profile in nonuniform-circuits', the route cannot force 'In refined regimes, lower bounds for Gap-MCSP slightly beyond known frontiers (e.g. N\textasciicircum{} 1+varepsilon -type circuit lower bounds for suitable promise gaps and input length N=2\textasciicircum{}n ) imply major consequences such as NP not in P/poly.'.
\item Proof strategy family: constructive-realization
\item Negative route: Assume 'MCSP obstruction family' cannot be realized while preserving 'MCSP complexity profile in nonuniform-circuits' strongly enough to activate the local target claim.
\item Failure diagnosis: The route fails if 'MCSP obstruction family' cannot be realized while preserving 'MCSP complexity profile in nonuniform-circuits'.
\item Proof obligations: For every admissible 'MCSP obstruction family', the first route-specific realization step preserves 'MCSP complexity profile in nonuniform-circuits' and prepares the passage from 'A hardness magnification theorem has the form:"A barely-nontrivial lower bound for an explicit problem Q in a weak model implies a major lower bound/separation in a strong model."Typically, Q is a meta-problem (e.g. a Gap-MCSP variant) and"major"means consequences such as NP not in P/poly.' to 'In refined regimes, lower bounds for Gap-MCSP slightly beyond known frontiers (e.g. N\textasciicircum{} 1+varepsilon -type circuit lower bounds for suitable promise gaps and input length N=2\textasciicircum{}n ) imply major consequences such as NP not in P/poly.'.
\item Candidate lemmas: Obstruction theorem for MCSP obstruction family
\item Independent value: A concrete special-case theorem for 'MCSP obstruction family' would be independently valuable because it converts a broad discovery programme into an inspectable mathematical subresult.
\item Expert gate: 6/6
\end{itemize}

\subsection*{conjpkg.disclayer-p-vs-np-pattern-hybridization-of-partial-bridges.12}
\begin{itemize}
\item Exact statement: There exists an admissible 'MCSP obstruction family' whose realization internalizes 'A hardness magnification theorem has the form:"A barely-nontrivial lower bound for an explicit problem Q in a weak model implies a major lower bound/separation in a strong model."Typically, Q is a meta-problem (e.g. a Gap-MCSP variant) and"major"means consequences such as NP not in P/poly.' while preserving 'MCSP complexity profile in nonuniform-circuits' strongly enough to support 'In refined regimes, lower bounds for Gap-MCSP slightly beyond known frontiers (e.g. N\textasciicircum{} 1+varepsilon -type circuit lower bounds for suitable promise gaps and input length N=2\textasciicircum{}n ) imply major consequences such as NP not in P/poly.'.
\item Restricted-regime mode: special-case
\item Bridge object: MCSP obstruction family (family of hard instances or obstructions forcing the target payoff)
\item Bridge definition: 'MCSP obstruction family' denotes the structured intermediate object required by this branch. It is meant to organize the source-side content 'A hardness magnification theorem has the form:"A barely-nontrivial lower bound for an explicit problem Q in a weak model implies a major lower bound/separation in a strong model."Typically, Q is a meta-problem (e.g. a Gap-MCSP variant) and"major"means consequences such as NP not in P/poly.'. It must preserve the invariant 'MCSP complexity profile in nonuniform-circuits'. Its operational role is 'Some blocked routes sit at the intersection of multiple breakthrough families, which is evidence that one bridge object may be insufficient and a composed route should be explored.'.
\item First nontrivial lemma: For every admissible 'MCSP obstruction family', the first route-specific realization step preserves 'MCSP complexity profile in nonuniform-circuits' and prepares the passage from 'A hardness magnification theorem has the form:"A barely-nontrivial lower bound for an explicit problem Q in a weak model implies a major lower bound/separation in a strong model."Typically, Q is a meta-problem (e.g. a Gap-MCSP variant) and"major"means consequences such as NP not in P/poly.' to 'In refined regimes, lower bounds for Gap-MCSP slightly beyond known frontiers (e.g. N\textasciicircum{} 1+varepsilon -type circuit lower bounds for suitable promise gaps and input length N=2\textasciicircum{}n ) imply major consequences such as NP not in P/poly.'.
\item Quantifier signature: For every admissible 'MCSP obstruction family' on the route from 'st.def.hardness-magnification-phenomenon-programme-leve' to 'st.thm.oliveira-pich-santhanam-magnification-near-state'.
\item LaTeX theorem: \textbackslash\{\}forall W,\textbackslash\{\} \textbackslash\{\}mathrm\{Adm\}(W) \textbackslash\{\}to \textbackslash\{\}exists w : \textbackslash\{\}mathsf\{BW\}(W),\textbackslash\{\} \textbackslash\{\}mathrm\{PresInv\}(W,w) \textbackslash\{\}wedge \textbackslash\{\}mathrm\{RealizesTarget\}(W,w).
\item Negative/barrier version: If 'MCSP obstruction family' cannot be realized with invariant 'MCSP complexity profile in nonuniform-circuits', the route cannot force 'In refined regimes, lower bounds for Gap-MCSP slightly beyond known frontiers (e.g. N\textasciicircum{} 1+varepsilon -type circuit lower bounds for suitable promise gaps and input length N=2\textasciicircum{}n ) imply major consequences such as NP not in P/poly.'.
\item Proof strategy family: constructive-realization
\item Negative route: Assume 'MCSP obstruction family' cannot be realized while preserving 'MCSP complexity profile in nonuniform-circuits' strongly enough to activate the local target claim.
\item Failure diagnosis: The route fails if 'MCSP obstruction family' cannot be realized while preserving 'MCSP complexity profile in nonuniform-circuits'.
\item Proof obligations: For every admissible 'MCSP obstruction family', the first route-specific realization step preserves 'MCSP complexity profile in nonuniform-circuits' and prepares the passage from 'A hardness magnification theorem has the form:"A barely-nontrivial lower bound for an explicit problem Q in a weak model implies a major lower bound/separation in a strong model."Typically, Q is a meta-problem (e.g. a Gap-MCSP variant) and"major"means consequences such as NP not in P/poly.' to 'In refined regimes, lower bounds for Gap-MCSP slightly beyond known frontiers (e.g. N\textasciicircum{} 1+varepsilon -type circuit lower bounds for suitable promise gaps and input length N=2\textasciicircum{}n ) imply major consequences such as NP not in P/poly.'.
\item Candidate lemmas: Realization theorem for MCSP obstruction family
\item Independent value: A concrete special-case theorem for 'MCSP obstruction family' would be independently valuable because it converts a broad discovery programme into an inspectable mathematical subresult.
\item Expert gate: 6/6
\end{itemize}

\subsection*{conjpkg.disclayer-p-vs-np-pattern-hybridization-of-partial-bridges.13}
\begin{itemize}
\item Exact statement: There exists an intermediate typed object built from 'MCSP obstruction family' such that the route from 'A hardness magnification theorem has the form:"A barely-nontrivial lower bound for an explicit problem Q in a weak model implies a major lower bound/separation in a strong model."Typically, Q is a meta-problem (e.g. a Gap-MCSP variant) and"major"means consequences such as NP not in P/poly.' to 'In refined regimes, lower bounds for Gap-MCSP slightly beyond known frontiers (e.g. N\textasciicircum{} 1+varepsilon -type circuit lower bounds for suitable promise gaps and input length N=2\textasciicircum{}n ) imply major consequences such as NP not in P/poly.' factors through theorem-sized maps preserving 'MCSP complexity profile in nonuniform-circuits'.
\item Restricted-regime mode: special-case
\item Bridge object: MCSP obstruction family (family of hard instances or obstructions forcing the target payoff)
\item Bridge definition: 'MCSP obstruction family' denotes the structured intermediate object required by this branch. It is meant to organize the source-side content 'A hardness magnification theorem has the form:"A barely-nontrivial lower bound for an explicit problem Q in a weak model implies a major lower bound/separation in a strong model."Typically, Q is a meta-problem (e.g. a Gap-MCSP variant) and"major"means consequences such as NP not in P/poly.'. It must preserve the invariant 'MCSP complexity profile in nonuniform-circuits'. Its operational role is 'Some blocked routes sit at the intersection of multiple breakthrough families, which is evidence that one bridge object may be insufficient and a composed route should be explored.'.
\item First nontrivial lemma: For every admissible 'MCSP obstruction family', the first route-specific realization step preserves 'MCSP complexity profile in nonuniform-circuits' and prepares the passage from 'A hardness magnification theorem has the form:"A barely-nontrivial lower bound for an explicit problem Q in a weak model implies a major lower bound/separation in a strong model."Typically, Q is a meta-problem (e.g. a Gap-MCSP variant) and"major"means consequences such as NP not in P/poly.' to 'In refined regimes, lower bounds for Gap-MCSP slightly beyond known frontiers (e.g. N\textasciicircum{} 1+varepsilon -type circuit lower bounds for suitable promise gaps and input length N=2\textasciicircum{}n ) imply major consequences such as NP not in P/poly.'.
\item Quantifier signature: For every admissible 'MCSP obstruction family' on the route from 'st.def.hardness-magnification-phenomenon-programme-leve' to 'st.thm.oliveira-pich-santhanam-magnification-near-state'.
\item LaTeX theorem: \textbackslash\{\}forall W,\textbackslash\{\} \textbackslash\{\}mathrm\{Adm\}(W) \textbackslash\{\}to \textbackslash\{\}exists w : \textbackslash\{\}mathsf\{BW\}(W),\textbackslash\{\} \textbackslash\{\}mathrm\{PresInv\}(W,w) \textbackslash\{\}wedge \textbackslash\{\}mathrm\{RealizesTarget\}(W,w).
\item Negative/barrier version: If 'MCSP obstruction family' cannot be realized with invariant 'MCSP complexity profile in nonuniform-circuits', the route cannot force 'In refined regimes, lower bounds for Gap-MCSP slightly beyond known frontiers (e.g. N\textasciicircum{} 1+varepsilon -type circuit lower bounds for suitable promise gaps and input length N=2\textasciicircum{}n ) imply major consequences such as NP not in P/poly.'.
\item Proof strategy family: constructive-realization
\item Negative route: Assume 'MCSP obstruction family' cannot be realized while preserving 'MCSP complexity profile in nonuniform-circuits' strongly enough to activate the local target claim.
\item Failure diagnosis: The route fails if 'MCSP obstruction family' cannot be realized while preserving 'MCSP complexity profile in nonuniform-circuits'.
\item Proof obligations: For every admissible 'MCSP obstruction family', the first route-specific realization step preserves 'MCSP complexity profile in nonuniform-circuits' and prepares the passage from 'A hardness magnification theorem has the form:"A barely-nontrivial lower bound for an explicit problem Q in a weak model implies a major lower bound/separation in a strong model."Typically, Q is a meta-problem (e.g. a Gap-MCSP variant) and"major"means consequences such as NP not in P/poly.' to 'In refined regimes, lower bounds for Gap-MCSP slightly beyond known frontiers (e.g. N\textasciicircum{} 1+varepsilon -type circuit lower bounds for suitable promise gaps and input length N=2\textasciicircum{}n ) imply major consequences such as NP not in P/poly.'.
\item Candidate lemmas: Factorization theorem for MCSP obstruction family
\item Independent value: A concrete special-case theorem for 'MCSP obstruction family' would be independently valuable because it converts a broad discovery programme into an inspectable mathematical subresult.
\item Expert gate: 6/6
\end{itemize}

\subsection*{conjpkg.disclayer-p-vs-np-pattern-regime-localization.01}
\begin{itemize}
\item Exact statement: There exists an admissible 'MCSP witness family' such that 'A hardness magnification theorem has the form:"A barely-nontrivial lower bound for an explicit problem Q in a weak model implies a major lower bound/separation in a strong model."Typically, Q is a meta-problem (e.g. a Gap-MCSP variant) and"major"means consequences such as NP not in P/poly.' can be organized into a route preserving 'MCSP complexity profile in nonuniform-circuits' and thereby supports 'In refined regimes, lower bounds for Gap-MCSP slightly beyond known frontiers (e.g. N\textasciicircum{} 1+varepsilon -type circuit lower bounds for suitable promise gaps and input length N=2\textasciicircum{}n ) imply major consequences such as NP not in P/poly.'.
\item Restricted-regime mode: threshold-version
\item Bridge object: MCSP witness family (family of instances equipped with a monotone witness measure)
\item Bridge definition: 'MCSP witness family' denotes the structured intermediate object required by this branch. It is meant to organize the source-side content 'A hardness magnification theorem has the form:"A barely-nontrivial lower bound for an explicit problem Q in a weak model implies a major lower bound/separation in a strong model."Typically, Q is a meta-problem (e.g. a Gap-MCSP variant) and"major"means consequences such as NP not in P/poly.'. It must preserve the invariant 'MCSP complexity profile in nonuniform-circuits'. Its operational role is 'The frontier repeatedly crosses model or threshold boundaries, so discovery should search for the smallest regime in which the route first becomes true.'.
\item First nontrivial lemma: There exists a threshold regime in which 'MCSP witness family' already preserves 'MCSP complexity profile in nonuniform-circuits' well enough to imply a restricted version of 'In refined regimes, lower bounds for Gap-MCSP slightly beyond known frontiers (e.g. N\textasciicircum{} 1+varepsilon -type circuit lower bounds for suitable promise gaps and input length N=2\textasciicircum{}n ) imply major consequences such as NP not in P/poly.'.
\item Quantifier signature: For every admissible 'MCSP witness family' on the route from 'st.def.hardness-magnification-phenomenon-programme-leve' to 'st.thm.oliveira-pich-santhanam-magnification-near-state'.
\item LaTeX theorem: \textbackslash\{\}exists \textbackslash\{\}theta,\textbackslash\{\} \textbackslash\{\}forall W,\textbackslash\{\} \textbackslash\{\}mathrm\{Adm\}\_\{\textbackslash\{\}theta\}(W) \textbackslash\{\}to \textbackslash\{\}exists w : \textbackslash\{\}mathsf\{BW\}(W),\textbackslash\{\} \textbackslash\{\}mathrm\{PresInv\}(W,w) \textbackslash\{\}wedge \textbackslash\{\}mathrm\{RealizesTarget\}(W,w).
\item Negative/barrier version: If 'MCSP witness family' cannot be realized with invariant 'MCSP complexity profile in nonuniform-circuits', the route cannot force 'In refined regimes, lower bounds for Gap-MCSP slightly beyond known frontiers (e.g. N\textasciicircum{} 1+varepsilon -type circuit lower bounds for suitable promise gaps and input length N=2\textasciicircum{}n ) imply major consequences such as NP not in P/poly.'.
\item Proof strategy family: threshold-localization
\item Negative route: Assume no threshold-localized version of 'MCSP witness family' preserves 'MCSP complexity profile in nonuniform-circuits' strongly enough for the restricted target.
\item Failure diagnosis: The route fails if no threshold-localized version of 'MCSP witness family' preserves 'MCSP complexity profile in nonuniform-circuits' strongly enough to force the restricted target.
\item Proof obligations: There exists a threshold regime in which 'MCSP witness family' already preserves 'MCSP complexity profile in nonuniform-circuits' well enough to imply a restricted version of 'In refined regimes, lower bounds for Gap-MCSP slightly beyond known frontiers (e.g. N\textasciicircum{} 1+varepsilon -type circuit lower bounds for suitable promise gaps and input length N=2\textasciicircum{}n ) imply major consequences such as NP not in P/poly.'.
\item Candidate lemmas: Prove that a hardness-magnification witness family turns 'A barely-nontrivial lower bound for an explicit problem Q in a weak model' into 'major consequences such as NP not in P/poly' while preserving 'MCSP complexity profile in nonuniform-circuits'.; Threshold localization lemma for MCSP witness family
\item Independent value: A threshold theorem for 'MCSP witness family' would be independently valuable because it localizes the first nontrivial regime where the route becomes mathematically real.
\item Expert gate: 6/6
\end{itemize}

\subsection*{conjpkg.disclayer-p-vs-np-pattern-regime-localization.02}
\begin{itemize}
\item Exact statement: There exists an admissible 'MCSP witness family' such that 'A hardness magnification theorem has the form:"A barely-nontrivial lower bound for an explicit problem Q in a weak model implies a major lower bound/separation in a strong model."Typically, Q is a meta-problem (e.g. a Gap-MCSP variant) and"major"means consequences such as NP not in P/poly.' can be organized into a route preserving 'MCSP complexity profile in nonuniform-circuits' and thereby supports 'There exist explicit"meta"problems (including MCSP variants) such that proving certain emph barely-superlinear lower bounds against restricted models (e.g. formulas) would imply major lower bounds for NP in much stronger models (including superpolynomial formula/circuit lower bounds, and in some parameter regimes consequences as strong as NP not in P/poly ).'.
\item Restricted-regime mode: threshold-version
\item Bridge object: MCSP witness family (family of instances equipped with a monotone witness measure)
\item Bridge definition: 'MCSP witness family' denotes the structured intermediate object required by this branch. It is meant to organize the source-side content 'A hardness magnification theorem has the form:"A barely-nontrivial lower bound for an explicit problem Q in a weak model implies a major lower bound/separation in a strong model."Typically, Q is a meta-problem (e.g. a Gap-MCSP variant) and"major"means consequences such as NP not in P/poly.'. It must preserve the invariant 'MCSP complexity profile in nonuniform-circuits'. Its operational role is 'The frontier repeatedly crosses model or threshold boundaries, so discovery should search for the smallest regime in which the route first becomes true.'.
\item First nontrivial lemma: There exists a threshold regime in which 'MCSP witness family' already preserves 'MCSP complexity profile in nonuniform-circuits' well enough to imply a restricted version of 'There exist explicit"meta"problems (including MCSP variants) such that proving certain emph barely-superlinear lower bounds against restricted models (e.g. formulas) would imply major lower bounds for NP in much stronger models (including superpolynomial formula/circuit lower bounds, and in some parameter regimes consequences as strong as NP not in P/poly ).'.
\item Quantifier signature: For every admissible 'MCSP witness family' on the route from 'st.def.hardness-magnification-phenomenon-programme-leve' to 'st.thm.oliveira-santhanam-hardness-magnification-for-na'.
\item LaTeX theorem: \textbackslash\{\}exists \textbackslash\{\}theta,\textbackslash\{\} \textbackslash\{\}forall W,\textbackslash\{\} \textbackslash\{\}mathrm\{Adm\}\_\{\textbackslash\{\}theta\}(W) \textbackslash\{\}to \textbackslash\{\}exists w : \textbackslash\{\}mathsf\{BW\}(W),\textbackslash\{\} \textbackslash\{\}mathrm\{PresInv\}(W,w) \textbackslash\{\}wedge \textbackslash\{\}mathrm\{RealizesTarget\}(W,w).
\item Negative/barrier version: If 'MCSP witness family' cannot be realized with invariant 'MCSP complexity profile in nonuniform-circuits', the route cannot force 'There exist explicit"meta"problems (including MCSP variants) such that proving certain emph barely-superlinear lower bounds against restricted models (e.g. formulas) would imply major lower bounds for NP in much stronger models (including superpolynomial formula/circuit lower bounds, and in some parameter regimes consequences as strong as NP not in P/poly ).'.
\item Proof strategy family: threshold-localization
\item Negative route: Assume no threshold-localized version of 'MCSP witness family' preserves 'MCSP complexity profile in nonuniform-circuits' strongly enough for the restricted target.
\item Failure diagnosis: The route fails if no threshold-localized version of 'MCSP witness family' preserves 'MCSP complexity profile in nonuniform-circuits' strongly enough to force the restricted target.
\item Proof obligations: There exists a threshold regime in which 'MCSP witness family' already preserves 'MCSP complexity profile in nonuniform-circuits' well enough to imply a restricted version of 'There exist explicit"meta"problems (including MCSP variants) such that proving certain emph barely-superlinear lower bounds against restricted models (e.g. formulas) would imply major lower bounds for NP in much stronger models (including superpolynomial formula/circuit lower bounds, and in some parameter regimes consequences as strong as NP not in P/poly ).'.
\item Candidate lemmas: Prove that a hardness-magnification witness family preserves 'MCSP complexity profile in nonuniform-circuits' while transporting the route from the nonuniform-circuits regime to the formula-size regime strongly enough to derive 'major lower bounds for NP in much stronger models (including superpolynomial formula/circuit lower bounds, and in some parameter...'.; Threshold localization lemma for MCSP witness family
\item Independent value: A threshold theorem for 'MCSP witness family' would be independently valuable because it localizes the first nontrivial regime where the route becomes mathematically real.
\item Expert gate: 6/6
\end{itemize}

\subsection*{conjpkg.disclayer-p-vs-np-pattern-regime-localization.03}
\begin{itemize}
\item Exact statement: There exists an admissible 'MCSP obstruction family' such that 'A hardness magnification theorem has the form:"A barely-nontrivial lower bound for an explicit problem Q in a weak model implies a major lower bound/separation in a strong model."Typically, Q is a meta-problem (e.g. a Gap-MCSP variant) and"major"means consequences such as NP not in P/poly.' can be organized into a route preserving 'MCSP complexity profile in nonuniform-circuits' and thereby supports 'In refined regimes, lower bounds for Gap-MCSP slightly beyond known frontiers (e.g. N\textasciicircum{} 1+varepsilon -type circuit lower bounds for suitable promise gaps and input length N=2\textasciicircum{}n ) imply major consequences such as NP not in P/poly.'.
\item Restricted-regime mode: threshold-version
\item Bridge object: MCSP obstruction family (family of hard instances or obstructions forcing the target payoff)
\item Bridge definition: 'MCSP obstruction family' denotes the structured intermediate object required by this branch. It is meant to organize the source-side content 'A hardness magnification theorem has the form:"A barely-nontrivial lower bound for an explicit problem Q in a weak model implies a major lower bound/separation in a strong model."Typically, Q is a meta-problem (e.g. a Gap-MCSP variant) and"major"means consequences such as NP not in P/poly.'. It must preserve the invariant 'MCSP complexity profile in nonuniform-circuits'. Its operational role is 'The frontier repeatedly crosses model or threshold boundaries, so discovery should search for the smallest regime in which the route first becomes true.'.
\item First nontrivial lemma: There exists a threshold regime in which 'MCSP obstruction family' already preserves 'MCSP complexity profile in nonuniform-circuits' well enough to imply a restricted version of 'In refined regimes, lower bounds for Gap-MCSP slightly beyond known frontiers (e.g. N\textasciicircum{} 1+varepsilon -type circuit lower bounds for suitable promise gaps and input length N=2\textasciicircum{}n ) imply major consequences such as NP not in P/poly.'.
\item Quantifier signature: For every admissible 'MCSP obstruction family' on the route from 'st.def.hardness-magnification-phenomenon-programme-leve' to 'st.thm.oliveira-pich-santhanam-magnification-near-state'.
\item LaTeX theorem: \textbackslash\{\}exists \textbackslash\{\}theta,\textbackslash\{\} \textbackslash\{\}forall W,\textbackslash\{\} \textbackslash\{\}mathrm\{Adm\}\_\{\textbackslash\{\}theta\}(W) \textbackslash\{\}to \textbackslash\{\}exists w : \textbackslash\{\}mathsf\{BW\}(W),\textbackslash\{\} \textbackslash\{\}mathrm\{PresInv\}(W,w) \textbackslash\{\}wedge \textbackslash\{\}mathrm\{RealizesTarget\}(W,w).
\item Negative/barrier version: If 'MCSP obstruction family' cannot be realized with invariant 'MCSP complexity profile in nonuniform-circuits', the route cannot force 'In refined regimes, lower bounds for Gap-MCSP slightly beyond known frontiers (e.g. N\textasciicircum{} 1+varepsilon -type circuit lower bounds for suitable promise gaps and input length N=2\textasciicircum{}n ) imply major consequences such as NP not in P/poly.'.
\item Proof strategy family: threshold-localization
\item Negative route: Assume no threshold-localized version of 'MCSP obstruction family' preserves 'MCSP complexity profile in nonuniform-circuits' strongly enough for the restricted target.
\item Failure diagnosis: The route fails if no threshold-localized version of 'MCSP obstruction family' preserves 'MCSP complexity profile in nonuniform-circuits' strongly enough to force the restricted target.
\item Proof obligations: There exists a threshold regime in which 'MCSP obstruction family' already preserves 'MCSP complexity profile in nonuniform-circuits' well enough to imply a restricted version of 'In refined regimes, lower bounds for Gap-MCSP slightly beyond known frontiers (e.g. N\textasciicircum{} 1+varepsilon -type circuit lower bounds for suitable promise gaps and input length N=2\textasciicircum{}n ) imply major consequences such as NP not in P/poly.'.
\item Candidate lemmas: Prove that a hardness-magnification witness family turns 'A barely-nontrivial lower bound for an explicit problem Q in a weak model' into 'major consequences such as NP not in P/poly' while preserving 'MCSP complexity profile in nonuniform-circuits'.; Threshold localization lemma for MCSP obstruction family
\item Independent value: A threshold theorem for 'MCSP obstruction family' would be independently valuable because it localizes the first nontrivial regime where the route becomes mathematically real.
\item Expert gate: 6/6
\end{itemize}

\subsection*{conjpkg.disclayer-p-vs-np-pattern-regime-localization.04}
\begin{itemize}
\item Exact statement: No admissible 'MCSP witness family' can simultaneously preserve 'MCSP complexity profile in nonuniform-circuits' and realize 'In refined regimes, lower bounds for Gap-MCSP slightly beyond known frontiers (e.g. N\textasciicircum{} 1+varepsilon -type circuit lower bounds for suitable promise gaps and input length N=2\textasciicircum{}n ) imply major consequences such as NP not in P/poly.' under the current route commitments.
\item Restricted-regime mode: threshold-version
\item Bridge object: MCSP witness family (family of instances equipped with a monotone witness measure)
\item Bridge definition: 'MCSP witness family' denotes the structured intermediate object required by this branch. It is meant to organize the source-side content 'A hardness magnification theorem has the form:"A barely-nontrivial lower bound for an explicit problem Q in a weak model implies a major lower bound/separation in a strong model."Typically, Q is a meta-problem (e.g. a Gap-MCSP variant) and"major"means consequences such as NP not in P/poly.'. It must preserve the invariant 'MCSP complexity profile in nonuniform-circuits'. Its operational role is 'The frontier repeatedly crosses model or threshold boundaries, so discovery should search for the smallest regime in which the route first becomes true.'.
\item First nontrivial lemma: There exists a threshold regime in which 'MCSP witness family' already preserves 'MCSP complexity profile in nonuniform-circuits' well enough to imply a restricted version of 'In refined regimes, lower bounds for Gap-MCSP slightly beyond known frontiers (e.g. N\textasciicircum{} 1+varepsilon -type circuit lower bounds for suitable promise gaps and input length N=2\textasciicircum{}n ) imply major consequences such as NP not in P/poly.'.
\item Quantifier signature: For every admissible 'MCSP witness family' on the route from 'st.def.hardness-magnification-phenomenon-programme-leve' to 'st.thm.oliveira-pich-santhanam-magnification-near-state'.
\item LaTeX theorem: \textbackslash\{\}forall W,\textbackslash\{\} \textbackslash\{\}mathrm\{Adm\}(W) \textbackslash\{\}to \textbackslash\{\}neg \textbackslash\{\}exists w : \textbackslash\{\}mathsf\{BW\}(W),\textbackslash\{\} \textbackslash\{\}mathrm\{PresInv\}(W,w) \textbackslash\{\}wedge \textbackslash\{\}mathrm\{RealizesTarget\}(W,w).
\item Negative/barrier version: If 'MCSP witness family' cannot be realized with invariant 'MCSP complexity profile in nonuniform-circuits', the route cannot force 'In refined regimes, lower bounds for Gap-MCSP slightly beyond known frontiers (e.g. N\textasciicircum{} 1+varepsilon -type circuit lower bounds for suitable promise gaps and input length N=2\textasciicircum{}n ) imply major consequences such as NP not in P/poly.'.
\item Proof strategy family: threshold-localization
\item Negative route: Assume no threshold-localized version of 'MCSP witness family' preserves 'MCSP complexity profile in nonuniform-circuits' strongly enough for the restricted target.
\item Failure diagnosis: The route fails if no threshold-localized version of 'MCSP witness family' preserves 'MCSP complexity profile in nonuniform-circuits' strongly enough to force the restricted target.
\item Proof obligations: There exists a threshold regime in which 'MCSP witness family' already preserves 'MCSP complexity profile in nonuniform-circuits' well enough to imply a restricted version of 'In refined regimes, lower bounds for Gap-MCSP slightly beyond known frontiers (e.g. N\textasciicircum{} 1+varepsilon -type circuit lower bounds for suitable promise gaps and input length N=2\textasciicircum{}n ) imply major consequences such as NP not in P/poly.'.
\item Candidate lemmas: Obstruction theorem for MCSP witness family; Threshold localization lemma for MCSP witness family
\item Independent value: A threshold theorem for 'MCSP witness family' would be independently valuable because it localizes the first nontrivial regime where the route becomes mathematically real.
\item Expert gate: 6/6
\end{itemize}

\subsection*{conjpkg.disclayer-p-vs-np-pattern-regime-localization.05}
\begin{itemize}
\item Exact statement: There exists an admissible 'MCSP witness family' whose realization internalizes 'A hardness magnification theorem has the form:"A barely-nontrivial lower bound for an explicit problem Q in a weak model implies a major lower bound/separation in a strong model."Typically, Q is a meta-problem (e.g. a Gap-MCSP variant) and"major"means consequences such as NP not in P/poly.' while preserving 'MCSP complexity profile in nonuniform-circuits' strongly enough to support 'In refined regimes, lower bounds for Gap-MCSP slightly beyond known frontiers (e.g. N\textasciicircum{} 1+varepsilon -type circuit lower bounds for suitable promise gaps and input length N=2\textasciicircum{}n ) imply major consequences such as NP not in P/poly.'.
\item Restricted-regime mode: threshold-version
\item Bridge object: MCSP witness family (family of instances equipped with a monotone witness measure)
\item Bridge definition: 'MCSP witness family' denotes the structured intermediate object required by this branch. It is meant to organize the source-side content 'A hardness magnification theorem has the form:"A barely-nontrivial lower bound for an explicit problem Q in a weak model implies a major lower bound/separation in a strong model."Typically, Q is a meta-problem (e.g. a Gap-MCSP variant) and"major"means consequences such as NP not in P/poly.'. It must preserve the invariant 'MCSP complexity profile in nonuniform-circuits'. Its operational role is 'The frontier repeatedly crosses model or threshold boundaries, so discovery should search for the smallest regime in which the route first becomes true.'.
\item First nontrivial lemma: There exists a threshold regime in which 'MCSP witness family' already preserves 'MCSP complexity profile in nonuniform-circuits' well enough to imply a restricted version of 'In refined regimes, lower bounds for Gap-MCSP slightly beyond known frontiers (e.g. N\textasciicircum{} 1+varepsilon -type circuit lower bounds for suitable promise gaps and input length N=2\textasciicircum{}n ) imply major consequences such as NP not in P/poly.'.
\item Quantifier signature: For every admissible 'MCSP witness family' on the route from 'st.def.hardness-magnification-phenomenon-programme-leve' to 'st.thm.oliveira-pich-santhanam-magnification-near-state'.
\item LaTeX theorem: \textbackslash\{\}forall W,\textbackslash\{\} \textbackslash\{\}mathrm\{Adm\}(W) \textbackslash\{\}to \textbackslash\{\}exists w : \textbackslash\{\}mathsf\{BW\}(W),\textbackslash\{\} \textbackslash\{\}mathrm\{PresInv\}(W,w) \textbackslash\{\}wedge \textbackslash\{\}mathrm\{RealizesTarget\}(W,w).
\item Negative/barrier version: If 'MCSP witness family' cannot be realized with invariant 'MCSP complexity profile in nonuniform-circuits', the route cannot force 'In refined regimes, lower bounds for Gap-MCSP slightly beyond known frontiers (e.g. N\textasciicircum{} 1+varepsilon -type circuit lower bounds for suitable promise gaps and input length N=2\textasciicircum{}n ) imply major consequences such as NP not in P/poly.'.
\item Proof strategy family: threshold-localization
\item Negative route: Assume no threshold-localized version of 'MCSP witness family' preserves 'MCSP complexity profile in nonuniform-circuits' strongly enough for the restricted target.
\item Failure diagnosis: The route fails if no threshold-localized version of 'MCSP witness family' preserves 'MCSP complexity profile in nonuniform-circuits' strongly enough to force the restricted target.
\item Proof obligations: There exists a threshold regime in which 'MCSP witness family' already preserves 'MCSP complexity profile in nonuniform-circuits' well enough to imply a restricted version of 'In refined regimes, lower bounds for Gap-MCSP slightly beyond known frontiers (e.g. N\textasciicircum{} 1+varepsilon -type circuit lower bounds for suitable promise gaps and input length N=2\textasciicircum{}n ) imply major consequences such as NP not in P/poly.'.
\item Candidate lemmas: Realization theorem for MCSP witness family; Threshold localization lemma for MCSP witness family
\item Independent value: A threshold theorem for 'MCSP witness family' would be independently valuable because it localizes the first nontrivial regime where the route becomes mathematically real.
\item Expert gate: 6/6
\end{itemize}

\subsection*{conjpkg.disclayer-p-vs-np-pattern-regime-localization.06}
\begin{itemize}
\item Exact statement: For every admissible 'MCSP witness family', there is a transport law from 'nonuniform-circuits' to 'formula-size' that preserves 'MCSP complexity profile in nonuniform-circuits' and retains enough strength to activate 'There exist explicit"meta"problems (including MCSP variants) such that proving certain emph barely-superlinear lower bounds against restricted models (e.g. formulas) would imply major lower bounds for NP in much stronger models (including superpolynomial formula/circuit lower bounds, and in some parameter regimes consequences as strong as NP not in P/poly ).'.
\item Restricted-regime mode: threshold-version
\item Bridge object: MCSP witness family (family of instances equipped with a monotone witness measure)
\item Bridge definition: 'MCSP witness family' denotes the structured intermediate object required by this branch. It is meant to organize the source-side content 'A hardness magnification theorem has the form:"A barely-nontrivial lower bound for an explicit problem Q in a weak model implies a major lower bound/separation in a strong model."Typically, Q is a meta-problem (e.g. a Gap-MCSP variant) and"major"means consequences such as NP not in P/poly.'. It must preserve the invariant 'MCSP complexity profile in nonuniform-circuits'. Its operational role is 'The frontier repeatedly crosses model or threshold boundaries, so discovery should search for the smallest regime in which the route first becomes true.'.
\item First nontrivial lemma: There exists a threshold regime in which 'MCSP witness family' already preserves 'MCSP complexity profile in nonuniform-circuits' well enough to imply a restricted version of 'There exist explicit"meta"problems (including MCSP variants) such that proving certain emph barely-superlinear lower bounds against restricted models (e.g. formulas) would imply major lower bounds for NP in much stronger models (including superpolynomial formula/circuit lower bounds, and in some parameter regimes consequences as strong as NP not in P/poly ).'.
\item Quantifier signature: For every admissible 'MCSP witness family' on the route from 'st.def.hardness-magnification-phenomenon-programme-leve' to 'st.thm.oliveira-santhanam-hardness-magnification-for-na'.
\item LaTeX theorem: \textbackslash\{\}forall W,\textbackslash\{\} \textbackslash\{\}mathrm\{Adm\}(W) \textbackslash\{\}to \textbackslash\{\}exists w : \textbackslash\{\}mathsf\{BW\}(W),\textbackslash\{\} \textbackslash\{\}mathrm\{PresInv\}(W,w) \textbackslash\{\}wedge \textbackslash\{\}mathrm\{RealizesTarget\}(W,w).
\item Negative/barrier version: If 'MCSP witness family' cannot be realized with invariant 'MCSP complexity profile in nonuniform-circuits', the route cannot force 'There exist explicit"meta"problems (including MCSP variants) such that proving certain emph barely-superlinear lower bounds against restricted models (e.g. formulas) would imply major lower bounds for NP in much stronger models (including superpolynomial formula/circuit lower bounds, and in some parameter regimes consequences as strong as NP not in P/poly ).'.
\item Proof strategy family: threshold-localization
\item Negative route: Assume no threshold-localized version of 'MCSP witness family' preserves 'MCSP complexity profile in nonuniform-circuits' strongly enough for the restricted target.
\item Failure diagnosis: The route fails if no threshold-localized version of 'MCSP witness family' preserves 'MCSP complexity profile in nonuniform-circuits' strongly enough to force the restricted target.
\item Proof obligations: There exists a threshold regime in which 'MCSP witness family' already preserves 'MCSP complexity profile in nonuniform-circuits' well enough to imply a restricted version of 'There exist explicit"meta"problems (including MCSP variants) such that proving certain emph barely-superlinear lower bounds against restricted models (e.g. formulas) would imply major lower bounds for NP in much stronger models (including superpolynomial formula/circuit lower bounds, and in some parameter regimes consequences as strong as NP not in P/poly ).'.
\item Candidate lemmas: Transport theorem for MCSP witness family; Threshold localization lemma for MCSP witness family
\item Independent value: A threshold theorem for 'MCSP witness family' would be independently valuable because it localizes the first nontrivial regime where the route becomes mathematically real.
\item Expert gate: 6/6
\end{itemize}

\subsection*{conjpkg.disclayer-p-vs-np-pattern-regime-localization.07}
\begin{itemize}
\item Exact statement: No admissible 'MCSP witness family' can simultaneously preserve 'MCSP complexity profile in nonuniform-circuits' and realize 'There exist explicit"meta"problems (including MCSP variants) such that proving certain emph barely-superlinear lower bounds against restricted models (e.g. formulas) would imply major lower bounds for NP in much stronger models (including superpolynomial formula/circuit lower bounds, and in some parameter regimes consequences as strong as NP not in P/poly ).' under the current route commitments.
\item Restricted-regime mode: threshold-version
\item Bridge object: MCSP witness family (family of instances equipped with a monotone witness measure)
\item Bridge definition: 'MCSP witness family' denotes the structured intermediate object required by this branch. It is meant to organize the source-side content 'A hardness magnification theorem has the form:"A barely-nontrivial lower bound for an explicit problem Q in a weak model implies a major lower bound/separation in a strong model."Typically, Q is a meta-problem (e.g. a Gap-MCSP variant) and"major"means consequences such as NP not in P/poly.'. It must preserve the invariant 'MCSP complexity profile in nonuniform-circuits'. Its operational role is 'The frontier repeatedly crosses model or threshold boundaries, so discovery should search for the smallest regime in which the route first becomes true.'.
\item First nontrivial lemma: There exists a threshold regime in which 'MCSP witness family' already preserves 'MCSP complexity profile in nonuniform-circuits' well enough to imply a restricted version of 'There exist explicit"meta"problems (including MCSP variants) such that proving certain emph barely-superlinear lower bounds against restricted models (e.g. formulas) would imply major lower bounds for NP in much stronger models (including superpolynomial formula/circuit lower bounds, and in some parameter regimes consequences as strong as NP not in P/poly ).'.
\item Quantifier signature: For every admissible 'MCSP witness family' on the route from 'st.def.hardness-magnification-phenomenon-programme-leve' to 'st.thm.oliveira-santhanam-hardness-magnification-for-na'.
\item LaTeX theorem: \textbackslash\{\}forall W,\textbackslash\{\} \textbackslash\{\}mathrm\{Adm\}(W) \textbackslash\{\}to \textbackslash\{\}neg \textbackslash\{\}exists w : \textbackslash\{\}mathsf\{BW\}(W),\textbackslash\{\} \textbackslash\{\}mathrm\{PresInv\}(W,w) \textbackslash\{\}wedge \textbackslash\{\}mathrm\{RealizesTarget\}(W,w).
\item Negative/barrier version: If 'MCSP witness family' cannot be realized with invariant 'MCSP complexity profile in nonuniform-circuits', the route cannot force 'There exist explicit"meta"problems (including MCSP variants) such that proving certain emph barely-superlinear lower bounds against restricted models (e.g. formulas) would imply major lower bounds for NP in much stronger models (including superpolynomial formula/circuit lower bounds, and in some parameter regimes consequences as strong as NP not in P/poly ).'.
\item Proof strategy family: threshold-localization
\item Negative route: Assume no threshold-localized version of 'MCSP witness family' preserves 'MCSP complexity profile in nonuniform-circuits' strongly enough for the restricted target.
\item Failure diagnosis: The route fails if no threshold-localized version of 'MCSP witness family' preserves 'MCSP complexity profile in nonuniform-circuits' strongly enough to force the restricted target.
\item Proof obligations: There exists a threshold regime in which 'MCSP witness family' already preserves 'MCSP complexity profile in nonuniform-circuits' well enough to imply a restricted version of 'There exist explicit"meta"problems (including MCSP variants) such that proving certain emph barely-superlinear lower bounds against restricted models (e.g. formulas) would imply major lower bounds for NP in much stronger models (including superpolynomial formula/circuit lower bounds, and in some parameter regimes consequences as strong as NP not in P/poly ).'.
\item Candidate lemmas: Obstruction theorem for MCSP witness family; Threshold localization lemma for MCSP witness family
\item Independent value: A threshold theorem for 'MCSP witness family' would be independently valuable because it localizes the first nontrivial regime where the route becomes mathematically real.
\item Expert gate: 6/6
\end{itemize}

\subsection*{conjpkg.disclayer-p-vs-np-pattern-regime-localization.08}
\begin{itemize}
\item Exact statement: There exists an admissible 'MCSP witness family' whose realization internalizes 'A hardness magnification theorem has the form:"A barely-nontrivial lower bound for an explicit problem Q in a weak model implies a major lower bound/separation in a strong model."Typically, Q is a meta-problem (e.g. a Gap-MCSP variant) and"major"means consequences such as NP not in P/poly.' while preserving 'MCSP complexity profile in nonuniform-circuits' strongly enough to support 'There exist explicit"meta"problems (including MCSP variants) such that proving certain emph barely-superlinear lower bounds against restricted models (e.g. formulas) would imply major lower bounds for NP in much stronger models (including superpolynomial formula/circuit lower bounds, and in some parameter regimes consequences as strong as NP not in P/poly ).'.
\item Restricted-regime mode: threshold-version
\item Bridge object: MCSP witness family (family of instances equipped with a monotone witness measure)
\item Bridge definition: 'MCSP witness family' denotes the structured intermediate object required by this branch. It is meant to organize the source-side content 'A hardness magnification theorem has the form:"A barely-nontrivial lower bound for an explicit problem Q in a weak model implies a major lower bound/separation in a strong model."Typically, Q is a meta-problem (e.g. a Gap-MCSP variant) and"major"means consequences such as NP not in P/poly.'. It must preserve the invariant 'MCSP complexity profile in nonuniform-circuits'. Its operational role is 'The frontier repeatedly crosses model or threshold boundaries, so discovery should search for the smallest regime in which the route first becomes true.'.
\item First nontrivial lemma: There exists a threshold regime in which 'MCSP witness family' already preserves 'MCSP complexity profile in nonuniform-circuits' well enough to imply a restricted version of 'There exist explicit"meta"problems (including MCSP variants) such that proving certain emph barely-superlinear lower bounds against restricted models (e.g. formulas) would imply major lower bounds for NP in much stronger models (including superpolynomial formula/circuit lower bounds, and in some parameter regimes consequences as strong as NP not in P/poly ).'.
\item Quantifier signature: For every admissible 'MCSP witness family' on the route from 'st.def.hardness-magnification-phenomenon-programme-leve' to 'st.thm.oliveira-santhanam-hardness-magnification-for-na'.
\item LaTeX theorem: \textbackslash\{\}forall W,\textbackslash\{\} \textbackslash\{\}mathrm\{Adm\}(W) \textbackslash\{\}to \textbackslash\{\}exists w : \textbackslash\{\}mathsf\{BW\}(W),\textbackslash\{\} \textbackslash\{\}mathrm\{PresInv\}(W,w) \textbackslash\{\}wedge \textbackslash\{\}mathrm\{RealizesTarget\}(W,w).
\item Negative/barrier version: If 'MCSP witness family' cannot be realized with invariant 'MCSP complexity profile in nonuniform-circuits', the route cannot force 'There exist explicit"meta"problems (including MCSP variants) such that proving certain emph barely-superlinear lower bounds against restricted models (e.g. formulas) would imply major lower bounds for NP in much stronger models (including superpolynomial formula/circuit lower bounds, and in some parameter regimes consequences as strong as NP not in P/poly ).'.
\item Proof strategy family: threshold-localization
\item Negative route: Assume no threshold-localized version of 'MCSP witness family' preserves 'MCSP complexity profile in nonuniform-circuits' strongly enough for the restricted target.
\item Failure diagnosis: The route fails if no threshold-localized version of 'MCSP witness family' preserves 'MCSP complexity profile in nonuniform-circuits' strongly enough to force the restricted target.
\item Proof obligations: There exists a threshold regime in which 'MCSP witness family' already preserves 'MCSP complexity profile in nonuniform-circuits' well enough to imply a restricted version of 'There exist explicit"meta"problems (including MCSP variants) such that proving certain emph barely-superlinear lower bounds against restricted models (e.g. formulas) would imply major lower bounds for NP in much stronger models (including superpolynomial formula/circuit lower bounds, and in some parameter regimes consequences as strong as NP not in P/poly ).'.
\item Candidate lemmas: Realization theorem for MCSP witness family; Threshold localization lemma for MCSP witness family
\item Independent value: A threshold theorem for 'MCSP witness family' would be independently valuable because it localizes the first nontrivial regime where the route becomes mathematically real.
\item Expert gate: 6/6
\end{itemize}

\subsection*{conjpkg.disclayer-p-vs-np-pattern-regime-localization.09}
\begin{itemize}
\item Exact statement: No admissible 'MCSP obstruction family' can simultaneously preserve 'MCSP complexity profile in nonuniform-circuits' and realize 'In refined regimes, lower bounds for Gap-MCSP slightly beyond known frontiers (e.g. N\textasciicircum{} 1+varepsilon -type circuit lower bounds for suitable promise gaps and input length N=2\textasciicircum{}n ) imply major consequences such as NP not in P/poly.' under the current route commitments.
\item Restricted-regime mode: threshold-version
\item Bridge object: MCSP obstruction family (family of hard instances or obstructions forcing the target payoff)
\item Bridge definition: 'MCSP obstruction family' denotes the structured intermediate object required by this branch. It is meant to organize the source-side content 'A hardness magnification theorem has the form:"A barely-nontrivial lower bound for an explicit problem Q in a weak model implies a major lower bound/separation in a strong model."Typically, Q is a meta-problem (e.g. a Gap-MCSP variant) and"major"means consequences such as NP not in P/poly.'. It must preserve the invariant 'MCSP complexity profile in nonuniform-circuits'. Its operational role is 'The frontier repeatedly crosses model or threshold boundaries, so discovery should search for the smallest regime in which the route first becomes true.'.
\item First nontrivial lemma: There exists a threshold regime in which 'MCSP obstruction family' already preserves 'MCSP complexity profile in nonuniform-circuits' well enough to imply a restricted version of 'In refined regimes, lower bounds for Gap-MCSP slightly beyond known frontiers (e.g. N\textasciicircum{} 1+varepsilon -type circuit lower bounds for suitable promise gaps and input length N=2\textasciicircum{}n ) imply major consequences such as NP not in P/poly.'.
\item Quantifier signature: For every admissible 'MCSP obstruction family' on the route from 'st.def.hardness-magnification-phenomenon-programme-leve' to 'st.thm.oliveira-pich-santhanam-magnification-near-state'.
\item LaTeX theorem: \textbackslash\{\}forall W,\textbackslash\{\} \textbackslash\{\}mathrm\{Adm\}(W) \textbackslash\{\}to \textbackslash\{\}neg \textbackslash\{\}exists w : \textbackslash\{\}mathsf\{BW\}(W),\textbackslash\{\} \textbackslash\{\}mathrm\{PresInv\}(W,w) \textbackslash\{\}wedge \textbackslash\{\}mathrm\{RealizesTarget\}(W,w).
\item Negative/barrier version: If 'MCSP obstruction family' cannot be realized with invariant 'MCSP complexity profile in nonuniform-circuits', the route cannot force 'In refined regimes, lower bounds for Gap-MCSP slightly beyond known frontiers (e.g. N\textasciicircum{} 1+varepsilon -type circuit lower bounds for suitable promise gaps and input length N=2\textasciicircum{}n ) imply major consequences such as NP not in P/poly.'.
\item Proof strategy family: threshold-localization
\item Negative route: Assume no threshold-localized version of 'MCSP obstruction family' preserves 'MCSP complexity profile in nonuniform-circuits' strongly enough for the restricted target.
\item Failure diagnosis: The route fails if no threshold-localized version of 'MCSP obstruction family' preserves 'MCSP complexity profile in nonuniform-circuits' strongly enough to force the restricted target.
\item Proof obligations: There exists a threshold regime in which 'MCSP obstruction family' already preserves 'MCSP complexity profile in nonuniform-circuits' well enough to imply a restricted version of 'In refined regimes, lower bounds for Gap-MCSP slightly beyond known frontiers (e.g. N\textasciicircum{} 1+varepsilon -type circuit lower bounds for suitable promise gaps and input length N=2\textasciicircum{}n ) imply major consequences such as NP not in P/poly.'.
\item Candidate lemmas: Obstruction theorem for MCSP obstruction family; Threshold localization lemma for MCSP obstruction family
\item Independent value: A threshold theorem for 'MCSP obstruction family' would be independently valuable because it localizes the first nontrivial regime where the route becomes mathematically real.
\item Expert gate: 6/6
\end{itemize}

\subsection*{conjpkg.disclayer-p-vs-np-pattern-regime-localization.10}
\begin{itemize}
\item Exact statement: There exists an admissible 'MCSP obstruction family' whose realization internalizes 'A hardness magnification theorem has the form:"A barely-nontrivial lower bound for an explicit problem Q in a weak model implies a major lower bound/separation in a strong model."Typically, Q is a meta-problem (e.g. a Gap-MCSP variant) and"major"means consequences such as NP not in P/poly.' while preserving 'MCSP complexity profile in nonuniform-circuits' strongly enough to support 'In refined regimes, lower bounds for Gap-MCSP slightly beyond known frontiers (e.g. N\textasciicircum{} 1+varepsilon -type circuit lower bounds for suitable promise gaps and input length N=2\textasciicircum{}n ) imply major consequences such as NP not in P/poly.'.
\item Restricted-regime mode: threshold-version
\item Bridge object: MCSP obstruction family (family of hard instances or obstructions forcing the target payoff)
\item Bridge definition: 'MCSP obstruction family' denotes the structured intermediate object required by this branch. It is meant to organize the source-side content 'A hardness magnification theorem has the form:"A barely-nontrivial lower bound for an explicit problem Q in a weak model implies a major lower bound/separation in a strong model."Typically, Q is a meta-problem (e.g. a Gap-MCSP variant) and"major"means consequences such as NP not in P/poly.'. It must preserve the invariant 'MCSP complexity profile in nonuniform-circuits'. Its operational role is 'The frontier repeatedly crosses model or threshold boundaries, so discovery should search for the smallest regime in which the route first becomes true.'.
\item First nontrivial lemma: There exists a threshold regime in which 'MCSP obstruction family' already preserves 'MCSP complexity profile in nonuniform-circuits' well enough to imply a restricted version of 'In refined regimes, lower bounds for Gap-MCSP slightly beyond known frontiers (e.g. N\textasciicircum{} 1+varepsilon -type circuit lower bounds for suitable promise gaps and input length N=2\textasciicircum{}n ) imply major consequences such as NP not in P/poly.'.
\item Quantifier signature: For every admissible 'MCSP obstruction family' on the route from 'st.def.hardness-magnification-phenomenon-programme-leve' to 'st.thm.oliveira-pich-santhanam-magnification-near-state'.
\item LaTeX theorem: \textbackslash\{\}forall W,\textbackslash\{\} \textbackslash\{\}mathrm\{Adm\}(W) \textbackslash\{\}to \textbackslash\{\}exists w : \textbackslash\{\}mathsf\{BW\}(W),\textbackslash\{\} \textbackslash\{\}mathrm\{PresInv\}(W,w) \textbackslash\{\}wedge \textbackslash\{\}mathrm\{RealizesTarget\}(W,w).
\item Negative/barrier version: If 'MCSP obstruction family' cannot be realized with invariant 'MCSP complexity profile in nonuniform-circuits', the route cannot force 'In refined regimes, lower bounds for Gap-MCSP slightly beyond known frontiers (e.g. N\textasciicircum{} 1+varepsilon -type circuit lower bounds for suitable promise gaps and input length N=2\textasciicircum{}n ) imply major consequences such as NP not in P/poly.'.
\item Proof strategy family: threshold-localization
\item Negative route: Assume no threshold-localized version of 'MCSP obstruction family' preserves 'MCSP complexity profile in nonuniform-circuits' strongly enough for the restricted target.
\item Failure diagnosis: The route fails if no threshold-localized version of 'MCSP obstruction family' preserves 'MCSP complexity profile in nonuniform-circuits' strongly enough to force the restricted target.
\item Proof obligations: There exists a threshold regime in which 'MCSP obstruction family' already preserves 'MCSP complexity profile in nonuniform-circuits' well enough to imply a restricted version of 'In refined regimes, lower bounds for Gap-MCSP slightly beyond known frontiers (e.g. N\textasciicircum{} 1+varepsilon -type circuit lower bounds for suitable promise gaps and input length N=2\textasciicircum{}n ) imply major consequences such as NP not in P/poly.'.
\item Candidate lemmas: Realization theorem for MCSP obstruction family; Threshold localization lemma for MCSP obstruction family
\item Independent value: A threshold theorem for 'MCSP obstruction family' would be independently valuable because it localizes the first nontrivial regime where the route becomes mathematically real.
\item Expert gate: 6/6
\end{itemize}

\subsection*{conjpkg.disclayer-p-vs-np-family-abstraction-search.01}
\begin{itemize}
\item Exact statement: There exists an admissible 'MCSP witness family' such that 'A hardness magnification theorem has the form:"A barely-nontrivial lower bound for an explicit problem Q in a weak model implies a major lower bound/separation in a strong model."Typically, Q is a meta-problem (e.g. a Gap-MCSP variant) and"major"means consequences such as NP not in P/poly.' can be organized into a route preserving 'MCSP complexity profile in nonuniform-circuits' and thereby supports 'In refined regimes, lower bounds for Gap-MCSP slightly beyond known frontiers (e.g. N\textasciicircum{} 1+varepsilon -type circuit lower bounds for suitable promise gaps and input length N=2\textasciicircum{}n ) imply major consequences such as NP not in P/poly.'.
\item Restricted-regime mode: special-case
\item Bridge object: MCSP witness family (family of instances equipped with a monotone witness measure)
\item Bridge definition: 'MCSP witness family' denotes the structured intermediate object required by this branch. It is meant to organize the source-side content 'A hardness magnification theorem has the form:"A barely-nontrivial lower bound for an explicit problem Q in a weak model implies a major lower bound/separation in a strong model."Typically, Q is a meta-problem (e.g. a Gap-MCSP variant) and"major"means consequences such as NP not in P/poly.'. It must preserve the invariant 'MCSP complexity profile in nonuniform-circuits'. Its operational role is 'Find a sharper abstraction that turns the blocked route into a typed, inspectable theorem target.'.
\item First nontrivial lemma: For every admissible 'MCSP witness family', the first route-specific realization step preserves 'MCSP complexity profile in nonuniform-circuits' and prepares the passage from 'A hardness magnification theorem has the form:"A barely-nontrivial lower bound for an explicit problem Q in a weak model implies a major lower bound/separation in a strong model."Typically, Q is a meta-problem (e.g. a Gap-MCSP variant) and"major"means consequences such as NP not in P/poly.' to 'In refined regimes, lower bounds for Gap-MCSP slightly beyond known frontiers (e.g. N\textasciicircum{} 1+varepsilon -type circuit lower bounds for suitable promise gaps and input length N=2\textasciicircum{}n ) imply major consequences such as NP not in P/poly.'.
\item Quantifier signature: For every admissible 'MCSP witness family' on the route from 'st.def.hardness-magnification-phenomenon-programme-leve' to 'st.thm.oliveira-pich-santhanam-magnification-near-state'.
\item LaTeX theorem: \textbackslash\{\}forall W,\textbackslash\{\} \textbackslash\{\}mathrm\{Adm\}(W) \textbackslash\{\}to \textbackslash\{\}exists w : \textbackslash\{\}mathsf\{BW\}(W),\textbackslash\{\} \textbackslash\{\}mathrm\{PresInv\}(W,w) \textbackslash\{\}wedge \textbackslash\{\}mathrm\{RealizesTarget\}(W,w).
\item Negative/barrier version: If 'MCSP witness family' cannot be realized with invariant 'MCSP complexity profile in nonuniform-circuits', the route cannot force 'In refined regimes, lower bounds for Gap-MCSP slightly beyond known frontiers (e.g. N\textasciicircum{} 1+varepsilon -type circuit lower bounds for suitable promise gaps and input length N=2\textasciicircum{}n ) imply major consequences such as NP not in P/poly.'.
\item Proof strategy family: constructive-realization
\item Negative route: Assume 'MCSP witness family' cannot be realized while preserving 'MCSP complexity profile in nonuniform-circuits' strongly enough to activate the local target claim.
\item Failure diagnosis: The route fails if 'MCSP witness family' cannot be realized while preserving 'MCSP complexity profile in nonuniform-circuits'.
\item Proof obligations: For every admissible 'MCSP witness family', the first route-specific realization step preserves 'MCSP complexity profile in nonuniform-circuits' and prepares the passage from 'A hardness magnification theorem has the form:"A barely-nontrivial lower bound for an explicit problem Q in a weak model implies a major lower bound/separation in a strong model."Typically, Q is a meta-problem (e.g. a Gap-MCSP variant) and"major"means consequences such as NP not in P/poly.' to 'In refined regimes, lower bounds for Gap-MCSP slightly beyond known frontiers (e.g. N\textasciicircum{} 1+varepsilon -type circuit lower bounds for suitable promise gaps and input length N=2\textasciicircum{}n ) imply major consequences such as NP not in P/poly.'.
\item Candidate lemmas: Prove that a hardness-magnification witness family turns 'A barely-nontrivial lower bound for an explicit problem Q in a weak model' into 'major consequences such as NP not in P/poly' while preserving 'MCSP complexity profile in nonuniform-circuits'.
\item Independent value: A concrete special-case theorem for 'MCSP witness family' would be independently valuable because it converts a broad discovery programme into an inspectable mathematical subresult.
\item Expert gate: 6/6
\end{itemize}

\subsection*{conjpkg.disclayer-p-vs-np-family-abstraction-search.02}
\begin{itemize}
\item Exact statement: There exists an admissible 'MCSP witness family' such that 'A hardness magnification theorem has the form:"A barely-nontrivial lower bound for an explicit problem Q in a weak model implies a major lower bound/separation in a strong model."Typically, Q is a meta-problem (e.g. a Gap-MCSP variant) and"major"means consequences such as NP not in P/poly.' can be organized into a route preserving 'MCSP complexity profile in nonuniform-circuits' and thereby supports 'There exist explicit"meta"problems (including MCSP variants) such that proving certain emph barely-superlinear lower bounds against restricted models (e.g. formulas) would imply major lower bounds for NP in much stronger models (including superpolynomial formula/circuit lower bounds, and in some parameter regimes consequences as strong as NP not in P/poly ).'.
\item Restricted-regime mode: special-case
\item Bridge object: MCSP witness family (family of instances equipped with a monotone witness measure)
\item Bridge definition: 'MCSP witness family' denotes the structured intermediate object required by this branch. It is meant to organize the source-side content 'A hardness magnification theorem has the form:"A barely-nontrivial lower bound for an explicit problem Q in a weak model implies a major lower bound/separation in a strong model."Typically, Q is a meta-problem (e.g. a Gap-MCSP variant) and"major"means consequences such as NP not in P/poly.'. It must preserve the invariant 'MCSP complexity profile in nonuniform-circuits'. Its operational role is 'Find a sharper abstraction that turns the blocked route into a typed, inspectable theorem target.'.
\item First nontrivial lemma: For every admissible 'MCSP witness family', the first route-specific realization step preserves 'MCSP complexity profile in nonuniform-circuits' and prepares the passage from 'A hardness magnification theorem has the form:"A barely-nontrivial lower bound for an explicit problem Q in a weak model implies a major lower bound/separation in a strong model."Typically, Q is a meta-problem (e.g. a Gap-MCSP variant) and"major"means consequences such as NP not in P/poly.' to 'There exist explicit"meta"problems (including MCSP variants) such that proving certain emph barely-superlinear lower bounds against restricted models (e.g. formulas) would imply major lower bounds for NP in much stronger models (including superpolynomial formula/circuit lower bounds, and in some parameter regimes consequences as strong as NP not in P/poly ).'.
\item Quantifier signature: For every admissible 'MCSP witness family' on the route from 'st.def.hardness-magnification-phenomenon-programme-leve' to 'st.thm.oliveira-santhanam-hardness-magnification-for-na'.
\item LaTeX theorem: \textbackslash\{\}forall W,\textbackslash\{\} \textbackslash\{\}mathrm\{Adm\}(W) \textbackslash\{\}to \textbackslash\{\}exists w : \textbackslash\{\}mathsf\{BW\}(W),\textbackslash\{\} \textbackslash\{\}mathrm\{PresInv\}(W,w) \textbackslash\{\}wedge \textbackslash\{\}mathrm\{RealizesTarget\}(W,w).
\item Negative/barrier version: If 'MCSP witness family' cannot be realized with invariant 'MCSP complexity profile in nonuniform-circuits', the route cannot force 'There exist explicit"meta"problems (including MCSP variants) such that proving certain emph barely-superlinear lower bounds against restricted models (e.g. formulas) would imply major lower bounds for NP in much stronger models (including superpolynomial formula/circuit lower bounds, and in some parameter regimes consequences as strong as NP not in P/poly ).'.
\item Proof strategy family: constructive-realization
\item Negative route: Assume 'MCSP witness family' cannot be realized while preserving 'MCSP complexity profile in nonuniform-circuits' strongly enough to activate the local target claim.
\item Failure diagnosis: The route fails if 'MCSP witness family' cannot be realized while preserving 'MCSP complexity profile in nonuniform-circuits'.
\item Proof obligations: For every admissible 'MCSP witness family', the first route-specific realization step preserves 'MCSP complexity profile in nonuniform-circuits' and prepares the passage from 'A hardness magnification theorem has the form:"A barely-nontrivial lower bound for an explicit problem Q in a weak model implies a major lower bound/separation in a strong model."Typically, Q is a meta-problem (e.g. a Gap-MCSP variant) and"major"means consequences such as NP not in P/poly.' to 'There exist explicit"meta"problems (including MCSP variants) such that proving certain emph barely-superlinear lower bounds against restricted models (e.g. formulas) would imply major lower bounds for NP in much stronger models (including superpolynomial formula/circuit lower bounds, and in some parameter regimes consequences as strong as NP not in P/poly ).'.
\item Candidate lemmas: Prove that a hardness-magnification witness family preserves 'MCSP complexity profile in nonuniform-circuits' while transporting the route from the nonuniform-circuits regime to the formula-size regime strongly enough to derive 'major lower bounds for NP in much stronger models (including superpolynomial formula/circuit lower bounds, and in some parameter...'.
\item Independent value: A concrete special-case theorem for 'MCSP witness family' would be independently valuable because it converts a broad discovery programme into an inspectable mathematical subresult.
\item Expert gate: 6/6
\end{itemize}

\subsection*{conjpkg.disclayer-p-vs-np-family-abstraction-search.03}
\begin{itemize}
\item Exact statement: There exists an admissible 'MCSP obstruction family' such that 'A hardness magnification theorem has the form:"A barely-nontrivial lower bound for an explicit problem Q in a weak model implies a major lower bound/separation in a strong model."Typically, Q is a meta-problem (e.g. a Gap-MCSP variant) and"major"means consequences such as NP not in P/poly.' can be organized into a route preserving 'MCSP complexity profile in nonuniform-circuits' and thereby supports 'In refined regimes, lower bounds for Gap-MCSP slightly beyond known frontiers (e.g. N\textasciicircum{} 1+varepsilon -type circuit lower bounds for suitable promise gaps and input length N=2\textasciicircum{}n ) imply major consequences such as NP not in P/poly.'.
\item Restricted-regime mode: special-case
\item Bridge object: MCSP obstruction family (family of hard instances or obstructions forcing the target payoff)
\item Bridge definition: 'MCSP obstruction family' denotes the structured intermediate object required by this branch. It is meant to organize the source-side content 'A hardness magnification theorem has the form:"A barely-nontrivial lower bound for an explicit problem Q in a weak model implies a major lower bound/separation in a strong model."Typically, Q is a meta-problem (e.g. a Gap-MCSP variant) and"major"means consequences such as NP not in P/poly.'. It must preserve the invariant 'MCSP complexity profile in nonuniform-circuits'. Its operational role is 'Find a sharper abstraction that turns the blocked route into a typed, inspectable theorem target.'.
\item First nontrivial lemma: For every admissible 'MCSP obstruction family', the first route-specific realization step preserves 'MCSP complexity profile in nonuniform-circuits' and prepares the passage from 'A hardness magnification theorem has the form:"A barely-nontrivial lower bound for an explicit problem Q in a weak model implies a major lower bound/separation in a strong model."Typically, Q is a meta-problem (e.g. a Gap-MCSP variant) and"major"means consequences such as NP not in P/poly.' to 'In refined regimes, lower bounds for Gap-MCSP slightly beyond known frontiers (e.g. N\textasciicircum{} 1+varepsilon -type circuit lower bounds for suitable promise gaps and input length N=2\textasciicircum{}n ) imply major consequences such as NP not in P/poly.'.
\item Quantifier signature: For every admissible 'MCSP obstruction family' on the route from 'st.def.hardness-magnification-phenomenon-programme-leve' to 'st.thm.oliveira-pich-santhanam-magnification-near-state'.
\item LaTeX theorem: \textbackslash\{\}forall W,\textbackslash\{\} \textbackslash\{\}mathrm\{Adm\}(W) \textbackslash\{\}to \textbackslash\{\}exists w : \textbackslash\{\}mathsf\{BW\}(W),\textbackslash\{\} \textbackslash\{\}mathrm\{PresInv\}(W,w) \textbackslash\{\}wedge \textbackslash\{\}mathrm\{RealizesTarget\}(W,w).
\item Negative/barrier version: If 'MCSP obstruction family' cannot be realized with invariant 'MCSP complexity profile in nonuniform-circuits', the route cannot force 'In refined regimes, lower bounds for Gap-MCSP slightly beyond known frontiers (e.g. N\textasciicircum{} 1+varepsilon -type circuit lower bounds for suitable promise gaps and input length N=2\textasciicircum{}n ) imply major consequences such as NP not in P/poly.'.
\item Proof strategy family: constructive-realization
\item Negative route: Assume 'MCSP obstruction family' cannot be realized while preserving 'MCSP complexity profile in nonuniform-circuits' strongly enough to activate the local target claim.
\item Failure diagnosis: The route fails if 'MCSP obstruction family' cannot be realized while preserving 'MCSP complexity profile in nonuniform-circuits'.
\item Proof obligations: For every admissible 'MCSP obstruction family', the first route-specific realization step preserves 'MCSP complexity profile in nonuniform-circuits' and prepares the passage from 'A hardness magnification theorem has the form:"A barely-nontrivial lower bound for an explicit problem Q in a weak model implies a major lower bound/separation in a strong model."Typically, Q is a meta-problem (e.g. a Gap-MCSP variant) and"major"means consequences such as NP not in P/poly.' to 'In refined regimes, lower bounds for Gap-MCSP slightly beyond known frontiers (e.g. N\textasciicircum{} 1+varepsilon -type circuit lower bounds for suitable promise gaps and input length N=2\textasciicircum{}n ) imply major consequences such as NP not in P/poly.'.
\item Candidate lemmas: Prove that a hardness-magnification witness family turns 'A barely-nontrivial lower bound for an explicit problem Q in a weak model' into 'major consequences such as NP not in P/poly' while preserving 'MCSP complexity profile in nonuniform-circuits'.
\item Independent value: A concrete special-case theorem for 'MCSP obstruction family' would be independently valuable because it converts a broad discovery programme into an inspectable mathematical subresult.
\item Expert gate: 6/6
\end{itemize}

\subsection*{conjpkg.disclayer-p-vs-np-family-abstraction-search.04}
\begin{itemize}
\item Exact statement: No admissible 'MCSP witness family' can simultaneously preserve 'MCSP complexity profile in nonuniform-circuits' and realize 'In refined regimes, lower bounds for Gap-MCSP slightly beyond known frontiers (e.g. N\textasciicircum{} 1+varepsilon -type circuit lower bounds for suitable promise gaps and input length N=2\textasciicircum{}n ) imply major consequences such as NP not in P/poly.' under the current route commitments.
\item Restricted-regime mode: special-case
\item Bridge object: MCSP witness family (family of instances equipped with a monotone witness measure)
\item Bridge definition: 'MCSP witness family' denotes the structured intermediate object required by this branch. It is meant to organize the source-side content 'A hardness magnification theorem has the form:"A barely-nontrivial lower bound for an explicit problem Q in a weak model implies a major lower bound/separation in a strong model."Typically, Q is a meta-problem (e.g. a Gap-MCSP variant) and"major"means consequences such as NP not in P/poly.'. It must preserve the invariant 'MCSP complexity profile in nonuniform-circuits'. Its operational role is 'Find a sharper abstraction that turns the blocked route into a typed, inspectable theorem target.'.
\item First nontrivial lemma: For every admissible 'MCSP witness family', the first route-specific realization step preserves 'MCSP complexity profile in nonuniform-circuits' and prepares the passage from 'A hardness magnification theorem has the form:"A barely-nontrivial lower bound for an explicit problem Q in a weak model implies a major lower bound/separation in a strong model."Typically, Q is a meta-problem (e.g. a Gap-MCSP variant) and"major"means consequences such as NP not in P/poly.' to 'In refined regimes, lower bounds for Gap-MCSP slightly beyond known frontiers (e.g. N\textasciicircum{} 1+varepsilon -type circuit lower bounds for suitable promise gaps and input length N=2\textasciicircum{}n ) imply major consequences such as NP not in P/poly.'.
\item Quantifier signature: For every admissible 'MCSP witness family' on the route from 'st.def.hardness-magnification-phenomenon-programme-leve' to 'st.thm.oliveira-pich-santhanam-magnification-near-state'.
\item LaTeX theorem: \textbackslash\{\}forall W,\textbackslash\{\} \textbackslash\{\}mathrm\{Adm\}(W) \textbackslash\{\}to \textbackslash\{\}neg \textbackslash\{\}exists w : \textbackslash\{\}mathsf\{BW\}(W),\textbackslash\{\} \textbackslash\{\}mathrm\{PresInv\}(W,w) \textbackslash\{\}wedge \textbackslash\{\}mathrm\{RealizesTarget\}(W,w).
\item Negative/barrier version: If 'MCSP witness family' cannot be realized with invariant 'MCSP complexity profile in nonuniform-circuits', the route cannot force 'In refined regimes, lower bounds for Gap-MCSP slightly beyond known frontiers (e.g. N\textasciicircum{} 1+varepsilon -type circuit lower bounds for suitable promise gaps and input length N=2\textasciicircum{}n ) imply major consequences such as NP not in P/poly.'.
\item Proof strategy family: constructive-realization
\item Negative route: Assume 'MCSP witness family' cannot be realized while preserving 'MCSP complexity profile in nonuniform-circuits' strongly enough to activate the local target claim.
\item Failure diagnosis: The route fails if 'MCSP witness family' cannot be realized while preserving 'MCSP complexity profile in nonuniform-circuits'.
\item Proof obligations: For every admissible 'MCSP witness family', the first route-specific realization step preserves 'MCSP complexity profile in nonuniform-circuits' and prepares the passage from 'A hardness magnification theorem has the form:"A barely-nontrivial lower bound for an explicit problem Q in a weak model implies a major lower bound/separation in a strong model."Typically, Q is a meta-problem (e.g. a Gap-MCSP variant) and"major"means consequences such as NP not in P/poly.' to 'In refined regimes, lower bounds for Gap-MCSP slightly beyond known frontiers (e.g. N\textasciicircum{} 1+varepsilon -type circuit lower bounds for suitable promise gaps and input length N=2\textasciicircum{}n ) imply major consequences such as NP not in P/poly.'.
\item Candidate lemmas: Obstruction theorem for MCSP witness family
\item Independent value: A concrete special-case theorem for 'MCSP witness family' would be independently valuable because it converts a broad discovery programme into an inspectable mathematical subresult.
\item Expert gate: 6/6
\end{itemize}

\subsection*{conjpkg.disclayer-p-vs-np-family-abstraction-search.05}
\begin{itemize}
\item Exact statement: There exists an admissible 'MCSP witness family' whose realization internalizes 'A hardness magnification theorem has the form:"A barely-nontrivial lower bound for an explicit problem Q in a weak model implies a major lower bound/separation in a strong model."Typically, Q is a meta-problem (e.g. a Gap-MCSP variant) and"major"means consequences such as NP not in P/poly.' while preserving 'MCSP complexity profile in nonuniform-circuits' strongly enough to support 'In refined regimes, lower bounds for Gap-MCSP slightly beyond known frontiers (e.g. N\textasciicircum{} 1+varepsilon -type circuit lower bounds for suitable promise gaps and input length N=2\textasciicircum{}n ) imply major consequences such as NP not in P/poly.'.
\item Restricted-regime mode: special-case
\item Bridge object: MCSP witness family (family of instances equipped with a monotone witness measure)
\item Bridge definition: 'MCSP witness family' denotes the structured intermediate object required by this branch. It is meant to organize the source-side content 'A hardness magnification theorem has the form:"A barely-nontrivial lower bound for an explicit problem Q in a weak model implies a major lower bound/separation in a strong model."Typically, Q is a meta-problem (e.g. a Gap-MCSP variant) and"major"means consequences such as NP not in P/poly.'. It must preserve the invariant 'MCSP complexity profile in nonuniform-circuits'. Its operational role is 'Find a sharper abstraction that turns the blocked route into a typed, inspectable theorem target.'.
\item First nontrivial lemma: For every admissible 'MCSP witness family', the first route-specific realization step preserves 'MCSP complexity profile in nonuniform-circuits' and prepares the passage from 'A hardness magnification theorem has the form:"A barely-nontrivial lower bound for an explicit problem Q in a weak model implies a major lower bound/separation in a strong model."Typically, Q is a meta-problem (e.g. a Gap-MCSP variant) and"major"means consequences such as NP not in P/poly.' to 'In refined regimes, lower bounds for Gap-MCSP slightly beyond known frontiers (e.g. N\textasciicircum{} 1+varepsilon -type circuit lower bounds for suitable promise gaps and input length N=2\textasciicircum{}n ) imply major consequences such as NP not in P/poly.'.
\item Quantifier signature: For every admissible 'MCSP witness family' on the route from 'st.def.hardness-magnification-phenomenon-programme-leve' to 'st.thm.oliveira-pich-santhanam-magnification-near-state'.
\item LaTeX theorem: \textbackslash\{\}forall W,\textbackslash\{\} \textbackslash\{\}mathrm\{Adm\}(W) \textbackslash\{\}to \textbackslash\{\}exists w : \textbackslash\{\}mathsf\{BW\}(W),\textbackslash\{\} \textbackslash\{\}mathrm\{PresInv\}(W,w) \textbackslash\{\}wedge \textbackslash\{\}mathrm\{RealizesTarget\}(W,w).
\item Negative/barrier version: If 'MCSP witness family' cannot be realized with invariant 'MCSP complexity profile in nonuniform-circuits', the route cannot force 'In refined regimes, lower bounds for Gap-MCSP slightly beyond known frontiers (e.g. N\textasciicircum{} 1+varepsilon -type circuit lower bounds for suitable promise gaps and input length N=2\textasciicircum{}n ) imply major consequences such as NP not in P/poly.'.
\item Proof strategy family: constructive-realization
\item Negative route: Assume 'MCSP witness family' cannot be realized while preserving 'MCSP complexity profile in nonuniform-circuits' strongly enough to activate the local target claim.
\item Failure diagnosis: The route fails if 'MCSP witness family' cannot be realized while preserving 'MCSP complexity profile in nonuniform-circuits'.
\item Proof obligations: For every admissible 'MCSP witness family', the first route-specific realization step preserves 'MCSP complexity profile in nonuniform-circuits' and prepares the passage from 'A hardness magnification theorem has the form:"A barely-nontrivial lower bound for an explicit problem Q in a weak model implies a major lower bound/separation in a strong model."Typically, Q is a meta-problem (e.g. a Gap-MCSP variant) and"major"means consequences such as NP not in P/poly.' to 'In refined regimes, lower bounds for Gap-MCSP slightly beyond known frontiers (e.g. N\textasciicircum{} 1+varepsilon -type circuit lower bounds for suitable promise gaps and input length N=2\textasciicircum{}n ) imply major consequences such as NP not in P/poly.'.
\item Candidate lemmas: Realization theorem for MCSP witness family
\item Independent value: A concrete special-case theorem for 'MCSP witness family' would be independently valuable because it converts a broad discovery programme into an inspectable mathematical subresult.
\item Expert gate: 6/6
\end{itemize}

\subsection*{conjpkg.disclayer-p-vs-np-family-abstraction-search.06}
\begin{itemize}
\item Exact statement: For every admissible 'MCSP witness family', there is a transport law from 'nonuniform-circuits' to 'formula-size' that preserves 'MCSP complexity profile in nonuniform-circuits' and retains enough strength to activate 'There exist explicit"meta"problems (including MCSP variants) such that proving certain emph barely-superlinear lower bounds against restricted models (e.g. formulas) would imply major lower bounds for NP in much stronger models (including superpolynomial formula/circuit lower bounds, and in some parameter regimes consequences as strong as NP not in P/poly ).'.
\item Restricted-regime mode: special-case
\item Bridge object: MCSP witness family (family of instances equipped with a monotone witness measure)
\item Bridge definition: 'MCSP witness family' denotes the structured intermediate object required by this branch. It is meant to organize the source-side content 'A hardness magnification theorem has the form:"A barely-nontrivial lower bound for an explicit problem Q in a weak model implies a major lower bound/separation in a strong model."Typically, Q is a meta-problem (e.g. a Gap-MCSP variant) and"major"means consequences such as NP not in P/poly.'. It must preserve the invariant 'MCSP complexity profile in nonuniform-circuits'. Its operational role is 'Find a sharper abstraction that turns the blocked route into a typed, inspectable theorem target.'.
\item First nontrivial lemma: For every admissible 'MCSP witness family', the first route-specific realization step preserves 'MCSP complexity profile in nonuniform-circuits' and prepares the passage from 'A hardness magnification theorem has the form:"A barely-nontrivial lower bound for an explicit problem Q in a weak model implies a major lower bound/separation in a strong model."Typically, Q is a meta-problem (e.g. a Gap-MCSP variant) and"major"means consequences such as NP not in P/poly.' to 'There exist explicit"meta"problems (including MCSP variants) such that proving certain emph barely-superlinear lower bounds against restricted models (e.g. formulas) would imply major lower bounds for NP in much stronger models (including superpolynomial formula/circuit lower bounds, and in some parameter regimes consequences as strong as NP not in P/poly ).'.
\item Quantifier signature: For every admissible 'MCSP witness family' on the route from 'st.def.hardness-magnification-phenomenon-programme-leve' to 'st.thm.oliveira-santhanam-hardness-magnification-for-na'.
\item LaTeX theorem: \textbackslash\{\}forall W,\textbackslash\{\} \textbackslash\{\}mathrm\{Adm\}(W) \textbackslash\{\}to \textbackslash\{\}exists w : \textbackslash\{\}mathsf\{BW\}(W),\textbackslash\{\} \textbackslash\{\}mathrm\{PresInv\}(W,w) \textbackslash\{\}wedge \textbackslash\{\}mathrm\{RealizesTarget\}(W,w).
\item Negative/barrier version: If 'MCSP witness family' cannot be realized with invariant 'MCSP complexity profile in nonuniform-circuits', the route cannot force 'There exist explicit"meta"problems (including MCSP variants) such that proving certain emph barely-superlinear lower bounds against restricted models (e.g. formulas) would imply major lower bounds for NP in much stronger models (including superpolynomial formula/circuit lower bounds, and in some parameter regimes consequences as strong as NP not in P/poly ).'.
\item Proof strategy family: constructive-realization
\item Negative route: Assume 'MCSP witness family' cannot be realized while preserving 'MCSP complexity profile in nonuniform-circuits' strongly enough to activate the local target claim.
\item Failure diagnosis: The route fails if 'MCSP witness family' cannot be realized while preserving 'MCSP complexity profile in nonuniform-circuits'.
\item Proof obligations: For every admissible 'MCSP witness family', the first route-specific realization step preserves 'MCSP complexity profile in nonuniform-circuits' and prepares the passage from 'A hardness magnification theorem has the form:"A barely-nontrivial lower bound for an explicit problem Q in a weak model implies a major lower bound/separation in a strong model."Typically, Q is a meta-problem (e.g. a Gap-MCSP variant) and"major"means consequences such as NP not in P/poly.' to 'There exist explicit"meta"problems (including MCSP variants) such that proving certain emph barely-superlinear lower bounds against restricted models (e.g. formulas) would imply major lower bounds for NP in much stronger models (including superpolynomial formula/circuit lower bounds, and in some parameter regimes consequences as strong as NP not in P/poly ).'.
\item Candidate lemmas: Transport theorem for MCSP witness family
\item Independent value: A concrete special-case theorem for 'MCSP witness family' would be independently valuable because it converts a broad discovery programme into an inspectable mathematical subresult.
\item Expert gate: 6/6
\end{itemize}

\subsection*{conjpkg.disclayer-p-vs-np-family-abstraction-search.07}
\begin{itemize}
\item Exact statement: No admissible 'MCSP witness family' can simultaneously preserve 'MCSP complexity profile in nonuniform-circuits' and realize 'There exist explicit"meta"problems (including MCSP variants) such that proving certain emph barely-superlinear lower bounds against restricted models (e.g. formulas) would imply major lower bounds for NP in much stronger models (including superpolynomial formula/circuit lower bounds, and in some parameter regimes consequences as strong as NP not in P/poly ).' under the current route commitments.
\item Restricted-regime mode: special-case
\item Bridge object: MCSP witness family (family of instances equipped with a monotone witness measure)
\item Bridge definition: 'MCSP witness family' denotes the structured intermediate object required by this branch. It is meant to organize the source-side content 'A hardness magnification theorem has the form:"A barely-nontrivial lower bound for an explicit problem Q in a weak model implies a major lower bound/separation in a strong model."Typically, Q is a meta-problem (e.g. a Gap-MCSP variant) and"major"means consequences such as NP not in P/poly.'. It must preserve the invariant 'MCSP complexity profile in nonuniform-circuits'. Its operational role is 'Find a sharper abstraction that turns the blocked route into a typed, inspectable theorem target.'.
\item First nontrivial lemma: For every admissible 'MCSP witness family', the first route-specific realization step preserves 'MCSP complexity profile in nonuniform-circuits' and prepares the passage from 'A hardness magnification theorem has the form:"A barely-nontrivial lower bound for an explicit problem Q in a weak model implies a major lower bound/separation in a strong model."Typically, Q is a meta-problem (e.g. a Gap-MCSP variant) and"major"means consequences such as NP not in P/poly.' to 'There exist explicit"meta"problems (including MCSP variants) such that proving certain emph barely-superlinear lower bounds against restricted models (e.g. formulas) would imply major lower bounds for NP in much stronger models (including superpolynomial formula/circuit lower bounds, and in some parameter regimes consequences as strong as NP not in P/poly ).'.
\item Quantifier signature: For every admissible 'MCSP witness family' on the route from 'st.def.hardness-magnification-phenomenon-programme-leve' to 'st.thm.oliveira-santhanam-hardness-magnification-for-na'.
\item LaTeX theorem: \textbackslash\{\}forall W,\textbackslash\{\} \textbackslash\{\}mathrm\{Adm\}(W) \textbackslash\{\}to \textbackslash\{\}neg \textbackslash\{\}exists w : \textbackslash\{\}mathsf\{BW\}(W),\textbackslash\{\} \textbackslash\{\}mathrm\{PresInv\}(W,w) \textbackslash\{\}wedge \textbackslash\{\}mathrm\{RealizesTarget\}(W,w).
\item Negative/barrier version: If 'MCSP witness family' cannot be realized with invariant 'MCSP complexity profile in nonuniform-circuits', the route cannot force 'There exist explicit"meta"problems (including MCSP variants) such that proving certain emph barely-superlinear lower bounds against restricted models (e.g. formulas) would imply major lower bounds for NP in much stronger models (including superpolynomial formula/circuit lower bounds, and in some parameter regimes consequences as strong as NP not in P/poly ).'.
\item Proof strategy family: constructive-realization
\item Negative route: Assume 'MCSP witness family' cannot be realized while preserving 'MCSP complexity profile in nonuniform-circuits' strongly enough to activate the local target claim.
\item Failure diagnosis: The route fails if 'MCSP witness family' cannot be realized while preserving 'MCSP complexity profile in nonuniform-circuits'.
\item Proof obligations: For every admissible 'MCSP witness family', the first route-specific realization step preserves 'MCSP complexity profile in nonuniform-circuits' and prepares the passage from 'A hardness magnification theorem has the form:"A barely-nontrivial lower bound for an explicit problem Q in a weak model implies a major lower bound/separation in a strong model."Typically, Q is a meta-problem (e.g. a Gap-MCSP variant) and"major"means consequences such as NP not in P/poly.' to 'There exist explicit"meta"problems (including MCSP variants) such that proving certain emph barely-superlinear lower bounds against restricted models (e.g. formulas) would imply major lower bounds for NP in much stronger models (including superpolynomial formula/circuit lower bounds, and in some parameter regimes consequences as strong as NP not in P/poly ).'.
\item Candidate lemmas: Obstruction theorem for MCSP witness family
\item Independent value: A concrete special-case theorem for 'MCSP witness family' would be independently valuable because it converts a broad discovery programme into an inspectable mathematical subresult.
\item Expert gate: 6/6
\end{itemize}

\subsection*{conjpkg.disclayer-p-vs-np-family-abstraction-search.08}
\begin{itemize}
\item Exact statement: There exists an admissible 'MCSP witness family' whose realization internalizes 'A hardness magnification theorem has the form:"A barely-nontrivial lower bound for an explicit problem Q in a weak model implies a major lower bound/separation in a strong model."Typically, Q is a meta-problem (e.g. a Gap-MCSP variant) and"major"means consequences such as NP not in P/poly.' while preserving 'MCSP complexity profile in nonuniform-circuits' strongly enough to support 'There exist explicit"meta"problems (including MCSP variants) such that proving certain emph barely-superlinear lower bounds against restricted models (e.g. formulas) would imply major lower bounds for NP in much stronger models (including superpolynomial formula/circuit lower bounds, and in some parameter regimes consequences as strong as NP not in P/poly ).'.
\item Restricted-regime mode: special-case
\item Bridge object: MCSP witness family (family of instances equipped with a monotone witness measure)
\item Bridge definition: 'MCSP witness family' denotes the structured intermediate object required by this branch. It is meant to organize the source-side content 'A hardness magnification theorem has the form:"A barely-nontrivial lower bound for an explicit problem Q in a weak model implies a major lower bound/separation in a strong model."Typically, Q is a meta-problem (e.g. a Gap-MCSP variant) and"major"means consequences such as NP not in P/poly.'. It must preserve the invariant 'MCSP complexity profile in nonuniform-circuits'. Its operational role is 'Find a sharper abstraction that turns the blocked route into a typed, inspectable theorem target.'.
\item First nontrivial lemma: For every admissible 'MCSP witness family', the first route-specific realization step preserves 'MCSP complexity profile in nonuniform-circuits' and prepares the passage from 'A hardness magnification theorem has the form:"A barely-nontrivial lower bound for an explicit problem Q in a weak model implies a major lower bound/separation in a strong model."Typically, Q is a meta-problem (e.g. a Gap-MCSP variant) and"major"means consequences such as NP not in P/poly.' to 'There exist explicit"meta"problems (including MCSP variants) such that proving certain emph barely-superlinear lower bounds against restricted models (e.g. formulas) would imply major lower bounds for NP in much stronger models (including superpolynomial formula/circuit lower bounds, and in some parameter regimes consequences as strong as NP not in P/poly ).'.
\item Quantifier signature: For every admissible 'MCSP witness family' on the route from 'st.def.hardness-magnification-phenomenon-programme-leve' to 'st.thm.oliveira-santhanam-hardness-magnification-for-na'.
\item LaTeX theorem: \textbackslash\{\}forall W,\textbackslash\{\} \textbackslash\{\}mathrm\{Adm\}(W) \textbackslash\{\}to \textbackslash\{\}exists w : \textbackslash\{\}mathsf\{BW\}(W),\textbackslash\{\} \textbackslash\{\}mathrm\{PresInv\}(W,w) \textbackslash\{\}wedge \textbackslash\{\}mathrm\{RealizesTarget\}(W,w).
\item Negative/barrier version: If 'MCSP witness family' cannot be realized with invariant 'MCSP complexity profile in nonuniform-circuits', the route cannot force 'There exist explicit"meta"problems (including MCSP variants) such that proving certain emph barely-superlinear lower bounds against restricted models (e.g. formulas) would imply major lower bounds for NP in much stronger models (including superpolynomial formula/circuit lower bounds, and in some parameter regimes consequences as strong as NP not in P/poly ).'.
\item Proof strategy family: constructive-realization
\item Negative route: Assume 'MCSP witness family' cannot be realized while preserving 'MCSP complexity profile in nonuniform-circuits' strongly enough to activate the local target claim.
\item Failure diagnosis: The route fails if 'MCSP witness family' cannot be realized while preserving 'MCSP complexity profile in nonuniform-circuits'.
\item Proof obligations: For every admissible 'MCSP witness family', the first route-specific realization step preserves 'MCSP complexity profile in nonuniform-circuits' and prepares the passage from 'A hardness magnification theorem has the form:"A barely-nontrivial lower bound for an explicit problem Q in a weak model implies a major lower bound/separation in a strong model."Typically, Q is a meta-problem (e.g. a Gap-MCSP variant) and"major"means consequences such as NP not in P/poly.' to 'There exist explicit"meta"problems (including MCSP variants) such that proving certain emph barely-superlinear lower bounds against restricted models (e.g. formulas) would imply major lower bounds for NP in much stronger models (including superpolynomial formula/circuit lower bounds, and in some parameter regimes consequences as strong as NP not in P/poly ).'.
\item Candidate lemmas: Realization theorem for MCSP witness family
\item Independent value: A concrete special-case theorem for 'MCSP witness family' would be independently valuable because it converts a broad discovery programme into an inspectable mathematical subresult.
\item Expert gate: 6/6
\end{itemize}

\subsection*{conjpkg.disclayer-p-vs-np-family-abstraction-search.09}
\begin{itemize}
\item Exact statement: No admissible 'MCSP obstruction family' can simultaneously preserve 'MCSP complexity profile in nonuniform-circuits' and realize 'In refined regimes, lower bounds for Gap-MCSP slightly beyond known frontiers (e.g. N\textasciicircum{} 1+varepsilon -type circuit lower bounds for suitable promise gaps and input length N=2\textasciicircum{}n ) imply major consequences such as NP not in P/poly.' under the current route commitments.
\item Restricted-regime mode: special-case
\item Bridge object: MCSP obstruction family (family of hard instances or obstructions forcing the target payoff)
\item Bridge definition: 'MCSP obstruction family' denotes the structured intermediate object required by this branch. It is meant to organize the source-side content 'A hardness magnification theorem has the form:"A barely-nontrivial lower bound for an explicit problem Q in a weak model implies a major lower bound/separation in a strong model."Typically, Q is a meta-problem (e.g. a Gap-MCSP variant) and"major"means consequences such as NP not in P/poly.'. It must preserve the invariant 'MCSP complexity profile in nonuniform-circuits'. Its operational role is 'Find a sharper abstraction that turns the blocked route into a typed, inspectable theorem target.'.
\item First nontrivial lemma: For every admissible 'MCSP obstruction family', the first route-specific realization step preserves 'MCSP complexity profile in nonuniform-circuits' and prepares the passage from 'A hardness magnification theorem has the form:"A barely-nontrivial lower bound for an explicit problem Q in a weak model implies a major lower bound/separation in a strong model."Typically, Q is a meta-problem (e.g. a Gap-MCSP variant) and"major"means consequences such as NP not in P/poly.' to 'In refined regimes, lower bounds for Gap-MCSP slightly beyond known frontiers (e.g. N\textasciicircum{} 1+varepsilon -type circuit lower bounds for suitable promise gaps and input length N=2\textasciicircum{}n ) imply major consequences such as NP not in P/poly.'.
\item Quantifier signature: For every admissible 'MCSP obstruction family' on the route from 'st.def.hardness-magnification-phenomenon-programme-leve' to 'st.thm.oliveira-pich-santhanam-magnification-near-state'.
\item LaTeX theorem: \textbackslash\{\}forall W,\textbackslash\{\} \textbackslash\{\}mathrm\{Adm\}(W) \textbackslash\{\}to \textbackslash\{\}neg \textbackslash\{\}exists w : \textbackslash\{\}mathsf\{BW\}(W),\textbackslash\{\} \textbackslash\{\}mathrm\{PresInv\}(W,w) \textbackslash\{\}wedge \textbackslash\{\}mathrm\{RealizesTarget\}(W,w).
\item Negative/barrier version: If 'MCSP obstruction family' cannot be realized with invariant 'MCSP complexity profile in nonuniform-circuits', the route cannot force 'In refined regimes, lower bounds for Gap-MCSP slightly beyond known frontiers (e.g. N\textasciicircum{} 1+varepsilon -type circuit lower bounds for suitable promise gaps and input length N=2\textasciicircum{}n ) imply major consequences such as NP not in P/poly.'.
\item Proof strategy family: constructive-realization
\item Negative route: Assume 'MCSP obstruction family' cannot be realized while preserving 'MCSP complexity profile in nonuniform-circuits' strongly enough to activate the local target claim.
\item Failure diagnosis: The route fails if 'MCSP obstruction family' cannot be realized while preserving 'MCSP complexity profile in nonuniform-circuits'.
\item Proof obligations: For every admissible 'MCSP obstruction family', the first route-specific realization step preserves 'MCSP complexity profile in nonuniform-circuits' and prepares the passage from 'A hardness magnification theorem has the form:"A barely-nontrivial lower bound for an explicit problem Q in a weak model implies a major lower bound/separation in a strong model."Typically, Q is a meta-problem (e.g. a Gap-MCSP variant) and"major"means consequences such as NP not in P/poly.' to 'In refined regimes, lower bounds for Gap-MCSP slightly beyond known frontiers (e.g. N\textasciicircum{} 1+varepsilon -type circuit lower bounds for suitable promise gaps and input length N=2\textasciicircum{}n ) imply major consequences such as NP not in P/poly.'.
\item Candidate lemmas: Obstruction theorem for MCSP obstruction family
\item Independent value: A concrete special-case theorem for 'MCSP obstruction family' would be independently valuable because it converts a broad discovery programme into an inspectable mathematical subresult.
\item Expert gate: 6/6
\end{itemize}

\subsection*{conjpkg.disclayer-p-vs-np-family-abstraction-search.10}
\begin{itemize}
\item Exact statement: There exists an admissible 'MCSP obstruction family' whose realization internalizes 'A hardness magnification theorem has the form:"A barely-nontrivial lower bound for an explicit problem Q in a weak model implies a major lower bound/separation in a strong model."Typically, Q is a meta-problem (e.g. a Gap-MCSP variant) and"major"means consequences such as NP not in P/poly.' while preserving 'MCSP complexity profile in nonuniform-circuits' strongly enough to support 'In refined regimes, lower bounds for Gap-MCSP slightly beyond known frontiers (e.g. N\textasciicircum{} 1+varepsilon -type circuit lower bounds for suitable promise gaps and input length N=2\textasciicircum{}n ) imply major consequences such as NP not in P/poly.'.
\item Restricted-regime mode: special-case
\item Bridge object: MCSP obstruction family (family of hard instances or obstructions forcing the target payoff)
\item Bridge definition: 'MCSP obstruction family' denotes the structured intermediate object required by this branch. It is meant to organize the source-side content 'A hardness magnification theorem has the form:"A barely-nontrivial lower bound for an explicit problem Q in a weak model implies a major lower bound/separation in a strong model."Typically, Q is a meta-problem (e.g. a Gap-MCSP variant) and"major"means consequences such as NP not in P/poly.'. It must preserve the invariant 'MCSP complexity profile in nonuniform-circuits'. Its operational role is 'Find a sharper abstraction that turns the blocked route into a typed, inspectable theorem target.'.
\item First nontrivial lemma: For every admissible 'MCSP obstruction family', the first route-specific realization step preserves 'MCSP complexity profile in nonuniform-circuits' and prepares the passage from 'A hardness magnification theorem has the form:"A barely-nontrivial lower bound for an explicit problem Q in a weak model implies a major lower bound/separation in a strong model."Typically, Q is a meta-problem (e.g. a Gap-MCSP variant) and"major"means consequences such as NP not in P/poly.' to 'In refined regimes, lower bounds for Gap-MCSP slightly beyond known frontiers (e.g. N\textasciicircum{} 1+varepsilon -type circuit lower bounds for suitable promise gaps and input length N=2\textasciicircum{}n ) imply major consequences such as NP not in P/poly.'.
\item Quantifier signature: For every admissible 'MCSP obstruction family' on the route from 'st.def.hardness-magnification-phenomenon-programme-leve' to 'st.thm.oliveira-pich-santhanam-magnification-near-state'.
\item LaTeX theorem: \textbackslash\{\}forall W,\textbackslash\{\} \textbackslash\{\}mathrm\{Adm\}(W) \textbackslash\{\}to \textbackslash\{\}exists w : \textbackslash\{\}mathsf\{BW\}(W),\textbackslash\{\} \textbackslash\{\}mathrm\{PresInv\}(W,w) \textbackslash\{\}wedge \textbackslash\{\}mathrm\{RealizesTarget\}(W,w).
\item Negative/barrier version: If 'MCSP obstruction family' cannot be realized with invariant 'MCSP complexity profile in nonuniform-circuits', the route cannot force 'In refined regimes, lower bounds for Gap-MCSP slightly beyond known frontiers (e.g. N\textasciicircum{} 1+varepsilon -type circuit lower bounds for suitable promise gaps and input length N=2\textasciicircum{}n ) imply major consequences such as NP not in P/poly.'.
\item Proof strategy family: constructive-realization
\item Negative route: Assume 'MCSP obstruction family' cannot be realized while preserving 'MCSP complexity profile in nonuniform-circuits' strongly enough to activate the local target claim.
\item Failure diagnosis: The route fails if 'MCSP obstruction family' cannot be realized while preserving 'MCSP complexity profile in nonuniform-circuits'.
\item Proof obligations: For every admissible 'MCSP obstruction family', the first route-specific realization step preserves 'MCSP complexity profile in nonuniform-circuits' and prepares the passage from 'A hardness magnification theorem has the form:"A barely-nontrivial lower bound for an explicit problem Q in a weak model implies a major lower bound/separation in a strong model."Typically, Q is a meta-problem (e.g. a Gap-MCSP variant) and"major"means consequences such as NP not in P/poly.' to 'In refined regimes, lower bounds for Gap-MCSP slightly beyond known frontiers (e.g. N\textasciicircum{} 1+varepsilon -type circuit lower bounds for suitable promise gaps and input length N=2\textasciicircum{}n ) imply major consequences such as NP not in P/poly.'.
\item Candidate lemmas: Realization theorem for MCSP obstruction family
\item Independent value: A concrete special-case theorem for 'MCSP obstruction family' would be independently valuable because it converts a broad discovery programme into an inspectable mathematical subresult.
\item Expert gate: 6/6
\end{itemize}

\subsection*{conjpkg.disclayer-p-vs-np-pattern-invariant-discovery.01}
\begin{itemize}
\item Exact statement: There exists an admissible 'MCSP witness family' such that 'A hardness magnification theorem has the form:"A barely-nontrivial lower bound for an explicit problem Q in a weak model implies a major lower bound/separation in a strong model."Typically, Q is a meta-problem (e.g. a Gap-MCSP variant) and"major"means consequences such as NP not in P/poly.' can be organized into a route preserving 'MCSP complexity profile in nonuniform-circuits' and thereby supports 'In refined regimes, lower bounds for Gap-MCSP slightly beyond known frontiers (e.g. N\textasciicircum{} 1+varepsilon -type circuit lower bounds for suitable promise gaps and input length N=2\textasciicircum{}n ) imply major consequences such as NP not in P/poly.'.
\item Restricted-regime mode: special-case
\item Bridge object: MCSP witness family (family of instances equipped with a monotone witness measure)
\item Bridge definition: 'MCSP witness family' denotes the structured intermediate object required by this branch. It is meant to organize the source-side content 'A hardness magnification theorem has the form:"A barely-nontrivial lower bound for an explicit problem Q in a weak model implies a major lower bound/separation in a strong model."Typically, Q is a meta-problem (e.g. a Gap-MCSP variant) and"major"means consequences such as NP not in P/poly.'. It must preserve the invariant 'MCSP complexity profile in nonuniform-circuits'. Its operational role is 'Several live routes are still searching for the right preserved quantity or witness measure, so invariant discovery remains a first-class discovery move.'.
\item First nontrivial lemma: For every admissible 'MCSP witness family', the first route-specific realization step preserves 'MCSP complexity profile in nonuniform-circuits' and prepares the passage from 'A hardness magnification theorem has the form:"A barely-nontrivial lower bound for an explicit problem Q in a weak model implies a major lower bound/separation in a strong model."Typically, Q is a meta-problem (e.g. a Gap-MCSP variant) and"major"means consequences such as NP not in P/poly.' to 'In refined regimes, lower bounds for Gap-MCSP slightly beyond known frontiers (e.g. N\textasciicircum{} 1+varepsilon -type circuit lower bounds for suitable promise gaps and input length N=2\textasciicircum{}n ) imply major consequences such as NP not in P/poly.'.
\item Quantifier signature: For every admissible 'MCSP witness family' on the route from 'st.def.hardness-magnification-phenomenon-programme-leve' to 'st.thm.oliveira-pich-santhanam-magnification-near-state'.
\item LaTeX theorem: \textbackslash\{\}forall W,\textbackslash\{\} \textbackslash\{\}mathrm\{Adm\}(W) \textbackslash\{\}to \textbackslash\{\}exists w : \textbackslash\{\}mathsf\{BW\}(W),\textbackslash\{\} \textbackslash\{\}mathrm\{PresInv\}(W,w) \textbackslash\{\}wedge \textbackslash\{\}mathrm\{RealizesTarget\}(W,w).
\item Negative/barrier version: If 'MCSP witness family' cannot be realized with invariant 'MCSP complexity profile in nonuniform-circuits', the route cannot force 'In refined regimes, lower bounds for Gap-MCSP slightly beyond known frontiers (e.g. N\textasciicircum{} 1+varepsilon -type circuit lower bounds for suitable promise gaps and input length N=2\textasciicircum{}n ) imply major consequences such as NP not in P/poly.'.
\item Proof strategy family: invariant-realization
\item Negative route: Assume 'MCSP witness family' cannot be realized while preserving 'MCSP complexity profile in nonuniform-circuits' strongly enough to activate the local target claim.
\item Failure diagnosis: The route fails if 'MCSP witness family' cannot be realized while preserving 'MCSP complexity profile in nonuniform-circuits'.
\item Proof obligations: For every admissible 'MCSP witness family', the first route-specific realization step preserves 'MCSP complexity profile in nonuniform-circuits' and prepares the passage from 'A hardness magnification theorem has the form:"A barely-nontrivial lower bound for an explicit problem Q in a weak model implies a major lower bound/separation in a strong model."Typically, Q is a meta-problem (e.g. a Gap-MCSP variant) and"major"means consequences such as NP not in P/poly.' to 'In refined regimes, lower bounds for Gap-MCSP slightly beyond known frontiers (e.g. N\textasciicircum{} 1+varepsilon -type circuit lower bounds for suitable promise gaps and input length N=2\textasciicircum{}n ) imply major consequences such as NP not in P/poly.'.
\item Candidate lemmas: Prove that a hardness-magnification witness family turns 'A barely-nontrivial lower bound for an explicit problem Q in a weak model' into 'major consequences such as NP not in P/poly' while preserving 'MCSP complexity profile in nonuniform-circuits'.
\item Independent value: A concrete special-case theorem for 'MCSP witness family' would be independently valuable because it converts a broad discovery programme into an inspectable mathematical subresult.
\item Expert gate: 6/6
\end{itemize}

\subsection*{conjpkg.disclayer-p-vs-np-pattern-invariant-discovery.02}
\begin{itemize}
\item Exact statement: There exists an admissible 'MCSP witness family' such that 'A hardness magnification theorem has the form:"A barely-nontrivial lower bound for an explicit problem Q in a weak model implies a major lower bound/separation in a strong model."Typically, Q is a meta-problem (e.g. a Gap-MCSP variant) and"major"means consequences such as NP not in P/poly.' can be organized into a route preserving 'MCSP complexity profile in nonuniform-circuits' and thereby supports 'There exist explicit"meta"problems (including MCSP variants) such that proving certain emph barely-superlinear lower bounds against restricted models (e.g. formulas) would imply major lower bounds for NP in much stronger models (including superpolynomial formula/circuit lower bounds, and in some parameter regimes consequences as strong as NP not in P/poly ).'.
\item Restricted-regime mode: special-case
\item Bridge object: MCSP witness family (family of instances equipped with a monotone witness measure)
\item Bridge definition: 'MCSP witness family' denotes the structured intermediate object required by this branch. It is meant to organize the source-side content 'A hardness magnification theorem has the form:"A barely-nontrivial lower bound for an explicit problem Q in a weak model implies a major lower bound/separation in a strong model."Typically, Q is a meta-problem (e.g. a Gap-MCSP variant) and"major"means consequences such as NP not in P/poly.'. It must preserve the invariant 'MCSP complexity profile in nonuniform-circuits'. Its operational role is 'Several live routes are still searching for the right preserved quantity or witness measure, so invariant discovery remains a first-class discovery move.'.
\item First nontrivial lemma: For every admissible 'MCSP witness family', the first route-specific realization step preserves 'MCSP complexity profile in nonuniform-circuits' and prepares the passage from 'A hardness magnification theorem has the form:"A barely-nontrivial lower bound for an explicit problem Q in a weak model implies a major lower bound/separation in a strong model."Typically, Q is a meta-problem (e.g. a Gap-MCSP variant) and"major"means consequences such as NP not in P/poly.' to 'There exist explicit"meta"problems (including MCSP variants) such that proving certain emph barely-superlinear lower bounds against restricted models (e.g. formulas) would imply major lower bounds for NP in much stronger models (including superpolynomial formula/circuit lower bounds, and in some parameter regimes consequences as strong as NP not in P/poly ).'.
\item Quantifier signature: For every admissible 'MCSP witness family' on the route from 'st.def.hardness-magnification-phenomenon-programme-leve' to 'st.thm.oliveira-santhanam-hardness-magnification-for-na'.
\item LaTeX theorem: \textbackslash\{\}forall W,\textbackslash\{\} \textbackslash\{\}mathrm\{Adm\}(W) \textbackslash\{\}to \textbackslash\{\}exists w : \textbackslash\{\}mathsf\{BW\}(W),\textbackslash\{\} \textbackslash\{\}mathrm\{PresInv\}(W,w) \textbackslash\{\}wedge \textbackslash\{\}mathrm\{RealizesTarget\}(W,w).
\item Negative/barrier version: If 'MCSP witness family' cannot be realized with invariant 'MCSP complexity profile in nonuniform-circuits', the route cannot force 'There exist explicit"meta"problems (including MCSP variants) such that proving certain emph barely-superlinear lower bounds against restricted models (e.g. formulas) would imply major lower bounds for NP in much stronger models (including superpolynomial formula/circuit lower bounds, and in some parameter regimes consequences as strong as NP not in P/poly ).'.
\item Proof strategy family: invariant-realization
\item Negative route: Assume 'MCSP witness family' cannot be realized while preserving 'MCSP complexity profile in nonuniform-circuits' strongly enough to activate the local target claim.
\item Failure diagnosis: The route fails if 'MCSP witness family' cannot be realized while preserving 'MCSP complexity profile in nonuniform-circuits'.
\item Proof obligations: For every admissible 'MCSP witness family', the first route-specific realization step preserves 'MCSP complexity profile in nonuniform-circuits' and prepares the passage from 'A hardness magnification theorem has the form:"A barely-nontrivial lower bound for an explicit problem Q in a weak model implies a major lower bound/separation in a strong model."Typically, Q is a meta-problem (e.g. a Gap-MCSP variant) and"major"means consequences such as NP not in P/poly.' to 'There exist explicit"meta"problems (including MCSP variants) such that proving certain emph barely-superlinear lower bounds against restricted models (e.g. formulas) would imply major lower bounds for NP in much stronger models (including superpolynomial formula/circuit lower bounds, and in some parameter regimes consequences as strong as NP not in P/poly ).'.
\item Candidate lemmas: Prove that a hardness-magnification witness family preserves 'MCSP complexity profile in nonuniform-circuits' while transporting the route from the nonuniform-circuits regime to the formula-size regime strongly enough to derive 'major lower bounds for NP in much stronger models (including superpolynomial formula/circuit lower bounds, and in some parameter...'.
\item Independent value: A concrete special-case theorem for 'MCSP witness family' would be independently valuable because it converts a broad discovery programme into an inspectable mathematical subresult.
\item Expert gate: 6/6
\end{itemize}

\subsection*{conjpkg.disclayer-p-vs-np-pattern-invariant-discovery.03}
\begin{itemize}
\item Exact statement: There exists an admissible 'MCSP obstruction family' such that 'A hardness magnification theorem has the form:"A barely-nontrivial lower bound for an explicit problem Q in a weak model implies a major lower bound/separation in a strong model."Typically, Q is a meta-problem (e.g. a Gap-MCSP variant) and"major"means consequences such as NP not in P/poly.' can be organized into a route preserving 'MCSP complexity profile in nonuniform-circuits' and thereby supports 'In refined regimes, lower bounds for Gap-MCSP slightly beyond known frontiers (e.g. N\textasciicircum{} 1+varepsilon -type circuit lower bounds for suitable promise gaps and input length N=2\textasciicircum{}n ) imply major consequences such as NP not in P/poly.'.
\item Restricted-regime mode: special-case
\item Bridge object: MCSP obstruction family (family of hard instances or obstructions forcing the target payoff)
\item Bridge definition: 'MCSP obstruction family' denotes the structured intermediate object required by this branch. It is meant to organize the source-side content 'A hardness magnification theorem has the form:"A barely-nontrivial lower bound for an explicit problem Q in a weak model implies a major lower bound/separation in a strong model."Typically, Q is a meta-problem (e.g. a Gap-MCSP variant) and"major"means consequences such as NP not in P/poly.'. It must preserve the invariant 'MCSP complexity profile in nonuniform-circuits'. Its operational role is 'Several live routes are still searching for the right preserved quantity or witness measure, so invariant discovery remains a first-class discovery move.'.
\item First nontrivial lemma: For every admissible 'MCSP obstruction family', the first route-specific realization step preserves 'MCSP complexity profile in nonuniform-circuits' and prepares the passage from 'A hardness magnification theorem has the form:"A barely-nontrivial lower bound for an explicit problem Q in a weak model implies a major lower bound/separation in a strong model."Typically, Q is a meta-problem (e.g. a Gap-MCSP variant) and"major"means consequences such as NP not in P/poly.' to 'In refined regimes, lower bounds for Gap-MCSP slightly beyond known frontiers (e.g. N\textasciicircum{} 1+varepsilon -type circuit lower bounds for suitable promise gaps and input length N=2\textasciicircum{}n ) imply major consequences such as NP not in P/poly.'.
\item Quantifier signature: For every admissible 'MCSP obstruction family' on the route from 'st.def.hardness-magnification-phenomenon-programme-leve' to 'st.thm.oliveira-pich-santhanam-magnification-near-state'.
\item LaTeX theorem: \textbackslash\{\}forall W,\textbackslash\{\} \textbackslash\{\}mathrm\{Adm\}(W) \textbackslash\{\}to \textbackslash\{\}exists w : \textbackslash\{\}mathsf\{BW\}(W),\textbackslash\{\} \textbackslash\{\}mathrm\{PresInv\}(W,w) \textbackslash\{\}wedge \textbackslash\{\}mathrm\{RealizesTarget\}(W,w).
\item Negative/barrier version: If 'MCSP obstruction family' cannot be realized with invariant 'MCSP complexity profile in nonuniform-circuits', the route cannot force 'In refined regimes, lower bounds for Gap-MCSP slightly beyond known frontiers (e.g. N\textasciicircum{} 1+varepsilon -type circuit lower bounds for suitable promise gaps and input length N=2\textasciicircum{}n ) imply major consequences such as NP not in P/poly.'.
\item Proof strategy family: invariant-realization
\item Negative route: Assume 'MCSP obstruction family' cannot be realized while preserving 'MCSP complexity profile in nonuniform-circuits' strongly enough to activate the local target claim.
\item Failure diagnosis: The route fails if 'MCSP obstruction family' cannot be realized while preserving 'MCSP complexity profile in nonuniform-circuits'.
\item Proof obligations: For every admissible 'MCSP obstruction family', the first route-specific realization step preserves 'MCSP complexity profile in nonuniform-circuits' and prepares the passage from 'A hardness magnification theorem has the form:"A barely-nontrivial lower bound for an explicit problem Q in a weak model implies a major lower bound/separation in a strong model."Typically, Q is a meta-problem (e.g. a Gap-MCSP variant) and"major"means consequences such as NP not in P/poly.' to 'In refined regimes, lower bounds for Gap-MCSP slightly beyond known frontiers (e.g. N\textasciicircum{} 1+varepsilon -type circuit lower bounds for suitable promise gaps and input length N=2\textasciicircum{}n ) imply major consequences such as NP not in P/poly.'.
\item Candidate lemmas: Prove that a hardness-magnification witness family turns 'A barely-nontrivial lower bound for an explicit problem Q in a weak model' into 'major consequences such as NP not in P/poly' while preserving 'MCSP complexity profile in nonuniform-circuits'.
\item Independent value: A concrete special-case theorem for 'MCSP obstruction family' would be independently valuable because it converts a broad discovery programme into an inspectable mathematical subresult.
\item Expert gate: 6/6
\end{itemize}

\subsection*{conjpkg.disclayer-p-vs-np-pattern-invariant-discovery.04}
\begin{itemize}
\item Exact statement: No admissible 'MCSP witness family' can simultaneously preserve 'MCSP complexity profile in nonuniform-circuits' and realize 'In refined regimes, lower bounds for Gap-MCSP slightly beyond known frontiers (e.g. N\textasciicircum{} 1+varepsilon -type circuit lower bounds for suitable promise gaps and input length N=2\textasciicircum{}n ) imply major consequences such as NP not in P/poly.' under the current route commitments.
\item Restricted-regime mode: special-case
\item Bridge object: MCSP witness family (family of instances equipped with a monotone witness measure)
\item Bridge definition: 'MCSP witness family' denotes the structured intermediate object required by this branch. It is meant to organize the source-side content 'A hardness magnification theorem has the form:"A barely-nontrivial lower bound for an explicit problem Q in a weak model implies a major lower bound/separation in a strong model."Typically, Q is a meta-problem (e.g. a Gap-MCSP variant) and"major"means consequences such as NP not in P/poly.'. It must preserve the invariant 'MCSP complexity profile in nonuniform-circuits'. Its operational role is 'Several live routes are still searching for the right preserved quantity or witness measure, so invariant discovery remains a first-class discovery move.'.
\item First nontrivial lemma: For every admissible 'MCSP witness family', the first route-specific realization step preserves 'MCSP complexity profile in nonuniform-circuits' and prepares the passage from 'A hardness magnification theorem has the form:"A barely-nontrivial lower bound for an explicit problem Q in a weak model implies a major lower bound/separation in a strong model."Typically, Q is a meta-problem (e.g. a Gap-MCSP variant) and"major"means consequences such as NP not in P/poly.' to 'In refined regimes, lower bounds for Gap-MCSP slightly beyond known frontiers (e.g. N\textasciicircum{} 1+varepsilon -type circuit lower bounds for suitable promise gaps and input length N=2\textasciicircum{}n ) imply major consequences such as NP not in P/poly.'.
\item Quantifier signature: For every admissible 'MCSP witness family' on the route from 'st.def.hardness-magnification-phenomenon-programme-leve' to 'st.thm.oliveira-pich-santhanam-magnification-near-state'.
\item LaTeX theorem: \textbackslash\{\}forall W,\textbackslash\{\} \textbackslash\{\}mathrm\{Adm\}(W) \textbackslash\{\}to \textbackslash\{\}neg \textbackslash\{\}exists w : \textbackslash\{\}mathsf\{BW\}(W),\textbackslash\{\} \textbackslash\{\}mathrm\{PresInv\}(W,w) \textbackslash\{\}wedge \textbackslash\{\}mathrm\{RealizesTarget\}(W,w).
\item Negative/barrier version: If 'MCSP witness family' cannot be realized with invariant 'MCSP complexity profile in nonuniform-circuits', the route cannot force 'In refined regimes, lower bounds for Gap-MCSP slightly beyond known frontiers (e.g. N\textasciicircum{} 1+varepsilon -type circuit lower bounds for suitable promise gaps and input length N=2\textasciicircum{}n ) imply major consequences such as NP not in P/poly.'.
\item Proof strategy family: invariant-realization
\item Negative route: Assume 'MCSP witness family' cannot be realized while preserving 'MCSP complexity profile in nonuniform-circuits' strongly enough to activate the local target claim.
\item Failure diagnosis: The route fails if 'MCSP witness family' cannot be realized while preserving 'MCSP complexity profile in nonuniform-circuits'.
\item Proof obligations: For every admissible 'MCSP witness family', the first route-specific realization step preserves 'MCSP complexity profile in nonuniform-circuits' and prepares the passage from 'A hardness magnification theorem has the form:"A barely-nontrivial lower bound for an explicit problem Q in a weak model implies a major lower bound/separation in a strong model."Typically, Q is a meta-problem (e.g. a Gap-MCSP variant) and"major"means consequences such as NP not in P/poly.' to 'In refined regimes, lower bounds for Gap-MCSP slightly beyond known frontiers (e.g. N\textasciicircum{} 1+varepsilon -type circuit lower bounds for suitable promise gaps and input length N=2\textasciicircum{}n ) imply major consequences such as NP not in P/poly.'.
\item Candidate lemmas: Obstruction theorem for MCSP witness family
\item Independent value: A concrete special-case theorem for 'MCSP witness family' would be independently valuable because it converts a broad discovery programme into an inspectable mathematical subresult.
\item Expert gate: 6/6
\end{itemize}

\subsection*{conjpkg.disclayer-p-vs-np-pattern-invariant-discovery.05}
\begin{itemize}
\item Exact statement: There exists an admissible 'MCSP witness family' whose realization internalizes 'A hardness magnification theorem has the form:"A barely-nontrivial lower bound for an explicit problem Q in a weak model implies a major lower bound/separation in a strong model."Typically, Q is a meta-problem (e.g. a Gap-MCSP variant) and"major"means consequences such as NP not in P/poly.' while preserving 'MCSP complexity profile in nonuniform-circuits' strongly enough to support 'In refined regimes, lower bounds for Gap-MCSP slightly beyond known frontiers (e.g. N\textasciicircum{} 1+varepsilon -type circuit lower bounds for suitable promise gaps and input length N=2\textasciicircum{}n ) imply major consequences such as NP not in P/poly.'.
\item Restricted-regime mode: special-case
\item Bridge object: MCSP witness family (family of instances equipped with a monotone witness measure)
\item Bridge definition: 'MCSP witness family' denotes the structured intermediate object required by this branch. It is meant to organize the source-side content 'A hardness magnification theorem has the form:"A barely-nontrivial lower bound for an explicit problem Q in a weak model implies a major lower bound/separation in a strong model."Typically, Q is a meta-problem (e.g. a Gap-MCSP variant) and"major"means consequences such as NP not in P/poly.'. It must preserve the invariant 'MCSP complexity profile in nonuniform-circuits'. Its operational role is 'Several live routes are still searching for the right preserved quantity or witness measure, so invariant discovery remains a first-class discovery move.'.
\item First nontrivial lemma: For every admissible 'MCSP witness family', the first route-specific realization step preserves 'MCSP complexity profile in nonuniform-circuits' and prepares the passage from 'A hardness magnification theorem has the form:"A barely-nontrivial lower bound for an explicit problem Q in a weak model implies a major lower bound/separation in a strong model."Typically, Q is a meta-problem (e.g. a Gap-MCSP variant) and"major"means consequences such as NP not in P/poly.' to 'In refined regimes, lower bounds for Gap-MCSP slightly beyond known frontiers (e.g. N\textasciicircum{} 1+varepsilon -type circuit lower bounds for suitable promise gaps and input length N=2\textasciicircum{}n ) imply major consequences such as NP not in P/poly.'.
\item Quantifier signature: For every admissible 'MCSP witness family' on the route from 'st.def.hardness-magnification-phenomenon-programme-leve' to 'st.thm.oliveira-pich-santhanam-magnification-near-state'.
\item LaTeX theorem: \textbackslash\{\}forall W,\textbackslash\{\} \textbackslash\{\}mathrm\{Adm\}(W) \textbackslash\{\}to \textbackslash\{\}exists w : \textbackslash\{\}mathsf\{BW\}(W),\textbackslash\{\} \textbackslash\{\}mathrm\{PresInv\}(W,w) \textbackslash\{\}wedge \textbackslash\{\}mathrm\{RealizesTarget\}(W,w).
\item Negative/barrier version: If 'MCSP witness family' cannot be realized with invariant 'MCSP complexity profile in nonuniform-circuits', the route cannot force 'In refined regimes, lower bounds for Gap-MCSP slightly beyond known frontiers (e.g. N\textasciicircum{} 1+varepsilon -type circuit lower bounds for suitable promise gaps and input length N=2\textasciicircum{}n ) imply major consequences such as NP not in P/poly.'.
\item Proof strategy family: invariant-realization
\item Negative route: Assume 'MCSP witness family' cannot be realized while preserving 'MCSP complexity profile in nonuniform-circuits' strongly enough to activate the local target claim.
\item Failure diagnosis: The route fails if 'MCSP witness family' cannot be realized while preserving 'MCSP complexity profile in nonuniform-circuits'.
\item Proof obligations: For every admissible 'MCSP witness family', the first route-specific realization step preserves 'MCSP complexity profile in nonuniform-circuits' and prepares the passage from 'A hardness magnification theorem has the form:"A barely-nontrivial lower bound for an explicit problem Q in a weak model implies a major lower bound/separation in a strong model."Typically, Q is a meta-problem (e.g. a Gap-MCSP variant) and"major"means consequences such as NP not in P/poly.' to 'In refined regimes, lower bounds for Gap-MCSP slightly beyond known frontiers (e.g. N\textasciicircum{} 1+varepsilon -type circuit lower bounds for suitable promise gaps and input length N=2\textasciicircum{}n ) imply major consequences such as NP not in P/poly.'.
\item Candidate lemmas: Realization theorem for MCSP witness family
\item Independent value: A concrete special-case theorem for 'MCSP witness family' would be independently valuable because it converts a broad discovery programme into an inspectable mathematical subresult.
\item Expert gate: 6/6
\end{itemize}

\subsection*{conjpkg.disclayer-p-vs-np-pattern-invariant-discovery.06}
\begin{itemize}
\item Exact statement: For every admissible 'MCSP witness family', there is a transport law from 'nonuniform-circuits' to 'formula-size' that preserves 'MCSP complexity profile in nonuniform-circuits' and retains enough strength to activate 'There exist explicit"meta"problems (including MCSP variants) such that proving certain emph barely-superlinear lower bounds against restricted models (e.g. formulas) would imply major lower bounds for NP in much stronger models (including superpolynomial formula/circuit lower bounds, and in some parameter regimes consequences as strong as NP not in P/poly ).'.
\item Restricted-regime mode: special-case
\item Bridge object: MCSP witness family (family of instances equipped with a monotone witness measure)
\item Bridge definition: 'MCSP witness family' denotes the structured intermediate object required by this branch. It is meant to organize the source-side content 'A hardness magnification theorem has the form:"A barely-nontrivial lower bound for an explicit problem Q in a weak model implies a major lower bound/separation in a strong model."Typically, Q is a meta-problem (e.g. a Gap-MCSP variant) and"major"means consequences such as NP not in P/poly.'. It must preserve the invariant 'MCSP complexity profile in nonuniform-circuits'. Its operational role is 'Several live routes are still searching for the right preserved quantity or witness measure, so invariant discovery remains a first-class discovery move.'.
\item First nontrivial lemma: For every admissible 'MCSP witness family', the first route-specific realization step preserves 'MCSP complexity profile in nonuniform-circuits' and prepares the passage from 'A hardness magnification theorem has the form:"A barely-nontrivial lower bound for an explicit problem Q in a weak model implies a major lower bound/separation in a strong model."Typically, Q is a meta-problem (e.g. a Gap-MCSP variant) and"major"means consequences such as NP not in P/poly.' to 'There exist explicit"meta"problems (including MCSP variants) such that proving certain emph barely-superlinear lower bounds against restricted models (e.g. formulas) would imply major lower bounds for NP in much stronger models (including superpolynomial formula/circuit lower bounds, and in some parameter regimes consequences as strong as NP not in P/poly ).'.
\item Quantifier signature: For every admissible 'MCSP witness family' on the route from 'st.def.hardness-magnification-phenomenon-programme-leve' to 'st.thm.oliveira-santhanam-hardness-magnification-for-na'.
\item LaTeX theorem: \textbackslash\{\}forall W,\textbackslash\{\} \textbackslash\{\}mathrm\{Adm\}(W) \textbackslash\{\}to \textbackslash\{\}exists w : \textbackslash\{\}mathsf\{BW\}(W),\textbackslash\{\} \textbackslash\{\}mathrm\{PresInv\}(W,w) \textbackslash\{\}wedge \textbackslash\{\}mathrm\{RealizesTarget\}(W,w).
\item Negative/barrier version: If 'MCSP witness family' cannot be realized with invariant 'MCSP complexity profile in nonuniform-circuits', the route cannot force 'There exist explicit"meta"problems (including MCSP variants) such that proving certain emph barely-superlinear lower bounds against restricted models (e.g. formulas) would imply major lower bounds for NP in much stronger models (including superpolynomial formula/circuit lower bounds, and in some parameter regimes consequences as strong as NP not in P/poly ).'.
\item Proof strategy family: invariant-realization
\item Negative route: Assume 'MCSP witness family' cannot be realized while preserving 'MCSP complexity profile in nonuniform-circuits' strongly enough to activate the local target claim.
\item Failure diagnosis: The route fails if 'MCSP witness family' cannot be realized while preserving 'MCSP complexity profile in nonuniform-circuits'.
\item Proof obligations: For every admissible 'MCSP witness family', the first route-specific realization step preserves 'MCSP complexity profile in nonuniform-circuits' and prepares the passage from 'A hardness magnification theorem has the form:"A barely-nontrivial lower bound for an explicit problem Q in a weak model implies a major lower bound/separation in a strong model."Typically, Q is a meta-problem (e.g. a Gap-MCSP variant) and"major"means consequences such as NP not in P/poly.' to 'There exist explicit"meta"problems (including MCSP variants) such that proving certain emph barely-superlinear lower bounds against restricted models (e.g. formulas) would imply major lower bounds for NP in much stronger models (including superpolynomial formula/circuit lower bounds, and in some parameter regimes consequences as strong as NP not in P/poly ).'.
\item Candidate lemmas: Transport theorem for MCSP witness family
\item Independent value: A concrete special-case theorem for 'MCSP witness family' would be independently valuable because it converts a broad discovery programme into an inspectable mathematical subresult.
\item Expert gate: 6/6
\end{itemize}

\subsection*{conjpkg.disclayer-p-vs-np-pattern-invariant-discovery.07}
\begin{itemize}
\item Exact statement: No admissible 'MCSP witness family' can simultaneously preserve 'MCSP complexity profile in nonuniform-circuits' and realize 'There exist explicit"meta"problems (including MCSP variants) such that proving certain emph barely-superlinear lower bounds against restricted models (e.g. formulas) would imply major lower bounds for NP in much stronger models (including superpolynomial formula/circuit lower bounds, and in some parameter regimes consequences as strong as NP not in P/poly ).' under the current route commitments.
\item Restricted-regime mode: special-case
\item Bridge object: MCSP witness family (family of instances equipped with a monotone witness measure)
\item Bridge definition: 'MCSP witness family' denotes the structured intermediate object required by this branch. It is meant to organize the source-side content 'A hardness magnification theorem has the form:"A barely-nontrivial lower bound for an explicit problem Q in a weak model implies a major lower bound/separation in a strong model."Typically, Q is a meta-problem (e.g. a Gap-MCSP variant) and"major"means consequences such as NP not in P/poly.'. It must preserve the invariant 'MCSP complexity profile in nonuniform-circuits'. Its operational role is 'Several live routes are still searching for the right preserved quantity or witness measure, so invariant discovery remains a first-class discovery move.'.
\item First nontrivial lemma: For every admissible 'MCSP witness family', the first route-specific realization step preserves 'MCSP complexity profile in nonuniform-circuits' and prepares the passage from 'A hardness magnification theorem has the form:"A barely-nontrivial lower bound for an explicit problem Q in a weak model implies a major lower bound/separation in a strong model."Typically, Q is a meta-problem (e.g. a Gap-MCSP variant) and"major"means consequences such as NP not in P/poly.' to 'There exist explicit"meta"problems (including MCSP variants) such that proving certain emph barely-superlinear lower bounds against restricted models (e.g. formulas) would imply major lower bounds for NP in much stronger models (including superpolynomial formula/circuit lower bounds, and in some parameter regimes consequences as strong as NP not in P/poly ).'.
\item Quantifier signature: For every admissible 'MCSP witness family' on the route from 'st.def.hardness-magnification-phenomenon-programme-leve' to 'st.thm.oliveira-santhanam-hardness-magnification-for-na'.
\item LaTeX theorem: \textbackslash\{\}forall W,\textbackslash\{\} \textbackslash\{\}mathrm\{Adm\}(W) \textbackslash\{\}to \textbackslash\{\}neg \textbackslash\{\}exists w : \textbackslash\{\}mathsf\{BW\}(W),\textbackslash\{\} \textbackslash\{\}mathrm\{PresInv\}(W,w) \textbackslash\{\}wedge \textbackslash\{\}mathrm\{RealizesTarget\}(W,w).
\item Negative/barrier version: If 'MCSP witness family' cannot be realized with invariant 'MCSP complexity profile in nonuniform-circuits', the route cannot force 'There exist explicit"meta"problems (including MCSP variants) such that proving certain emph barely-superlinear lower bounds against restricted models (e.g. formulas) would imply major lower bounds for NP in much stronger models (including superpolynomial formula/circuit lower bounds, and in some parameter regimes consequences as strong as NP not in P/poly ).'.
\item Proof strategy family: invariant-realization
\item Negative route: Assume 'MCSP witness family' cannot be realized while preserving 'MCSP complexity profile in nonuniform-circuits' strongly enough to activate the local target claim.
\item Failure diagnosis: The route fails if 'MCSP witness family' cannot be realized while preserving 'MCSP complexity profile in nonuniform-circuits'.
\item Proof obligations: For every admissible 'MCSP witness family', the first route-specific realization step preserves 'MCSP complexity profile in nonuniform-circuits' and prepares the passage from 'A hardness magnification theorem has the form:"A barely-nontrivial lower bound for an explicit problem Q in a weak model implies a major lower bound/separation in a strong model."Typically, Q is a meta-problem (e.g. a Gap-MCSP variant) and"major"means consequences such as NP not in P/poly.' to 'There exist explicit"meta"problems (including MCSP variants) such that proving certain emph barely-superlinear lower bounds against restricted models (e.g. formulas) would imply major lower bounds for NP in much stronger models (including superpolynomial formula/circuit lower bounds, and in some parameter regimes consequences as strong as NP not in P/poly ).'.
\item Candidate lemmas: Obstruction theorem for MCSP witness family
\item Independent value: A concrete special-case theorem for 'MCSP witness family' would be independently valuable because it converts a broad discovery programme into an inspectable mathematical subresult.
\item Expert gate: 6/6
\end{itemize}

\subsection*{conjpkg.disclayer-p-vs-np-pattern-invariant-discovery.08}
\begin{itemize}
\item Exact statement: There exists an admissible 'MCSP witness family' whose realization internalizes 'A hardness magnification theorem has the form:"A barely-nontrivial lower bound for an explicit problem Q in a weak model implies a major lower bound/separation in a strong model."Typically, Q is a meta-problem (e.g. a Gap-MCSP variant) and"major"means consequences such as NP not in P/poly.' while preserving 'MCSP complexity profile in nonuniform-circuits' strongly enough to support 'There exist explicit"meta"problems (including MCSP variants) such that proving certain emph barely-superlinear lower bounds against restricted models (e.g. formulas) would imply major lower bounds for NP in much stronger models (including superpolynomial formula/circuit lower bounds, and in some parameter regimes consequences as strong as NP not in P/poly ).'.
\item Restricted-regime mode: special-case
\item Bridge object: MCSP witness family (family of instances equipped with a monotone witness measure)
\item Bridge definition: 'MCSP witness family' denotes the structured intermediate object required by this branch. It is meant to organize the source-side content 'A hardness magnification theorem has the form:"A barely-nontrivial lower bound for an explicit problem Q in a weak model implies a major lower bound/separation in a strong model."Typically, Q is a meta-problem (e.g. a Gap-MCSP variant) and"major"means consequences such as NP not in P/poly.'. It must preserve the invariant 'MCSP complexity profile in nonuniform-circuits'. Its operational role is 'Several live routes are still searching for the right preserved quantity or witness measure, so invariant discovery remains a first-class discovery move.'.
\item First nontrivial lemma: For every admissible 'MCSP witness family', the first route-specific realization step preserves 'MCSP complexity profile in nonuniform-circuits' and prepares the passage from 'A hardness magnification theorem has the form:"A barely-nontrivial lower bound for an explicit problem Q in a weak model implies a major lower bound/separation in a strong model."Typically, Q is a meta-problem (e.g. a Gap-MCSP variant) and"major"means consequences such as NP not in P/poly.' to 'There exist explicit"meta"problems (including MCSP variants) such that proving certain emph barely-superlinear lower bounds against restricted models (e.g. formulas) would imply major lower bounds for NP in much stronger models (including superpolynomial formula/circuit lower bounds, and in some parameter regimes consequences as strong as NP not in P/poly ).'.
\item Quantifier signature: For every admissible 'MCSP witness family' on the route from 'st.def.hardness-magnification-phenomenon-programme-leve' to 'st.thm.oliveira-santhanam-hardness-magnification-for-na'.
\item LaTeX theorem: \textbackslash\{\}forall W,\textbackslash\{\} \textbackslash\{\}mathrm\{Adm\}(W) \textbackslash\{\}to \textbackslash\{\}exists w : \textbackslash\{\}mathsf\{BW\}(W),\textbackslash\{\} \textbackslash\{\}mathrm\{PresInv\}(W,w) \textbackslash\{\}wedge \textbackslash\{\}mathrm\{RealizesTarget\}(W,w).
\item Negative/barrier version: If 'MCSP witness family' cannot be realized with invariant 'MCSP complexity profile in nonuniform-circuits', the route cannot force 'There exist explicit"meta"problems (including MCSP variants) such that proving certain emph barely-superlinear lower bounds against restricted models (e.g. formulas) would imply major lower bounds for NP in much stronger models (including superpolynomial formula/circuit lower bounds, and in some parameter regimes consequences as strong as NP not in P/poly ).'.
\item Proof strategy family: invariant-realization
\item Negative route: Assume 'MCSP witness family' cannot be realized while preserving 'MCSP complexity profile in nonuniform-circuits' strongly enough to activate the local target claim.
\item Failure diagnosis: The route fails if 'MCSP witness family' cannot be realized while preserving 'MCSP complexity profile in nonuniform-circuits'.
\item Proof obligations: For every admissible 'MCSP witness family', the first route-specific realization step preserves 'MCSP complexity profile in nonuniform-circuits' and prepares the passage from 'A hardness magnification theorem has the form:"A barely-nontrivial lower bound for an explicit problem Q in a weak model implies a major lower bound/separation in a strong model."Typically, Q is a meta-problem (e.g. a Gap-MCSP variant) and"major"means consequences such as NP not in P/poly.' to 'There exist explicit"meta"problems (including MCSP variants) such that proving certain emph barely-superlinear lower bounds against restricted models (e.g. formulas) would imply major lower bounds for NP in much stronger models (including superpolynomial formula/circuit lower bounds, and in some parameter regimes consequences as strong as NP not in P/poly ).'.
\item Candidate lemmas: Realization theorem for MCSP witness family
\item Independent value: A concrete special-case theorem for 'MCSP witness family' would be independently valuable because it converts a broad discovery programme into an inspectable mathematical subresult.
\item Expert gate: 6/6
\end{itemize}

\subsection*{conjpkg.disclayer-p-vs-np-pattern-invariant-discovery.09}
\begin{itemize}
\item Exact statement: No admissible 'MCSP obstruction family' can simultaneously preserve 'MCSP complexity profile in nonuniform-circuits' and realize 'In refined regimes, lower bounds for Gap-MCSP slightly beyond known frontiers (e.g. N\textasciicircum{} 1+varepsilon -type circuit lower bounds for suitable promise gaps and input length N=2\textasciicircum{}n ) imply major consequences such as NP not in P/poly.' under the current route commitments.
\item Restricted-regime mode: special-case
\item Bridge object: MCSP obstruction family (family of hard instances or obstructions forcing the target payoff)
\item Bridge definition: 'MCSP obstruction family' denotes the structured intermediate object required by this branch. It is meant to organize the source-side content 'A hardness magnification theorem has the form:"A barely-nontrivial lower bound for an explicit problem Q in a weak model implies a major lower bound/separation in a strong model."Typically, Q is a meta-problem (e.g. a Gap-MCSP variant) and"major"means consequences such as NP not in P/poly.'. It must preserve the invariant 'MCSP complexity profile in nonuniform-circuits'. Its operational role is 'Several live routes are still searching for the right preserved quantity or witness measure, so invariant discovery remains a first-class discovery move.'.
\item First nontrivial lemma: For every admissible 'MCSP obstruction family', the first route-specific realization step preserves 'MCSP complexity profile in nonuniform-circuits' and prepares the passage from 'A hardness magnification theorem has the form:"A barely-nontrivial lower bound for an explicit problem Q in a weak model implies a major lower bound/separation in a strong model."Typically, Q is a meta-problem (e.g. a Gap-MCSP variant) and"major"means consequences such as NP not in P/poly.' to 'In refined regimes, lower bounds for Gap-MCSP slightly beyond known frontiers (e.g. N\textasciicircum{} 1+varepsilon -type circuit lower bounds for suitable promise gaps and input length N=2\textasciicircum{}n ) imply major consequences such as NP not in P/poly.'.
\item Quantifier signature: For every admissible 'MCSP obstruction family' on the route from 'st.def.hardness-magnification-phenomenon-programme-leve' to 'st.thm.oliveira-pich-santhanam-magnification-near-state'.
\item LaTeX theorem: \textbackslash\{\}forall W,\textbackslash\{\} \textbackslash\{\}mathrm\{Adm\}(W) \textbackslash\{\}to \textbackslash\{\}neg \textbackslash\{\}exists w : \textbackslash\{\}mathsf\{BW\}(W),\textbackslash\{\} \textbackslash\{\}mathrm\{PresInv\}(W,w) \textbackslash\{\}wedge \textbackslash\{\}mathrm\{RealizesTarget\}(W,w).
\item Negative/barrier version: If 'MCSP obstruction family' cannot be realized with invariant 'MCSP complexity profile in nonuniform-circuits', the route cannot force 'In refined regimes, lower bounds for Gap-MCSP slightly beyond known frontiers (e.g. N\textasciicircum{} 1+varepsilon -type circuit lower bounds for suitable promise gaps and input length N=2\textasciicircum{}n ) imply major consequences such as NP not in P/poly.'.
\item Proof strategy family: invariant-realization
\item Negative route: Assume 'MCSP obstruction family' cannot be realized while preserving 'MCSP complexity profile in nonuniform-circuits' strongly enough to activate the local target claim.
\item Failure diagnosis: The route fails if 'MCSP obstruction family' cannot be realized while preserving 'MCSP complexity profile in nonuniform-circuits'.
\item Proof obligations: For every admissible 'MCSP obstruction family', the first route-specific realization step preserves 'MCSP complexity profile in nonuniform-circuits' and prepares the passage from 'A hardness magnification theorem has the form:"A barely-nontrivial lower bound for an explicit problem Q in a weak model implies a major lower bound/separation in a strong model."Typically, Q is a meta-problem (e.g. a Gap-MCSP variant) and"major"means consequences such as NP not in P/poly.' to 'In refined regimes, lower bounds for Gap-MCSP slightly beyond known frontiers (e.g. N\textasciicircum{} 1+varepsilon -type circuit lower bounds for suitable promise gaps and input length N=2\textasciicircum{}n ) imply major consequences such as NP not in P/poly.'.
\item Candidate lemmas: Obstruction theorem for MCSP obstruction family
\item Independent value: A concrete special-case theorem for 'MCSP obstruction family' would be independently valuable because it converts a broad discovery programme into an inspectable mathematical subresult.
\item Expert gate: 6/6
\end{itemize}

\subsection*{conjpkg.disclayer-p-vs-np-pattern-invariant-discovery.10}
\begin{itemize}
\item Exact statement: There exists an admissible 'MCSP obstruction family' whose realization internalizes 'A hardness magnification theorem has the form:"A barely-nontrivial lower bound for an explicit problem Q in a weak model implies a major lower bound/separation in a strong model."Typically, Q is a meta-problem (e.g. a Gap-MCSP variant) and"major"means consequences such as NP not in P/poly.' while preserving 'MCSP complexity profile in nonuniform-circuits' strongly enough to support 'In refined regimes, lower bounds for Gap-MCSP slightly beyond known frontiers (e.g. N\textasciicircum{} 1+varepsilon -type circuit lower bounds for suitable promise gaps and input length N=2\textasciicircum{}n ) imply major consequences such as NP not in P/poly.'.
\item Restricted-regime mode: special-case
\item Bridge object: MCSP obstruction family (family of hard instances or obstructions forcing the target payoff)
\item Bridge definition: 'MCSP obstruction family' denotes the structured intermediate object required by this branch. It is meant to organize the source-side content 'A hardness magnification theorem has the form:"A barely-nontrivial lower bound for an explicit problem Q in a weak model implies a major lower bound/separation in a strong model."Typically, Q is a meta-problem (e.g. a Gap-MCSP variant) and"major"means consequences such as NP not in P/poly.'. It must preserve the invariant 'MCSP complexity profile in nonuniform-circuits'. Its operational role is 'Several live routes are still searching for the right preserved quantity or witness measure, so invariant discovery remains a first-class discovery move.'.
\item First nontrivial lemma: For every admissible 'MCSP obstruction family', the first route-specific realization step preserves 'MCSP complexity profile in nonuniform-circuits' and prepares the passage from 'A hardness magnification theorem has the form:"A barely-nontrivial lower bound for an explicit problem Q in a weak model implies a major lower bound/separation in a strong model."Typically, Q is a meta-problem (e.g. a Gap-MCSP variant) and"major"means consequences such as NP not in P/poly.' to 'In refined regimes, lower bounds for Gap-MCSP slightly beyond known frontiers (e.g. N\textasciicircum{} 1+varepsilon -type circuit lower bounds for suitable promise gaps and input length N=2\textasciicircum{}n ) imply major consequences such as NP not in P/poly.'.
\item Quantifier signature: For every admissible 'MCSP obstruction family' on the route from 'st.def.hardness-magnification-phenomenon-programme-leve' to 'st.thm.oliveira-pich-santhanam-magnification-near-state'.
\item LaTeX theorem: \textbackslash\{\}forall W,\textbackslash\{\} \textbackslash\{\}mathrm\{Adm\}(W) \textbackslash\{\}to \textbackslash\{\}exists w : \textbackslash\{\}mathsf\{BW\}(W),\textbackslash\{\} \textbackslash\{\}mathrm\{PresInv\}(W,w) \textbackslash\{\}wedge \textbackslash\{\}mathrm\{RealizesTarget\}(W,w).
\item Negative/barrier version: If 'MCSP obstruction family' cannot be realized with invariant 'MCSP complexity profile in nonuniform-circuits', the route cannot force 'In refined regimes, lower bounds for Gap-MCSP slightly beyond known frontiers (e.g. N\textasciicircum{} 1+varepsilon -type circuit lower bounds for suitable promise gaps and input length N=2\textasciicircum{}n ) imply major consequences such as NP not in P/poly.'.
\item Proof strategy family: invariant-realization
\item Negative route: Assume 'MCSP obstruction family' cannot be realized while preserving 'MCSP complexity profile in nonuniform-circuits' strongly enough to activate the local target claim.
\item Failure diagnosis: The route fails if 'MCSP obstruction family' cannot be realized while preserving 'MCSP complexity profile in nonuniform-circuits'.
\item Proof obligations: For every admissible 'MCSP obstruction family', the first route-specific realization step preserves 'MCSP complexity profile in nonuniform-circuits' and prepares the passage from 'A hardness magnification theorem has the form:"A barely-nontrivial lower bound for an explicit problem Q in a weak model implies a major lower bound/separation in a strong model."Typically, Q is a meta-problem (e.g. a Gap-MCSP variant) and"major"means consequences such as NP not in P/poly.' to 'In refined regimes, lower bounds for Gap-MCSP slightly beyond known frontiers (e.g. N\textasciicircum{} 1+varepsilon -type circuit lower bounds for suitable promise gaps and input length N=2\textasciicircum{}n ) imply major consequences such as NP not in P/poly.'.
\item Candidate lemmas: Realization theorem for MCSP obstruction family
\item Independent value: A concrete special-case theorem for 'MCSP obstruction family' would be independently valuable because it converts a broad discovery programme into an inspectable mathematical subresult.
\item Expert gate: 6/6
\end{itemize}

\subsection*{conjpkg.disclayer-p-vs-np-family-equivalence-coherence-reframing.01}
\begin{itemize}
\item Exact statement: There exists an equivalence witness identifying 'MCSP witness family' with the target-side object class relevant to 'In refined regimes, lower bounds for Gap-MCSP slightly beyond known frontiers (e.g. N\textasciicircum{} 1+varepsilon -type circuit lower bounds for suitable promise gaps and input length N=2\textasciicircum{}n ) imply major consequences such as NP not in P/poly.', and 'MCSP complexity profile in nonuniform-circuits' transports coherently along that witness.
\item Restricted-regime mode: equivalence-theorem
\item Bridge object: MCSP witness family (family of instances equipped with a monotone witness measure)
\item Bridge definition: 'MCSP witness family' denotes the structured intermediate object required by this branch. It is meant to organize the source-side content 'A hardness magnification theorem has the form:"A barely-nontrivial lower bound for an explicit problem Q in a weak model implies a major lower bound/separation in a strong model."Typically, Q is a meta-problem (e.g. a Gap-MCSP variant) and"major"means consequences such as NP not in P/poly.'. It must preserve the invariant 'MCSP complexity profile in nonuniform-circuits'. Its operational role is 'Recast the blocked frontier as an identity, equivalence, transport, or coherence problem rather than as a missing brute-force bridge.'.
\item First nontrivial lemma: For every admissible 'MCSP witness family', the target-side operationalization is equivalent to the route-side operationalization in the sense needed to transport 'MCSP complexity profile in nonuniform-circuits' toward 'In refined regimes, lower bounds for Gap-MCSP slightly beyond known frontiers (e.g. N\textasciicircum{} 1+varepsilon -type circuit lower bounds for suitable promise gaps and input length N=2\textasciicircum{}n ) imply major consequences such as NP not in P/poly.'.
\item Quantifier signature: For every admissible 'MCSP witness family' on the route from 'st.def.hardness-magnification-phenomenon-programme-leve' to 'st.thm.oliveira-pich-santhanam-magnification-near-state'.
\item LaTeX theorem: \textbackslash\{\}forall W,\textbackslash\{\} \textbackslash\{\}mathrm\{Adm\}(W) \textbackslash\{\}to \textbackslash\{\}exists w : \textbackslash\{\}mathsf\{BW\}(W),\textbackslash\{\} \textbackslash\{\}mathrm\{PresInv\}(W,w) \textbackslash\{\}wedge \textbackslash\{\}mathrm\{RealizesTarget\}(W,w).
\item Negative/barrier version: Absent an explicit equivalence witness for 'MCSP witness family', transport toward 'In refined regimes, lower bounds for Gap-MCSP slightly beyond known frontiers (e.g. N\textasciicircum{} 1+varepsilon -type circuit lower bounds for suitable promise gaps and input length N=2\textasciicircum{}n ) imply major consequences such as NP not in P/poly.' is not licensed. no\_equivalence\_witness
\item Proof strategy family: equivalence-transport
\item Negative route: Assume no equivalence witness for 'MCSP witness family' aligns the route-side and target-side objects.
\item Failure diagnosis: The route fails if no equivalence witness for 'MCSP witness family' transports 'MCSP complexity profile in nonuniform-circuits' to the target-side operationalization. Local negative probe: no\_equivalence\_witness.
\item Proof obligations: For every admissible 'MCSP witness family', the target-side operationalization is equivalent to the route-side operationalization in the sense needed to transport 'MCSP complexity profile in nonuniform-circuits' toward 'In refined regimes, lower bounds for Gap-MCSP slightly beyond known frontiers (e.g. N\textasciicircum{} 1+varepsilon -type circuit lower bounds for suitable promise gaps and input length N=2\textasciicircum{}n ) imply major consequences such as NP not in P/poly.'.; Transport obligation: move the invariant across the proposed equivalence for 'carrier-realization' without changing the target claim.; Transport obligation: move the invariant across the proposed equivalence for 'formula-size-transport' without changing the target claim.; Equivalence claim: determine whether the live carrier for 'carrier-realization' is identical to, equivalent to, or only analogous with the target-side object class.
\item Candidate lemmas: Equivalence and coherent transport theorem for MCSP witness family; Equivalence transport lemma for MCSP witness family
\item Independent value: An equivalence theorem for 'MCSP witness family' would independently matter because it identifies when two live operationalizations are mathematically interchangeable rather than merely analogous.
\item Expert gate: 6/6
\end{itemize}

\subsection*{conjpkg.disclayer-p-vs-np-family-equivalence-coherence-reframing.02}
\begin{itemize}
\item Exact statement: There exists an equivalence witness identifying 'MCSP witness family' with the target-side object class relevant to 'There exist explicit"meta"problems (including MCSP variants) such that proving certain emph barely-superlinear lower bounds against restricted models (e.g. formulas) would imply major lower bounds for NP in much stronger models (including superpolynomial formula/circuit lower bounds, and in some parameter regimes consequences as strong as NP not in P/poly ).', and 'MCSP complexity profile in nonuniform-circuits' transports coherently along that witness.
\item Restricted-regime mode: equivalence-theorem
\item Bridge object: MCSP witness family (family of instances equipped with a monotone witness measure)
\item Bridge definition: 'MCSP witness family' denotes the structured intermediate object required by this branch. It is meant to organize the source-side content 'A hardness magnification theorem has the form:"A barely-nontrivial lower bound for an explicit problem Q in a weak model implies a major lower bound/separation in a strong model."Typically, Q is a meta-problem (e.g. a Gap-MCSP variant) and"major"means consequences such as NP not in P/poly.'. It must preserve the invariant 'MCSP complexity profile in nonuniform-circuits'. Its operational role is 'Recast the blocked frontier as an identity, equivalence, transport, or coherence problem rather than as a missing brute-force bridge.'.
\item First nontrivial lemma: For every admissible 'MCSP witness family', the target-side operationalization is equivalent to the route-side operationalization in the sense needed to transport 'MCSP complexity profile in nonuniform-circuits' toward 'There exist explicit"meta"problems (including MCSP variants) such that proving certain emph barely-superlinear lower bounds against restricted models (e.g. formulas) would imply major lower bounds for NP in much stronger models (including superpolynomial formula/circuit lower bounds, and in some parameter regimes consequences as strong as NP not in P/poly ).'.
\item Quantifier signature: For every admissible 'MCSP witness family' on the route from 'st.def.hardness-magnification-phenomenon-programme-leve' to 'st.thm.oliveira-santhanam-hardness-magnification-for-na'.
\item LaTeX theorem: \textbackslash\{\}forall W,\textbackslash\{\} \textbackslash\{\}mathrm\{Adm\}(W) \textbackslash\{\}to \textbackslash\{\}exists w : \textbackslash\{\}mathsf\{BW\}(W),\textbackslash\{\} \textbackslash\{\}mathrm\{PresInv\}(W,w) \textbackslash\{\}wedge \textbackslash\{\}mathrm\{RealizesTarget\}(W,w).
\item Negative/barrier version: Absent an explicit equivalence witness for 'MCSP witness family', transport toward 'There exist explicit"meta"problems (including MCSP variants) such that proving certain emph barely-superlinear lower bounds against restricted models (e.g. formulas) would imply major lower bounds for NP in much stronger models (including superpolynomial formula/circuit lower bounds, and in some parameter regimes consequences as strong as NP not in P/poly ).' is not licensed. no\_equivalence\_witness
\item Proof strategy family: equivalence-transport
\item Negative route: Assume no equivalence witness for 'MCSP witness family' aligns the route-side and target-side objects.
\item Failure diagnosis: The route fails if no equivalence witness for 'MCSP witness family' transports 'MCSP complexity profile in nonuniform-circuits' to the target-side operationalization. Local negative probe: no\_equivalence\_witness.
\item Proof obligations: For every admissible 'MCSP witness family', the target-side operationalization is equivalent to the route-side operationalization in the sense needed to transport 'MCSP complexity profile in nonuniform-circuits' toward 'There exist explicit"meta"problems (including MCSP variants) such that proving certain emph barely-superlinear lower bounds against restricted models (e.g. formulas) would imply major lower bounds for NP in much stronger models (including superpolynomial formula/circuit lower bounds, and in some parameter regimes consequences as strong as NP not in P/poly ).'.; Transport obligation: move the invariant across the proposed equivalence for 'carrier-realization' without changing the target claim.; Transport obligation: move the invariant across the proposed equivalence for 'formula-size-transport' without changing the target claim.; Equivalence claim: determine whether the live carrier for 'carrier-realization' is identical to, equivalent to, or only analogous with the target-side object class.
\item Candidate lemmas: Equivalence and coherent transport theorem for MCSP witness family; Equivalence transport lemma for MCSP witness family
\item Independent value: An equivalence theorem for 'MCSP witness family' would independently matter because it identifies when two live operationalizations are mathematically interchangeable rather than merely analogous.
\item Expert gate: 6/6
\end{itemize}

\subsection*{conjpkg.disclayer-p-vs-np-family-equivalence-coherence-reframing.03}
\begin{itemize}
\item Exact statement: There exists an equivalence witness identifying 'MCSP witness family' with the target-side object class relevant to 'There exist general-purpose magnification theorems (not problem-specific) showing that certain mild lower bounds for distinguishers/sparse-code objects can be amplified into stronger circuit/formula lower bounds, achieving magnification thresholds below previously known lower bounds in some settings.', and 'exponential-lower-bound witness complexity' transports coherently along that witness.
\item Restricted-regime mode: equivalence-theorem
\item Bridge object: MCSP witness family (family of instances equipped with a monotone witness measure)
\item Bridge definition: 'MCSP witness family' denotes the structured intermediate object required by this branch. It is meant to organize the source-side content 'A hardness magnification theorem has the form:"A barely-nontrivial lower bound for an explicit problem Q in a weak model implies a major lower bound/separation in a strong model."Typically, Q is a meta-problem (e.g. a Gap-MCSP variant) and"major"means consequences such as NP not in P/poly.'. It must preserve the invariant 'exponential-lower-bound witness complexity'. Its operational role is 'Recast the blocked frontier as an identity, equivalence, transport, or coherence problem rather than as a missing brute-force bridge.'.
\item First nontrivial lemma: For every admissible 'MCSP witness family', the target-side operationalization is equivalent to the route-side operationalization in the sense needed to transport 'exponential-lower-bound witness complexity' toward 'There exist general-purpose magnification theorems (not problem-specific) showing that certain mild lower bounds for distinguishers/sparse-code objects can be amplified into stronger circuit/formula lower bounds, achieving magnification thresholds below previously known lower bounds in some settings.'.
\item Quantifier signature: For every admissible 'MCSP witness family' on the route from 'st.def.hardness-magnification-phenomenon-programme-leve' to 'st.thm.atserias-m-uller-simple-general-magnification-fr'.
\item LaTeX theorem: \textbackslash\{\}forall W,\textbackslash\{\} \textbackslash\{\}mathrm\{Adm\}(W) \textbackslash\{\}to \textbackslash\{\}exists w : \textbackslash\{\}mathsf\{BW\}(W),\textbackslash\{\} \textbackslash\{\}mathrm\{PresInv\}(W,w) \textbackslash\{\}wedge \textbackslash\{\}mathrm\{RealizesTarget\}(W,w).
\item Negative/barrier version: Absent an explicit equivalence witness for 'MCSP witness family', transport toward 'There exist general-purpose magnification theorems (not problem-specific) showing that certain mild lower bounds for distinguishers/sparse-code objects can be amplified into stronger circuit/formula lower bounds, achieving magnification thresholds below previously known lower bounds in some settings.' is not licensed. no\_equivalence\_witness
\item Proof strategy family: equivalence-transport
\item Negative route: Assume no equivalence witness for 'MCSP witness family' aligns the route-side and target-side objects.
\item Failure diagnosis: The route fails if no equivalence witness for 'MCSP witness family' transports 'exponential-lower-bound witness complexity' to the target-side operationalization. Local negative probe: no\_equivalence\_witness.
\item Proof obligations: For every admissible 'MCSP witness family', the target-side operationalization is equivalent to the route-side operationalization in the sense needed to transport 'exponential-lower-bound witness complexity' toward 'There exist general-purpose magnification theorems (not problem-specific) showing that certain mild lower bounds for distinguishers/sparse-code objects can be amplified into stronger circuit/formula lower bounds, achieving magnification thresholds below previously known lower bounds in some settings.'.; Transport obligation: move the invariant across the proposed equivalence for 'carrier-realization' without changing the target claim.; Transport obligation: move the invariant across the proposed equivalence for 'formula-size-transport' without changing the target claim.; Equivalence claim: determine whether the live carrier for 'carrier-realization' is identical to, equivalent to, or only analogous with the target-side object class.
\item Candidate lemmas: Equivalence and coherent transport theorem for MCSP witness family; Equivalence transport lemma for MCSP witness family
\item Independent value: An equivalence theorem for 'MCSP witness family' would independently matter because it identifies when two live operationalizations are mathematically interchangeable rather than merely analogous.
\item Expert gate: 6/6
\end{itemize}

\subsection*{conjpkg.disclayer-p-vs-np-family-categorical-reframing.01}
\begin{itemize}
\item Exact statement: There exists a mediating object built from 'MCSP witness family' such that the route from 'A hardness magnification theorem has the form:"A barely-nontrivial lower bound for an explicit problem Q in a weak model implies a major lower bound/separation in a strong model."Typically, Q is a meta-problem (e.g. a Gap-MCSP variant) and"major"means consequences such as NP not in P/poly.' to 'In refined regimes, lower bounds for Gap-MCSP slightly beyond known frontiers (e.g. N\textasciicircum{} 1+varepsilon -type circuit lower bounds for suitable promise gaps and input length N=2\textasciicircum{}n ) imply major consequences such as NP not in P/poly.' factors through typed morphisms preserving 'MCSP complexity profile in nonuniform-circuits' and making the relevant square commute.
\item Restricted-regime mode: special-case
\item Bridge object: MCSP witness family (family of instances equipped with a monotone witness measure)
\item Bridge definition: 'MCSP witness family' denotes the structured intermediate object required by this branch. It is meant to organize the source-side content 'A hardness magnification theorem has the form:"A barely-nontrivial lower bound for an explicit problem Q in a weak model implies a major lower bound/separation in a strong model."Typically, Q is a meta-problem (e.g. a Gap-MCSP variant) and"major"means consequences such as NP not in P/poly.'. It must preserve the invariant 'MCSP complexity profile in nonuniform-circuits'. Its operational role is 'Recast the blocked frontier as a diagrammatic composition problem: identify objects, morphisms, mediating objects, and the exact commutativity or factorization obligation.'.
\item First nontrivial lemma: For every admissible 'MCSP witness family', the first route-specific realization step preserves 'MCSP complexity profile in nonuniform-circuits' and prepares the passage from 'A hardness magnification theorem has the form:"A barely-nontrivial lower bound for an explicit problem Q in a weak model implies a major lower bound/separation in a strong model."Typically, Q is a meta-problem (e.g. a Gap-MCSP variant) and"major"means consequences such as NP not in P/poly.' to 'In refined regimes, lower bounds for Gap-MCSP slightly beyond known frontiers (e.g. N\textasciicircum{} 1+varepsilon -type circuit lower bounds for suitable promise gaps and input length N=2\textasciicircum{}n ) imply major consequences such as NP not in P/poly.'.
\item Quantifier signature: For every admissible 'MCSP witness family' on the route from 'st.def.hardness-magnification-phenomenon-programme-leve' to 'st.thm.oliveira-pich-santhanam-magnification-near-state'.
\item LaTeX theorem: \textbackslash\{\}forall W,\textbackslash\{\} \textbackslash\{\}mathrm\{Adm\}(W) \textbackslash\{\}to \textbackslash\{\}exists w : \textbackslash\{\}mathsf\{BW\}(W),\textbackslash\{\} \textbackslash\{\}mathrm\{PresInv\}(W,w) \textbackslash\{\}wedge \textbackslash\{\}mathrm\{RealizesTarget\}(W,w).
\item Negative/barrier version: If 'MCSP witness family' cannot be realized with invariant 'MCSP complexity profile in nonuniform-circuits', the route cannot force 'In refined regimes, lower bounds for Gap-MCSP slightly beyond known frontiers (e.g. N\textasciicircum{} 1+varepsilon -type circuit lower bounds for suitable promise gaps and input length N=2\textasciicircum{}n ) imply major consequences such as NP not in P/poly.'. non\_commuting\_route
\item Proof strategy family: constructive-realization
\item Negative route: Assume 'MCSP witness family' cannot be realized while preserving 'MCSP complexity profile in nonuniform-circuits' strongly enough to activate the local target claim.
\item Failure diagnosis: The route fails if 'MCSP witness family' cannot be realized while preserving 'MCSP complexity profile in nonuniform-circuits'. Local negative probe: non\_commuting\_route.
\item Proof obligations: For every admissible 'MCSP witness family', the first route-specific realization step preserves 'MCSP complexity profile in nonuniform-circuits' and prepares the passage from 'A hardness magnification theorem has the form:"A barely-nontrivial lower bound for an explicit problem Q in a weak model implies a major lower bound/separation in a strong model."Typically, Q is a meta-problem (e.g. a Gap-MCSP variant) and"major"means consequences such as NP not in P/poly.' to 'In refined regimes, lower bounds for Gap-MCSP slightly beyond known frontiers (e.g. N\textasciicircum{} 1+varepsilon -type circuit lower bounds for suitable promise gaps and input length N=2\textasciicircum{}n ) imply major consequences such as NP not in P/poly.'.; Commutativity obligation: the route for 'carrier-realization' must commute between source theorem hypotheses, the mediating carrier, and the target theorem hypotheses.
\item Candidate lemmas: Commutative factorization theorem for MCSP witness family; Coherence diagram lemma for MCSP witness family
\item Independent value: A concrete special-case theorem for 'MCSP witness family' would be independently valuable because it converts a broad discovery programme into an inspectable mathematical subresult.
\item Expert gate: 6/6
\end{itemize}

\subsection*{conjpkg.disclayer-p-vs-np-family-categorical-reframing.02}
\begin{itemize}
\item Exact statement: There exists a mediating object built from 'MCSP witness family' such that the route from 'A hardness magnification theorem has the form:"A barely-nontrivial lower bound for an explicit problem Q in a weak model implies a major lower bound/separation in a strong model."Typically, Q is a meta-problem (e.g. a Gap-MCSP variant) and"major"means consequences such as NP not in P/poly.' to 'There exist explicit"meta"problems (including MCSP variants) such that proving certain emph barely-superlinear lower bounds against restricted models (e.g. formulas) would imply major lower bounds for NP in much stronger models (including superpolynomial formula/circuit lower bounds, and in some parameter regimes consequences as strong as NP not in P/poly ).' factors through typed morphisms preserving 'MCSP complexity profile in nonuniform-circuits' and making the relevant square commute.
\item Restricted-regime mode: special-case
\item Bridge object: MCSP witness family (family of instances equipped with a monotone witness measure)
\item Bridge definition: 'MCSP witness family' denotes the structured intermediate object required by this branch. It is meant to organize the source-side content 'A hardness magnification theorem has the form:"A barely-nontrivial lower bound for an explicit problem Q in a weak model implies a major lower bound/separation in a strong model."Typically, Q is a meta-problem (e.g. a Gap-MCSP variant) and"major"means consequences such as NP not in P/poly.'. It must preserve the invariant 'MCSP complexity profile in nonuniform-circuits'. Its operational role is 'Recast the blocked frontier as a diagrammatic composition problem: identify objects, morphisms, mediating objects, and the exact commutativity or factorization obligation.'.
\item First nontrivial lemma: For every admissible 'MCSP witness family', the first route-specific realization step preserves 'MCSP complexity profile in nonuniform-circuits' and prepares the passage from 'A hardness magnification theorem has the form:"A barely-nontrivial lower bound for an explicit problem Q in a weak model implies a major lower bound/separation in a strong model."Typically, Q is a meta-problem (e.g. a Gap-MCSP variant) and"major"means consequences such as NP not in P/poly.' to 'There exist explicit"meta"problems (including MCSP variants) such that proving certain emph barely-superlinear lower bounds against restricted models (e.g. formulas) would imply major lower bounds for NP in much stronger models (including superpolynomial formula/circuit lower bounds, and in some parameter regimes consequences as strong as NP not in P/poly ).'.
\item Quantifier signature: For every admissible 'MCSP witness family' on the route from 'st.def.hardness-magnification-phenomenon-programme-leve' to 'st.thm.oliveira-santhanam-hardness-magnification-for-na'.
\item LaTeX theorem: \textbackslash\{\}forall W,\textbackslash\{\} \textbackslash\{\}mathrm\{Adm\}(W) \textbackslash\{\}to \textbackslash\{\}exists w : \textbackslash\{\}mathsf\{BW\}(W),\textbackslash\{\} \textbackslash\{\}mathrm\{PresInv\}(W,w) \textbackslash\{\}wedge \textbackslash\{\}mathrm\{RealizesTarget\}(W,w).
\item Negative/barrier version: If 'MCSP witness family' cannot be realized with invariant 'MCSP complexity profile in nonuniform-circuits', the route cannot force 'There exist explicit"meta"problems (including MCSP variants) such that proving certain emph barely-superlinear lower bounds against restricted models (e.g. formulas) would imply major lower bounds for NP in much stronger models (including superpolynomial formula/circuit lower bounds, and in some parameter regimes consequences as strong as NP not in P/poly ).'. non\_commuting\_route
\item Proof strategy family: constructive-realization
\item Negative route: Assume 'MCSP witness family' cannot be realized while preserving 'MCSP complexity profile in nonuniform-circuits' strongly enough to activate the local target claim.
\item Failure diagnosis: The route fails if 'MCSP witness family' cannot be realized while preserving 'MCSP complexity profile in nonuniform-circuits'. Local negative probe: non\_commuting\_route.
\item Proof obligations: For every admissible 'MCSP witness family', the first route-specific realization step preserves 'MCSP complexity profile in nonuniform-circuits' and prepares the passage from 'A hardness magnification theorem has the form:"A barely-nontrivial lower bound for an explicit problem Q in a weak model implies a major lower bound/separation in a strong model."Typically, Q is a meta-problem (e.g. a Gap-MCSP variant) and"major"means consequences such as NP not in P/poly.' to 'There exist explicit"meta"problems (including MCSP variants) such that proving certain emph barely-superlinear lower bounds against restricted models (e.g. formulas) would imply major lower bounds for NP in much stronger models (including superpolynomial formula/circuit lower bounds, and in some parameter regimes consequences as strong as NP not in P/poly ).'.; Commutativity obligation: the route for 'carrier-realization' must commute between source theorem hypotheses, the mediating carrier, and the target theorem hypotheses.
\item Candidate lemmas: Commutative factorization theorem for MCSP witness family; Coherence diagram lemma for MCSP witness family
\item Independent value: A concrete special-case theorem for 'MCSP witness family' would be independently valuable because it converts a broad discovery programme into an inspectable mathematical subresult.
\item Expert gate: 6/6
\end{itemize}

\subsection*{conjpkg.disclayer-p-vs-np-family-categorical-reframing.03}
\begin{itemize}
\item Exact statement: There exists a mediating object built from 'MCSP obstruction family' such that the route from 'A hardness magnification theorem has the form:"A barely-nontrivial lower bound for an explicit problem Q in a weak model implies a major lower bound/separation in a strong model."Typically, Q is a meta-problem (e.g. a Gap-MCSP variant) and"major"means consequences such as NP not in P/poly.' to 'In refined regimes, lower bounds for Gap-MCSP slightly beyond known frontiers (e.g. N\textasciicircum{} 1+varepsilon -type circuit lower bounds for suitable promise gaps and input length N=2\textasciicircum{}n ) imply major consequences such as NP not in P/poly.' factors through typed morphisms preserving 'MCSP complexity profile in nonuniform-circuits' and making the relevant square commute.
\item Restricted-regime mode: special-case
\item Bridge object: MCSP obstruction family (family of hard instances or obstructions forcing the target payoff)
\item Bridge definition: 'MCSP obstruction family' denotes the structured intermediate object required by this branch. It is meant to organize the source-side content 'A hardness magnification theorem has the form:"A barely-nontrivial lower bound for an explicit problem Q in a weak model implies a major lower bound/separation in a strong model."Typically, Q is a meta-problem (e.g. a Gap-MCSP variant) and"major"means consequences such as NP not in P/poly.'. It must preserve the invariant 'MCSP complexity profile in nonuniform-circuits'. Its operational role is 'Recast the blocked frontier as a diagrammatic composition problem: identify objects, morphisms, mediating objects, and the exact commutativity or factorization obligation.'.
\item First nontrivial lemma: For every admissible 'MCSP obstruction family', the first route-specific realization step preserves 'MCSP complexity profile in nonuniform-circuits' and prepares the passage from 'A hardness magnification theorem has the form:"A barely-nontrivial lower bound for an explicit problem Q in a weak model implies a major lower bound/separation in a strong model."Typically, Q is a meta-problem (e.g. a Gap-MCSP variant) and"major"means consequences such as NP not in P/poly.' to 'In refined regimes, lower bounds for Gap-MCSP slightly beyond known frontiers (e.g. N\textasciicircum{} 1+varepsilon -type circuit lower bounds for suitable promise gaps and input length N=2\textasciicircum{}n ) imply major consequences such as NP not in P/poly.'.
\item Quantifier signature: For every admissible 'MCSP obstruction family' on the route from 'st.def.hardness-magnification-phenomenon-programme-leve' to 'st.thm.oliveira-pich-santhanam-magnification-near-state'.
\item LaTeX theorem: \textbackslash\{\}forall W,\textbackslash\{\} \textbackslash\{\}mathrm\{Adm\}(W) \textbackslash\{\}to \textbackslash\{\}exists w : \textbackslash\{\}mathsf\{BW\}(W),\textbackslash\{\} \textbackslash\{\}mathrm\{PresInv\}(W,w) \textbackslash\{\}wedge \textbackslash\{\}mathrm\{RealizesTarget\}(W,w).
\item Negative/barrier version: If 'MCSP obstruction family' cannot be realized with invariant 'MCSP complexity profile in nonuniform-circuits', the route cannot force 'In refined regimes, lower bounds for Gap-MCSP slightly beyond known frontiers (e.g. N\textasciicircum{} 1+varepsilon -type circuit lower bounds for suitable promise gaps and input length N=2\textasciicircum{}n ) imply major consequences such as NP not in P/poly.'. non\_commuting\_route
\item Proof strategy family: constructive-realization
\item Negative route: Assume 'MCSP obstruction family' cannot be realized while preserving 'MCSP complexity profile in nonuniform-circuits' strongly enough to activate the local target claim.
\item Failure diagnosis: The route fails if 'MCSP obstruction family' cannot be realized while preserving 'MCSP complexity profile in nonuniform-circuits'. Local negative probe: non\_commuting\_route.
\item Proof obligations: For every admissible 'MCSP obstruction family', the first route-specific realization step preserves 'MCSP complexity profile in nonuniform-circuits' and prepares the passage from 'A hardness magnification theorem has the form:"A barely-nontrivial lower bound for an explicit problem Q in a weak model implies a major lower bound/separation in a strong model."Typically, Q is a meta-problem (e.g. a Gap-MCSP variant) and"major"means consequences such as NP not in P/poly.' to 'In refined regimes, lower bounds for Gap-MCSP slightly beyond known frontiers (e.g. N\textasciicircum{} 1+varepsilon -type circuit lower bounds for suitable promise gaps and input length N=2\textasciicircum{}n ) imply major consequences such as NP not in P/poly.'.; Commutativity obligation: the route for 'carrier-realization' must commute between source theorem hypotheses, the mediating carrier, and the target theorem hypotheses.
\item Candidate lemmas: Commutative factorization theorem for MCSP obstruction family; Coherence diagram lemma for MCSP obstruction family
\item Independent value: A concrete special-case theorem for 'MCSP obstruction family' would be independently valuable because it converts a broad discovery programme into an inspectable mathematical subresult.
\item Expert gate: 6/6
\end{itemize}

\end{document}
